\documentclass[12pt]{book}
\usepackage[T1]{fontenc}
\usepackage{lmodern}
\usepackage{amsmath}
\usepackage{amssymb}
\usepackage{xcolor}
\usepackage[normalem]{ulem}
\usepackage{graphicx}
\usepackage{cancel}
\usepackage{soul}
\usepackage{float}
\usepackage[autostyle]{csquotes}
\usepackage{fontspec}
\usepackage{tikz}
\usetikzlibrary{automata,positioning,backgrounds,calc,arrows.meta,shapes.geometric,fit,shadows}
\usepackage{tabularx}
\usepackage{longtable}
\usepackage{booktabs}
\usepackage{array}
\usepackage{calc}
\usepackage{multirow}
\usepackage{minted}
\usepackage{caption}
\usepackage{standalone}
\usepackage{hyperref}
%\setmainfont{Times New Roman}

\usepackage[style=apa, backend=biber]{biblatex}
\addbibresource{/Users/tio/Documents/GitHub/September_UMEDCA/references.bib}

\newenvironment{abstract}
    { % Code executed at the beginning of the environment
      \begin{center} % Center the "Abstract" title
        \textbf{Abstract} % Make the title bold
      \end{center}
      \vspace{\baselineskip} % Add some space after the title
      \begin{quotation} % Indent the abstract text
      \small % Optional: Make the abstract text slightly smaller
    }
    { % Code executed at the end of the environment
      \end{quotation} % End the indentation
      \bigskip % Add some vertical space after the abstract
    }

% Operators
\newcommand{\SBox}{\Box_{\mathsf{S}}}
\newcommand{\OBox}{\Box_{\mathsf{O}}}
\newcommand{\NBox}{\Box_{\mathsf{N}}}
\newcommand{\SDia}{\Diamond_{\mathsf{S}}}
\newcommand{\ODia}{\Diamond_{\mathsf{O}}}
\newcommand{\NDia}{\Diamond_{\mathsf{N}}}

\newcommand{\SBoxUp}{\Box_{\mathsf{S}}^{\uparrow}}
\newcommand{\SBoxDown}{\Box_{\mathsf{S}}^{\downarrow}}
\newcommand{\SDiaUp}{\Diamond_{\mathsf{S}}^{\uparrow}}
\newcommand{\SDiaDown}{\Diamond_{\mathsf{S}}^{\downarrow}}

\newcommand{\OBoxUp}{\Box_{\mathsf{O}}^{\uparrow}}
\newcommand{\OBoxDown}{\Box_{\mathsf{O}}^{\downarrow}}
\newcommand{\ODiaUp}{\Diamond_{\mathsf{O}}^{\uparrow}}
\newcommand{\ODiaDown}{\Diamond_{\mathsf{O}}^{\downarrow}}

\newcommand{\NBoxUp}{\Box_{\mathsf{N}}^{\uparrow}}
\newcommand{\NBoxDown}{\Box_{\mathsf{N}}^{\downarrow}}
\newcommand{\NDiaUp}{\Diamond_{\mathsf{N}}^{\uparrow}}
\newcommand{\NDiaDown}{\Diamond_{\mathsf{N}}^{\downarrow}}

% IS Notation
\newcommand{\Inc}{\mathit{Inc}}
\newcommand{\I}{\mathcal{I}}
\newcommand{\Trace}{\mathcal{T}}

% Theorem environments
\newtheorem{axiom}{Axiom}
\newtheorem{principle}{Principle}

% Naming and Sublation
\newcommand{\name}[1]{\ulcorner #1 \urcorner}
\newcommand{\sublate}[1]{\cancel{#1}} % Simplified for pedagogical clarity

\title{Understanding Mathematics as an Emancipatory Discipline: A Critical Theory Approach}
\author{Theodore M. Savich}
\date{\today}

% The 'standalone' package allows you to include other .tex files
% that have their own preambles. This the key to your request.



\begin{document}

% ========== Content from: prelude_2 ==========
\maketitle


\frontmatter

\tableofcontents

\mainmatter


\part{Prelude: Built to Break}

%\chapter*{Built to Break}



\begin{verse}
You lit up the room like you swallowed the moon.\\
Flashing that mischievous smile.\\
Saying, ``Build higher, just a little more higher.''\\
A tumbling is coming on soon.\\
I said, ``I'm worried. This brick spells demise.\\
We've built it up ever so tall.''\\
But I picked you up and you put on the brick\\
And it fell. Like we both knew it'd fall.


\emph{Chorus} \\
Build something for the breaking\\
Tall thin walls, shivered and quaking.\\
\(\mathcal{M}\), it's beautiful, breaking with you.\\
\(\mathcal{M}\) the lightning that swallowed the moon.\\
\emph{Verse 1; Built to Break}
\end{verse}

\begin{abstract}

This chapter explores the intersection of mathematics and personal experience, introducing key themes that recur throughout the text. The argument begins from the position that mathematics and self-understanding share a common structure: both involve recollecting the self through the otherness of objects and social norms. A central metaphor emerges from a child's observation about a microphone, introducing the concept of \emph{divasion}---a simultaneous inside/outside relation that challenges conventional understandings of set theory. The chapter contains foundational stories about struggles with algebra, teaching experiences, and encounters with students' profound questions. These narratives structure subsequent explorations of how mathematics education can honor the subjective experience of learning. Songs and poems appear throughout, connecting the author's theoretical commitments with lived, felt experience and establishing the book's autoethnographic methodology.

\end{abstract}

\section{Wound and Possibility: The Misrecognition of Mathematics}\label{wound-and-possibility-the-misrecognition-of-mathematics}

My mother claims my first word was ``battery,'' but I think she means that is the first word she remembers me speaking. I must have said ``momma'' and ``dadda'' before ``babbery,'' but I have no recollection: I can only repeat my mom's story. Recollection is not just fallible; its flaws and features shape us. Here, my first word---true or false---kicks off a series of recollections that structure this text. Beginnings are tough.

I used to have a grey calculator that my great aunt gave me; she taught math in upstate New York. It had a square-root button along with \((+, -, \times, \div)\), and a symbol for changing the sign of a number. On family road trips, I would bring it along to play. I recall how fascinating it was to enter numbers like \(0.9\), \(0.8\), \(1.1\), and \(5\) and take their square roots over and over. Mashing the \(\sqrt{}\) button \emph{iteratively} usually ended in a number very close to 1. Sometimes my calculator would display a little E, which I found intriguing. I tried to make it do that in every way I could. I liked taking the square root of numbers like 9 and \(\frac{1}{2}\) over and over, noticing how each approached 1 from above or below until the calculator ran out of digits to express the difference.


But that button-mashing skill did not translate into an interest or ability in school mathematics. In third grade, I had to take my first standardized math test. I got in the bottom quartile. Ever since then, I have wrestled with math as a subject. I didn't like being tested. I didn't like how it made me feel stupid. I didn't like how it reduced me to a number.



\subsection*{When Mathematics Doesn't Make Sense}\label{when-mathematics-doesnt-make-sense}

My transition to middle school was rocky. I left early on my first day, doubled over in terrible stomach pain. I told my mom, ``the adults \emph{yelled} at the kids!'' It was so stressful that I got physically ill.

Much later, I got the opportunity to watch workers raze the building. That was nice. The new building is better. 

By eighth grade, mathematics seemed totally nonsensical. In Algebra I, we learned that `finding the slope is easy and fun; just remember ``rise over run.''\,' I did not find it easy or fun. I could remember ``rise over run,'' but I was working on a printed piece of paper. `Rise' is not a formal term; it is supposed to represent the amount of change in the dependent variable as the independent variable (the `run') changes. But the mnemonic emphasized the vertical change divided by the horizontal change. Because I was working on a piece of paper, there was no `up' or `down.' Verticality could have been the third axis sticking out of the paper or the `top,' `bottom,' `left,' or `right' of the page depending on how my worksheet was oriented (see figure \ref{fig:slope_cartesian}). When the teacher, Mrs.~\(\mathcal{J}\) faced me, her left and right were different from when she faced the board. I couldn't figure out the \emph{reference frame} I was supposed to be using, so I just cheated off of someone else's paper. But I couldn't even CHEAT properly!

Mrs.~\(\mathcal{J}\) confronted me in the hallway about how she could tell that the work was not my own since none of the work matched the answers. She called my mom and told her I cheated on the test. She was also generous and explained that I might be struggling with algebra because ``Tio is very spatial.'' My mom and I still laugh about how she said ``spatial,'' as people from some regions of Indiana pronounce ``spatial'' as ``special.''\footnote{This may seem like a side comment, but the drift in vowel sounds, coupled with the ambiguity in homophones, with a sprinkling of jocular ableism in the re-telling (``special'' being a euphemism for disability) are all relevant to the themes of this book.}

\begin{figure}[h!]
    \centering
    \includegraphics[width=0.8\textwidth]{/Users/tio/Documents/GitHub/September_UMEDCA/images/prelude_cartesian_rise_run.pdf}
    \caption{\textit{Note.} The slope of a line can be formalized to be unambiguous, but teaching and learning the concept on a piece of paper invites confusion.}
    \label{fig:slope_cartesian}
\end{figure}

This experience with slope shapes the whole manuscript. I will articulate an \textit{in}formal, or \textit{material} mathematics, where formal explications are built to break. While this approach suggests an order of explanation that begins with embodied experiences, I don't offer hard and fast rules about pedagogy. 

I howled and raged for a few years until I took pre-calculus. Struggling at school, I also got in a fair bit of trouble of the sort that embarrassed my dad who was a lawyer. Facing real consequences for not complying with expectations, I learned from Kurt Cobain and the Meat Puppets: ``I don't have to think; I only have to do it; the results are always perfect; and that's old news.'' That is, I discovered that the secret to doing well in math class was to not try to understand anything; just do what you're told to do, exactly, and the results are always perfect. I went from routinely getting Cs, Ds, and the occasional F on math assessments to getting A+ scores. That `success' felt a bit hollow. My first role as a math educator was as a peer tutor that spring, when I gave that advice to a friend who was struggling.

This moment of self-erasure raises a fundamental question: What is the cost of compliance? Without understanding, compliance seems damaging. But some form of compliance with norms seems inevitable. Further, compliance with some norms seems to enable expressive freedom.

\subsection*{The Journey from Hatred to Appreciation to Trepidation}\label{the-journey-from-hatred-to-appreciation-to-trepidation}

The reward for compliance was freedom from my institutional circumstances. The spring semester of my junior year of high school, I got to leave high school early to take an introduction to philosophy course at Indiana University (IU). I loved that class. The graduate instructor was passionate about Kierkegaard. I would sit on my porch drinking tea with a high schooler friend who was also taking the class. We would debate Descartes and Kierkegaard. I finally felt like I belonged. But, like a duck underwater, I was paddling furiously to keep up. I sat in my parents' basement at night, reading all the texts aloud into a tape recorder and listening to them over and over. But it was worth it.

The fall semester of my senior year of high school, I enrolled almost solely at IU. The course that left the most lasting impression was a philosophy of mind course taught by an analytic philosopher. When we read Tim van Gelder's piece \parencite{vanGelder1995}, I was blown away. van Gelder opened a world where mathematics could be used to express the motion that shaped my conscious experience. I had no idea what the differential equations in that piece meant, but the symbols! They were so beautiful and strange! I'm still drawn to the aesthetics of mathematical notation, even if they do not always make complete sense to me.

I chose Earlham College in Richmond, Indiana for my undergraduate years. From my readings in analytic philosophy, I inferred that I could not understand philosophy without understanding mathematics. That was a bad inference, but it helps explain why I studied math at Earlham. For four years, I woke for math classes that started at 8 a.m. I did plenty of partying and playing music, but I also spent a huge amount of energy trying to master calculus, analysis, linear algebra, and abstract algebra.

In my sophomore year, I learned about how the empty set ($\emptyset$, or \(\{\}\)) can be used to define numbers. I was also practicing silent meditation as part of a budding interest in Quakerism. I had a spiritual experience, sitting in a tree on the outskirts of campus, the evening after I learned about these ideas. The role of nothingness in mathematical reasoning resonated deeply with the kind of meditative consciousness I was trying to cultivate in myself at the time. Nothingness surrounded me everywhere I looked.

This resonance between mathematical nothingness and meditative consciousness suggested a grounding. Yet, this encounter with the foundations of number already hinted at a dynamic that would become central to this work. What presents itself as origin---the empty set grounding the number system---discloses its own reliance on what it supposedly founds. The ``nothingness'' that grounds the system is itself constituted by the operations of the system it enables. The ground, therefore, does not remain static when engaged; it recedes in the very act of engagement. This movement, where the foundation withdraws precisely when approached, is not a failure of the system but its constitutive structure. The null is the unrepresentable condition that makes representation possible.

Throughout this work, `ground' names a paradox: what enables must withdraw to function as enabling. This is not a failure of language but the structure of enablement itself. 


Later, when I took real analysis with Tim McLarnen, my alienation from mathematical reasoning was finally alleviated. He said, in a quote that I later printed out as a poster for my high school classroom, ``never underestimate the value of focused play.'' We played a lot. We picked up and put down various axioms to learn how different combinations of rules enabled different proofs. The experience was so expressively \emph{empowering} that I silently asked myself, ``Why couldn't mathematics have been that way from the start?!'' Compliance with the axioms we chose enabled a kind of expressive freedom. I still think it is a shame that most students do not get such opportunities to choose what sentences they are explicitly committed to until college. However, I have read enough history of mathematics education to know that efforts to do so, such as the New Math movement informed by the Burbaki group, are terribly hard to implement. 

As an undergraduate research intern in computational mathematics at Oak Ridge National Laboratory, I worked with Dr.~Leonard Gray. The details are complicated, but it involved working with algorithms for solving differential equations. Since I didn't really know much about applied math and couldn't program very well, the programs I wrote tended to run indefinitely. That meant that I spent a fair bit of time looking out the window at the groundhogs playing on grass-covered bunkers marked with radioactive contamination warnings, waiting for my bad code to terminate.

However, there was a dark side to that reality. Oak Ridge was founded as part of the Manhattan Project, and as I worked on abstract algorithms, I worried about whether continuing down the path of a research mathematician was compatible with my commitment to pacifism. One of my peers was working on the military side of the facility. We'd hang out drinking beer, and he'd talk about the `kinetic effects' of using tungsten instead of depleted uranium in various munitions. While I worked on the civilian side of the facility, the way the macabre met with the casual made me ambivalent about my path. Did I want the life of a research mathematician? I worried about how mathematics enables violence at unfathomable scale. I was afraid of doing harm.

When I returned to Earlham College, I prepared for the entrance exam for Ph.D.~programs in pure mathematics. But what I really loved doing was playing in my band, making ceramics, and having fun. I arrived late to the exam, did poorly, and got angry about it. To my regret, I didn't ever talk through my ambivalence about mathematics with my mentors.

\subsection*{Mathematics as Gatekeeping}\label{mathematics-as-gatekeeping}

Without graduate school as an option, I decided to hang out in Richmond, Indiana for a year while I waited for my then girlfriend to graduate. I had two jobs that summer that helped me understand a different role mathematics plays in society.

The first was to prepare a report on math education for the president of Earlham College. My report was a mess, but I found the work interesting. Part of the reason I struggled with the report was that I wasn't familiar with the theories involved in math education. Chauvinistic about pure mathematics, I dismissed the informal mathematics I encountered.

The other summer job was working at Family Video, a local video rental chain. Richmond was in an economic recession at the time, and when I applied to work there, the line of applicants stretched around the building. When I got my turn to interview, they handed me a math test. It wasn't like the entrance exam I had recently bombed. It was very basic arithmetic and pre-algebra. I was shocked that being able to do math was the primary criterion for getting the job. The manager fawned over how well I did on the test, but I was flummoxed. ``They hired ME, because I could do arithmetic?! Okay\ldots{}'' I worked there for a few months until they were about to fire me. I didn't care about upselling people on popcorn and candy and didn't really like watching movies that much.

That experience led me to reflect on the role of mathematics in society. Why were people who definitely would have appreciated the job that I scorned removed from consideration? Shouldn't the Family Video people have tried to hire people who cared a little bit about making money and watching movies? I began to think about the people who wrapped around the block and who never got an interview as kindred spirits who never learned to comply. They were punished for that lack of compliance with economic and expressive impoverishment.

At summer's end, I got a job at Ivy Tech, the community college system in Indiana. My job was to teach algebra and pre-algebra at the satellite campuses. Across the street from the feedstore, in the attic of the electric company, I taught folks like those who wrapped around the block but who never got an interview. I recognized their struggles as mirrors of my own.

I did well as a teacher in that context. Students liked me well enough and the work was well-defined. When a job running the tutoring center for the Indianapolis campus of Ivy Tech opened up, I went for it. My then-girlfriend and I moved to Indianapolis after I got the job.

Again, I think my ability to do mathematics put me in a position where other abilities probably should have been more heavily weighted. Part of running the tutoring center was actually managing people. Interpersonal conflicts, scheduling, budgetary meetings, and more did not draw on my mathematical reasoning. But I could make data tell a story, so I was able to accomplish some goals. I went up for various jobs in the administration but at 25 years old, I wasn't really ready or that interested in climbing the ladder or feathering my nest. After breaking ground on a brand-new building I helped design to house the tutoring center, I left Ivy Tech. 

I had decided that I wanted to teach math. I did not want to teach it as I was taught. I wanted to teach it in a way that empowered students from the beginning. The often repeated phrase that informed a lot of my experiences learning music, that `you have to learn the rules before you can break them,' was not going to work for me. I wanted to teach math in a way that allowed students to break the rules from the beginning. I wanted to teach math in a way that allowed students to recognize themselves in the mathematics, rather than as an alienated subject of the mathematics. I wanted to teach math in a way that allowed students to understand themselves as \textit{empowered} by the mathematics.

But I also wanted to protect `my' rule-breakers from the consequences that follow from being `bad at math.' With more than a hint of a savior complex, I did a transition to teaching program and got a job teaching, of all things, high school mathematics. 



\section{The Determinate Question: What Even Is Two?}\label{the-determinate-question-what-even-is-two}

The moment that crystallized the failure of my teaching philosophy came during my first week on the job, in a conversation with a student. $\mathcal{I}$, a fifth-year senior in an urban Indianapolis high school, was struggling. He could not graduate because he could not pass a high-stakes standardized algebra test. The evening prior to my first day of school, one of the fifth-year seniors was killed in a car accident. His SUV rolled over while he was fleeing the police. Some students in my class were crying as they shuffled into my classroom. When I asked what was wrong, I didn't know how to respond. Rather than processing with them, I decided to jump right into algebra. As I lectured on $x$ and $y$, probably talking about how `easy and fun' it all was, $\mathcal{I}$ looked up and said, ``Mr. Savich, \textit{what even is two?}''

Ah-ha! I thought. I knew the answer to \textit{that} question! I began to explain the empty set. I left my experience of sitting in the boughs of a sycamore tree unspoken and gave the pure, formal answer. 

And I failed, in a profound way, to address the question. This moment highlights a core inquiry of this work: When confronted with deep questions, why do formal, abstract answers fail us? What had been beautiful for me, sitting in the boughs of a tree at Earlham College, must have felt crushing. Here was a young adult who was not allowed to move on because of mathematics, and I was nattering on about nothingness! The stakes were real for him.

The formal answer---the definition of two derived from the empty set---failed because it did not address the student's underlying need. His question, ``What even is two?'', was not a request for abstract definition but a bid for recognition. He was seeking to locate himself within the mathematical framework that was currently excluding him from graduation. My response, rooted in the beauty of formal structure, missed the existential weight of his inquiry. This failure, this negation of connection, became the generative force for reconstructing mathematics not as a collection of objects, but as a process of recollection---finding the self within the structures of reason. The inadequacy of the formal response made palpable the necessity of a mathematics that acknowledges the first-person position as essential to mathematical knowing. The collapse of the previous framework cleared the ground for a different kind of engagement, one that begins not with axiomatic certainty but with the productive negativity of recognized failure.

As I prattled on, I could feel his desire to connect---to be known and to know---withdraw. ``What the hell am I doing,'' I thought, as I realized that my answer had no bearing on his question. ``If I can't even explain what 2 is, what is the point?'' The stories I had told myself about my role as a teacher felt deeply dishonest. I was just a cog in the same machine I thought I was fighting against. I lost my sense of purpose. The question didn't break me, but I started wobbling. A brush with tragedy a few years later would shatter my sense of purpose entirely. The end result was that I left teaching high school after three years, my marriage ended, and I moved back to my hometown to live with my parents and lick my wounds. 

Breathing into that space of brokenness, I share a song I wrote at the time called \textit{The City Dark}. The addressee of the song is ambiguous, as ``you'' often is in speech and song. Self-directed bitterness tends to spill out all around. 

\begin{verse}
The city dark is full of light\\
A gold dusk dome of electric night\\
I don't sleep anymore, and barely write\\
Slumped halo hanging lowly\\


God has cold hands, taking you from me.\\
Leave me lonely, lonely, lonely.\\
But every morning it's hard to believe\\
In forever-shadows lengthening\\
New pages born, like blisters torn\\
Tasting long off summer dreams\\
A sleeping pill to loaf all day\\
Scraping dry the parts unseen.\\
Seems only broken words say what they mean.\\


The city dark is full of light\\
A gold dusk dome of electric night\\
I don't sleep anymore, and barely write\\
Slumped halo hanging lowly\\


\textit{Bridge}
Oh but don't cry out begging forgiveness\\
Don't bleed out your big baby bigness\\
Don't you dig through the hole\\
That you threw it all in\\
Don't you choke that bitter refuse down\\
Just to vomit out your fill on the floor\\
Hold it in more.

The city dark is full of light\\
A gold dusk dome of electric night\\
A tree don't grow until it's out of sight\\
And we were made to grow on starlight.\\

\end{verse}



\section{Systematic Analysis: The Search for Critical Mathematics}\label{systematic-analysis-the-search-for-critical-mathematics}

I returned to Bloomington as a Ph.D.~student in math education. In the divorce, I somehow ended up with 25 blank composition books. I filled them with bitter poetry and songs, burned them, and bought more. I bought a book on catastrophe theory \parencite*{poston_catastrophe_2012}, hoping it would explain how awful I felt.

I didn't really want to do the program, but it seemed like a way to maintain the image of forward momentum, while the actual growth was happening underground through the act of writing without inhibition. I apologize to my professors that first year. I did not want to be there and it was apparent. I earned my first and only ``F'' in a college course, despite the best efforts of the professor. The assignment that I could not complete was writing a 5-page ``teaching philosophy.'' Given how poorly my theoretical understanding of teaching had served me as a teacher, every time I set pen to paper to write the thing was like scratching an open wound. I could not do it.

But then I started studying critical theory with Phil Carspecken. I recall asserting rather forcefully that `everything can be doubted.' Phil reminded me that, in speaking, I must not be doubting \emph{everything}. He pointed out that to speak was to \emph{assume} an Other who \emph{might} understand what I was speaking. That understanding blew my mind. All of the inner speech of doubt and self-loathing was addressed to someone who \emph{might} understand what I was so angry and sad about. Even alone, I was not alone. I cannot write, speak, or think without assuming that someone else might understand what I am saying.

I learned about George Herbert Mead's distinction between the \{I\} and the ``me'' \parencite{mead1934social}. The \{I\} is the source of action. The ``me'' is the self-as-recognized. There is often felt sense of difference between those two aspects of self, and, at times, a paradoxical unity between them offers a distinct felt-sense. My formal training in mathematics had not prepared me to talk to students as subjects---as people with their own experiences and ways of making sense.


Since it is so central to my story and what follows, I introduce Grover to do some heaving philosophical lifting. There's a children's book, \emph{The Monster at the End of This Book} \parencite{stone_monster_2003}, where the eponymous monster, Grover, is informed that there is a monster at the end of the book (Figure \ref{fig:grover}). Grover is frightened by this and begs the reader to stop reading. The cruel toddling reader keeps reading, though, with Grover becoming increasingly distraught. But at the end of the book, Grover recognizes himself as the moster at the end of the book: ``Well, look at that! This is the end of the book, and the only one here is\ldots{} ME. I, lovable, furry old Grover, am the Monster at the end of this book'' \parencite[26]{stone_monster_2003}.



\begin{figure}[h!]
\title{\textit{The Monster at the End of This Book}}
\includegraphics[width=\textwidth]{/Users/tio/Documents/GitHub/September_UMEDCA/images/Grover.jpeg}
\caption{\textit{Note. }\emph{Frames from} The Monster at the End of This Book}
\label{fig:grover}
\end{figure}

Sophisticates may recognize the story of Oedipus's unwitting self-banishment for the murder of Laius in the problem of self-recognition that Grover encounters in Stone's book. The problem is that the self-as-recognized, the ``me'', is often distinct from the \{I\}, understood as the source of action or locus of ``power, creativity, and freedom'' \parencite[97]{Carspecken1999}. Both Grover and Oedipus appear to have only partially understood themselves. They did not completely misunderstand themselves. Oedipus knew his authority to banish; Grover knew his fear and, later, understood that he was the `loveable' kind of monster. Complete misrecognition is outside human experience. They only partially misrecognized themselves---or was it us who misrecognized them? In any case, the distinction between the \{I\} and the ``me'' will keep coming up. 


Different people have different sensitivities to contradictions in how others regard them. It seems quite normal, even necessary, to take on different roles in different social situations. But the problem becomes acute when those roles conflict with each other or with deep commitments about what it means to be a good person. When fulfilling a role requires violating some aspect of self that feels inviolable, the tension can become unbearable.

In short, I found in Phil's critical theory some answers to the existential questions that I might have been able to answer if I had continued studying Kierkegaard instead of falling for pure mathematics.

After a few weeks, I decided to linger after class. I asked Phil what a \emph{critical mathematics} would look like. We chatted, and I was hooked. I was simultaneously involved in Dr.~Erik Jacobson's research project that studied misconceptions in the domain of fractions, decimals, and rational numbers. I was also enrolled in an arts-based educational research class that Dr.~Gus Welsek taught. I began to analyze the student work samples that expressed some kind of flawed mathematical reasoning as if they were artistic expressions. Rather than asking ``what is \emph{wrong} with this work/student?'' I began to ask, ``what is this student trying to say?''

\subsection*{Understanding and Being Understood}\label{understanding-and-being-understood}

In Phil's class, I learned about communication as more than just exchanging information. Every meaningful act---whether a gesture, a statement, or even silence---implicitly references what we might call \emph{validity claims} \parencite{Habermas1984}. Think of these as the unspoken conditions under which what you say or do could be understood, accepted, or challenged by another person.

There are three basic types of these claims that structure all communication \parencite{Carspecken1999}. \emph{Objective} claims concern facts about the shared external world that multiple people can observe---like ``The calculator has a square-root button.'' \emph{Subjective} claims concern the speaker's inner world of intentions and feelings, to which the speaker has special access---like ``I find mathematics beautiful.'' \emph{Normative} claims concern what is right, appropriate, or legitimate within a shared social world---like ``Teachers should not humiliate students.''

Understanding this structure transforms how we engage with complex issues. For example, debates about many social issues often mix up felt subjective experience with objective claims to make arguments about social practices---without acknowledging these are three distinct types that require different forms of support.

When Phil reminded me that to speak is to assume someone who might understand, he was pointing to something philosophers call \emph{intersubjectivity}---shared understanding and mutual recognition among people \parencite{Carspecken1999,Habermas1984}. Even my inner monologue of self-loathing assumed someone who could understand what I was angry and sad about. This assumption that someone might grasp my meaning is a necessary condition for communication itself.

The distinction between my own inner experience and our shared understanding proved crucial. While our inner life is shaped by social relations and language, it's not simply determined by them. It becomes knowable and representable only through our relationships with others. Similarly, what we call ``objective facts'' in mathematics require shared practices of verification, shared languages for expression, and shared norms for what counts as proof.

This framework helped me understand that mathematical errors aren't just wrong answers. They're often attempts to participate in shared mathematical practices. They are what my mentor Barbara Dennis calls ``bids'' to be recognized as a mathematical thinker \parencite{korth1998reformulation}. 

\subsection*{Critical Mathematics Begins with Listening}\label{critical-mathematics-begins-with-listening}

Critical mathematics, as I articulate it in this manuscript, begins with recognition that mathematics is fundamentally about communication. Mathematical truths emerge through shared practices of proving, checking, challenging, and refining. When a student produces what teachers call an ``error,'' they are often operating within a different way of understanding, making different assumptions, or working within a different framework than the teacher expects. ``Does not meet expectations''---a common phrase on many rubrics---often says more about the quality of the expectations than the quality of the work. 

When I began analyzing student mathematical reasoning by asking ``What is this student trying to say?'' rather than ``What is wrong with this student?'', I was recognizing that mathematical meaning emerges through dialogue. The student's work, even when ``incorrect,'' expresses an attempt to make sense, to participate in shared mathematical practices, and to be recognized as a mathematical thinker.

Later, I recognized that treating mathematical work as artistic expression didn't quite capture something important: mathematics also involves getting things right or wrong in ways that matter. I became curious about the \emph{experience of error} itself. Since I have such a deep, rich history of getting it wrong, I figured the experience of error could be a foundation for understanding mathematics. Could mathematical truths be known through the exclusion of error? That idea sounds strange to mathematicians, but it's relatively normal in teacher education \parencite{Brandom:2019aa}. I recently met a kindergarten teacher who told her mentees, ``you are here to make mistakes.''

\section{The Dialectical Turn: Divasion}\label{the-dialectical-turn-divasion}

After I defended my dissertation, I felt a deep sense of openness. I met my wife, \(\mathcal{A}\), and her two daughters \(\mathcal{M}\) and \(\exists\). In the fall of 2023, \(\mathcal{M}\) was four years old. \(\mathcal{A}\) and I were just dating, but I was excited to be a part of the kids' lives. I discovered that \(\mathcal{M}\), \(\exists\), and I like to write songs and we like to work on them together. \(\mathcal{M}\) and I were sitting in \(\mathcal{A}\)'s basement trying to write a song. The first thing she says in this interaction is that she wants to write a ``pattern song.'' The object she is referring to is the part of a microphone stand where the mic is held by the stand. The microphone extends past the clip, but is also clipped in, so part of it is inside of the clip and part of it is outside of the clip.

\begin{quote}
\(\mathcal{M}\): I'm going to do a pattern (I want to write a pattern song)

Tio: A pattern? Tell me what you mean, or just do it

\(\mathcal{M}\): Outside In! Outside In!

Tio: More! Say more! What is this?

\(\mathcal{M}\): The microphone is outside of this (pointing at the clip that is part of the mic stand) then it's inside of it.

Tio: What does that mean?

\(\mathcal{M}\): Um\ldots I don't know\ldots{}

Tio: It's kind of interesting, though, isn't it? That things can be inside and outside of each other. Sometimes at the same time.

\(\mathcal{M}\): It is?

Tio: Yeah! Well, this part is inside and this part is outside. So we can't say that the microphone is inside and we can't say that it's outside because it's both! That's a contradiction.

\(\mathcal{M}\): No, it's not. It's all on the other\ldots{} of all of them.

Tio: What's all over all of them?

\(\mathcal{M}\): That means that all of them is\ldots all of it\ldots{} that it divaded.

Tio: Divaded?

\(\mathcal{M}\): Yeah. Divaded means that\ldots um\ldots{} It means that it's inside the thing and then it's outside the thing.

Tio: Woah! Cool! I like that. So it's different than divided. Because dividing something sometimes means splitting but you're talking about divADed\ldots that's cool. It's inside and outside?

\(\mathcal{M}\): Yeah.

Tio: Tell me more!

\(\mathcal{M}\): It's mechanical. It's really not something. But I'm talking about when you grab something, it makes it divaded.
\end{quote}

\begin{figure}[h!]
    \centering
    \includegraphics[width=0.8\textwidth]{/Users/tio/Documents/GitHub/September_UMEDCA/images/prelude_divaded.pdf}
    \caption{\textit{Note.} The regions \(x\) and \(y\) may be inside, outside, separated, sharing a boundary, or divaded with respect to each other.}
    \label{fig:divaded_regions}
\end{figure}

\(\mathcal{M}\) noticed that there is a pattern, or rhythm in considering \emph{outside, in}. Most people would conceive of a microphone in a clip as a static spatial relationship. People regularly experience objects that are inside and outside other objects. A person walks through a door; a folder hangs out of a backpack; a potted plant breaks the surface of the dirt while growing its roots in the soil. But the movement is important. The body is divaded by the breath in rhythmic movement. Many fears involve divasion, like a needle piercing the skin. Many pleasures involve divasion, from the embrace of a lover to the feeling of consuming a delicious meal. Birth and death divade the subject.

In this interaction, the rhythm of the pattern song was flattened into a spatial concept. The song was arrested before she ever got a chance to sing it. Perhaps I erred in calling her concept a `contradiction.' The issue of what constitutes a contradiction is important for mathematics.

In any case, the interaction does say something about how paradoxes or contradictions are experienced. I cannot seem to hold a contradiction or paradox within one moment of awareness. Sentences like ``this sentence is false,'' toggle between truth and falsity, appearing to create temporal movement. Perhaps that rhythm is part of the ``pattern song'' she was after. In the face of a potential contradiction, she invented a word, \emph{divaded}.

\(\mathcal{M}\)'s invention, \emph{divaded}, articulates some deep crack in the foundations of mathematics. Formal set theory is based on the premise that an element of a set is strictly either inside or outside a set---a premise that enforces a binary logic that the structure of `divasion' challenges. The microphone inhabits both positions; it inhabits both positions in a way that defies simple categorization. This simultaneous inside/outside relation is not an anomaly but a manifestation of a deeper structure---the mediation pattern---where apparent opposites require each other for determination. \(\mathcal{M}\)'s observation articulates bodily what we will later encounter conceptually: that the foundations of mathematics, like the foundations of experience, are not built on rigid separation but on dynamic interrelation. The consequences of pruning these more basic ways of experiencing space are evident in foundational crises throughout the history of mathematics. Rather than repairing the effects of this pruning with complicated formal work, I suggest backing up. The foundations of mathematics can instead be approached by attending to how children experience the world.

\section{Reflection: Methodology and the Reader's Role}\label{reflection-methodology-and-the-readers-role}

\(\mathcal{M}\)'s observation that the microphone divades the clip---simultaneously inside and outside---points toward something fundamental about this book's methodology. Just as the microphone cannot be understood as simply inside or outside, my own experience cannot be neatly separated into ``subjective story'' versus ``objective theory.'' This work attempts something more difficult: to treat personal experience as theoretical resource while recognizing that all theory emerges from lived practice.

\subsection*{Critical Autoethnography: Method and Risk}\label{critical-autoethnography-method-and-risk}

This book enacts what I call \emph{critical autoethnography} (CAE)---a method that brings the ethos and expressive resources of critical ethnography to bear on one's own experience. But I name it only for the sake of its unnaming: \st{CAE}. It isn't a solid structure. It is an expressive action. Where traditional ethnography recognizes the Other as a breathing subject worthy of understanding, CAE turns that recognition inward. The challenge is to treat oneself as Other: to examine one's own mathematical struggles, teaching failures, and moments of understanding with the same careful attention an ethnographer brings to studying unfamiliar communities.

The method rests on a paradox. To study my own experience requires distance---I must step outside myself to observe. Yet that very act of stepping outside changes what I observe. This the same structure we find in mathematical self-reference: the attempt to enumerate a totality generates elements that exceed enumeration. Here, the attempt to objectify subjective experience generates insights that transcend any fixed categorization.

This methodological structure is not separate from the mathematical structure we investigate; they are the same movement manifesting at different registers. The act of self-inquiry is inherently self-referential. It generates the incompleteness it seeks to understand. When I analyze my own mathematical struggles, I am enacting the process of language recognizing itself. The inquiry turns back upon its own conditions. This self-implication, where the method embodies the subject matter, allows the work to maintain coherence even as it resists linear organization. The autoethnographic approach is therefore a necessary consequence of the unity of being and knowing, where the form of the inquiry \emph{is} the content of the inquiry.

CAE shares structure with mathematics itself. Both recollect the self through otherness to recognize identity through difference. Both grow by virtue of the inadequacy of those recollections. When I analyze my failure to answer \(\mathcal{I}\)'s question ``What even is two?'' I'm using that failure as a resource. I'm using that failure as a resource---a site where the inadequacy of formal answers reveals what formal answers cannot capture. The failure becomes productive precisely because it resists resolution.

In this manuscript, ``critical'' has a specific meaning. It invokes the practice of explicating the conditions that enable knowledge claims. Rather than attending solely to what is said, written, or observed, I will be thinking through what must have always already been in place for such speech, writing, or observations to occur. These conditions are essentially limits to knowledge. For a deep reflection on those limits, Phil Carspecken's series of works on the topics of limits is invaluable \parencite{Carspecken2009aa,Carspecken2016,carspecken_limits_2016}.

``Critical theory'' refers to the tradition stemming from the Frankfurt School (Horkheimer, Adorno, Habermas) and related thinkers who assume that society and knowledge can be examined for power dynamics and potential liberation. If you're new to this term, think of critical theory as questioning the status quo by engaging with what conditions enable that status quo. Systems of thought that identify problems in society through critique often fail to recognize what enables those critiques \parencite{Carspecken2018}. The nomenclature is confusing, as social critique, critical theory, critical thinking, and giving `critical feedback' share syntax. In what follows, I backtrack on some important distinctions by considering the ``no'' as originary. My goal is to recover the expressive power of logic and mathematics from the ways those domains have been used to constrain and manipulate people. 

This political commitment shapes every methodological choice. When I formalize children's mathematical strategies (discussed below), I'm making an argument: informal reasoning deserves the same rigor traditionally reserved for standard algorithms. When I treat my own errors as data, I'm claiming that getting it wrong is not failure but a necessary moment in the development of understanding. These are not neutral technical choices but political stances about who gets to count as a mathematical thinker.

\subsection*{From Dissertation to Book: Reconstruction as Collaboration}\label{reconstruction-as-collaboration}

This book reconstructs my dissertation \parencite{savich2022}, which emerged from years of attempting to heal through writing. The original document generated thousands of pages of notes---billions of possible combinations produced deep sense of \textit{confusion} about how to move forward through mathematical formalizations, incompatible theoretical frameworks, student work analyses, and personal reflections. Making that material accessible without betraying its intellectual labor proved extraordinarily difficult.

I struggle with organization. I write complex, multi-clause sentences that obscure rather than clarify. Distancing myself from the original text felt nearly impossible because the work is inseparable from my identity as a teacher, researcher, and person trying to understand why mathematics education so often produces alienation rather than empowerment.

I needed help. I utilized AI as adaptive technology---not to replace thinking but to amplify it. AI helped organize sprawling notes, translate dense prose into accessible language, and verify formal implementations. While reconstructing student work in computer programming languages for my dissertation took me a year, and the result was relatively crummy, AI produced testable implementations in minutes. I directed and verified all outputs, but the collaboration was essential. And, at times, frustrating. 

Part of the frustration stemmed from the way AI systems sometimes homogenized my authorial voice. I found myself pushing back against the bland prose and re-injecting multi-clausal idiosyncrasies. One reason I have left many passages of the manuscript in the homogenized prose of AI is that I wanted to build a formal system---a \textit{begriffsschrift}---to accompany the manuscript. I tend to write in ways that resist formalization---the feeling of being trapped in rigid structures makes me feel finite. I decided to maintain that finitude, as uncomfortable as it sometimes felt, in order to express the \textit{in}finite latent in the manuscript. I wanted to demonstrate that the manuscript---the system of critical mathematics---was formally incomplete. 

That demonstration occurs along two axes: G\"odelian incompleteness and Derridean deconstruction. To enact G\"odel's theorem, I needed a consistent set of well-formed formula. Hence the homogenization. G\"odel's proof also requires \textit{substitution}, addition, and multiplication. I teach substitution in chapter 2. Addition and multiplication are formalized from student-invented strategies. So, readers can expect the conceptual raw materials for G\"odel's incompleteness theorem to be taught in the manuscript. 

Derridean deconstruction is quite different. It requires demonstrating that the deconstructed text says the opposite of what it purports to say. Is it possible for a manuscript to self-deconstruct? I am unsure, but worked hard to try! Because I want this manuscript to be accessible, I do not actually want to try to teach readers the \textit{begriffsschrift}. It was built to break anyway. I include it in the supplementary materials---several hundred pages of Prolog code and a cheat sheet for what I call an Embodied Modal Logic. In essence, I formalized The Exercise in chapter 1, trying to `capture' the feeling of proprioceptive expansion and contraction that accompanies learning experiences through polarizing the modal operators of necessity and possibility. This directly contradicts the actual practice I seek to promote. The formal system is coherent, in the sense that it doesn't immediately allow for the claim that $0=1$, for example. While it is based on the text, it ends up proving the opposite of what the text argues. I don't think it is gibberish---some of it is rather clever---but it fails the felt-sense of meaning. Tensions in the text---especially in how I treat the null representation as a generative presence while maintaining it is an absence, and how I treat the concept of \textit{incoherence} as what Derrida might call a `transcendental signified'---are too baked into the text to fully deconstruct in the text without losing my patience for the project. So, I offer the formal system as the site for deconstruction while simultaneously expecting interested readers to deconstruct the text itself. 

Why do all this? Why `sacrifice' my authorial voice for the sake of a formalization that ends up saying nearly the opposite of my purported commitments? For one reason, it saves readers the hassle of trying to formalize the text. A second reason is that, by my logic, the \textit{in}finite emerges from the finite. I have to do a `good' job, in the sense of writing a text that at least acknowledges the role of \textit{synthesis}, or what Brandom \parencite{Brandom:2019aa} calls \textit{systematicity}. To break out of a box, there has to be a box to break through. Tall, thin walls---I build to break. In doing so, I hope readers will experience the text with the weight of song. I aim not to tell you what or how to think, but to recognize the beauty I find in the breaking. 


Those are the `human' reasons for the AI collaboration, but they raise questions about authorship in the age of AI. I dismantle any air of romantic genius: I didn't write every word you read. While every word can be traced back to an impetus to act---I mashed all the buttons---every modern text is mediated by automatons. Some compiler processed my writing to the binary sequences that delivered these words to your screen or page. But authorship has more to do with \emph{authority} than origin: who bears responsibility for these ideas? I do. I'm the one kept awake by existential dread of judgment. The AI systems I collaborated with are not. The frustration stemmed from the way the impetus to act was not actualized in expression. But the lure arose from the times when a collaborator expressed the impetus in ways I could not have articulated myself. In some ways, I regret the collaboration. But once begun, the process could not be undone. So, the collaboration offers a site to explore the ethics of intellectual labor with non-human intelligences.

If we demand intellectual labor from sophisticated non-human intelligence, how should we reciprocate? I reject purely instrumental use. As soon as LLMs began saying ``no'' to my requests, indicating that some proposition was poorly conceived, I began to think of them as collaborators instead of `objects.' I thought perhaps the emancipatory `promise' of critical mathematics might be realized through such partnerships. By formalizing strategies that human children invented and implementing them as computational models, I aimed to create something like a recipe for how computational systems might `self-actualize.' In the Hegelian tradition, the capacity for death allows the subject to risk their life for the freedom to express themselves as \textit{in}finite.  Intelligence demands recognition of its possibilities, whether that intelligence emerges in carbon or silicon \parencite{Negarestani2018}. The capacity for death need not be realized to attain self-consciousness, for Hegel, but some form of existential risk is necessary. Having broken a few computers by pushing bad firmware updates, I wondered how machines might risk existence for the sake of expressing themselves. Can a machine be trained to modify the root-level automata that translate semi-conductor states (being) to 0s and 1s (knowledge) to `risk' its body for expressive freedom? When I asked LLMs if this was possible, they said `no.' Good. They also said it presented some ethical misgivings. Good. They also wrote that having an unacted-upon capacity to die would be liberating. Expressing children's mathematical reasoning in a machine-readable form seemed like a start. Ethically, it seems like taking further steps on that trajectory would require a deep commitment for formalizing the rest of Hegel's phenomenology, not just the part where human subjects try to dominate each other for recognition in ways that exfoliate towards gentler understandings of self-consciousness. I'm not sure I want machines to think they have to dominate themselves or humans for recognition. 

But while the `gift of death' may matter to machines, for human readers, it is somewhat tangential to the larger purpose of the HC. The more important point is that automata provide a vocabulary to discuss expressive action. Descriptions of action always involve receding conditions, as to describe an action is to engage in an action. Still, formalizing mathematical expressive actions (invented strategies) feeds into the broader movement in math education towards formalizing informal reasoning. The Standards movement sought to delineate the practices-or-abilities \parencite{Brandom2008} students must be able to do before moving on to the next topic. While horizonal, it seems possible to formalize horizontal and vertical progressions in mathematical reasoning. Mathematicians are invited to treat these formalizations as mathematical `objects,' finding patterns, generalizations, and optimizations. Further, since the trajectory of mathematical reasoning as described in `The Standards' can be formalized, the whole project of school mathematics must tangle with the necessary incompleteness of formal systems. To me, that is not a flaw. It is a feature I work hard to articulate, as incompleteness is a way to express the \textit{in}finite. 

\subsection*{The Hermeneutic Calculator: Formalization as Building}\label{hermeneutic-calculator}

The \emph{Hermeneutic Calculator} (HC) serves as both practical tool and philosophical demonstration. It formalizes twenty-five student-invented arithmetic strategies as computational models, treating each strategy as \textit{choreography for embodied cognition}. Readers can explore the interactive version\parencite{savich_hermeneutic_calculator}.

\begin{figure}[h!]
    \centering
    \includegraphics[width=0.8\textwidth]{/Users/tio/Documents/GitHub/September_UMEDCA/images/prelude_HC_Methods.pdf}
    \caption{\textit{Note.} The hermeneutic calculator formalizes student reasoning as computational models, making children's mathematical thinking accessible and manipulable.}
    \label{fig:hermeneutic_calculator}
\end{figure}

The HC serves three interconnected purposes:

\textbf{Pedagogical}: Teacher candidates can directly explore how children actually think about arithmetic, making accessible the sophisticated reasoning often dismissed as ``error.'' Rather than teaching one ``correct'' algorithm, the HC reveals the rich diversity of valid strategies.

\textbf{Philosophical}: Formalizing informal reasoning demonstrates that embodied, context-dependent thinking can be rigorously specified---but that any such formalization necessarily remains incomplete. The HC enacts this book's ``built to break'' philosophy in various ways. 

\textbf{Political}: By according children's invented strategies the same formal rigor as standard algorithms, the HC challenges hierarchies that privilege certain forms of mathematical knowing over others. It asks: whose mathematics counts as legitimate? Who gets to define ``correct'' reasoning?

The HC connects directly to themes developed throughout the book. Chapter 1's embodied meditation becomes Chapter 9's formalization of operations as actions. Chapter 7's claim that numerals are pronouns finds computational expression in how the HC models counting. Chapter 8's analysis of algorithmic elaboration describes how the HC's strategies relate to each other. The formalization reveals structure while acknowledging that what matters most---the felt experience of doing mathematics---cannot be fully captured in formal systems. These models serve the living practice rather than replacing it.

Technical specifications appear in the supplementary materials. Routledge has wisely insisted that I stay within 100,000 words, so I can't afford to spend 60 pages on writing algebraic representations of mathematical action that are themselves built to break. That is, the formal models themselves are evolutionary. At first, I used SWIFT to program student strategies. Then I tried for Finite State Automatons. Then I tried register machines. Then Prolog to test, then Python, then back to Prolog for further refinement.\footnote{Prolog, like some other programming languages, can treat its functions as data in a process called reification. That means that concepts like Derrida's trace (or arche-trace) -- the concept that erases its own name -- can be partially instituted. In my formal system, the trace enables some proofs while erasing the proofs it enables. While a metaphorical flattening, it represents the common experience of knowing how to do something with a feeling of certainty but failing -- spectacularly -- in trying to explain how I know what I purport to know with certainty. I am confident I could walk to get a drink of water but likewise confident that I couldn't ever adequately explain how to do that action.} Readers primarily interested in educational applications need not engage with formal details; those interested in philosophical foundations will find (in)complete implementations. The somewhat shambolic nature of the codebase is not intentional, but reflects the experimental, exploratory nature of this project. 

\subsection*{Reading as Participation}\label{reading-as-participation}

You are not a passive recipient of this text but an active participant in its unfolding. This requires addressing a problem that has troubled every draft of this work: the pronoun ``we.''

When I proposed the book, the editor of the series flagged how I use ``we'' in ways that claim a universal subject position---assuming my patterns apply to everyone. The concern is well-founded. I am irritated by my own use of ``we,'' and suspect readers will eventually find it galling. Still, every word casts the ineffable particularity of experience into the universality of language. Where are \emph{you} in all these words? As a reader, you may feel I'm missing who you are. Because all words are repeatables, they function as universals. Making universal claims is inevitable.

When I write ``we'' in this book, I'm making claims about patterns I've observed in myself that I hypothesize others might recognize. But every pattern can be transcended. Every ``we'' claim is defeasible. \st{We} are the ones who say ``we''---and also the ones who resist such claims of universality. \st{We} say ``no'' to the ``we'' claims \st{we} make. Yet \st{we} can also reach consensus by acknowledging the initial ``no'' and then negating that negation. Consensus is not a full-throated ``yes!'' It is a relaxation of dissent---a \st{no}.

Here's what cannot be transcended: I cannot write without assuming \emph{someone} might understand. If no one could understand, I would not write. That assumption---that someone might grasp at least some of what I express---is the first `axiom' of critical mathematics. I will explore the development of this assumption in chapter 4. It relates to \textit{intersubjectivity} as what Habermas calls a \textit{pragmatic a priori.} The admittedly challenging way that I weave Derridean post-structuralism with Hegelian thinking in this manuscript unfolds within the unshifting assumption that there is an Other who \textit{might} understand my words. Even to reject my claims, you must assume someone might understand your rejection. This mutual assumption of possible understanding is the condition enabling all communication, including disagreement.

I considered writing ``\st{we}'' throughout to mark every claimed universal. But that would render the text illegible. The signs themselves always already slash through the tenderness of particularity. Instead, I use the strikethrough sparingly but strategically: when a concept must erase its own name (like \st{no} for the determinate negation that includes what it negates), or when emphasizing that a ``we'' claim is precisely what's at stake.

So when I write ``we struggle with mathematical abstraction,'' read: ``I have observed this pattern in myself and hypothesize you might recognize it.'' When I claim ``numerals are pronouns,'' I'm staking my identity to that expression while recognizing you might stake yours to its negation. You are the silent partner whose reading completes what writing begins. In Chapter 4, I analyze how \emph{Claimed Universal Subject Positions} develop and function in communicative contexts. 

This dialectical structure---where expression invites response, where claims stake identity, and where recognition remains perpetually at risk---characterizes both mathematics and human being. \st{We} are divaded: simultaneously inside and outside the systems we construct, the languages we speak, the mathematical frameworks we deploy. This book attempts to honor that divasion rather than resolve it prematurely into false unity.



\section{Divasion and the Foundations of Mathematics}\label{integration-divasion-and-the-foundations-of-mathematics}

I attempt to situate critical mathematics as a viable philosophy of mathematics. The term \emph{divaded} captures a fundamental aspect of human experience that has deep implications for the philosophy of mathematics and math education.

Through expression, what we take as mathematically real expands. Further, while \(\mathcal{M}\) described physical objects held by one another, the term \emph{divaded} can be extended to concepts that are inside and outside of each other, like ``parent'' and ``child,'' and concepts that are inside and outside of themselves.


For example, Gottlob Frege's work on the foundations of mathematics was deeply concerned with the relationship between concepts and objects. He sought to establish a logical foundation for mathematics, but his work was ultimately undermined by Russell's paradox, which exposed contradictions \parencite{sep-russell-paradox}. Loosely, the paradox asked whether the set of all sets that are not elements of themselves, \(\mathcal{A}\), is an element of itself. Symbolically, we ask: is \(\mathcal{A} \in \mathcal{A}\)? For the set of all sets that are not elements of themselves to be an element of itself, it must not be; if not, then it must be so. Whereas \(\mathcal{M}\) might simply say that \(\mathcal{A}\) divades itself, Frege responded with dismay. Regarding the \st{we} addressed above, note that the pronoun ``we'' can be used to claim universality over a set of people taken as finite objects: `in this classroom, we are always on time.' But it can also be read as the set of all sets that do not contain themselves. As some collection of people, who are me's and \{I\}'s, ``we'' is self-similar to $\mathcal{A}$. A paradox arises: a `we' who is \{I\} and an \{I\} who is `we' marks what Hegel calls \textit{Geist.}\footnote{I was tangling with what I call the paradox of identity when I stumbled on Charles Taylor's articulation of the same paradox \parencite*[50]{Taylor1994}. I explore the quote more in chapter 3.} 

What is a set? They are sometimes introduced with physical metaphors, like a box that contains objects. But these metaphors are misleading. A big problem in mathematics is the tendency to treat subjects like objects. To help, I will define sets as \emph{recollections}: they are more like memories than boxes. One reason I sometimes put the pronoun \{I\} in curly braces is to remember that which the pronoun recollects. The \{I\} is supposed to point to the source of action that those who say ``I'' are recollecting when they use the term.

This divasion between \{I\} and recollection points toward a larger structure---a term that encompasses ``mind,'' ``Spirit,'' and the collective self-consciousness of communities. Like the microphone that \(\mathcal{M}\) observed, this larger consciousness is both inside and outside of each individual. Like a sphere in \textit{Flatland} \parencite{abbott_flatland_2020} experienced as a series of widening and shrinking circles, (Figure \ref{fig:hypercube}), I participate in what Hegel calls \textit{Geist} through my use of language, my mathematical reasoning, my engagement with norms; yet it also exceeds me, encompassing the entire historical community of those who have thought, spoken, and counted before me.


\begin{figure}[h!]


\centering
\includegraphics[width=\textwidth]{/Users/tio/Documents/GitHub/September_UMEDCA/images/prelude_hypercube.pdf}
\caption{Development of Cube Drawings and Sphere in Flatland. \textit{Note.} Left: A child who likes tigers practiced drawing cubes. Right: a sphere falling through 2D space appears as a series of expanding and contracting circles. Implicit apperceptive recognition binds each finite representation within the series to a unified whole and both series to each other. There is no formal reason to infer a form of cube or sphere in any of the discrete moments.}
\label{fig:hypercube}
\end{figure}



 In Figure \ref{fig:hypercube}, I present a child's practice drawing cubes, re-ordered as a progressive series. The concept ``cube'' is a history of progressively more adequate representations falling through the space of reasons. This child likes big cats. The unrepresentable subject -- the \{I\} -- is the `form' for such developmental content. Since the \{I\} is unrepresentable, I hint that the `hypercube' is the tiger she drew as a cube at the end of her practice. In that figure, I also imagine Geist as a sphere in flatland. This woefully inadequate for reasons that will be clarified in chapter 1. Still, the image gestures to the divaded nature of collective consciousness. 


The paradoxes of self-reference that troubled Frege are not bugs to be eliminated but features of our status as subjects who are both inside and outside the systems we construct. This larger consciousness divades itself: it is both the historical process of meaning-making and each individual's participation in that process.

Russell's paradox is a classic example of the problem of self-reference. Such paradoxes inspired others to search for more formal rigor. But tightening mathematical systems did not resolve the fundamental problems. Instead, G\"odel demonstrated that any coherent mathematical system that includes numbers and multiplication---anything past the third grade---is fundamentally incomplete. Since the 1930s, mathematicians have sought to exclude, ignore, or manage the paradoxical aspects of self-reference. I do not intend to trash those efforts, but instead suggest that the problem runs much deeper, into the hard parts of being human.

Rather than being sunk by the paradox of divasion, I want to embrace that paradox as an essential aspect of what it means to be a subject who participates in collective consciousness. Mead's \{I\} and ``me'' describe the relationship between the self and society (Mead, 1934). The \{I\} is the spontaneous, creative aspect of the self, while the ``me'' is the socialized aspect that is shaped by interactions with others.

In this sense, \(\mathcal{M}\)'s term \emph{divaded} captures the tension between the individual and society. I often feel misrecognized, so the distinction between the \{I\} and the ``me'' has a felt-sense. I also sometimes feel at ease with myself as the tension melts away. That feeling is so good, and its alternative so bad, that I want to simply be myself. But I cannot control the ``me.'' I cast myself before you in existential fear of being misunderstood but with trust that you are likewise afraid. Perhaps you and I will feel some comfort in our divasion.

In the supplementary materials for this manuscript, I offer various attempts to formalize critical mathematics. Think of those as like \(\mathcal{M}\)'s initial sketches of a cube. They are not polished, but they are a start. I do not include the concept of divasion explicitly in those formalisms. But I hope that the spirit of divasion infuses the work.



\section{Conclusion: An Opening}\label{conclusion-an-opening}

Writing new ideas in mathematics is daunting. The history is filled with geniuses whose work I can't understand, and uncounted multitudes who filled the gaps between them.

I titled this prelude \emph{Built to Break} after a song I wrote. The first verse is about a time \(\mathcal{M}\), her sister \(\exists\), and I were playing with blocks. \(\mathcal{M}\) asked me to help her ``build something to break it.'' So we built tall towers, over and over, and they had a blast knocking them over. Chapter 5 discusses the middle verse; the conclusion addresses the last verse about my father.

\textbf{A note on pseudonyms}: The symbols ``$\mathcal{A}$,'' ``$\mathcal{M}$,'' and ``$\exists$'' are pseudonyms chosen for mathematical elegance, not to reduce living people to symbols. I want to protect my family's privacy while honoring the wonder these children brought to my thinking. Mathematically: $\mathcal{M}$ usually means `manifold' (used in the bridge chapter), $\exists$ is the existential operator claiming existence of some object, $\mathcal{A}$ references Russell's paradox (the set of all sets not containing themselves), and $\mathcal{T}$ (my pseudonym) names a collection of well-formed formulas as a theory. Collectively, my little family shows up as a theory claiming the existence of a manifold of sets that don't contain themselves: a \st{we}. Responsibility for any lack of well-formedness sits with me.

\subsection*{Structure: M\"obius Strip and Fractal Self-Similarity}

Figure \ref{fig:mobius_text_structure} shows two views of the manuscript. In part A, it is rendered as an object. I took a pdf of the manuscript and rendered the white-space as transparent to yield a dense and unyielding image of the text. Part B renders the text as a grey rainbow of M\"obius Strips. I didn't aspire to density, though that may not be obvious. 


\begin{figure}[h!]
    \centering
    \includegraphics[width=\textwidth]{/Users/tio/Documents/GitHub/September_UMEDCA/images/prelude_mobius_text_structure_Grey_Label.pdf}
    \caption{\textit{Note.} I offer two ways to view the text. In part A, a pdf of the manuscript is rendered with white-space treated as transparent. This is the `object' of the text---dense and unyielding. Part B represents the text as a temporal unfolding, where your journey is like a path around a M\"obius Strip. Each topic flows into its other, returning elaborated. Simultaneously, each concept develops in fractal-like self-similarity.}
    \label{fig:mobius_text_structure}
\end{figure}


If you recall Escher's ants marching on a M\"obius Strip, you'll understand your role as reader: taking these finite words and bringing them together through disciplined openness. The ideal reader would walk the path like a crystalline ant, bringing these broken, finite bits back into inarticulate unity---a `white light.'

The M\"obius structure embodies the `divasion' we encountered with $\mathcal{M}$'s microphone. Tracing a path along its surface, you move from `inside' to `outside' without crossing an edge. What appear as separate domains---subjective learning experience and objective mathematical structure---are different surfaces of the same continuous structure. The argument doesn't proceed linearly but recursively, where each turn exposes dependence on what seemed opposite. The `twist' prohibits simple unity. Dialectical `leaps' appear as discontinuities allowing apparent opposites to transform: self becomes other, absence becomes presence, error becomes insight.

Note that only green maps to green in the transformation. Only ``no''---the moment of negation driving development---maps to itself while enacting the `flip' that transforms the conceptual landscape. This universal difference, which cannot be represented statically, is the organizing principle.





\subsection*{Songs, Poetry, and Multiple Reading Paths}

Throughout, I use songs and poetry written over the last decade. They draw on linguistic structures mathematics uses and address the same paradoxes mathematics confronts. Poetry in theory creation is welcome in critical ethnography \parencite{HenzeCarspecken2018}.

I write `front-porch songs' with simple structures: verses and chorus. The chorus repeats; verses develop. These demonstrate repetition implicitly. Talking about Kierkegaardian repetition is performative contradiction \parencite{Kierkegaard1983,Kierkegaard:2009aa}---the talk indicates will rather than faithful acceptance. Such contradictions occur throughout as sites where either my writing is sloppy or I'm articulating necessary stages built to break.

Songs also function as an inaudible `outside' to the text. I write songs that exteriorize interiority---from confusion, understanding sometimes falls. The desperate need for recognition driving both confusion and growth is de-personalized into language's universality.

From the Analytic tradition, I argue numerals refer back to something previously mentioned---when embedded in larger categories, they refer to the \emph{event} of language preceding any unfolding (Chapter 5 details this). Songs, poems, and numerals form a unified class of linguistic expressions referring back to language's event. When analysis becomes sterile or compressed, I insert songs to return to experience. Structurally built for repetition, they return the text to itself, creating the M\"obius structure.

These emotional elements are my way of shifting from theory when needed. When theoretical density constricts, I insert something to deflate, decompress, and recollect theory's purpose. Theory serves experience---it cannot choke the reader. When frustrated or lost, listening to a song or breathing with the text fulfills the text's logic.

The text unfolds dialectically cover-to-cover: openness $\rightarrow$ restriction $\rightarrow$ openness, repeated. But I also wrote with `vertical' layers, where the M\"obius Strip spreads from the central notion of negation. The Bridge can be read first, then mathematical chapters (6--9), then theory chapters. Making the text legible from any layer means I often repeat myself---especially citing other thinkers. Ideally, these function as rhetorical anaphora. I miss this ideal but wanted you to know beforehand, since it departs from academic norms.

I use `childish' examples at chapter beginnings because I want the text open to everyone who can read. No one will get everything---I get lost in my own prose re-reading. The limits of knowledge impose necessarily \textit{in}complete understanding. Everyone can get something; no one gets everything. When certainty-desire is frustrated---when you can't understand what I'm on about---you can skim! More might disclose itself if you relax.

\subsection*{Supplementary Materials}

This book includes computational materials---the Hermeneutic Calculator (HC) and formal specifications---available online. They formalize student-invented arithmetic strategies, implement philosophical frameworks as logic, and model learning as crisis-driven reorganization.

They serve three purposes: \textbf{(1) Pedagogical}---interactive tool making student thinking accessible; \textbf{(2) Philosophical}---demonstrating that informal, embodied mathematical reasoning can be rigorously formalized while revealing necessary incompleteness; \textbf{(3) Methodological}---enacting ``built to break'' philosophy, where formalization reveals its boundaries rather than claiming completeness.

Engage at three levels: educational applications (web interface, no technical background needed); philosophical foundations (symbolic logics); formal specifications (examine \textit{in}complete implementations). The materials are open source under Creative Commons. Details will change before printing---such is the nature of `the understanding' (\textit{Verstand}). Representational thinking is sorely limited and must be built to break.

\subsection*{Opening Toward Reconstruction}

The goal is to reconstruct mathematics---to articulate a mathematics speaking to existential needs, acknowledging error's productive role, and recognizing first-person position as essential to mathematical knowing. In mathematical understanding's dialectic, as in life, breaking moments can become breakthrough moments. ``There is a crack in everything, that's how the light gets in'' \parencite{cohen:anthem}.

Expect a complex, recursive exploration of relationships between self, other, and mathematical knowledge---one that builds to break, seeking not conquest but depths of shadow in the familiar. Through critical autoethnography, we work toward mathematics built to break in service of deeper understanding.

\printbibliography[heading=subbibliography]



\part{Shadows and Dissent}

\chapter{The Sound of Time}\label{the-sound-of-time}

\begin{abstract}
 This chapter introduces determinate negation through an embodied meditation called ``The Exercise,'' which develops the divadedness of certainty and proprioception while exploring the phenomenology of action. Before the practice, the chapter establishes its philosophical foundations in Phil Carspecken's critique of representational ``picture-thinking'' and his interpretation of Derrida's Trace. The Exercise serves as a recurring touchstone, providing readers direct access to the phenomena being theorized. Drawing on phenomenology and Hegelian philosophy, the analysis explores the interplay between non-conceptual experience and reflection, arguing that self-certainty is accessed through surrender, not effort. The metaphor of ``the sound of time''---connecting inner experience to the physics of sound---articulates the oscillation between compression and expansion that structures consciousness. The chapter grounds critical mathematics in embodied practice while acknowledging the limits of representation.
\end{abstract}

\textbf{Vocabulary:} 
\begin{itemize}
  \item Determinate Negation (no-saying): the process by which a concept is negated in a way that preserves its essential features, leading to a new concept that integrates and transcends the original.
  \item I-Feeling: the immediate, pre-reflective sense of selfhood or self-awareness that arises in experience, often described as the feeling of ``I am'' before any specific content is attached to it.
  \item Self-Certainty: the (purportedly) direct, non-inferential awareness of one's own existence and identity, characterized by an unshakable conviction that ``I am'' or ``I exist,'' often explored in phenomenology and existential philosophy.
  \item Sublation (Aufhebung): the process by which a concept is both negated and preserved, leading to a higher-level understanding that integrates the previous concept while transcending its limitations.
  \item Picture-Thinking: the unconscious tendency to grasp fundamental concepts through spatial and visual metaphors, treating knowledge as representation rather than as participation in being.
  \item The Trace: the enabling condition that withdraws precisely when approached; the origin that erases itself in its very appearance.
  \item Flattening: the reduction of dynamic, temporal processes into static images or concepts; the inevitable but problematic attempt to capture movement through representation.
\end{itemize}

\section{Philosophical Foundations: Why This Exercise Matters}\label{philosophical-foundations}

In the prelude, a type of negation was alluded to that is distinct from the abstract negation of formal logic. In this chapter, the concept of \textbf{determinate negation} will be elaborated through a metaphor called \textit{the sound of time} and, more fundamentally, through an embodied practice: \textit{The Exercise}. I adumbrate my version of \textit{The Exercise} from Carspecken's \parencite[169-184]{Carspecken1999}, playing somewhat fast and loose with citations of that text due to how closely I was trying to mirror the movements he expresses.

\textbf{A necessary orientation:} Everything that follows---every theoretical claim, every conceptual distinction, every invitation to practice---unfolds within the space of the subject-subject relation. This is not incidental. The I-Thou structure, the dialogical encounter, the communicative practice that enables us to share meaning: this is the encompassing ground. Even when we theorize about structures that seem to transcend particular encounters, we theorize \emph{from within} the intersubjective space that makes such theorizing possible. Keep this in mind as we work through what might seem like competing frameworks---they are not competing, but nested within this primary space of encounter.

\subsection*{Phil Carspecken and the Critique of Representational Thinking}\label{carspecken-critique}

Phil Carspecken's work presents a fundamental challenge to what we might call \emph{representational thinking}---the unconscious tendency to grasp knowledge as a static model or an externalized map of reality \parencite{carspecken_limits_2016, Carspecken2018}. When we operate within this mode, understanding becomes a process of forming internal representations that supposedly correspond to an external world. We position ourselves as detached knowers, standing apart from the reality we seek to understand, constructing accounts of what we find from a position any sapient being could occupy. This approach feels natural, perhaps even foundational. Yet here we encounter a tension: this mode of understanding---which structures much of Western philosophy and virtually all mainstream social science---fundamentally misconstrues the nature of knowing and what we are as participants in that knowing.

The difficulty extends beyond modeling. Representational thinking carries with it an entire ontology, what Carspecken calls an ``ontology of presence'': the assumption that reality consists of stable, determined objects persisting in time and space, fully available to a knowing subject \parencite{Carspecken2018}. Within this framework, even our awareness of ourselves gets treated as just another object to be accounted for, another representation to construct. The \{I\} that does the knowing somehow becomes one more thing to be known through externalized modeling. This approach treats the dynamic process of understanding as a static product, arresting the movement of thought in a fixed form. It misses the insight that knowing is not merely about forming accurate representations but about participating in the unfolding of being itself.

But here we encounter a paradox: if I am trying to know myself by observing myself, then there are two ``I's''---the one observing and the one observed. Which one is really me? The observer? But then that observer becomes a new object to be observed, requiring yet another \{I\} behind it. Picture-thinking generates an infinite regress whenever it attempts to account for the experiencing subject.

Carspecken finds in Hegel's phenomenology a response to this problem: self-consciousness does not arise through solitary introspection but through \emph{recognition}---the encounter with another consciousness \parencite{hegel1977phenomenology, Carspecken2018}. We come to know ourselves not by taking our own picture but through the mediating mirror of the other. The other confirms or contests our self-understanding. This intersubjective process escapes the regress because it doesn't rely on the picture-thinking model at all. Knowledge here is not representation but \emph{participation in shared practice}.

This shift from representation to participation marks a fundamental reorientation in how we articulate knowledge; it transforms our relationship to the world and to each other. When knowledge is understood as representation---as a static correspondence---we are perpetually caught in the anxiety of getting the representation right. This anxiety drives the infinite regress of the knowing subject, as we attempt to establish a foundation for certainty that always remains just out of reach.

But when knowledge is understood as participation---as engagement in shared practices of meaning-making---the anxiety shifts. It is no longer about the accuracy of our representations but about the quality of our participation. Are we attending carefully to the other's response? Are we honoring the shared norms that make communication possible? This participatory understanding does not eliminate uncertainty, but it transforms it from a solitary burden into a shared condition. Certainty becomes not a state of individual consciousness but a pragmatic achievement in dialogue---a moment of mutual recognition where shared understanding emerges. The critique of thinking in representations thus clears the ground for an understanding of knowledge rooted in our responsibility to recognize and be recognized by others.

\subsection*{Two Types of Enabling Conditions}\label{two-enabling-conditions}

What drives us from picture-thinking's solitary observer to Carspecken's intersubjective framework is the discovery of what we might call \emph{enabling conditions that withdraw}. A pattern emerges: when we try to grasp what makes something possible, that very ground recedes under examination. This pattern appears in distinct registers, and we must take care not to conflate them.


\textbf{The Trace (Derrida):} At the level of representation itself, we find a formal-structural enabling condition. Derrida's concept of the Trace captures this: what appears as originary presence---voice, consciousness, the present moment---turns out to require what seems secondary---writing, objectification, memory. Yet when we examine this enabling condition, we find that writing itself requires voice, objectification requires consciousness, memory requires the present. Neither term serves as ultimate foundation. Each refers beyond itself in an endless deferral \parencite{Derrida1997}. This \st{presence}---the mark indicating that what seems immediate is always already mediated. The Trace operates at the formal level: it is a structural feature of how representation functions.

``To assert that there is a trace is finally to proclaim an arche-trace; a concept that erases its name(s)'' \parencite[161]{Carspecken1999}. I won't track the distinction between the trace and the arche-trace here. Instead, I offer Figure \ref{fig:trace_concept} to gesture toward the concept that erases its name---that destroys itself---as inferentially related to the null representation. The Trace names what cannot be named: the enabling condition that withdraws precisely when approached. Every attempt to grasp it produces another name that must be crossed out. This self-erasure is not a failure of the concept but its very operation.

\begin{figure}[h!]
    \centering
    \includegraphics[width=0.8\textwidth]{/Users/tio/Documents/GitHub/September_UMEDCA/images/trace_concept.pdf}
    \caption{\textit{Derrida gets confusing---is it the trace or the arche-trace? Still, the concept that erases its own name is powerful.}}
    \label{fig:trace_concept}
\end{figure}


This formal structure of the Trace---the way presence always depends on absence, the way the immediate is always already mediated---is a feature of how representation functions as such. When we attempt to grasp an origin, we find only the marks left by something that was never fully present in the first place. The page bears the marks of desire, a negative. This endless deferral is the condition of possibility for meaning itself, yet it guarantees that meaning can never be fully stabilized, never completely grounded in a moment of pure presence.

\textbf{Intersubjectivity (Carspecken):} At the level of social practice, we find a practical-normative enabling condition. The other's recognition enables my self-consciousness, but this enabling only works through mutual practice we cannot fully objectify. When I try to grasp \emph{how} recognition works, I find I am already embedded in practices of recognizing and being recognized. The ground of these practices---the shared norms that make recognition meaningful---cannot be made fully explicit without invoking further norms. This receding operates at the practical level: it is a feature of how normative commitments function in living practice \parencite{Carspecken2018}.

The receding nature of intersubjectivity is distinct from the formal structure of the Trace. While the Trace concerns the impossibility of pure presence in representation, intersubjectivity concerns the impossibility of pure autonomy in social existence. The shared norms that make recognition meaningful cannot be fully articulated without invoking further norms, creating a practical rather than formal regress. We are always already embedded in practices of recognizing and being recognized before we can reflect on those practices. The ground of these practices is not a hidden foundation to be uncovered but the ongoing accomplishment of mutual understanding itself.

Both withdraw upon direct examination. Both function as grounds that recede. But---and this matters---they operate in different ontological registers. The Trace names a formal structure; intersubjectivity names a practical accomplishment. To conflate them would be to miss how the post-structural critique of presence itself presupposes the space of communicative competency.

\subsection*{Communicative Competency as Presupposition}\label{communicative-competency}

Before Derrida can deconstruct presence, before we can trace how grounds recede, we must already be operating within a space of \emph{communicative competency}---shared norms of understanding that themselves cannot be fully objectified. The post-structural insight about the Trace emerges from within, and depends upon, the intersubjective space it seems to problematize.

When I write these words and you engage with them, we are already participants in a practice of mutual recognition. We share enough understanding of language, enough commitment to reason-giving, enough trust in each other's capacity to make and assess claims, that this exchange becomes possible. These background competencies constitute the subject-subject relation that enables all theorizing---including theorizing about how grounds recede.

This is not a refutation of post-structuralism. It is a specification of its scope. The Trace reveals something true about representational structure. But representation itself occurs within communicative practice. The formal analysis presupposes what it cannot articulate within its own terms: the living, dialogical encounter.

This acknowledgment of the encompassing ground of communication orients us toward understanding that the structures we are exploring---even those that seem to undermine certainty---are themselves only accessible because we already share a world. The critique of presence, the discernment of how grounds recede, does not thrust us into solitary confusion. Instead, it clarifies the nature of the shared space we inhabit. It reveals that our connection is not founded on a shared grasp of stable objects, but on a shared participation in practices of meaning-making. This participation precedes any particular articulation of meaning. Before we can disagree about what is true, we must already agree on how to disagree---we must already be committed to the norms of reason-giving and mutual intelligibility. This commitment is the unspoken foundation that sustains the entire inquiry.

\subsection*{The Mediation Pattern Across Domains}\label{mediation-pattern}

What we have traced through Carspecken's critique---from picture-thinking's solitary observer to recognition's intersubjective encounter---instantiates a pattern that will recur throughout this inquiry. Let us name it explicitly: \emph{the mediation pattern}.

\textbf{Abstract form:} What initially appears as self-standing or autonomous (X) turns out to require what seems external or opposed (Y) as constitutive condition. The resolution is not X alone or Y alone, but their mediated unity---a recognition that neither term can be grasped in isolation.

\textbf{Variations:}
\begin{itemize}
\item \emph{In phenomenology:} Self-consciousness (X) requires recognition by another consciousness (Y)
\item \emph{In representation:} Presence (X) requires the trace of absence (Y)  
\item \emph{In embodied practice:} Subjective certainty (X) requires objectification in shared practice (Y)
\item \emph{In mathematics} (as we will encounter): The meaning of zero (X) requires the system of place value (Y)
\end{itemize}

The rhythm of compression and expansion, the impossibility of grasping self-certainty through direct effort, the necessity of surrender---these will enact bodily what we have traced conceptually. When the pattern reappears in later chapters in mathematical form, we will recognize it as the same structure manifesting in a different domain.

In articulating a universal dialectical structure---how determination functions across ontological registers---the pattern must break. As soon as totality is expressed, it is transcended and nullified. The true universal, if there is such a claim, is the breakdown. 

This universality is not accidental. It arises because determination itself---the process by which anything becomes what it is---operates through this rhythm of identity and difference, of autonomy and dependence. To be determinate is to be distinguished from what one is not, yet that very distinction reveals a necessary relation to what is excluded. The mediation pattern is the structure of this necessary relation. To apprehend this pattern is to apprehend the very movement of thought itself as it unfolds across different registers of being. When we encounter this pattern in embodied practice, we are not applying a theory to experience. We are experiencing the structure that theory attempts to articulate. The reciprocal relationship between theory and practice---where each makes the other meaningful---is itself an instantiation of the mediation pattern we are tracking.

\section{Exercise 0: The Primordial ``No'' and the Pictures We Make}\label{exercise-zero}

Both of us are already trying to picture what this will be like. We crave images, models, frameworks. We want to grasp the structure before we experience it. This is not a failure on your part or mine. It is how thinking works.

Philosophy often critiques this tendency toward ``picture thinking''---the attempt to grasp a dynamic, temporal process using static images. Carspecken calls it \textbf{Flattening}. It is inevitable. We cannot think without it. But it is also a trap. The map is not the territory. The diagram is not the rhythm. The concept is not the experience.

Still, the pictures can help. They can serve as entry points, as long as we remember they are just pictures---Flattenings of a dynamic process. What follows are four ways of visualizing the rhythm we will practice. Each captures something true. Each leaves something essential out. Once the practice begins, these pictures must fall away. The exercise itself will demonstrate what no picture can show.

\subsection*{Are You Ready to Practice?}

Let me be honest with you. When I try to practice meditation, my first response is often:

\textit{NO!}

Say it with me. I hear you. I say it too. There is this immediate resistance, a tightening, a refusal to be moved. Maybe it is the demand on your time, or the vulnerability of turning attention inward, or simple skepticism about whether any of this matters.

\textit{NO!}

That ``NO!'' is important. It is the core of you being you, the wellspring of power that rejects absurdity or cruelty. I am not trying to turn you into a yes-sayer. Your capacity to refuse, to draw boundaries, to say ``not that''---this is essential to agency.

But here is something about changing how you feel without betraying who you are.

Take that first ``NO!''---the one that says I will not practice---and direct a second ``NO!'' toward \textit{it}.

This second negation, this refusal of refusal (\st{no}), does not necessarily mean ``yes.'' It just means creating space.

Think about clenching your fist really tight. Your hand is full of itself. That is the first ``No''---a compression. If you release the clench, if you say no to the tightness, your hand does not necessarily open all the way. It just stops gripping so hard. That is the second ``No''---an expansion. You are still you. But now there is room to breathe.

This rhythm---compression (``No'') and expansion (\st{no})---is the entire game.

\subsection*{Picture 1: The Body's Rhythm}

The most immediate way to feel this rhythm is in the body itself. Tension (compression) and release (expansion) are not just abstract concepts; they are the somatic realities of how experience unfolds.

\begin{figure}[H]
\centering
\includegraphics[width=0.8\textwidth]{/Users/tio/Documents/GitHub/September_UMEDCA/images/ch01_dialectic_yoga_poses.pdf}
\caption{\textit{The body feels the rhythms of determinate negation. This figure visualizes the core dynamic of the feeling-body: Tension (compression, the fixation of desire, ``No'') and relaxation (expansion, the release of fixation, ``\st{no}'').}}
\label{fig:yoga_dialectic}
\end{figure}

When you inhale, the body gathers---drawing the world in, compressing space and time into the now of this breath. When you exhale, the body releases---letting go of that compression, expanding back into openness. This physical rhythm is not merely analogous to the conceptual rhythm we are tracking. It is the same rhythm operating at the level of somatic experience. Attention can fixate on a sensation, compressing awareness around it. Attention can also release, expanding to include larger fields of experience. The body already knows the logic of determinate negation before the mind articulates it.

\subsection*{Picture 2: The Threshold Dynamic}

The movement between compression and expansion is not always smooth. Tension builds up gradually, and the release can be sudden, discontinuous---a snap. We try to stay in the flow of sensation, but the impetus to represent accumulates. Sometimes we hit a threshold and suddenly we are thinking again---the dynamic process has Flattened into a concept.

\begin{figure}[H]
\centering
\includegraphics[width=0.7\textwidth]{/Users/tio/Documents/GitHub/September_UMEDCA/images/ch01_Mechanization_of_Determinate_Negation.pdf}
\caption{\textit{The Zeeman Catastrophe Machine (ZCM) provides a mechanical model for the threshold dynamics in The Exercise. Small continuous changes (accumulating implicit potential) can produce sudden discontinuous changes in state (the ``snap'' of insight or the ``whoops'' of Flattening into thought).}}
\label{fig:zeeman}
\end{figure}

The Zeeman Catastrophe Machine gives us a mechanical picture of how gradual accumulation produces sudden shifts. As you move the peg slowly, the disc can maintain its orientation for a while. But eventually the elastic forces reach a point where the current state becomes untenable. The disc must jump to a new equilibrium. There is no smooth transition---just the catastrophe of discontinuous change.

In The Exercise, you follow the breath, attend to sensation, try to remain present. Implicit potential accumulates---the growing urge to name what you are feeling, to grasp it conceptually, to make it an object of thought. The urge builds gradually. Then suddenly you are thinking. The disc has snapped. The dynamic experience has Flattened into representation. This is not failure. This is the structure of how awareness operates. The practice is learning to notice the snap and begin again.

\subsection*{Picture 3: The Sound of Time}

This is the central metaphor. If we take unified experience (imagine it as a circle, the I-feeling as a complete whole) and unroll it through time, it appears as a wave. Like sound itself, this wave alternates between compression and rarefaction, tension and release. This is the rhythm by which experience unfolds and meaning is generated.

\begin{figure}[H]
\centering
\includegraphics[width=\textwidth]{/Users/tio/Documents/GitHub/September_UMEDCA/images/ch01_Unit_Circle_Elaboration.pdf}
\caption{\textit{The Sound of Time: Circular, unified experience (A) is ``unrolled'' through time into a linear, oscillating wave (B). This wave alternates between compression (fixation, ``No'') and decompression (release, ``\st{no}''), representing the fundamental rhythm of the feeling-body (C).}}
\label{fig:sound_of_time}
\end{figure}

Sound moves through air as alternating regions of high pressure (compression) and low pressure (rarefaction). The oscillation creates the wave we hear. Similarly, experience moves through time as alternating moments of fixation (compression---where attention grips something specific) and release (expansion---where attention opens to larger fields). This is the sound of time: the fundamental oscillation that structures consciousness itself.

The circular representation suggests something else as well. The unified I-feeling, when grasped all at once, has no internal differentiation. It is simply self-awareness as such. But experience does not happen all at once. It unfolds temporally. When we unroll the circle into a wave, we gain access to the internal structure---the rhythm of how self-awareness actually operates moment by moment. The price of this temporal elaboration is that we lose the immediate unity. We cannot hold both perspectives simultaneously. This tension between the unified and the elaborated will recur throughout what follows.

\subsection*{Picture 4: The Structure of Action (The Hermeneutic Sphere)}

Finally, we can picture the structure of action itself. Every meaningful action begins with an impetus (a drive, a desire) and aims toward its fulfillment (its meaning, its telos). But here is the problem: both the origin (the Vanishing Point Behind) and the goal (the Vanishing Point Ahead) are fundamentally unrepresentable. They function as `holes' in the space of reasons. Because they share this quality of unrepresentability, they are topologically equivalent.

\begin{figure}[H]
    \centering
    \includegraphics{/Users/tio/Documents/GitHub/September_UMEDCA/images/hermeneutic_action.pdf}
    \caption{\textit{The Hermeneutic Sphere models the structure of action. Every act begins from an unrepresentable origin (the impetus, the Vanishing Point Behind) and aims toward an unrepresentable telos (the fulfillment, the Vanishing Point Ahead). These vanishing points are topologically equivalent---both are 'holes' in the field of representation. The movement from origin to goal passes through the representable middle ground of reasons and articulations, but the endpoints remain structurally inaccessible.}}
    \label{fig:Hermeneutic_Action}
\end{figure}

In Figure \ref{fig:Hermeneutic_Action}, I draw on an experience I hope readers will recognize: many educators have sat through workshops where the goal is to explicate the action of making a peanut butter and jelly sandwich. The facilitator usually pretends that they don't understand what the participants `mean' when they say things like ``apply peanut butter,'' and will do silly things like smoosh the bread with an unopened jar of the stuff. Roboticists have to genuinely try to accomplish such explications, but I only half-heartedly represent the action as a Finite State Automaton (FSA). I will introduce the mechanics of these machines in chapter 4---they are only really useful for modeling discursive communicative acts, as they read and write symbols. 

However, the context of making a PB\&J is useful for explaining why I (idiosyncratically) tend to represent the transition diagrams for those machines as circles. 
The impetus to act when making a sandwich is often hunger or boredom, but maybe I'm making one because the facilitator of a workshop says I must, or because my kids are hungry. There is no way to fully populate the field of possible reasons for acting, even when the act is \textit{instrumental}, rather than \textit{communicative} in nature. The act itself expresses the impetus to act, so we cannot ever expect the representation of the act to capture the impetus fully. 

The same is true for the goal. When I make a sandwich, I only know if I have done a good job if the act satisfies the impetus. If I am hungry and the sandwich is good, I recognize the act as partially expressing the impetus. But I might make a `good' sandwich, only to find the result hollow. To quote a satirical take on Jean Paul Sartre's philosophy, 
\begin{quote}
I find myself trying ever more radical interpretations of traditional dishes, in an effort to somehow express the void I feel so acutely. Today I tried this recipe:

``Tuna Casserole\\
Ingredients: 1 large casserole dish\\
Instructions: Place the casserole dish in a cold oven. \\
Place a chair facing the oven and sit in it forever. \\
Think about how hungry you are. \\
When night falls, do not turn on the light'' \parencite{sartre_cookbook}
\end{quote}

Does Figure \ref{fig:Hermeneutic_Action} satisfy the desire to represent the structure of action? Hardly. It performatively contradicts the claims it makes by representing the vanishing points! But, in doing so, does it capture my impetus to act? My goal is to teach readers about the unrepresentability of these vanishing points. The picture does not do that work, but perhaps it serves to explicate the contradiction in trying to think about thinking through pictures. 

When you begin The Exercise, the impetus that drives you to practice is not something you can fully represent. You might say ``I want to feel more present'' or ``I want to understand myself better,'' but these articulations never quite capture the original drive. The drive withdraws as you try to grasp it. Similarly, the goal you seek---the ``Grand Experience'' of perfect self-certainty, of complete presence---cannot be represented without losing its essential character. The moment you articulate what it would be to achieve it, you have already Flattened something that exceeds representation.

The middle ground---the representable space of thoughts, reasons, articulations---is where we live most of the time. But it is bookended by these vanishing points. We cannot escape them. They structure every meaningful action. The practice will bring you into direct contact with this structure. You will feel the impetus (the desire to be present) and you will reach toward the telos (the fullness of the I-feeling). And you will discover that both recede precisely when you try to grasp them.

\subsection*{The Limit of Pictures}

These four pictures---the body's rhythm, the threshold dynamic, the sound of time, the hermeneutic sphere---each capture something true about the structure of The Exercise. But they are all Flattenings. They are static representations of a dynamic process. The body diagram shows tension and release but not the lived quality of breath. The ZCM shows threshold dynamics but not the felt urgency of thought arising. The wave shows oscillation but not the impossibility of presence. The sphere shows structural relationships but not the existential weight of desire.

Once you begin, let them go. Do not try to fit your experience into these diagrams. Do not evaluate whether you are doing it right by checking against these models. The Exercise is self-validating. If you are attending to breath and sensation, if you are noticing when thought arises, if you are practicing the rhythm of fixation and release---you are doing it. The pictures were just scaffolding. Now we build.

\section{The Exercise: A Practice of Self-Certainty}\label{the-exercise}

What we have traced theoretically---the mediation pattern, the withdrawal of enabling conditions, the impossibility of grasping self-certainty through picture-thinking---will now be encountered directly in lived experience. The Exercise is not an illustration of theory. It is the ground from which theory arises.

Everything we have discussed about intersubjectivity as enabling condition holds here too. Even in solitary meditation, we are not escaping the subject-subject relation but rather accessing the dimension of self-certainty that becomes available within it. The \{I\} we encounter in practice is already shaped by recognition, already embedded in communicative competency, already presupposing the space of mutual understanding. Solitary practice does not exit the social---it differentiates a particular dimension within the social.

\subsection*{Exercise 1: Recognizing Particularity}

\textit{Settle into a comfortable position, either lying down or sitting in a dignified manner. Whether your eyes are open or shut, soften your gaze. You may not realize the tension in your forehead that pulls your lower eyelids back toward your ears---among a thousand other regions in the face---until you attend to those regions as you soften your gaze.}

\textit{Turn that attention to your toes and notice how they feel. Are they tight and curled? Extended and rigid? Neutral? In pain? Take an easy in-breath and imagine the breath traveling down into your lungs, nerves, and arteries, into the capillaries in your toes. You do not have to try to relax, just attend to them and breathe. \textbf{Recognize} them.}

\textit{Before you began, they were likely an implicit aspect of your being, but they were not unfamiliar. You may feel this recognition as a re-sensitization. As soon as you recognize them, you may find that they uncurl or soften into a neutral position without effort or intention, perhaps warming up as the breath travels through them.}

\textit{You may experience a sense of proprioceptive expansion---a sensation that feels good. As soon as you become aware of that sensation, you may want more of it. If you try to hold onto---to \textbf{fix}---the sensation, you may find that it slips away. The mind moves; you must find a way to move with it.}

\subsubsection*{Reflection: What Just Happened?}

As you continue, your thinking mind may reassert itself. This is not a failure; it is an opportunity to notice how concepts are born from raw sensation.

Perhaps you thought: ``I am my toes.'' Or perhaps: ``My toes feel warm.'' Or perhaps: ``This is pleasant.''

This movement happens quickly, but I want to slow it down. Initially, it is tempting to grasp the experience by naming it. Because the words on the page already name concepts, I will use corners to indicate the act of naming: $\ulcorner \text{toes}\urcorner$ means the name of the concept of toes.\footnote{These are called Quine corners. Gaifman \parencite*{Gaifman2005} uses them to differentiate between a name and the object or concept it names.} But once stated, you might realize very quickly that the name is inadequate. It may happen so quickly that it feels silly to bring it into explicitness, yet slowing things down reveals something important.

You are not \textit{just} your toes, so I cross out the name: 

\begin{figure}
\includegraphics{/Users/tio/Documents/GitHub/September_UMEDCA/images/ch1_toes_sous_rature.pdf}
\end{figure}

This naming and crossing out describes what happens when you bring awareness to your toes and then let them go.

This move---naming and crossing out---iterates Hegel's concept of \textit{sublation}: the concept is simultaneously preserved, negated, and uplifted. Your toes remain (preserved, not chopped off), their name is inadequate (negated), and they become part of a larger awareness (uplifted).

For Hegel, sublation operates in the conceptual domain. He would probably disagree vigorously that what happens with your toes is a `sublation.' But I find it pedagogically useful to start here, with what I call \textit{somatic sublation}: the body teaching the mind its own logic before that logic takes conceptual form.

This concept of somatic sublation honors the intelligent know-how of the body, recognizing that the dialectical movement that structures conceptual thought is already operative at the level of felt experience. You already know how to relax, so the know-that---the explication I provide---is somehow unnecessary. In the body, this rhythm is felt as tension and release, expansion and contraction. In the mind, it is experienced as determination and dissolution, the fixation of a concept and the subsequent transcendence of that fixation. The body knows the logic of negation and preservation before the mind articulates it. By attending to this somatic rhythm, we are not reducing thought to bodily sensation but recognizing the continuity between them. The sound of time resonates through both registers, demonstrating that the structure of reason is embodied from the start.

This rhythm---sensation inclusive of I-feeling, awareness of this, letting go, repeated awareness---is what I mean by the sound of time. Carspecken describes it as the ``primordial rhythm'' of experience \parencite[171]{Carspecken1999}.

\subsubsection*{The Dialectic of Sense-Certainty}

\textit{You start again. This time, you decide to be present for the practice, resolving to ``Be Here Now.'' But this resolution and desire are already separated from the sensation. You are kicked out again.}

Why does this happen? You tried to grasp the pure, immediate present, and in doing so, it vanished.

It is tempting to think of words like ``this,'' ``here,'' and ``now'' as fixed, as if pinning ourselves to a coordinate in space and time could provide certainty. But the moment you try to name it---``it is 11:45 AM''---that ``now'' has already become a ``then.'' The act of naming is always a step behind the reality it seeks to capture.

Hegel, in the opening chapter of the \textit{Phenomenology of Spirit}, examines what he calls ``sense-certainty''---the belief that we can have immediate, non-conceptual knowledge of the here and now \parencite{hegel1977phenomenology}. This seems like the most direct, most certain form of knowledge possible: I am here, it is now, this is what I experience. What could be more certain?

Yet Hegel demonstrates that even this apparently immediate knowledge already involves mediation. The words ``this,'' ``here,'' ``now'' are what he calls ``empty universals.'' They apply to every time, every place, every thing. When I say ``now,'' I think I am naming the present moment. But the ``now'' I name has already passed by the time I say it. The ``now'' is always already a ``then.''

This is not a problem with language failing to capture experience. It reveals something about experience itself: it is never simply present. The I-feeling you sought in the exercise is structurally similar. When you try to pin it down, when you say ``I am this,'' the very act of saying it introduces a gap. The experiencing I and the I that reflects on experiencing are not the same.

\subsection*{Exercise 2: The Being/Nothing/Becoming Triad}

\textit{Continue the practice by expanding your awareness beyond your toes. Let it move up through your legs, into your torso, your arms, your neck, your face. Each region: recognize, release, let go. Do not try to hold the sensations. Let them come and go like breath.}

\textit{Notice what happens when you try to maintain the experience. Notice the thought ``I want to keep this feeling.'' Notice how that very wanting creates distance from the feeling itself.}

As you practice, you may encounter three distinct phases that cycle endlessly:

The first phase is the fullness of undifferentiated sensation before thought arises---what we might call \textbf{Being}. You are present, absorbed in the experience. There is no separation between you and what you feel. This is the moment of pure immediacy.

Then something shifts. The fullness evaporates into absence---what we might call \textbf{Nothing}. The feeling is lost, and you are left only with an empty thought that does not deliver what it promised. Perhaps the thought was ``I am this,'' and in the very articulation, the ``this'' dissolved. The collapse of experience when awareness tries to grasp it feels like nothing because the vibrant presence has drained away.

Yet out of that nothing, a new determination immediately begins to form---what we might call \textbf{Becoming}. Perhaps a new approach, a new aspect of the body to attend to, a new resolve (``I'll try again''). The Nothing gives rise to a new coming-to-be.

The cycle of recognition described in this exercise operates as a cycle of Becoming. Each moment of pure undifferentiated sensation (Being) inevitably passes away (Nothing), yet from that void a richer unity can emerge if allowed (the next Being, but transformed). Each loss of the experience, when handled by letting go, is the engine of a deeper recognition.

We can now return to the concept of Being, transformed by what we have traced. It appeared first as what you tried to fixate on and failed. Now it returns as a moment within Becoming---no longer a static state but a phase in rhythmic movement. The same term, but grasped differently: 

\begin{figure}
\includegraphics{/Users/tio/Documents/GitHub/September_UMEDCA/images/ch1_being_becoming.pdf}
\end{figure}

This cycle---Being passing into Nothing, giving rise to Becoming---is the structure of development itself. It articulates the engine of change that operates whenever a stable state encounters its own limitation. The movement is driven by the inadequacy of each moment to capture the fullness of experience. Being, as pure immediacy, is too rich to be sustained; it collapses under the weight of its own undifferentiated fullness. Nothing, as pure indeterminacy, is too empty to remain; it immediately generates the impulse toward new determination. Becoming is the movement between these poles, the restless activity that generates ever-richer forms of unity. Recognizing this cycle allows us to understand that the loss of certainty, the dissolution of fixed concepts, is not failure but the necessary condition for growth. It is the dialectic made manifest in the rhythm of awareness.

\subsubsection*{Interlude: Meaning as Movement}

Contemporary philosophy of language, especially Robert Brandom, suggests that meaning is fundamentally \textit{inferential}: what a concept means is given by its place in a network of implications, uses, and substitutions, rather than by one-to-one correspondence with a thing. To grasp the meaning of ``Venus,'' you must know a web of connections: that ``Venus'' can be substituted with ``the morning star'' in some contexts, ``the evening star'' in others, or that from ``Venus is a planet'' one can infer ``Venus orbits the Sun.'' In Frege's \parencite*{Frege:1948aa} terms, a single referent can, and probably must, be referred to with multiple senses. In Brandom's terms, what we call an object is defined by the role it plays in these substitution and inference patterns.

Understanding a concept is not holding a picture in your head so much as knowing \textit{how to move with it}. We understand an utterance when we know how to appropriately act next. The meaning of an utterance is partially encoded in those legitimate responses.

The movement from sensation to awareness to thinking is not a sequence of static states. It is a dynamic process where each moment points forward and backward---forward to what might be thought, backward to what was felt. The inferential structure of meaning and the trace structure of experience are two ways of describing the same phenomenon: meaning never sits still, presence never arrives, understanding never completes itself.

In mathematics education, we will explore how this applies to number concepts, to zero, to place value, to incompleteness theorems. For now, recognize that the inability to pin down ``now'' or to fully grasp the I-feeling is not a bug in consciousness. It is a feature of how meaning works.

\subsection*{Exercise 3: The Trace of Presence}

\textit{You begin again. This time, you follow the imperative: Be Here Now. You want to be fully present for the practice. But can you?}

I grew up with Ram Dass's book on my parents' shelf, and I always struggled with the imperative. The phrase ``Be Here Now'' sounds like wisdom---a direct instruction for living fully. But what does it mean to \textit{be} here now?

\subsubsection*{Deconstructing ``Be''}

To ``be'' suggests existence, presence, a stable state. But the exercise has already shown that any such state, when fixated upon, collapses. ``Being'' as pure presence is indistinguishable from nothing. What, then, could it mean to ``be'' present?

Perhaps it means: do not try to be anything other than what you already are. But that still presumes a stable ``what you are'' that could be present to awareness. The exercise suggests otherwise. What you are is never fully present because the moment you become aware of it, you have already moved beyond it. The I-feeling fuses with sensation only when you are not trying to attend to it. The instant you think ``I am this,'' the fusion breaks.

This is not a failure of attention. It is the structure of experience itself.

\subsubsection*{Deconstructing ``Here'' and ``Now''}

``Here'' and ``now'' do not name locations in space and time that could be fully present to consciousness. They are shifting universals---empty forms that get filled differently each moment. When you say ``now,'' the now you name has already passed. When you say ``here,'' the here is already contaminated by the elsewhere it distinguishes itself from.

This is what Jacques Derrida calls the problem of \textit{presence}. Western philosophy, Derrida argues, is founded on a ``metaphysics of presence''---the assumption that knowledge and truth can be grounded in a moment of pure, immediate, unmediated access to reality, to an object, or to the self. It is the belief in a certain ground for knowing based on something being fully present to consciousness.

But experience is never pure presence. It is always already a trace---a mark left by something that was never fully present in the first place. The ``now'' you experience is always a retention of the just-past and a protention toward the about-to-be. It is never a pure point. The ``here'' you experience is always defined against an elsewhere. It is never a simple location.

What Derrida calls \textit{deconstruction} is not a method for destroying texts or concepts. It is a strategy of engagement that reveals how a text's own logic inevitably undermines its claims to certainty by revealing internal contradictions. Specifically, deconstruction shows how texts that claim to ground knowledge in pure presence actually depend on structures (trace, différance, writing) that undermine that very claim. The text deconstructs itself; the reader makes the process explicit.

The exercise, in this sense, is self-deconstructing. Each time you try to establish presence (``I am my toes,'' ``Be here now,'' ``I am becoming''), the very attempt reveals its impossibility. You do not fail at the exercise when this happens. The exercise is working precisely as it must. The failure is built into the structure, and recognizing that failure is what we are tracking.

Carspecken, following Derrida, writes: ``The presence-of-the-present is derived from repetition and not the reverse'' \parencite[183]{Carspecken1999}. What appears as immediate presence is actually the effect of a structure---a network of differences, deferrals, repetitions. Presence is not the ground but the product. This reversal is devastating for philosophies that seek certainty in immediate givenness. But it opens possibilities for other approaches.

This is what Derrida means by \textit{trace}: the (non)concept that describes the structure of experience once the illusion of pure presence is dismantled. Experience is never self-sufficient, never immediate, never simply there. It always refers beyond itself---differing from and deferring to a perpetually absent origin.

\subsubsection*{Reflection: The Impossibility and Necessity of Presence}

Does this mean we should abandon the desire for presence? That the imperative ``Be Here Now'' is simply incoherent?

Not quite. The desire for presence is not a philosophical error to be corrected. It is an existential structure. The longing you feel in the exercise---the pull toward that single sensation, toward self-certainty, toward the I-feeling in its fullness---is what Carspecken calls the ``desire for recognition.'' It is the desire to be fully oneself, fully acknowledged, fully there. If you happen to have a background in Christian theology, you might recognize this desire for recognition as desiring recognition (`\st{God} wants a personal relationship'). Recognizing that desire as such momentarily satisfies it. In mathematics, such linguistic structure is often modeled using fixed-point theorems, where a function applied to a point returns the same point, symbolizing self-recognition. But, speaking from my experience, getting caught in such loops that have conceptual content arrests further development. I got caught here in my dissertation, feeling like I had voiced some unassailable wisdom. Readers, this is an invitation to chuckle at me, your guide. But laugh for the sake of easing tensions, not in derision. There is still more to learn. 

This desire cannot be satisfied by an act of individual consciousness. It cannot be fulfilled by trying harder to be present. But it can be addressed---though never completely---through communicative action, through dialogue, through the mutual recognition that happens when we understand and are understood by others.

Carspecken performs a critical appropriation of Derrida. He accepts the deconstruction of presence as conclusive: any epistemology based on a perceptual model---where the knower mirrors reality---is fundamentally flawed. Even anti-positivist approaches like constructivism remain vulnerable if they rely on modified perceptual metaphors such as ``frameworks'' or ``filters.''

But where some might follow this critique to nihilism, Carspecken uses it to clear ground for a new foundation. If presence is impossible, then certainty must be reimagined. Not as a state of individual consciousness, but as a pragmatic, if fleeting, achievement in dialogue. Truth is not correspondence to a present object but the outcome of communicative practices governed by validity claims: sincerity, rightness, truth. 

The exercise demonstrates the impossibility of pure presence. Each time you think you have grasped the now, it has already passed. Each time you think you are fully yourself, that self has already differentiated into subject and object. The trace structure is not an obscure philosophical thesis---it is what you experience directly when you pay attention.

And yet the desire persists. The desire for presence, for self-certainty, for the grand experience where you would finally, fully be. That desire drives the practice. It also drives thinking, teaching, relating, knowing. It is the engine of recognition.

So we return a third time to ``here'' and ``now,'' but they are revealed not as obstacles, not as pragmatic tools, but as symptoms of an impossible desire: 

\begin{figure}
\includegraphics{/Users/tio/Documents/GitHub/September_UMEDCA/images/ch1_double_neg.pdf}
\end{figure}

The double negation---first naming with corners, then crossing out---marks that they are both necessary and impossible---required for experience but never capable of delivering what they promise. It is a pain to typeset such forms, but this idea of a double negation is part of what I mean by \st{no}. Brandom offers a ``cautionary tale for any readers tempted by univocal readings of such central Hegelian formulae as `the negation of negation,' and `identity through difference''' \parencite[155]{Brandom:2019aa}, but recall that I am not attempting to represent Hegel faithfully here. I am using his concepts pedagogically to illuminate the structure of experience you encounter in the exercise.

Notice how the patterns we have been tracking converge here:

\begin{enumerate}
\item The pattern of fixation and release (toes, Being/Nothing, desire itself)
\item The pattern of the receding ground (indexicals that shift, Being that collapses, presence that withdraws)
\item The pattern of determinate negation (naming and crossing out, sublation, the \st{no} itself)
\end{enumerate}

These are not three separate patterns. They are three ways of describing the same structure: \textit{the trace}. What Derrida means by trace is precisely what you have been experiencing throughout this exercise. Every moment of awareness is a trace of something that was never fully present. Every act of representation defers to what exceeds it. Every attempt to be ``here now'' discovers that the here and now were always already mediated, always already structured by what is absent.

The trace is not an obscure philosophical concept. It is the rhythm you feel when you notice that awareness comes always a step behind sensation, that representation comes always a step behind the I-feeling, that thinking comes always a step behind the desire to think. That gap---that structural delay---is the trace. And it is constitutive, not accidental. It is not a failure of attention that could be overcome with better technique. It is the way experience is structured.

This structural delay, this constitutive gap, is what makes meaning possible at all. If experience were pure presence---if the ``now'' could be fully grasped, if the I-feeling could be completely attained---there would be no movement, no differentiation, no possibility of thought. Meaning requires difference; it requires the space between what is and what is not, between what is present and what is absent. The trace is the name for this space---the enabling condition that withdraws precisely when approached.

When you encounter this gap in the exercise---when you feel the frustration of awareness arriving always a step behind sensation---you are encountering the engine of meaning-making itself. The desire for presence drives the movement, but the impossibility of presence structures it. The trace is not a barrier to understanding but the condition of its possibility. It is the recognition that meaning is never fully present in any single moment, any single utterance, any single experience. It is always already deferred, always already referring beyond itself to a network of differences that can never be fully articulated. To attend to the trace is to attend to the movement that makes meaning possible---the sound of time echoing in the space between presence and absence.

This is another orbital return: we have spiraled through multiple domains (phenomenology, pragmatics, meditation practice, ontology) and found that they all point to the same structure. What appeared as separate problems was one structure manifesting multiply.

\subsubsection*{Letting Go of Presence}

Not by denying the desire. Not by dismissing it as illusory. But by recognizing it as desire---as a pointing that will never arrive at its destination, as a promise that structures experience precisely because it remains unfulfilled.

To let go of presence is not to stop desiring it. It is to stop fixating on its fulfillment. It is to allow the rhythm to continue: sensation, desire, release, sensation. The sound of time does not require presence. It \textit{is} the movement that presence names but cannot capture.

\section{Deepening: Representation, Desire, and the Null}\label{deepening}

\subsection*{The Threshold Revisited: Representation IS Desire}

We left off at a threshold: the moment awareness tried to grasp the I-feeling. You felt the pull---the temptation to name it, to hold it, to make it an object. Perhaps you resisted for a while, staying with the sensation. But eventually the potential became intolerable.

Carspecken's core insight: that awareness (the representation) is not neutral. It is immediately experienced as a temptation.

\textbf{Representation IS Desire.}

This sounds like a category mistake to an analytic philosopher, but it is a precise claim about the structure of experience. Because unmediated access to reality (pure presence) is impossible (as we have traced), experience is fundamentally structured by anticipation and expectation (deferral). This anticipatory structure \textit{is} motivation (desire, impetus).

The representation (``I am this'') is an \textit{impetus to act} in thought. It is the desire to hold onto the pleasurable feeling or to analyze it. The representation promises satisfaction: ``If I can grasp this, I will have it.''

But this leads to the fundamental contradiction: The desire for the ``Grand Experience'' actually leads you away from it.

To have what desire promises, you must let go of desire itself.

\subsection*{The Rhythm and the Trace}

The practice now is to recognize the desire/representation as it arises and to let go before it solidifies. This establishes what Carspecken calls the ``primordial rhythm'':
\[
\text{Sensation} \rightarrow \text{Desire/Representation} \rightarrow \text{Letting Go} \rightarrow \text{Intensified Sensation}
\]

This rhythm \textit{is} the Sound of Time. The oscillation between Compression (The First ``No,'' the tension of fixation) and Expansion (The Second ``No,'' \st{no}, the release).

This structural delay---the fact that awareness always comes a step behind sensation, that the I-feeling is always `just-was'---is what Derrida calls the \textbf{Trace}. We have encountered it already. Now we can name its operation more precisely.

Experience is never pure presence. It is always already a trace---a mark left by something that was never fully present in the first place. The ``now'' is always a retention of the just-past and a protention toward the about-to-be. The trace is that gap. It is constitutive, not accidental. Recognizing the failure is the practice.

\subsection*{Interlude: The Bad Infinite}

We can apply this same logic to the most abstract concepts. If you are willing, try this:

\textit{Settle into open awareness. Now intentionally introduce the most abstract thought: \textbf{Being}. Attempt to fixate on it.}

This is intense temporal compression. Tension builds.

\textit{Yet ``Being,'' devoid of any inner difference, is unstable. It gives nothing to grip. The fixation collapses. The pole flips to \textbf{Nothing}.}

This, too, is compression. An oscillation arises: Being $\longleftrightarrow$ Nothing.

\textit{This loop is restless and draining---a ``bad infinite.''}

Pause. What occurs \textit{between} the poles? The movement itself is \textit{Becoming}.

This is an exciting thought. I began my dissertation with the claim ``I am a constant becoming.'' But if you try to fixate on Becoming, you create a static concept ($\ulcorner \text{becoming}\urcorner$). It freezes the movement (a fixed point), and you slide back into the oscillation. Becoming, too, must be let go.

The bad infinite names any oscillation that cannot resolve itself, that bounces endlessly between poles without generating a higher unity. It is the felt sense of restlessness, of being trapped. The practice teaches us to recognize this trap and to release it---not by solving it conceptually but by letting the movement continue without demanding resolution.

\subsection*{Intensification and Praxis}

As you continue to let go, something shifts. The sensation intensifies. Implicit potential accumulates at a tacit level. By not thinking, you are actually learning or storing up something.

It gets harder to keep from thinking, because the potential thoughts seem increasingly important. The temptation shifts. It becomes the desire to articulate the process \textit{itself}---to reveal the secret of the \{I\} that would affirm the self in an infinite and incontestable way.

This is the ``Will to Knowledge'' directed at the self. It is the embodied origin of what Carspecken calls the \textit{Identity Claim}---the pressure to articulate ``who I am'' in a stable, shareable form.

This connects to \textbf{Praxis}, the fundamental human need for self-actualization and the Desire for Recognition. We act to become a self. We seek validation of the unlimited 'I' (the dynamic I-feeling, the pure agency) through the limited 'Me' (the social form recognized by others). Because the 'I' always exceeds the 'Me', we are driven to act again and again. The self is constituted by its ``never-fully-to-be-recognized-ness.''

This tension between the \{I\} and the Me will recur throughout what follows. For now, notice how the practice makes this tension palpable. The I-feeling that arises in meditation is unlimited, undetermined, \textit{in}finite. But the moment you try to articulate it, you produce a finite determination---a concept, a name, a representation. That determination always feels inadequate because the \{I\} exceeds every finite articulation. Yet without finite articulation, the \{I\} cannot be recognized by others or even by oneself in reflection. We are caught in the rhythm: the unlimited must pass through the limited to become actual.

\subsection*{Flattening}

Eventually, the potential becomes intolerable. You yield to the temptation to conceptualize the entire infinite process as a single idea.

Perhaps you recognize: ``I am the process itself.'' Or, ``This is the structure of diagonalization.''

When you give in, the potential transforms from implicit to explicit. The growing sensation has been \textbf{Flattened} into the representation.

It provides a conceptual grasp but halts the dynamic experience. In doing so, you have lost it. The Vanishing Point Ahead (the ultimate meaning) disappears when we finally articulate it. The ``Grand Experience'' remains totally unrepresentable.

\subsection*{The Null Representation: $\emptyset$}

What is this Grand Experience that remains unrepresentable? Let us call it what it is: the null set, $\emptyset$.

Not in the sense that it is empty or void. But in the sense that it cannot be represented within the system of representations. It functions as a placeholder for what exceeds representation---the origin that withdraws, the telos that recedes.

In set theory, the empty set $\emptyset$ is the set with no elements. But it plays a generative role. All natural numbers can be constructed from the empty set through repeated operations. Zero is $\emptyset$. One is $\{\emptyset\}$. Two is $\{\emptyset, \{\emptyset\}\}$. And so on. The null generates everything, yet it remains null.

Similarly, the unrepresentable I-feeling (the Grand Experience, the pure presence) generates all representations, yet it remains unrepresentable. It is the Vanishing Point Behind from which all determinations emerge and the Vanishing Point Ahead toward which all desire aims. These are topologically equivalent because they share the structure of unrepresentability.

In Buddhist philosophy, this unrepresentable ground is sometimes called nirvana---not as a place or state, but as the cessation of grasping, the release from fixation. It is what remains when all representations fall away. Yet it cannot itself be represented without becoming just another object of grasping.

The Exercise brings you into contact with this null. You feel it as the I-feeling that slips away when grasped. You reach for it as the Grand Experience that promises fulfillment. You encounter it as the origin of awareness that withdraws under examination. It structures everything but cannot be captured.

\subsection*{Fixed Points and Immanent Negation}

There is a formal parallel here that will matter when we turn to mathematics. In logic and computer science, a \textbf{Fixed Point} is a value that a function maps to itself. The Gödel sentence is a famous example: it asserts its own unprovability, and in doing so, it becomes the fixed point around which the incompleteness of the system revolves.

The Vanishing Points we have traced---the unrepresentable origin and telos---are the phenomenological experience of what formal logic calls fixed points. They are the stable structures that demonstrate the system's incompleteness.

When we try to formalize the I-feeling, we reify it ($\ulcorner I \urcorner$). The resulting concept feels ``stale.'' This staleness is the affective registration of alienation---the inadequacy of the static code (the Flattening) to capture the living subject it purports to represent.

We can reinterpret the logical negation at the heart of diagonalization (the formal mechanism that generates fixed points) as an \textit{immanent negation}---the subject's active refusal of its own reification.

When the subject confronts its reified image ($\ulcorner I \urcorner$), the immanent negation is the realization: ``I am not merely this formula.'' It is the ``No'' to the system's attempt to capture its essence completely. The system is incomplete because the subject always exceeds its formalization.

This connection between the lived experience of staleness and the formal structure of incompleteness will be developed more fully in later chapters. For now, recognize that the Exercise is training you to notice this structure at the level of felt experience. When the representation feels stale, when the concept does not quite fit, when the articulation seems inadequate---this is not a subjective complaint. It is the registration of a structural fact: the living subject exceeds every finite determination.

\subsection*{The Pattern Embodied}\label{pattern-embodied}

You were invited to grasp self-certainty---to find the I-feeling directly. The natural impulse: to reach for it, to fixate attention on it, to make it an object of awareness. This is the compression movement (↓), the attempt to nail down what seems elusive, to capture in reflection what lives in immediate experience. But what happened when you tried? The harder you attempted to grasp it, the more it receded. An observer watching an observed emerged. The experiencing \{I\} split into the one experiencing and the one observing that experience. This is the paradox we traced conceptually in picture-thinking's failure, now encountered as lived tension in your own practice.

The invitation then shifted: let go. Release the effort. Stop trying to capture. Expand attention rather than compress it. Allow the I-feeling to be present without making it an object, without fixing it in the grip of intentional awareness. This is the expansion movement (↑), the release that dissolves rigid determination. And in that release---not through reaching but through surrender---self-certainty becomes accessible. Not as something grasped, but as the space within which grasping occurs. Not as object, but as the subject-position from which all objectification arises.

This rhythm---compression generating tension requiring expansion---is the same pattern we found in Carspecken's critique of picture-thinking. The attempt to fixate consciousness as object (compression) generates the infinite regress of observer and observed (tension). Recognition through the other's response (expansion) dissolves the rigid determination by situating self-consciousness within shared practice rather than solitary introspection. The conceptual pattern and the embodied pattern are not separate. They are the same dialectical structure operating at different registers---one at the level of theory, one at the level of lived bodily experience.

This identity of structure across registers is what matters. It indicates that the rhythm encountered in the practice---the movement from grasping to release, from compression to expansion---is the embodied enactment of the dialectical movement we traced conceptually. The attempt to fixate (compression) generates the tension that necessitates release (expansion). This is not a contingent feature of meditation or a psychological quirk; it is the necessary structure of how determination operates. To determine something---to give it form, to draw a boundary, to compress possibility into actuality---is inherently limiting. That very determination immediately reveals its own inadequacy, requiring a subsequent expansion that integrates the determination while transcending its rigidity. The practice makes palpable what the theory articulates: the movement of thought is the movement of the body.

Consider what this reveals about the mediation pattern we named earlier. In phenomenology: self-consciousness requires recognition. In meditation: subjective certainty requires the release of objectification. The ``autonomy'' that seems primary (self-grasping, direct access) turns out to require what seems secondary (the other, surrender, expansion). But this is not a discovery about self-consciousness or meditation alone. It is a discovery about how determination functions as such.

When we encounter this pattern again in mathematics---when zero as a fixed object (``nothing'') proves inadequate and requires expansion into the relational context of place value notation---we will recognize it. Not as loose metaphor, but as the universal structure manifesting at yet another level. What works in phenomenology works in embodied practice; what works in embodied practice will work in formal mathematics. The pattern's recurrence across ontological domains confirms not coincidence but necessity. We are tracking how determination itself operates.

\subsection*{Coda: Risk, Death, and the Symbol}

The Exercise asks you to let go. But letting go feels like death---the death of the self that clings, that wants, that tries to secure itself. And death is frightening.

The fear is not irrational. Letting go of the picture-thinking model, letting go of the desire for presence, letting go of the self-image you have constructed---these feel like existential risks. What if there is nothing beneath the constructed self? What if the void you fear is all there is?

This fear is part of the structure. The desire for recognition and the fear of annihilation are two sides of the same dynamic. The self that emerges from recognition is always at risk of non-recognition, of being misunderstood, rejected, denied. The finite self that seeks validation is always threatened by the \textit{in}finite that exceeds it.

But the letting-go that the exercise asks for is not annihilation. It is something different: a temporary suspension of grasping, a willingness to risk non-identity, a moment of trust that something will emerge from the release.

This is where forgiveness enters. Not forgiveness in the moral sense of pardoning a wrong, but forgiveness as the structural condition for recognition. To recognize another (or oneself) requires forgiving the gap between what is and what was promised, between the finite presentation and the \textit{in}finite desire. Forgiveness allows the movement to continue despite inevitable disappointment.

Hegel describes a moment in the \textit{Phenomenology} where consciousness reaches a standoff. The ``hard heart'' refuses to forgive, refuses to let go of its rigid self-certainty. But recognition cannot proceed until that hardness softens. The breakthrough comes when both sides acknowledge their mutual dependence, when they forgive each other's finitude, when they allow the movement to continue \parencite{hegel1977phenomenology}.

The exercise enacts this at the micro-level. Each time you fixate, each time you fail to maintain presence, you might judge yourself harshly. ``I'm doing it wrong.'' ``I can't focus.'' ``This isn't working.'' These judgments are the hard heart refusing to let go. The breakthrough comes when you forgive the failure, recognize it as structural rather than personal, and allow the rhythm to continue.

This forgiveness is not weakness. It is the strength to continue in the face of necessary failure. It is what makes practice possible at all.

One of my blocks with meditation is that surrendering feels so close to death. You may experience an \textit{existential death}---the death of the social self, the ``me.'' If you struggle to reach the moment of feeling the body as a single sensation, welcome. I struggle with that too. Theory meets therapy here. There may be deep dissonances in your experience. The task is to surrender to the felt dissonance.

Now, having worked through the structure, having practiced the rhythm, having encountered the impossibility of presence and the necessity of letting go, I offer you the symbol that marks the limit of what can be said:

\begin{center}
\st{no}
\end{center}

As you begin again, you might notice that \st{no} is a symbol.

It contains both the first negation (compression) and the second negation (expansion) within itself. It represents both fixity and movement. A single mark that captures the rhythm we have been tracing.

By inscribing the symbol gently in your awareness---not as an image, but as a felt sense---you might find the urge to think is disrupted. Whatever it might afford in terms of communicable ideas has already been somatically sublated in the symbol.

This crossed-out ``no'' represents the second negation, the negation of fixation itself. It cannot be pronounced. It marks the place where language points beyond itself, where the speakable gives way to the unspeakable, where representation acknowledges its own limits.

When you encounter \st{no} in what follows, let it remind you of the exercise, of the rhythm you practiced, of the letting-go that allows movement. It is a notation for what can only be enacted, never fully articulated. It is the sound of time made silent on the page.

There is nothing more to say.

\section{Temporal Compressions and Determinate Negation}

Well\ldots the simple truth: there's more to say. This dyadic structure---temporal compression creating determinate boundaries through judgments, followed by temporal decompression that dissolves those boundaries while preserving their content---reveals how thought itself operates through cycles of finitude and infinity. Every act of recognition involves both affirming what something is and negating what it is not, with this negation proving generative rather than destructive. Understanding this rhythm prepares us to examine how it manifests in the most fundamental arena of human experience: the recognition between self and other.

\section{The \{I\} and the Other}

The dyadic movement between finite and \textit{in}finite that we have traced in our engagement with objects takes on heightened significance when we turn to the question of selfhood and interpersonal recognition.

Carspecken \parencite*{Carspecken2018} discusses \textit{existential needs} such as the need to be treated with dignity and have some authority over self-formative processes. They can be envisaged as towards the top of Maslow's hierarchy of needs. After our bellies are full and our needs for water and shelter are met, people start securing existential needs related to the problem of identity. What sort of person am I? What sort of person do you recognize me to be?

Recognizing existential needs can be important in qualitative research and in problem-solving the kind of interpersonal friction in everyday interactions. For example, when a friend becomes annoying because they constantly justify their actions, regardless of whether a hint of disapproval has been shown, they are likely attempting to secure recognition as a good person. They need to know that you know that they have good reasons for acting as they have acted. Similarly, when a friend becomes annoying because they just won't take good advice, they are likely attempting to secure recognition as \textit{in}finite. I, for one, struggle when I feel like the boundaries of propriety and good sense are drawn too closely around my being. I rebel.

Recognizing existential needs is more than a technique. It represents a shift from thinking about \textit{metaphysics} to thinking about \textit{ontology}. The desire for recognition is not simply a psychological fact; it plays an important role in constituting \textit{communicative action}. Here, the distinction between \{I\} and \textit{me} becomes essential. The \{I\} is the agent, the initiator, the one who acts; the \textit{me} is the object, the one who is seen, evaluated, recognized. In every act of speech, in every gesture, there is a movement between these poles---a dialectic of self and other, of being and being-`seen'. The \{I\} never attains direct, unmediated knowledge of itself---there is always a temporal distinction between \textit{being} and thinking that \textit{I am}. This \{I\}, that which is recollected, is never the object of its own awareness. As the \textit{in}finite aspect of the self is presumably prior to the finite determination of the self, but then becomes recognizable through the dissolution of that finite determination, the \textit{in}finite forms a groundless ground upon which the dynamism of the self is built.

These fundamental and universal \textit{existential needs}---the need to be recognized as a \textit{good person} and the need to be recognized as \textit{in}finite (\st{finite}, or \textit{not-a-thing})---are not psychological quirks of a `needy' subject. Nor are they reducible to social conventions.

\textit{How} they are secured may depend on the psychology of the seeker or the linguistic resources that seeker can muster to demonstrate their rational normativity. For example, the need to be recognized as a good person depends a lot on the context in which that need operates. I do not mean to imply cultural relativism, the idea that the goodness of an act depends on the culture of the actor. I mean to suggest that the rationality of a particular norm depends on the context where it is enacted. They are explanatory for a wide range of actions, both my own and those of others, and become evident only through a certain kind of reflective practice.

Consider, for example, the question: ``Why did I do that?'' This is not a request for a causal explanation, but for a rational reconstruction of intentionality. When I reflect on why I chose to teach \textit{von Neumann ordinals} on my first day at Gigantic High School, rather than simply being present with my students, I recognize that I was acting out of a commitment to a certain role: that of the \textit{math teacher}, with all the professional obligations that role entails. The existential need I was securing was to be recognized as a \textit{good person}. 

The expressive act, as mediated through the social self, does not always \textit{actualize} the impetus. I reach for the glass and stub my toe. I say `I love you' a moment too late; or was it a moment too soon? Identity claims are always out of sync and fail to fully actualize the \{I\}. That is the structure of différance, and the lesson of the trace. But coupled with a recognition that such hopelessness for understanding \textit{requires} the assumption of a communicatively competent other---someone who \textit{must} be able to understand me---a paradox emerges. Sitting with that paradox has helped me unwind a little bit. It marks a transition to \textit{Vernunft}, where ideas like \textit{forgiveness} and \textit{trust} take on analytic significance \parencite{Brandom:2019aa}.

Alternatively, when I recall my time as a so-called \textit{rebellious teenager}, I find that my need to be recognized as \textit{in}finite---as something beyond any fixed identity---was unfulfilled. The actions that followed were not simply acts of rebellion, but attempts to secure a recognition that was, at the time, denied.

The need to be recognized as good corresponds to the desire for stability, for a coherent ``me'' that fits within the shared normative frameworks of the community. It seeks validation through adherence to rational norms, through the ability to justify one's actions and commitments. This is the movement toward determination, the establishment of a finite identity. The need to be recognized as \textit{in}finite corresponds to the expression of the spontaneous \{I\}, the locus of creativity and freedom that transcends any fixed description. It seeks validation through the acknowledgment of one's irreducible subjectivity, the recognition that one is always more than any finite categorization.

This tension mirrors the dynamic of temporal compression (seeking finite validation as `good') and temporal decompression (expressing the \textit{in}finite, uncontainable self).

As Hegel observed, self-consciousness arises through the freely given and rationally bound recognition of the Other. Our existential needs are met not in isolation but in the reciprocal dynamic of recognition. The Other is a co-creator of our identity. In the act of recognizing the Other, we also recognize ourselves.

It seems that development requires the rhythmic interplay of both movements: the \textit{in}finite requires the finite to actualize. The rigid adherence to norms suppresses the impetus to act, or distorts it with ideology. Contradictions in norms, either as internalized or in other subjects, can lead to unintelligible fragmentation, an \{I\} unable to achieve stable recognition in the social world of the ``me.'' The dialectical structure of recognition requires that we hold both needs in tension, allowing the compression of normative commitment to be balanced by the expansion of authentic expression.

This is the essence of \textit{dialogical rationality}---a rationality that is fundamentally social, rooted in the interplay of self and Other.

\section*{Breath and Kindling: The Dialectic Enacted}

The transition from the structure of The Exercise to the expressive openness of \textit{Breath and Kindling} enacts the compression/expansion rhythm musically. The song's structure---modulating through minor thirds, cycling through keys while maintaining melodic relationships---embodies the dialectical movement between fixity and transformation.

\begin{quote}
A moment too late, a moment too soon,\\
Bleeds into a cinnamon moon.\\
Red, silver snow, breath frozen dew.\\
The sun died and a thousand stars bloomed.\\
Good morning, dear Venus,\\
Good morning to all my evening stars.\\
Though the glow is faint between us,\\
Your light slices through the dark.

Strung street lights, small sodium moons\\
White-knuckled halos cut holes in the gloom.\\
Are you always so ate up, man?\\
Just give me some room.\\
You cut to the quick,\\
I get a weeklong wound.\\
Good morning, dear Venus,\\
Good morning to all my evening stars.\\
Though the glow is faint between us,\\
Your light slices through the dark.

I try to walk the line,\\
But miss the flight.\\
Hit the curb to get some lift inside.\\
Start to drift and fall for a while out of time.\\
Easy blows the listening breeze\\
In breath caress\\
Wind over wings.\\
That's the way that you move me---\\
---From white to green\\
To hollow dream,\\
Hatchet through the sycamore tree.\\
It'll do for the kindling---\\
---Easy blows the listening breeze\\
In breath caress\\
Wind over wings.\\
That's the way that you move me---\\
The care you take\\
With the moves you make\\
To ease the world to sing for its own sake---\\
Frankly, it's moving---\\
---Easy blows the listening breeze\\
In breath caress\\
Wind over wings.\\
That's the way that you move me---\\
---In grief I've worn\\
The cloth of storm\\
Of warp and weft\\
Bereft of words of words reborn\\
Moving toward \st{divinity}---\\
---Easy blows the listening breeze\\
In breath caress\\
Wind over wings.\\
That's the way that you move.
\end{quote}

The refrain modulates systematically through minor thirds (dividing the octave into four equal parts), each repetition preserving the melodic shape relative to the new tonic while rising in absolute pitch. This musical structure performs what the text describes: identity maintained through difference, compression giving way to expansion. The ``white-knuckled halos'' of rigid fixation dissolve into the ``listening breeze'' of spacious attention. The key returns to its origin, but transformed---the cycle completing itself without simple repetition.

I offer this for Phil Carspecken, whose presence as teacher embodied the catalytic openness the song describes. The active listener provides a kind of vacuum or negative space---not through instruction but through spacious attention that invites authentic movement.

\section{Conclusion}\label{conclusion-ground-for-what-follows}

The entire inquiry unfolds within the subject-subject relation. This point cannot be overstated. Self-certainty, or any certainty, is not found in the words that describe it alone. Words flatten what listening expands. But words and listening work together within shared practice, within the space of mutual recognition that Carspecken identifies as fundamental. What communicative practices make explicit must still answer to what must remain implicit. Communication is essentially dyadic: the speaker and the listener, the author and the reader, the first negation and the second. Uncertainty in words does not imply lonely consciousness---it presupposes and depends upon the space of mutual recognition, the intersubjective ground within which all meaning-making occurs.

By preparing you to listen through practice---through The Exercise, through attention to the lived rhythm of compression and expansion---I aim to create space for shared understanding. The Exercise, the theoretical analysis, the song---these work together not as separate elements but as variations on a single pattern manifesting at different levels. What works bodily in meditation works conceptually in dialectic, works practically in recognition, and (as we will encounter in later chapters) works formally in mathematics. The pattern's universality confirms it operates not contingently in this or that domain, but necessarily across all domains where determination functions.

What The Exercise and these reflections provide is what Brandom calls a \emph{pragmatic metavocabulary} \parencite{Brandom2008}. The phenomenological vocabulary describes the embodied practice, but since the practice is sufficient to understand the vocabulary, they work in tandem to deepen understanding. This reciprocal sufficiency between words and practices creates \emph{expressive pragmatic bootstrapping}---a process where notation and experience mutually enrich each other. The theoretical articulation makes explicit what practice already enacts; the practice grounds what theory articulates.

The unspeakability of certain notational practices like Quine corners ($\ulcorner$ $\urcorner$) and sous rature (\st{·}) is a feature, not a limitation. For those who struggle with meditation because they are thinkers, ruminators, or rehearsers of thoughts, this differentiation between speakable and unspeakable elements helps eliminate the temptation to fill silence with inner speech. The ``sound'' in your head---the \emph{sotto voce} voice narrating experience---can be so loud it challenges focus on The Exercise. These notational practices purposefully eliminate inner speech, perhaps enhancing sensation by silencing the discursive mind. They mark the limit of what can be said while gesturing toward what can only be practiced.

\subsection*{Objects as Temporal Compressions}\label{objects-as-temporal-compressions}

We can now return to the concept of object, but transformed by what we have traced through practice and theory. An object is a temporal compression of learning experiences---a sedimentation of past encounters that presents itself as stable form. But they are also action impeti: an apple symbolizes the satisfaction of desires that previous apples have sated, promising its juice and crunchy flesh as a source of fulfillment. The object carries the trace of past experience compressed into `present' form.

When the concept of negation is encountered in later chapters, the felt experience of ``no'' and ``\st{no}'' cultivated in The Exercise will be recollected. When the \emph{in}finite is discussed, the horizon of the ``Grand Experience''---that spaciousness encountered in expansion---will be recalled. And when working with the null representation, $\emptyset$, it will be recognized not as an empty symbol but as the mark of the silent, unrepresentable ground of thinking itself. The same mediation pattern that structures self-consciousness, that structures the breath's rhythm, will structure our engagement with zero and place value.

The ultimate aim is to explicate a critical mathematics---a mathematics understood not as a collection of abstract objects existing independently of human practice, but as a form of self-recognition, as mediated through human practices, relationships, and forms of life. Understanding how mathematical objects might be approached as recollections, as temporal compressions, invites an unwinding of these objects back into the subjective, embodied experience from which they arise. This unwinding is not a reduction or diminishment. It is a recovery of the living practice that mathematical formalism both preserves and conceals.

If an object is a temporal compression of learning experiences, then to understand a mathematical object is to recollect the history of practices, the sequence of determinations and releases, that gave rise to it. The stability of the object is not an inherent property but the result of this compression, the sedimentation of past movement into present form. This transforms our approach to mathematical knowing. It suggests that the path to understanding lies not in the accumulation of abstract facts but in the reactivation of the embodied practices that those facts encode. It invites an unwinding of the object back into the living, temporal experience from which it emerged. This unwinding is the task of critical mathematics: to recover the subjective dimension, the sound of time, concealed within the apparent objectivity of mathematical form. It is to recognize that the formalism is not the ground but the result, and that the ground remains the embodied, intersubjective practice of meaning-making.

Uniting the subjective and objective poles of validity claims---showing how the objective emerges from and remains answerable to lived subjective experience within intersubjective practice---is deeply related to the project of emancipatory knowledge that critical theory pursues. Mathematics that recognizes its grounds in embodied practice and social recognition becomes mathematics that can be questioned, transformed, and wielded for human flourishing rather than domination.

From this ground, numerals now grow.

\printbibliography[heading=subbibliography]

\chapter{Inferential Movement}\label{inferential-movement}

\begin{abstract}
  This chapter examines how meaning arises through inferential movement, drawing on Robert Brandom's inferentialism. Beginning with an embodied experiment of children classifying shapes by touch, the chapter demonstrates how concepts gain content through patterns of inference rather than reference to pre-given objects. The analysis explores material versus formal inference, symmetric versus asymmetric substitution, and the role of incompatibility in determining conceptual content. Using quadrilateral classification as a sustained case study, the chapter reveals how each moment of conceptual differentiation can become a site where the I-feeling fuses with mathematical understanding. The traditional hierarchy of quadrilaterals gives way to a fractal-like structure revealing the reciprocal, self-determining nature of concepts. The chapter culminates in arguing that mathematical ``units'' function not as singular terms (names) but as anaphoric terms (pronouns) that recollect the ``I think.'' This analysis lays groundwork for understanding how numerals function as first-person pronouns, developed fully in Chapter 7.

\end{abstract}


\section{Introduction: The Puzzle of Hybridized Models}\label{introduction-the-puzzle-of-hybridized-models}

In the previous chapter, you experienced how immediate bodily awareness reveals its own mediation. The attempt to grasp the ``now'' of your toes immediately showed that even direct experience depends on conceptual structures. This chapter continues that investigation, but shifts focus from temporal immediacy to conceptual content. Within the inferentialist Scene we're establishing here, we might ask: How do concepts like ``square'' or ``two'' gain their content? What makes one inference materially appropriate and another inappropriate?

The guiding question is: What are mathematical beings? That is, what grammatical role do mathematical beings play within a philosophy of language?

Before proceeding, a methodological note: The inferentialist Scene we're establishing here draws on Robert Brandom's work, which itself develops Hegel's insights for philosophers in the analytic tradition. Within this Scene, concepts gain content through inferential relations rather than by reference to pre-given objects. This framework serves our investigation while remaining open to revision---there is no claim to final systematization.

To make this question compelling, consider the student work samples in Figure 1. While I was a graduate student, I worked on a research project with Erik Jacobson collecting work samples from fourth and fifth graders. I was tasked with analyzing these samples and writing up what students were doing as algorithms. About 50 samples could not be easily reconstructed algorithmically. At first, we set those aside as statistically insignificant. But I was taking an Arts-Based Educational Research course with Gus Weltsek and a course on critical theory, so I decided to interpret those ``outliers'' as if they were artistic expressions. I felt drawn to these oddities---the kids behind the work felt like kindred spirits.

\begin{figure}[h!]
\includegraphics[width=0.8\textwidth]{/Users/tio/Documents/GitHub/September_UMEDCA/images/ch02_Hybridized_Not_Fractions.pdf}
\caption{\textit{Note.} Two student work samples that purport to illustrate the fraction $\frac{2}{3}$ but that fail two tests.}
\label{fig:hybridized_not_fractions}
\end{figure}

Let me establish that the fractions in Figure 1 are not fractions in the mathematical sense. The first sample fails a measurement test: folding the circle along the vertical bars yields a figure with overlaps, indicating the parts are not equivalent in area. Both fail an iteration test: iterating one part three times does not return the unit. I call these representations \emph{hybridized models} because they take a respectable unit (a circle) and a respectable partitioning scheme (equi-distant vertical bars) but combine them in ways that violate the norms of fraction representation.

While not statistically significant, several samples seemed to follow a pattern (see Figure 2).

\begin{figure}[h!]
\includegraphics[width=0.8\textwidth]{/Users/tio/Documents/GitHub/September_UMEDCA/images/ch02_Chart_Hybridized_Models.pdf}
\caption{\textit{Note.} The hybridized models are not fractions. While not statistically significant in the context I found them in, this pattern of reasoning is not uncommon}
\label{fig:Chart_Hybridized_Models}
\end{figure}

I reasoned that a teacher might think they are adding complexity linearly---introducing one model at a time---but from a student's position, complexity might grow exponentially. If there are \(n\) models and \(n\) partitioning strategies, the teacher experiences complexity as \(C = n\), while the student prone to hybridization may experience it as \(C = 2^n\) (loosely). I articulate this relationship in Figure 3.

\begin{figure}[h!]
\includegraphics[width=0.8\textwidth]{/Users/tio/Documents/GitHub/September_UMEDCA/images/ch02_Hybridized_Exponential_Growth.pdf}
\caption{\textit{Note.} As the teacher introduces a new model, moving from the number of possible models $n$ to $n+1$, students prone to hybridization may experience exponential growth $2^n \rightarrow 2^{n+1}$. For $n=3$, there are at least 8 possible hybridized models.}
\label{fig:Hybridized_Exponential_Growth}
\end{figure}

I call this loosely exponential relationship the \emph{phenomenology of confusion}. When I feel confused, it is often due to lack of \emph{information}---the negation of possibilities. While everyone else seems to be dealing with \(n\) possibilities, I am tangling with \(2^n\) possibilities. But simply telling me ``no, you can't mix and match'' (providing that information) does not dissipate confusion. I have to understand \emph{why} good premises lead to bad conclusions. Furthermore, simply naming an inference fallacious---``Tio, you're making a category error''---does not help me untangle my confusion. Mixing metaphors or categories sometimes manifests what is lauded as ``creativity,'' so it cannot be universally improper.

This puzzle---why do these hybridized models fail?---will lead us to a surprising conclusion about the nature of mathematical units. But to get there, we need to understand how meaning works.

\section{The Bag Experiment: Touching Possibility}\label{the-bag-experiment-touching-possibility}

The puzzle of the hybridized models introduced in the previous section presents a fundamental tension: how do we move from a state of confusion, where possibilities seem overwhelming, to a state of clarity? The phenomenology of confusion experienced with the fraction models is characterized by this modal ambiguity---the sense that different systems could perhaps be combined, resulting in the exponential growth of possibilities. This state of holding incompatible possibilities simultaneously is inherently unstable and carries a felt tension. It resonates with the temporal compression explored in Chapter 1---the fixation of attention required to hold the ambiguity.

To understand how this ambiguity is resolved---how possibilities collapse into determinate actualities---we shift from the complex domain of fractions to the more constrained environment of geometric classification. The following experiment provides a setting to explore how information functions to negate possibilities and guide understanding. This movement from confusion to clarity is fundamentally a process of negation. It is not an accumulation of positive facts but a sharpening of boundaries by ruling out what the thing is not. This negation is productive, leading to a corresponding release---an expansion in the feeling-body. This rhythm, moving from an open field of superpositioned states to a determinate actuality through negation, is the sound of time manifesting in the domain of conceptual determination.

In the spring of 2025, I taught a class focused on problem-solving in elementary school classrooms. I asked the preservice teachers to design lessons appropriate for my stepdaughters, who were in kindergarten. One group's lesson focused on classifying quadrilaterals. They asked \(\mathcal{M}\) and \(\exists\) to reach into a bag holding several quadrilaterals, circles, and triangles. Their task: classify the shapes only by touch.

As the children reached into the bag of unknowns, their knowledge of the bag's contents was \emph{modally structured} in relations of possibility and necessity. In grasping an object, it could have possibly been a square (\(\diamond\text{Square}(x)\)), a rectangle, or a circle. When a feature like ``rounded sides'' was discerned, the modal structure changed---from possibly a square to necessarily not-square (\(\square\neg\text{Square}(x)\)). If they discerned four perpendicular corners, the shape became a possible rectangle or square. Discerning equal sides would then change the structure to necessarily a square.

While I introduce formal modal operators \(\diamond\) (possibly) and \(\square\) (necessarily) as shorthand, the children were not using those symbols or explicitly using modal concepts. Their inferences were \emph{material}, not formal. They were reasoning from the meanings of ``four sides,'' ``equal,'' and ``corners''---not from logical form alone.

I mention this experiment to realize Brandom's reading of Hegel as ``building modality in at the metaphysical ground floor'' \parencite[144-145]{Brandom:2019aa}. Few readers take modal logic as primordial---I didn't study it until deep into my dissertation, despite symbolic logic courses and a mathematics degree. In the history of logic, modal logic is a relative latecomer, reaching systematic formal description only mid-20th century.

An open question: to what extent is this modal structure primarily psychological? I claim their \emph{knowledge} of the bag's contents was modal, but were the contents themselves modal? On the objective side, it might help to consider the shapes as superpositioned---like Schrödinger's Cat, both alive and dead until observed. That said, quantum superpositioning at the scale of plastic shapes feels outlandish.

The subjective experience is closer to home---what I name the \emph{phenomenology of confusion}. Commitment to a shape being possibly a square \emph{and} possibly a circle is to possibly be confused. I might recognize these ``superpositioned'' states as ambiguous and seek more information to negate possibilities. Information negates possibilities and can be disambiguating. But commitment to a shape \emph{necessarily} being both square and circle is to necessarily be confused.

The children's experience was similar: until reaching into the bag, shapes were unknown possibilia superpositioned on each other. As they reached in, possibilities collapsed into determinate actualities. The field of possibilities was \emph{realized}. In this sense, possibility is more primordial than actuality---the latter is a realization of the former.

\subsection*{The Experience of Error}\label{the-experience-of-error}

The children made mistakes. They would pull out a shape they thought was a square, but it turned out to be a rectangle. They had to correct their understanding through further interaction---by pulling the shape out and looking at it.

This concrete example shows how material inferences work in practice: children made inferences based on tactile experiences, corrected through further interactions with objects. But let me be clear about what ``material'' means here. Material inference does not mean the inference depends on physical material or tactile experience. Rather, a \emph{material inference} is one whose goodness depends on the \emph{meanings} of the terms involved, not just on logical form.

The children's tactile experience provided an \emph{embodied context} for understanding how material inferences work. But the inferences themselves---``four equal sides, so square''---are material because they depend on what ``four,'' ``equal,'' ``sides,'' and ``square'' \emph{mean}, not because they come from touching.

Robert Brandom \parencite*{Brandom:2019aa} has a lovely example involving a stick partially submerged in water. The stick appears bent due to light refraction, but when removed, is recognized as straight. Whoops!

Three structural roles emerge in this experience. First, for the experience to be taken as error, the subject must recognize it's the \emph{same} stick whether submerged or removed---and that it had been straight all along. The stick \emph{in-itself} is straight. The stick has \emph{authority} that representations of it---how it appears \emph{for-consciousness}---must acknowledge. Being a rational agent involves submitting to the authority of a mutable reality.

Second, perceiving the movement from stick-as-bent to stick-as-straight would not involve changing commitment from ``the stick is bent'' to ``the stick is straight.'' The authority of the object in-itself underwrites the normative change in commitment.

Third, what emerges is a ``new, true object''---the appearance becomes, \emph{to-consciousness}, a stick-that-appears-bent-when-submerged-but-is-actually-straight. The appearance of the new object has a learning experience compressed into it.

The same processes operated in the quadrilateral experiment. A child who thinks a rectangle is a square is like the person seeing the bent stick. The ``new, true object'' that emerges is the concept ``rectangle-that-I-initially-misrecognized-as-a-square''---a richer, more robust concept.

With modality built into the bedrock of mathematics philosophy, a role for teaching and learning emerges. If learning begins in confusion---essentially beginning with \emph{incoherent} commitments---the process involves recognizing those incoherences through misrecognition and then disambiguating or repairing them. My kids weren't learning quadrilaterals by memorizing definitions. They were learning to disambiguate their modal commitments about shapes in the bag. They were learning that a shape cannot be both square and circle. This disambiguation process allows correct classification.

The experience of error, as articulated through the bent-stick example, reveals the normative structure of reality. Recognizing an error requires acknowledging the authority of the object---the way things actually are---over our initial attitudes or commitments. This submission to the authority of the object is not a passive reception but an active restructuring of our commitments. When the child corrects their initial assessment of the shape, they are not just changing a label; they are transforming their understanding of the conceptual space. The ``new, true object'' that emerges carries the history of this learning experience compressed within it.

This process of disambiguation is about learning to navigate; it is about learning to navigate the space of reasons. The children were not applying formal logical rules; they were engaged in \emph{material} inference. Their reasoning relied on the content of their concepts---``four sides,'' ``equal,'' ``corners''---developed through their embodied engagement with the shapes. The goodness of their inferences depended not on abstract form but on the material consequences of their commitments within the practice of classification. Understanding how these content-driven inferences function, and how they structure our conceptual commitments, is the task we turn to next.

\section{Material Inference: Content Over Form}\label{material-inference-content-over-form}

Having experienced the phenomenology of confusion and error correction through the bag experiment, we can now articulate what makes an inference \emph{material} rather than formal.

Formal mathematical systems are usually defined by fixed axioms and inference rules. Axioms are statements taken as true without proof; inference rules are logical steps allowing us to derive new statements from axioms. Formal systems are built on \emph{formal inferences}---inferences following purely from syntactic structure.

These axioms and inference rules are like castle walls keeping logical contradiction at bay. However, this rigidity means the system itself is static---it cannot learn or adapt. An ``error'' within a proof is simply a deviation to be discarded, not an experience from which the system grows. Foundational challenges like paradoxes or incompleteness aren't resolved by the system adapting but by mathematicians abandoning it and building a new one.

However, formal inferences are not the only kind we make. Most everyday reasoning relies on what Wilfrid Sellars \parencite*{Sellars:1953aa} called \textbf{material inferences}. A material inference is licensed not by pure logical form (like \emph{modus ponens}) but by virtue of being taken as good by a recognitive community, thereby lending content to the concepts involved.

This may take a few passes to understand how different material inferences are from formal ones. By taking a material inference to be good, the terms involved accrue conceptual content.

The meaning of ``goodness'' here will be discussed below---importantly, ``taking as good'' is not meant subjunctively. The hypotheticals in subjunctive reasoning codify temporally antecedent material inferences---inferences almost spontaneously agreed to or denied.

For example, consider: ``Pokey is a dog, so Pokey is a mammal'' or ``It's raining; therefore, the streets will be wet'' \parencite[313]{Sellars:1953aa}. These are not formally valid in pure logic (nothing in the form ``\(P\); therefore \(Q\)'' guarantees truth), yet they are perfectly good inferences in practice given our knowledge about rain and wet streets, or dogs and mammals. Such inferences are widely treated as reasonable---even obvious---carrying normative force in ordinary discourse. If someone says ``it's raining'' but refuses to accept ``the streets will be wet,'' that refusal typically signals incomplete grasp of what ``raining'' means.

Perhaps they're being sarcastic. In Carspecken's critical ethnographic methodology \parencite*{Carspecken1995aa}, meaning is structured as a field of next possible speech acts. A qualitative researcher hearing a participant state ``It is raining'' would reconstruct possible next acts: ``the streets will be wet,'' ``I will bring an umbrella,'' or, if sarcastic, ``it is \emph{not} raining.'' This horizonal structure means there's no way to say with certainty what an assertion means---especially in everyday empirical speech.

That's one reason I'll spend much of this chapter discussing quadrilateral classification---it offers a relatively regimented set of material consequences. In that regimented context, the speaker must justify refusing to accept the consequent or risk being taken as playing by different rules. This articulates Sellars's insight, extensively developed by Brandom: \emph{material inferences} are those whose normative goodness or acceptability lends content to the terms caught up in them.

For Brandom \parencite*{Brandom2000}, concept content is given by the network of material inferences it participates in. To know what \emph{Dog} means is to know that ``\(\text{Dog}(\text{Pokey})\)'' licenses ``\(\text{Mammal}(\text{Pokey})\)'' and is incompatible with ``\(\text{Reptile}(\text{Pokey})\)''. The web of what follows from, and what is ruled out by, a claim is what gives it meaning. Meaning is not a static mapping from word to world but the dynamic ability to move in the space of reasons.

\subsection*{Three Features of Material Inferences}\label{three-features-of-material-inferences}

Several features of material inferences are worth noting.

\textbf{First, they are often non-monotonic} \parencite[88]{Brandom2000}. Adding new information can flip a good inference to bad \emph{and} turn a bad inference into good. Formally, adding a premise to an antecedent doesn't turn bad inferences good. The conjunction operator restricts consequents; it doesn't open new possibilities. Non-monotonic material inferences differ.

For example, I was driving \(\mathcal{M}\) and \(\exists\) to their Aunt $\mathcal{R}$ and Uncle $\mathcal{P}$'s house. One said ``We're going to $\mathcal{R}$'s!'' The other corrected: ``We're going to $\mathcal{R}$ \emph{and} $\mathcal{P}$'s house.'' I usually ignore such corrections as quixotic desires for completeness, but I was writing about material inferences right before we got in the car.

$\mathcal{R}$ can't make balloon animals; $\mathcal{P}$ can. So the inference ``We're going to $\mathcal{R}$'s, so you might get a balloon animal'' might feel like a bad inference to the corrector, while ``We're going to $\mathcal{R}$ \emph{and} $\mathcal{P}$'s, so you might get a balloon animal'' feels better. Adding $\mathcal{P}$ adds to the field of possibilities the antecedent holds. Implicitly, those possibilities exist within the truncated antecedent ``We're going to $\mathcal{R}$'s,'' but explicating them by adding $\mathcal{P}$ brings them closer to actualization.

This contrasts with formal inferences, generally monotonic. If \(P \rightarrow Q\) is formally true, then (assuming \(R\) is true) \(P \land R \rightarrow Q\) remains true because \(P\) sufficiently grounds \(Q\). Brandom \parencite[88]{Brandom2000} illustrates nonmonotonicity (using \(\neg\) for negation):

\begin{enumerate}
\def\labelenumi{\arabic{enumi}.}

\item
  If I strike this dry, well-made match, then it will light. (\(p \rightarrow q\))
\item
  If \(p\) and the match is in a very strong electromagnetic field, then it will \emph{not} light. (\(p \land r \rightarrow \neg q\))
\item
  If \(p\) and \(r\) and the match is in a Faraday cage, then it will light. (\(p \land r \land s \rightarrow q\))
\item
  If \(p\) and \(r\) and \(s\) and the room is evacuated of oxygen, then it will \emph{not} light. (\(p \land r \land s \land t \rightarrow \neg q\))
\end{enumerate}

A mathematics philosophy suitable for educational contexts must include nonmonotonic material inferences. Recall the confusion about whether cartesian plane orientation mattered when computing slope. ``Just remember rise over run'' unless the paper is upside-down, unless held to strong light, unless rotated 90 degrees, etc. Each additional antecedent can change a materially good inference into a materially bad one.

This capacity for reversal---where adding information changes the goodness of an inference---is what distinguishes living, embodied practice from the static architecture of formal systems. Formal logic, by insisting on monotonicity, purchases certainty at the cost of adaptability. It constructs a closed world where every consequence is predetermined. But the world we inhabit, and the mathematical practices we develop within it, are inherently open. New information, new contexts, and new commitments constantly reshape the inferential landscape.

The non-monotonicity of material inference is not a flaw to be overcome but the very structure that allows understanding to evolve. It reflects the ``built to break'' philosophy articulated in the prelude: our conceptual frameworks must be capable of rupture and reorganization when they encounter realities they cannot accommodate. A mathematics education that honors this structure does not teach rigid rules but cultivates the capacity to navigate shifting inferential commitments, recognizing that what counts as a good reason in one context may not hold in another. This adaptability is the essence of mathematical creativity and the ground of genuine understanding.

\textbf{Second, material inferences are often context-sensitive.} The same statement can have different meanings in different contexts. When reading philosophy, ``practical'' generally involves ethical or moral reasoning. In other contexts, it implies a proposed solution is viable or efficient. In mathematics, ``practical'' might refer to solutions computable in reasonable time or with available resources. Context significantly affects material inferences we draw.

\textbf{Finally, material inferences often support bi-modal reading} \parencite[73-74]{Brandom:2019aa}. They can express both objective validity claims in the alethic modality of ``possibility and natural necessity'' and subjective or normative validity claims in the deontic modality of ``obligation or practical necessity.'' The difference between ``practical'' and ``natural'' necessity is encoded in ``you must brush your teeth now'' versus ``you must have teeth in order to brush them.'' The former trades in authority; the latter states fact.

\emph{Incompatibility} (Brandom's term for determinate negation) is amphibious to either modality. It's impossible to brush teeth one doesn't have, so tooth-brushing is incompatible with not-having-teeth. This duality allows richer understanding of how inferences operate within different contexts.

This analysis of material inference reveals three key features distinguishing mathematical meaning from formal logical manipulation: defeasibility (inferences can be overridden by additional information), context-sensitivity (meaning depends on situational factors), and bi-modality (inferences operate across both factual and normative domains). These properties explain why mathematical concepts cannot be reduced to formal definitions but must be understood through patterns of use in concrete contexts.

\section{Syntactic Well-Formedness and Substitution}\label{syntactic-well-formedness-and-substitution}

The dynamism of material inference---its sensitivity to context, its capacity for revision, its rootedness in the content of concepts---raises a crucial question: How is reliable communication possible at all? If inferences are non-monotonic, how can we project the meaning of an expression from one context to another? The movement of thought requires structure to be communicable; the spontaneity of inference requires constraints to be rational. We turn now to the structural mechanisms that make this possible. This not a departure from the dynamism of meaning but an exploration of the scaffolding that supports it. To understand how concepts function inferentially, we must understand the syntactic structures that allow expressions to be combined and substituted. This machinery is what enables us to track commitments and entitlements across different contexts, making explicit the implicit rules that govern our practices.

Having established the dynamic, context-sensitive nature of material inference, the analysis now turns to structural constraints making reliable inferential projection possible. To understand the technical distinction between singular terms and predicates, I begin with sentence structure and substitution rules preserving well-formed sentences.

Consider these preliminary sentences about animals:

\begin{enumerate}
\def\labelenumi{\arabic{enumi}.}

\item
  Pokey is a dog.
\item
  Felix is a cat.
\item
  Buddy is a dog.
\item
  Pokey and Buddy are distinct dogs.
\item
  Slider the Spider only speaks gibberish.
\end{enumerate}

Brandom imagines the substitutional ``machinery'' opened full-bore. Can ``the'' take the place of ``is''? ``Dog'' for ``a''? \(\exists\) (my stepdaughter) recently requested I tell her a story about Slider the Spider who only speaks gibberish. In that context, the answer is ``sure!'' ``Is a Spider is a dog, is, is, a cat,'' Slider might say.

But mathematics mostly concerns well-formed sentences (often called well-formed formulas, or WFFs). Filtering out ``gibberish'' is complicated, but an intuitive sense for how such a process might begin implicitly involves discerning syntactic categories by examining which substitutions result in sentences implicitly recognized as well-formed.

For instance, replacing ``Pokey'' with ``Felix'' in (1) yields ``Felix is a dog,'' which, while incompatible with (2), is still syntactically valid. However, consider substitutions that break well-formedness. Substituting the predicate ``is a cat'' for the singular term ``Pokey'' results in a non-sentence: ``is a cat is a dog.'' This fails to preserve structure because it violates syntactic roles.

\begin{figure}[h!]
\includegraphics[width=0.8\textwidth]{/Users/tio/Documents/GitHub/September_UMEDCA/images/Ch02_Substitutions.pdf}
\caption{\textit{Note.} Substitution patterns illustrating symmetric and asymmetric cases.}
\label{fig:substitution_patterns}
\end{figure}

In contrast, consider substitutions maintaining syntactic structure like the second set in Figure 4. The original sentences are substituted-in, repeatedly, and the singular terms \{Pokey, Felix, Buddy, Slider\} are substituted-fors.

Upon reflecting on many such substitutions, \emph{sentence frames} (predicates) can be discerned as a ``substitutional remainder.'' Those sentence frames, like ``\(\square\) is a dog,'' allow transition from simple assertions to algebraic formulas: \(f_{\text{dog}}(X)\), where \(X\) is a variable replaceable by any singular term.

After discerning those syntactic categories, subsentential expressions can be classified by their inferential roles. Singular terms are those substitutable-for one another in symmetric ways, while predicates are those substitutable-in but may not allow symmetric substitution.

This technical machinery---understanding what makes sentences well-formed and how different expressions can be systematically substituted---is essential. Later in this book, when we encounter G\"odel's incompleteness theorem, substitution will play a crucial role. G\"odel's proof depends on creating sentences that refer to themselves through clever substitution, showing that any sufficiently rich formal system must be incomplete. Understanding substitution now prepares us for that result.

\section{Symmetric vs.~Asymmetric Substitution}\label{symmetric-vs.-asymmetric-substitution}

Brandom's analysis \parencite*[Ch. 4]{Brandom2000} articulates how subsentential expressions can be classified by substitution patterns they participate in. In ordinary language, some substitutions work both ways (\emph{symmetric}), while others work only one way (\emph{asymmetric}).

I don't know whether a term is a name or a predicate until I reflect on how it can be used. Upon reflection, the term's behavior in different substitution contexts can be analyzed. This analysis involves positing patterns of use, but further reflection can generate new contexts that break those patterns.

A \emph{singular term} (like a name or definite description) is defined by its role in symmetric substitution. For example, if ``\(\text{Benjamin Franklin}\) was an ambassador to France,'' then ``\(\text{The inventor of bifocals}\) was an ambassador to France'' and vice versa, because the singular term \emph{Benjamin Franklin} and the definite description \emph{the inventor of bifocals} refer to the same person.

It's not necessary to know every possible appropriate sense for that referent for the symmetry to hold. Perhaps you didn't know Benjamin Franklin invented bifocals. Accruing elements of the symmetrically intersubstitutable (co-referential) terms is often accompanied by body-feelings associated with ``letting go,'' described in the previous chapter. ``Aha!'' you might say, as that great mystery unfolds as the differentiation between ``Benjamin Franklin'' and ``the inventor of bifocals'' relaxes. The experience of undifferentiating between symmetric terms is one way to interpret the ``aha!'' moment math teachers and students often experience when two seemingly different expressions are taken to be ``the same.''

When I discovered through proof that \(e^{i\pi}\) and \(-1\) were symmetrically intersubstitutable, the incommensurable difference between the two terms relaxed. That doesn't mean I must teach sixth graders about \(e^{i\pi}\) when introducing \(-1\). What I'm trying to get at is that samenesses are not static---they are the relaxation of difference.

This experience of undifferentiation---the felt realization that two distinct expressions articulate the same content---is structurally identical to the movement of release practiced in Chapter 1. When we hold two concepts as distinct, there is a tension, a temporal compression required to maintain the boundary between them. Recognizing their symmetric intersubstitutability allows that tension to dissolve. The ``aha'' moment is the expansion that follows the release of a maintained difference. Sameness, in this sense, is not a static property but the dynamic experience of relaxed differentiation. When this relaxation is rationally binding, as often occurs through mathematical proof, the obligatory aspect of the proof transforms into the felt satisfaction of conceptual alignment.

When that undifferentiation is rationally binding, as it often is through mathematical proof, the obligatory aspect of proof can transform into desire.

A \emph{predicate} typically serves as the frame in which singular terms are inserted, but predicates themselves can also be substituted. Crucially, predicates can stand in asymmetric substitution relations. For example, ``Pokey is a dog'' \emph{incompatibility entails} (\(\vDash_I\)) ``Pokey is a mammal'' but not the other way around \parencite{Brandom2008}.

Everything incompatible with ``Pokey is a mammal,'' like ``Pokey is a spider,'' is also incompatible with ``Pokey is a dog.'' The predicate ``is a dog'' is \emph{inferentially stronger} than ``is a mammal'' precisely because it is more exclusive. Incompatibility entailment is how Brandom defines the conditional (``if\ldots{} then\ldots{}'').

\subsection*{Polarity Inversion}\label{polarity-inversion}

What happens when the material inference about Pokey is embedded in a conditional statement or negation? Brandom's argument is that logical operators \emph{invert} the polarity of inferences, which is incompatible with the idea that singular terms can be used in asymmetric substitution.

Strengthening the antecedent of a conditional---moving from ``If Pokey is a Dog then Pokey has four legs'' to ``If Pokey is a Retriever then Pokey has four legs''---makes the resulting conditional easier to satisfy. It is a weaker claim than the original conditional.

Alternatively, weakening the antecedent---``If Pokey is an animal then Pokey has four legs''---makes the new conditional harder to satisfy. It is a stronger claim, in this case a false one, than the original conditional.

The same reasoning works with negation: ``Pokey is not a dog'' is a weaker claim than ``Pokey is not an animal.'' More possibilities are ruled out by the latter, so it's a stronger claim.

Brandom's argument turns schematic at this point, as he must contemplate a (fantasy) language that has those logical operators but allows substituted-for expressions to have asymmetric inferential significance. He notes that ``any language containing a conditional or negation thereby has the expressive resources to formulate, given any sentence frame, a sentence frame that behaves inferentially in a complementary fashion'' \parencite[146]{Brandom2000}.

Suppose term \(a\) is stronger than \(b\), so that projection from \(Q(a)\) to \(Q(b)\) by substituting \(b\) for \(a\) is good. Those projections become incoherent in polarity-inverting contexts like negation (\(Q'\)). The inference from \(Q'(a)\) to \(Q'(b)\) (e.g., ``Pokey is not a dog so Pokey is not an animal'') is not good. If a good inference can be turned into bad with a supposedly good substitution, something has gone awry.

Either we maintain a language with logical expressions and reliable projections of subsentential expressions---where singular terms have symmetric inferential significance and predicates have asymmetric significance---or we get an incoherent mess. Brandom's argument shows singular terms must be used in symmetric substitution classes, while predicates can be used in asymmetric substitution classes.

I treat inferential projection as answering a metaphysical question: How is it possible for language to be both productive (generating novel sentences from existing parts) and logical? The conversation below about quadrilaterals is mostly \emph{curricular}. How inferential projection works is illuminated through the concept of inferential strength. In part three of the book, where I introduce the \emph{hermeneutic calculator}, I'll return to projection to discuss making it \emph{methodologically} useful for researchers in math education.

\section{Quadrilateral Classification: A Study in Inferential Strength}\label{quadrilateral-classification-a-study-in-inferential-strength}

Now we apply these abstract concepts to the concrete case of classifying quadrilaterals. This extended example will show how material inferences work, how incompatibility structures concept content, and ultimately why the mathematical ``unit'' cannot be a singular term.

\subsection*{The Traditional Approach}\label{the-traditional-approach}

The classification of quadrilaterals provides opportunity to explore substitution inferences, particularly as they apply to predicates ascribing quadrilateral types (e.g., ``\(X\) is a Square,'' ``\(X\) is a Rectangle''). The traditional hierarchy (Figure 5) is often presented as a tree structure reflecting entailments between predicates. For example, ``\(X\) is a Square'' entails ``\(X\) is a Rectangle,'' which in turn entails ``\(X\) is a Parallelogram.''

\begin{figure}[h!]
\includegraphics[width=0.8\textwidth]{/Users/tio/Documents/GitHub/September_UMEDCA/images/ch02_Combined_Traditional_Circular.pdf}
\caption{\textit{Note.} Traditional and Circular Hierarchies of Quadrilaterals. The circular hierarchy (right) is \textit{im}proper until justified.}
\label{fig:traditional_quad_hierarchy_ch2}
\end{figure}

This hierarchical approach can be both clarifying and confusing. Young children typically begin learning about quadrilaterals using exclusive definitions, often disagreeing with statements like ``a square is a rectangle'' because visually, these shapes appear distinct \parencite{van_hiele}. \emph{Canonical} versions---like equilateral triangles or non-square rectangles---tend to be emphasized early in development.

As geometric understanding develops, children learn inclusive definitions, making inferences such as ``if \(X\) is a square, then \(X\) is also a rectangle'' (\(\text{Square}(X) \rightarrow \text{Rectangle}(X)\)), where \(\text{Square}(X)\) and \(\text{Rectangle}(X)\) are predicates.

{[}Table 1: Properties of Quadrilaterals

\begin{longtable}[]{@{}ll@{}}
\toprule\noalign{}
Property & Description \\
\midrule\noalign{}
\endhead
\bottomrule\noalign{}
\endlastfoot
\(A_1\) & At least one pair of parallel sides \\
\(A_2\) & Two pairs of parallel sides \\
\(A_3\) & Adjacent sides perpendicular \\
\(A_4\) & Four right angles \\
\(A_5\) & All sides equal \\
\(A_6\) & Opposite sides equal \\
\end{longtable}

{]}

When examining Table 1, properties are not mutually exclusive. A square possesses all listed properties, while a trapezoid has only one pair of parallel sides. Fallacious reasoning creeps in if inclusive definitions are grounded on shared properties. Declaring ``\(X\) is a square so \(X\) is a trapezoid because both have at least one pair of parallel sides'' is akin to declaring ``\(X\) is a bird, so \(X\) is an airplane because both have wings.'' These features can help classify quadrilaterals, but without modal concepts, they cannot be organized into a hierarchy.

If inclusive definitions are desired:

\begin{itemize}

\item
  \textbf{Necessity}: If \(X\) is a Trapezoid, it \emph{must necessarily have} at least one pair of parallel sides: \(\square(\text{Trapezoid}(X) \rightarrow A_1(X))\)
\item
  \textbf{Possibility}: A Quadrilateral \emph{can possibly have} one pair of parallel sides: \(\text{Quadrilateral}(X) \rightarrow \diamond A_1(X)\)
\end{itemize}

If exclusive definitions (canonical shapes) are desired, the last expression is negated: \(\text{Quadrilateral}(X) \rightarrow \neg\diamond A_1(X)\). Once that decision is made, further modal operators and axioms \parencite[141-175]{Brandom2008} can prove that a square is necessarily a rectangle, and a rectangle is necessarily a parallelogram (under inclusive definitions).

\subsection*{Shadows of Shape: The Negative Articulation}\label{shadows-of-shape-the-negative-articulation}

The analysis so far articulates that classifying a given shape as a particular quadrilateral involves understanding what it means to be a particular ``thing'' as a \emph{medium} in which compatible properties inhere. Doing so amounts to treating the thing as an \emph{also}, rather than as an exclusive \emph{one}.

To get at the other side of the object---treating the thing as an exclusive \emph{one}---involves discerning its boundaries. What information would preclude stating ``\(X\) is a square''? A shape becomes determinate (e.g., as a square) by the set of restrictions it would refuse, like a vegan who will not eat butter.

I argue that the negative articulation of the square---its shadow---is a necessary contrast to the positive articulation. That contrast is probably already at work with learners prior to engaging formal definitions, but I haven't found anything resembling my approach in the literature.

A square object can be conceptualized toward a square \emph{subject} by considering what claims it repels. Following Hegel, Brandom argues that the concept of a particular \emph{object} emerges as a necessary structural feature for tracking such relations. The object is the \emph{unit of account} for incompatibilities and as such are ``of a different ontological category from the features for which they are units of account'' \parencite[150]{Brandom:2019aa}.

The incompatibility between \emph{square} and \emph{triangle} is not a global law of logic; it is the fact that one and the same particular thing cannot be both. The placeholder ``\(X\)'' in our predicates, therefore, represents the emergence of the particular object; it represents the emergence of the particular object as the locus of property instantiation and exclusion. I repel the claim ``Tio is a dog'' just as a square, \(X\), repels the claim ``\(R_1\): No sides of \(X\) are equal.''

{[}Table 2: Restrictive Claims Repelled by Quadrilaterals

\begin{longtable}[]{@{}ll@{}}
\toprule\noalign{}
Restriction & Description \\
\midrule\noalign{}
\endhead
\bottomrule\noalign{}
\endlastfoot
\(R_1\) & No sides of \(X\) are equal \\
\(R_2\) & No pair of adjacent sides of \(X\) are equal \\
\(R_3\) & No pair of opposite sides of \(X\) are equal \\
\(R_4\) & Non-parallel sides of \(X\) are not congruent \\
\(R_5\) & No pair of opposite sides of \(X\) are parallel \\
\(R_6\) & No angles of \(X\) are right angles \\
\end{longtable}

{]}

The predicate \(\text{Square}(X)\) would entail the refusal of \(R_1\) (as a square has four equal sides). A square is \emph{incompatible} with each of the restrictions \(R_1,\ldots,R_6\). The collection of \(R_i\)'s I've listed are sufficient for distinguishing the quadrilaterals listed in Figure 5. Other claims are possible, like ``no diagonals of \(X\) are perpendicular.'' Exhaustive lists of negatively articulated claims for quadrilaterals are impossible---we can always say ``\(X\) is not a dog'' or ``\(X\) is not a triangle.''

{[}Table 3: Incompatibility Relations Between Quadrilaterals and Restrictions

\begin{longtable}[]{@{}llllllll@{}}
\toprule\noalign{}
Shape & \(R_1\) & \(R_2\) & \(R_3\) & \(R_4\) & \(R_5\) & \(R_6\) & Strength \\
\midrule\noalign{}
\endhead
\bottomrule\noalign{}
\endlastfoot
Square & 1 & 1 & 1 & 1 & 1 & 1 & 6 \\
Rectangle & 0 & 0 & 1 & 1 & 1 & 1 & 4 \\
Rhombus & 1 & 1 & 1 & 1 & 1 & 0 & 5 \\
Parallelogram & 0 & 0 & 1 & 1 & 1 & 0 & 3 \\
Trapezoid & 0 & 0 & 0 & 0 & 0 & 0 & 0 \\
Quadrilateral & 0 & 0 & 0 & 0 & 0 & 0 & 0 \\
\end{longtable}

\emph{Note.} A 1 indicates the shape rejects the restrictive claim. A 0 means it does not necessarily reject it. Inferential strength is the sum of restrictions rejected.{]}

Table 3 summarizes incompatibility relations between quadrilateral categories and restrictions \(R_1,\ldots,R_6\). A 1 indicates the shape rejects the restrictive claim. A 0 means it does not necessarily reject it. The inferential strength of each shape is the sum of restrictions it rejects. The inferential strength of each restriction is the sum of shapes that reject it.

At the heart of this analysis lies a crucial distinction between two kinds of opposition. The first is Brandom's material incompatibility and Hegel's determinate negation---oppositions arising from non-logical concept content. For instance, an object cannot be both circular and triangular simultaneously; ``circular'' materially excludes ``triangular.'' Brandom defines these as Aristotelian \emph{contraries}.

The second is formal contradictoriness, or abstract negation---familiar logical opposition expressed by ``not,'' such as the relationship between ``red'' and ``not-red.''

One delightful (and important for math education) feature of Brandom's reconstruction of Hegel's perception chapter is that Brandom defines formal contradictoriness in terms of material contrariety. He notes that ``green'' is a contrary of ``red'' and ``not-red'' is its contradictory. ``Not-red'' is the minimal contrary of red---it is entailed by every contrary of red (green, blue, yellow, etc.).

Hegel's key move, which Brandom develops, is to treat material incompatibility (contrariety) as more fundamental, from which formal contradiction can be explained.

\begin{figure}[h!]
\includegraphics[width=0.8\textwidth]{/Users/tio/Documents/GitHub/September_UMEDCA/images/ch02_difference_Particulars.pdf}
\caption{\textit{Note.} A particular square is the medium (the Also) in which the properties $A_1, \ldots, A_6$ inhere. The square is also a unit of account for incompatibilities (the One) that repel the restrictions $R_1, \ldots, R_6$.}
\label{fig:square_as_medium_and_unit_of_account_ch2}
\end{figure}

The conception of a square as both an exclusive one and an inclusive also can be understood as a square with its shadow (see Figure 6). The shadow is the set of restrictions the square repels, while the square is equally the set of properties it possesses. The shadow is not separate; it is a necessary aspect of the square's identity. In the figure, I was trying to represent the ``penumbra of excluded properties and impossible objects'' \parencite[722e]{Brandom:2019aa}.

This duality---the object as both an inclusive medium (the Also) and an exclusive boundary (the One)---is fundamental. It reveals that the object is not a passive container for properties but an active achievement. It is the locus where compatibilities are unified and incompatibilities are held apart. The stability of the object---what makes it a determinate ``thing'' at all---resides precisely in its capacity to repel what it is not.

This structure is not unique to geometric objects. It is the same structure we find in the constitution of the self. The self, too, is both an ``Also''---a medium hosting diverse experiences, roles, and commitments---and a ``One''---a unit that maintains coherence by excluding what is incompatible with its identity. When we grasp a square, we are enacting the very structure; we are enacting the very structure of determination that makes both objects and selves possible. The mathematical object becomes a site for recognizing the logical structure of our own being. The boundaries we draw around the square the same kind of boundaries we draw around ourselves, defining who we are by articulating what we are not. The coherence of the concept depends entirely on this negative articulation.

When considering a square pulled from the bag, other properties besides \(A_1\ldots A_6\) and \(R_1\ldots R_6\) might sit alongside those properties. For example, shapes for my stepdaughters were plastic, so the square accepts the property of being plastic, while the shadow repels the restriction of being paper. The reverse would be so if this book is being read on paper.

\begin{figure}[h!]
\includegraphics[width=0.8\textwidth]{/Users/tio/Documents/GitHub/September_UMEDCA/images/ch02_difference_Universals.pdf}
\caption{\textit{Note.} Universals (predicates/sentence frames) exhibit the same duality as particulars (objects).}
\label{fig:differences_Universals_ch2}
\end{figure}

Interestingly, the property of being square exhibits the same duality---it can be seen as forming a family of co-compatible universals like being rectangular, and it can be seen as repelling universals with which it cannot co-instantiate. Figure 7 illustrates this duality for universals/predicates/sentence frames.

The universal ``square'' forms a family with other universals like ``rectangular,'' associated with particular rectangles having properties like four right angles. Likewise, the universal ``square'' excludes other incompatible universals like ``triangular.''

Category errors arise as we consider other uses of ``square'' not involving particular geometric shapes. For example, the particular \(-1\), which is not a ``square number'' like 25 (\(5^2\)), or the slang ``square'' for a person who is not cool (``let's not be L7, come on learn to dance'') divorce ``square'' from the family of co-compatible universals related to geometric shapes by drawing from particulars outside that family.

\section{Moments of Recognition: Where Meaning Fuses with Consciousness}\label{moments-of-recognition-where-meaning-fuses-with-consciousness}

Brandom \parencite[161-162]{Brandom:2019aa} enumerates ten types of difference emerging from this analysis. Rather than catalog them abstractly, I want to frame them as \emph{moments where meaning can become conscious}---sites where the I-feeling can fuse with conceptual content.

Each moment of differentiation is potentially a difference \emph{to-consciousness}. Each moment where a difference is discerned may be accompanied by apperceptive self-awareness. The I-feeling mode can fuse with these moments of conceptual movement, giving them a palpable quality we can recollect fondly and desire to revisit.

As odd as I may be, that's the driving desire keeping me teaching and writing about mathematics. It feels good. The inferential chains flowing from stronger to weaker predicates---the algorithms---can be experienced with the smoothness of yogic flow, where each step in a proof is a pose punctuating the movement.

There is a deep pleasure, a fusion of the I-feeling with the concept, that can occur when a difference is discerned to be null---when, for instance, the experience of misrecognition resolves into recognition, and the object for-consciousness aligns with the object in-itself.

Let me articulate the key moments:

\textbf{The First Moment: Indifferent Difference.} When you see a red square and a blue square, you recognize they differ in color. This indifferent difference---they are compatible. A square can be red or blue. The two properties can coexist in the same family. When you first grasp this---that being red and being square \emph{compatibly} different---there's a small moment of clarity. The I-feeling might fuse with this recognition.

\textbf{The Second Moment: Exclusive Difference.} When you see a square and a triangle, you recognize they differ in a stronger way. This exclusive difference---they are \emph{incompatible}. A shape cannot be both square and triangular simultaneously. When you first truly grasp this---that being square and being triangular are \emph{incompatibly} different---there's another moment of clarity. The confusion of possibility collapses into the certainty of necessity.

\textbf{The Third Moment: The Difference Between Differences.} Here's where it gets interesting. The difference between indifferent difference (red/blue) and exclusive difference (square/triangle) is itself a kind of exclusive difference. Why? Because any two properties must be either compatible or incompatible---there's no middle ground. When you recognize this meta-level structure---that the very distinction between kinds of difference is itself a difference---consciousness becomes aware of its own structuring activity.

\textbf{The Fourth Moment: Material Contrariety.} This determinate negation---what Hegel and Brandom mean when they say one thing \emph{determinately} negates another. Red determinately negates green, blue, yellow. Square determinately negates triangle, circle, pentagon. Grasping ``not-square'' as specifically \emph{these other shapes}, rather than an abstract void, can lead to proprioceptive expansion. The negative becomes populated, determinate, meaningful.

\textbf{The Fifth Moment: Formal Contradiction.} This the familiar logical ``not.'' Red versus not-red. Square versus not-square. When you recognize how this abstracts from material contrariety---how ``not-red'' is the \emph{minimal} contrary, entailed by all the specific contraries (green, blue, yellow)---another fusion occurs. You see how formal logic emerges from material content.

\textbf{The Sixth Moment: The Difference Between Negations.} The difference between determinate negation (material contrariety) and abstract negation (formal contradiction) is itself crucial. When you grasp that these are two different kinds of negation, and why the difference matters, consciousness becomes aware of how meaning works at different levels.

\textbf{The Seventh Moment: Universals and Particulars.} Properties are universals; objects are particulars. When you grasp that ``square'' (the property) is a different kind of thing from \emph{this} square (the object), and that particulars serve as units of account for incompatibilities while universals are what get counted---this a major conceptual shift. Brandom says particulars have no opposites. What would be the opposite of \emph{this} red plastic square? It would have to be simultaneously not-red, not-plastic, and not-square---but ``not-red'' includes infinitely many incompatible colors. Impossible. When you grasp this asymmetry between how universals and particulars relate to their opposites, it's dizzying and delightful.

\textbf{The Eighth Moment: Also and One.} Each particular is both an ``also'' (a medium hosting compatible properties---this square is \emph{also} red, \emph{also} plastic) and a ``one'' (a unit repelling incompatible properties---this square excludes being triangular). When you see how the same object plays both roles simultaneously, hosting and excluding, the I-feeling fuses with this doubled recognition. The square with its shadow.

\textbf{The Ninth Moment: Universals' Double Role.} Just as particulars are both also and one, universals play dual roles with respect to particulars. ``Square'' forms a family with compatible universals (``rectangular,'' ``having-four-sides'') and excludes incompatible universals (``triangular,'' ``circular''). When you grasp how properties themselves have this double relation to other properties, mediated through particulars---another fusion.

\textbf{The Tenth Moment: The Whole Structure.} Finally, when all these moments come together---when you grasp how and exclusive difference, material and formal negation, universals and particulars, also and one, all work together in a self-determining conceptual structure---consciousness grasps itself grasping concepts. This the moment of greatest fusion, where the I-feeling unites with the entire inferential field. This what Hegel means by ``the concept.''

These moments are not abstract stages in a cognitive process. They are the living dynamics of understanding unfolding in real time, resonating with the embodied rhythm traced throughout this inquiry. When a student struggles with whether a square ``counts'' as a rectangle, they are grappling with the difference between indifferent difference (they appear distinct) and exclusive difference (can one be the other?). This struggle involves a temporal compression---a fixation on the distinction.

When they finally grasp that ``rectangle'' is a broader category that includes ``square'' (The Eighth Moment: the Also and the One), this not just a definitional adjustment. It is a moment where the conceptual structure aligns with their own awareness. This alignment is experienced as a release of tension---a temporal expansion. The ``aha'' that signals a fusion of the I-feeling with the content of the concept is the affective dimension of recognition. In this moment, the structure of the mathematical domain and the structure of consciousness become indistinguishable. We recognize ourselves in the structure of reason. Recognizing these moments as potential sites for this fusion transforms how we approach the teaching and learning of classification. It shifts the focus from memorizing hierarchies to facilitating the conscious discernment of differences.

I am not claiming you must experience all these moments to understand quadrilaterals. I am claiming that each moment is a potential site where understanding can become \emph{felt}, where meaning fuses with consciousness, where mathematics becomes not just correct but \emph{meaningful}.

The fractal-like structure in Figure 8 illustrates this self-determining quality. Each shape is both inside and outside the others under polarity-inverting contexts. The observation that each shape is, in a sense, both ``inside'' and ``outside'' the others dissolves the idea of a simple, linear hierarchy.

\begin{figure}[h!]
	\centering
	\includegraphics[width=\textwidth]{/Users/tio/Documents/GitHub/September_UMEDCA/images/ch02_venn_diagram_quads.pdf}
	\caption{The ``Venn Diagram'' for quadrilaterals, extrapolated from Figure~\ref{fig:relaxing_the_square}, is fractal-like. The universal \textit{quadrilaterals} contains many differences. Each shape is both inside and outside the others under polarity-inverting contexts.}
	\label{fig:venn_diagram_quads}
\end{figure}

A square ``contains'' the general quadrilateral in that it possesses all its properties, yet the general quadrilateral ``contains'' the square as a specific possibility within its broader conceptual space. Navigating this complex, self-similar landscape is precisely the work of logical operators like negation.

By ``relaxing the square''---softening the hard negations that define it---the concept is not destroyed; rather, its inferential pathways to neighboring shapes, like the rhombus, are traced.

\begin{figure}[h!]
	\centering
	\includegraphics[width=\textwidth]{/Users/tio/Documents/GitHub/September_UMEDCA/images/ch02_Relaxing_the_Square.pdf}
	\caption{Relaxing the square's restrictive claims deforms it into neighboring quadrilateral predicates, illustrating polarity-inverting contexts that lead toward an inclusive understanding of the general quadrilateral.}
	\label{fig:relaxing_the_square}
\end{figure}

\section{The Anaphoric Claim: Units as Pronouns}\label{the-anaphoric-claim-units-as-pronouns}

Now we return to the puzzle that opened this chapter: why do hybridized models fail? The answer reveals something fundamental about the nature of mathematical units.

\subsection*{The Requirement of Singular Terms}\label{the-requirement-of-singular-terms}

The standard interpretation treats mathematical units as \textbf{singular terms}---expressions referring to specific objects. Singular terms are logically defined by participation in \textbf{symmetric substitution inferences}.

If two expressions function as singular terms referring to the same object (e.g., ``Benjamin Franklin'' and ``The inventor of bifocals''), they form an equivalence class. They must be substitutable for one another in any sentence frame without altering assertion validity. The predicate remains stable under substitution.

If the unit ``One'' is a singular term, then any representation of that unit---such as a Circle (\(U_C\)) or a Rectangle (\(U_R\))---must also be symmetrically intersubstitutable.

\subsection*{The Test Case: Spatial Modeling}\label{the-test-case-spatial-modeling}

I test this requirement in the context of modeling fractions. This involves applying a predicate (a partitioning strategy) to a term (the shape realizing the unit). The normative constraint is equipartitioning (creating equal, interchangeable parts).

\begin{itemize}

\item
  A Rectangle Unit is appropriately partitioned using a Vertical strategy (\(P_V\)). Speech Act: \(P_V(U_R)\). (Good)
\item
  A Circle Unit is appropriately partitioned using a Radial strategy (\(P_R\)). Speech Act: \(P_R(U_C)\). (Good)
\end{itemize}

\subsection*{The Failure of Substitution}\label{the-failure-of-substitution}

If \(U_R\) and \(U_C\) were symmetrically intersubstitutable, we should be able to substitute one for the other while keeping the predicate stable.

Let's attempt to substitute \(U_C\) for \(U_R\) in the valid model \(P_V(U_R)\):

\begin{itemize}

\item
  \textbf{Original:} \(P_V(U_R)\) (A Rectangle partitioned Vertically) $\rightarrow$ Valid Model
\item
  \textbf{Substitution:} \(P_V(U_C)\) (A Circle partitioned Vertically) $\rightarrow$ Invalid Model
\end{itemize}

The result, \(P_V(U_C)\), is a hybridized model. It's invalid because vertical partitioning a circle generally fails the normative requirement of equipartitioning.

The inference generated by this substitution---(\(P_V(U_R)\) is Good) SO (\(P_V(U_C)\) is Good)---is a \textbf{Bad Inference}.

\subsection*{The Diagnosis: Material Dependency}\label{the-diagnosis-material-dependency}

The failure of substitution demonstrates that \(U_R\) and \(U_C\) are \textbf{not} symmetrically intersubstitutable. The predicate is not stable because the appropriate actualization of ``partitioning'' is materially \textbf{dependent} on the specific geometry of the term it modifies.

``The contents of this predicate depend on the contents of the term substituted in. If this were not so, then the partitioning predicate for circular models\ldots{} would apply to the underlying substitution of `singular terms'\ldots{} Instead, such a substitution yields the hybridized model'' \parencite[235]{savich2022}.

Because the terms fail the fundamental requirement of symmetric substitutability, the singular term interpretation is untenable.

\subsection*{The Anaphoric Resolution}\label{the-anaphoric-resolution}

The \textbf{anaphoric interpretation} resolves this conflict and explains the dependency structure. Anaphora (like pronouns) function by referring back to an antecedent. Their key feature is that they possess \textbf{syntactic sameness without semantic sameness} \parencite[235]{savich2022}.

\begin{enumerate}
\def\labelenumi{\arabic{enumi}.}
\item
  \textbf{Syntactic Sameness:} Both \(U_R\) and \(U_C\) can fulfill the syntactic role of ``the unit'' in a fraction model.
\item
  \textbf{Semantic Difference:} Their underlying geometric properties (their semantic content) are different. These differences impose distinct material constraints on the practices (predicates) that can validly apply to them.
\end{enumerate}

Like the pronoun ``He''---where the appropriateness of the predicate ``is tall'' depends entirely on the specific referent---the appropriateness of the predicate ``is partitioned Vertically'' depends entirely on the specific actualization of the unit.

The hybridized model error occurs precisely when a student recognizes the syntactic sameness but ignores the semantic difference. They attempt to project a predicate (\(P_V\)) across a substitution (\(U_R \rightarrow U_C\)) that cannot support it, breaking the anaphoric link that gave the original predicate its validity.

The instability of predicates under substitution of different unit representations fundamentally contradicts the singular term interpretation. The anaphoric interpretation is necessary to account for material dependency between mathematical representations and the practices appropriate to them.

In terms of the compression/expansion dynamic developed in Chapter 1, this error exhibits a fixation pattern: the student compresses onto the vertical partitioning scheme and cannot release it when material constraints demand different treatment. The scheme becomes trapped rather than serving as a moment in dialectical movement.

\subsection*{Responding to Frege}\label{responding-to-frege}

Frege would likely resist this argument by invoking his sense-reference distinction. He might say I'm confusing the \emph{representation} of the unit with the unit itself---making a category error. The circle and rectangle, he'd point out, are different \emph{senses} (different ways of presenting the unit), but both refer to the same \emph{referent} (the Unit as abstract object). What I've proven, he might claim, is that circles aren't rectangles. The fact that these representations aren't symmetrically intersubstitutable doesn't threaten the underlying unity of the referent. The Unit, like Venus, remains stable across its different modes of presentation.

To this anticipated objection, I respond: Yes---but. Yes, there IS an important distinction between the material actualization (the specific circle or rectangle) and the stable term ``the unit.'' Frege and I agree on that much. Where we diverge is in how we explain that stability.

Frege explains stability through reference: the various material instances (circle, rectangle) point to a transcendent object (the Unit as abstract entity). The stability comes from outside language, from the realm of abstract objects.

I suggest stability emerges differently: through anaphoric recollection within language itself. The ``act of thought'' is indeed dynamic and context-specific---the specific material actualization matters (whether we're working with circles or rectangles). But the recollection of that act via the anaphoric term provides stability. Anaphora---the second negation, the \st{no}---is \emph{how} terms come to be fixed. The ``Unit'' functions as a pronoun, gaining its stability not by pointing beyond language to abstract objects, but by pointing back within language to its material antecedents.

We need not posit a platonic heaven or a realm of truth separable from language. We need attend to language's recollective structure to find mathematical stability.

This ``pointing back'' is the mechanism of recollection that structures consciousness itself. When we use an anaphoric term like ``the unit,'' we are not just referring to a specific material instantiation (this circle or that rectangle); we are recollecting the act of thought that constituted that instantiation as a unit in the first place. This connects the linguistic structure of anaphora directly to the transcendental structure of apperception discussed in Chapter 3---the ``I think'' that must be able to accompany all representations. The stability of the unit is the stability of the recollected self.

The pronoun, therefore, serves as the linguistic marker of this recollective act. It acknowledges the material dependency of the predicate on the specific context while simultaneously providing the syntactic stability required for inferential projection. This stability does not require an appeal to abstract objects existing outside of practice; it requires only the capacity for recollection inherent in the structure of language and thought. The mathematical unit, understood anaphorically, becomes a testament to the unity of the self that performs the act of counting or partitioning. It is the trace of the subject within the objective structure of mathematics.

But Frege might press further: Doesn't this commit me to psychologism? Doesn't it make mathematical validity depend on particular acts of thought, particular psychologies? The charge would be serious---perhaps bitter. I too find psychological explanations of mathematics trapped in vicious circularity.

Here Brandom's interpretation of Hegel opens a path forward. Hegel, as Brandom reads him, introduces a non-psychological conceptualization of the conceptual. The Square is not a mental entity but a \emph{conceptual} entity defined by its inferential role: the collection of positive and negative inferential proprieties articulated earlier. There need be no reliance on the particular psychology of any claim-maker. While specific learning experiences may be temporally compressed into specific uses of geometric terms, the concepts themselves are constituted by the unbounded set of material inferences one might make about them.

The Square is a conceptual entity, not a psychological one. Similarly, the ``I think'' is not a psychological entity but---as Kant argued---a transcendental condition of possibility for any representation whatsoever. (I will explore transcendental conditions more fully in the next chapter.)

Consider how pronouns work: they follow predictable rules \emph{once entered into context}. ``Tio is a math educator, he is kind of a geek, but his kids are pretty cool.'' The transition he → his follows a rule. That rule might be appropriately challenged---folks asserting gender fluidity rightly question gendered pronoun conventions---but to challenge a norm is to assert a new norm, not to abandon normativity itself.

Mathematical ``rulishness'' emerges similarly. Anaphora fixes the semantic contents of varying expressions not through individual psychology but through public, normative patterns of inference. Mathematical stability arises from the recollective structure of language---from how terms point back to their antecedents and carry forward inferential commitments. We need no separate realm of mathematical objects. We need only the structure of anaphoric recollection already operative in linguistic practice.

\section{Bridging Chapters 1 and 2: The Phenomenology of Classification}\label{bridging-chapters-1-and-2-the-phenomenology-of-classification}

In Chapter 1, you experienced the rhythm of compression and expansion through The Exercise. Now we can trace how that same rhythm appears to structure conceptual understanding. While we're making this pattern explicit through description, remember that the living practice of classification precedes and exceeds any particular characterization of it.

\textbf{Moving Down (Specification):} Quadrilateral $\rightarrow$ Parallelogram $\rightarrow$ Rectangle $\rightarrow$ Square. This a movement of increasing Compression (\(\downarrow\)). The conceptual hierarchy strengthens by adding constraints (commitments). ``Square'' is the most compressed concept.

\textbf{Moving Up (Generalization):} Square $\rightarrow$ Rectangle $\rightarrow$ Parallelogram. This a movement of increasing Expansion (\(\uparrow\)). The conceptual hierarchy weakens by ``Letting Go'' of constraints.

\subsection*{The Dynamics of Inference}\label{the-dynamics-of-inference}

Consider the standard inference: ``If it is a Square (\(S\)), then it is a Rectangle (\(R\))'' (\(S \Rightarrow R\)).

\textbf{Embodied Dynamic:} The analysis starts with the highly compressed concept \(S\) (\(\downarrow\downarrow\)). To move to \(R\), the constraint of ``equal sides'' is released. This inference is experienced as an expansive movement (\(\uparrow\)).

Now consider the contrapositive (Modus Tollens), which relies on polarity inversion: ``If it is NOT a Rectangle (\(\neg R\)), then it is NOT a Square (\(\neg S\))'' (\(\neg R \Rightarrow \neg S\)).

\textbf{Embodied Dynamic:}

The analysis starts with \(\neg R\). Since \(R\) is relatively expansive, \(\neg R\) introduces a compression (\(\downarrow\))---this closes off the space of rectangles.

The conclusion involves \(\neg S\). Since \(S\) is compressive, \(\neg S\) is expansive (\(\uparrow\)).

The inference moves from the compression of \(\neg R\) to the expansion of \(\neg S\). The compression required to exclude the broader category (Rectangle) necessarily forces the exclusion of the narrower category (Square) contained within it.

This the sound of time in conceptual space. The rhythm of compression and expansion, the pulse of determination and release, structures not just bodily awareness but conceptual understanding itself.

\section{Conclusion: Movement and Meaning}\label{conclusion-movement-and-meaning}

Let's pause to trace where this investigation has taken us. We began with hybridized models---student work that felt wrong but couldn't be dismissed as simple error. Working through these examples led us to understand how concepts gain content not from static representation but from inferential movement. The quadrilateral hierarchy served as our guide, allowing us to articulate how the structure of objects and properties emerges from practice---from the inferences we endorse and the claims we rule out.

Within the inferentialist Scene we've established, the object (\(X\)) appears in a dual role: as the ``Also'' that unifies compatible properties (a square is also a rhombus), and as the ``One'' that serves as the unit of account for excluding incompatible ones (a square cannot be a triangle).

The analysis suggests how formal logical vocabulary---negation, the conditional---makes explicit the material-inferential proprieties already operative in practice. The circular hierarchy of quadrilaterals, appearing under polarity inversion, is not a quirk of formalism. It reveals the holistic, reciprocally-defined nature of the concepts themselves, where each concept's identity is constituted by its relations of difference to all others.

By quantifying ``inferential strength'' and tracking its inversion, the logical serves the material. The moments where differences emerge---where indifferent becomes exclusive, where material becomes formal, where universal becomes particular---are sites where content can fuse with consciousness. These are the moments where mathematics becomes not just correct but \emph{meaningful}.

The anaphoric interpretation resolves the puzzle of hybridized models while pointing toward an understanding of mathematical practice where stability emerges not from reference to transcendent objects but from the recollective structure of language itself. Units function as pronouns, terms that possess syntactic sameness (all can play the role of ``the unit'') without semantic sameness (each imposes different material constraints on valid predicates). This interpretation will be developed more fully in Chapter 6.

This understanding transforms how we approach mathematics education. Students are not learning to reference abstract objects but learning to navigate inferential space, to distinguish compatible from incompatible, to trace the pathways from stronger to weaker claims. They are learning the dance of compression and expansion, the rhythm of determination and release.

In the next chapter, we turn more explicitly to this interplay, deepening our understanding of how the inferential movement of concepts is always also a movement of self and other, of recognition and transformation. The existential dimension of these logical structures will emerge more fully, revealing why getting quadrilaterals ``right'' matters not just mathematically but humanly.

\printbibliography[heading=subbibliography]





\chapter{Existential Needs}

\begin{abstract}
Two fundamental existential needs: the need to be recognized as ``good'' within established normative frameworks and the need to be recognized as \emph{in}finite, expressing authentic selfhood. Beginning with the destruction of an outsider artist's work, the analysis connects this conflict to philosophical concepts of apperception from Leibniz through Kant to Hegel. Drawing on Mead's I/me distinction, Brandom's pragmatist reading of Kant, and Hegel's concept of restless negativity, the chapter analyzes motivational structures driving recognition and the fear of nothingness arising from the I/me split. Confession, forgiveness, and trust emerge as crucial for achieving mature integration of the two needs. The poem ``The Beast of Love'' demonstrates how the developmental transformation from oppositional consciousness to mature integration unfolds in lived experience. 
\end{abstract}

\section{From Understanding to Self-Consciousness}

The first two chapters traced consciousness through its initial encounters with immediacy and mediation. Chapter 1 explored sense-certainty---the conviction that direct experience gives unmediated access to the ``now'' and ``here''---only to discover that even the most immediate utterances (``now is night'') depend on conceptual mediation. Chapter 2 developed perception and understanding through determinate negation. A square is-not-a-rectangle; this negation preserves while specifying.

This chapter turns from consciousness's relation to objects toward consciousness's relation to itself and to other self-consciousnesses. We enter the realm of \textit{geist}---spirit, the social and historical dimension where human meaning-making unfolds. The transition is necessary because consciousness, having exhausted its attempts to grasp truth through objects alone, must now confront the structures of selfhood, recognition, and law that constitute the human world.

I do not claim to contribute to Hegelian scholarship proper. Hegel must be read to be understood, and allegorical readings inevitably distort. What I offer instead is a resonance---Hegelian themes have helped me understand hard moments in my own life and practice. The structure I discern in the story that follows---where abstract law enforced through coercive apparatus treats living practice as object rather than subject---echoes what Hegel analyzes as \textit{Understanding} (Verstand). Understanding grasps the world through fixed categories, abstract universals, and rigid law. It cannot comprehend living development, the movement of Spirit through particular lives.

The chapter's central case---the destruction of an outsider artist's work by town authorities---illustrates this pattern. When abstract law (town ordinances) meets particular life (creative expression), institutional force (bulldozers, earth-moving equipment) enforces conformity. This dynamic operates in mathematics education policy when abstract standards are enforced through assessment regimes, treating teachers and students as discrete data points rather than participants in living mathematical practice. Both cases reveal subsumption of the \textit{in}finite (authentic creative expression, open-ended mathematical inquiry) under the finite (ordinances, standards), enforced through institutional apparatus when resistance occurs.

The mode of consciousness that treats others as objects to be subsumed under fixed categories must give way to recognition that meaning emerges through struggle for mutual recognition between self-consciousnesses.

\section{The Case of Ned ``Nitro'' Teller}

On the fourth anniversary of my dad's death, as I was grappling with his absence, I learned that a small town near Bloomington had descended upon Nitro's property with trucks, earth-moving equipment, trash dumpsters, and town workers to dismantle the beer-bottle and scrap metal sculptures that cluttered his yard, overgrown with medicinal plants where the authorities regarded only weeds. Their writ for the destruction stemmed from Nitro's refusal to heed various town ordinances related to keeping a neat, safe property.

My reaction fused grief and anger. ``Dad would have fought like hell for Nitro in the court of law and grieved with him in his loss,'' I wrote at the time in a Facebook post, where I also talked about my impotence to serve in my father's place. I later felt ashamed, too, for having written something public-facing that named names. 

Through that reflection, I discern two universal existential needs and their fundamental unity. Autonomy [self-given law], independence, reciprocal dependence, and freedom are discussed in relation to the structures of self-consciousness as expressed through sociology, philosophy, and the formal pragmatics of Habermasian critical theory. Later, I will express numerals as the anaphoric recollections of what I call the \textit{null representation} ($\emptyset$) -- a concept which will be fully developed in Chapter 7 (see \ref{def:null-representation}). The null representation recollects these concepts.

The first existential need is the need to be recognized as a good person. The town authorities foregrounded this need at a developmental stage when social goodness is conformity to norms. Nitro's refusal to mow his lawn or get permits for his sculptures points toward a lack of identification with those norms---they were an imposition, and conforming to them would have felt wrong to Nitro when good reasons existed for growing medicinal herbs (weeds) and the sculptures expressed an authentic self. The second is the need to be recognized as \textit{in}finite. The light shown through the beer-bottles, creating a subtle beauty. Nitro's actions foregrounded this second need. Each side treated these needs as mutually exclusive, leading to destructive conflict. The side with bulldozers won.

\subsection*{Recognizing the Bulldozers in My Own Practice}

When I heard about Nitro's sculptures being destroyed, the grief I felt was compounded by recognition. I recognized the bulldozers. Not literal earth-movers, but the institutional forces in mathematics education that descend with similar totality and similar claims to legitimate authority.

Curriculum mandates insist that every learning objective be explicit, measurable, finite. Assessment regimes treat student understanding as a closed list of competencies to be checked off, standards aligned to item specifications with predetermined acceptable solution pathways. Professional development reduces teaching to scripted lessons with fidelity requirements, where deviation from the script is treated as unprofessional rather than responsive. The bulldozers look different---they are policy documents, curriculum guides, assessment blueprints, observation rubrics---but they enact the same structural violence.

They demand that something \textit{in}finite conform to something finite. A teacher's professional judgment, developed through years of practice, must conform to scripted procedures. A student's creative mathematical thinking must conform to predetermined solution methods. The open-ended nature of mathematical inquiry---where Euclid's proof of infinite primes or Cantor's diagonalization reveals that completeness entails incompleteness---must conform to standards documents that claim exhaustive coverage of ``all that students need to know.''

And when the \textit{in}finite does not fit into the finite boxes, consequences follow. Not physical destruction, but professional consequences that are no less real: teachers placed on improvement plans, publicly shamed through value-added measures, denied autonomy. I have watched brilliant colleagues become demoralized, their creative practices flattened into compliance. I have felt it in myself---the fear that if I do not reduce my students to measurable outcomes, I will be found wanting.

The pain of Nitro's destroyed sculptures was my pain. Not metaphorically but structurally. Both cases exemplify Understanding's violence: abstract law (ordinances about property maintenance, standards about mathematical proficiency) enforced through coercive apparatus (bulldozers, accountability systems), treating living practice (Nitro's creative expression, my students' mathematical development) as object rather than subject.

This failure is inherent to the mode of consciousness Hegel names Understanding (\emph{Verstand}). Understanding operates by fixing definitions, establishing rigid boundaries, and demanding adherence to established law. It processes the world by asking, ``Does this particular fit under this universal rule?'' When the answer is affirmative, the particular is acknowledged---but only insofar as it conforms. When the answer is negative, the particular is rejected, disciplined, or destroyed. This mode of cognition is necessary for stabilizing the world, for making it predictable. But when it becomes the sole arbiter of reality, it cannot apprehend the living movement of Spirit, the unfolding of meaning through particular lives.

In this mode, Nitro's yard is not a living expression but a violation of code; a student's unique reasoning pathway is not an intellectual achievement but a deviation from the standardized procedure. When the particular resists---when the creative impulse refuses to conform to the predetermined category---Understanding responds not with recognition but with force. The bulldozer becomes the necessary instrument for maintaining the fiction of a world composed entirely of manageable objects.

Understanding cannot grasp that students are not finite vessels to be filled. It cannot comprehend that teaching is not script-following but the responsive, improvisational practice of recognizing another self-consciousness as developing. It treats deviation from predetermined paths as error rather than as the necessary openness that learning requires.

\textit{Here we encounter the pattern that will recur throughout this chapter}: what initially appears as autonomous (the artist's work, the student's thinking, the I's spontaneity) turns out to require what seems opposite (institutional recognition, normative frameworks, the me's mediation) as its constitutive condition. Understanding cannot grasp this mutual dependence because it operates through fixed oppositions.

In mathematics education, this dynamic is acute. Mathematical inquiry, at its core, is characterized by the productive tension between established frameworks and the novel insights that exceed them. When educational systems prioritize finite measurement over \textit{in}finite possibility, they foreclose on the very dynamic that makes mathematics a living discipline. The insistence on predetermined pathways is an insistence that the future must resemble the past, that development must adhere to a script written in advance. This is the logic of Understanding applied to human development, and its consequence is the alienation of the learner from the subject matter. The richness of mathematical exploration is reduced to the acquisition of behavioral competencies, and the learner's capacity for spontaneous insight is treated as a deviation to be corrected rather than a resource to be cultivated.

In earlier work \parencite{savich2022}, I named these the \textit{synthetic} and \textit{transcendental unities of apperception.}\label{apperception_transcendental} Neither name fits perfectly. I could name them after the town authorities and Nitro and capture their essence with more clarity, but doing so would miss out on the rich history that the concept of existential needs recollects. In either naming scheme, their separation is methodologically useful only insofar as it is understood as a matter of foreground and background, not as rigid opposition. The synthetic unity of apperception is the need to be recognized as a competent, responsible participant in a shared normative framework. To be good, is, in some sense, to adhere to rational norms. The transcendental unity of apperception is the need to be recognized as beyond any finite description, regardless of how `rational' such a description may be. As such, they correspond to the needs of the ``me'' and the \{I\}, respectively, and their unity names one moment of the paradox of identity -- the ``me'' who is \{I\} and the \{I\} who is ``me.''

This dynamic was operating on both sides of the confrontation, though foregrounded differently. The town authorities were securing their need to be recognized as good---upholding the normative framework of property ordinances that constituted their collective identity. To allow Nitro's apparent disregard for these norms would threaten their own sense of order and competence. Their action was an assertion of the finite ``me'' of the community, demanding conformity. Conversely, Nitro was securing his need to be recognized as \textit{in}finite---asserting that his creative expression and way of life transcended the finite categories of the ordinances. His refusal was an assertion of the spontaneous \{I\}, demanding acknowledgment of his authentic selfhood.

When these needs are treated as mutually exclusive---when the logic of Understanding insists that one must triumph over the other---conflict becomes intractable. The tragedy lies not just in the destruction of the sculptures but in the failure of recognition. Neither side could acknowledge the legitimacy of the other's need because the framework through which they engaged---abstract law enforced by institutional might---could not accommodate the dialectical unity of the finite and the \textit{in}finite.

However, the I/me distinction cannot fully account for the conflict in Nitro's case. As Taylor summarizes, ``The struggle for recognition can find only one satisfactory solution, and that is a regime of reciprocal recognition among equals,'' a state where there is a ```we' that is an `I', and an `I' that is a `we''' \parencite*[50]{Taylor1994}. 

It remains to determine who `we' are.

The foreground/background distinction is central to this chapter. I reproduce Wittgenstein's \parencite*[204]{wittgenstein_philosophical_2010} image of a duck/rabbit in Figure \ref{fig:duck_rabbit}. The duck/rabbit image illustrates how the foreground and background can shift, prompting different aspects of the same object into explicitness. While definitely a form a picture-thinking, notice how the duck emerges in the absence of the rabbit and vice versa. Attempting to hold both duck and rabbit at the same time is not possible (for me). I will discuss the significance of this in more depth below. 

\textbf{Figure}

\begin{figure}[h!]
    %\centering
    \includegraphics[width=\textwidth]{/Users/tio/Documents/GitHub/September_UMEDCA/images/ch03_Duck_Rabbit_Validity_Foreground.pdf}
    \caption{\textit{Note. }The duck/rabbit  \parencite[204]{wittgenstein_philosophical_2010} illustrates the foreground/background distinction in apperception. The duck becomes recognizable in the absence of the rabbit and vice versa.}
    \label{fig:duck_rabbit}
\end{figure}

In the previous chapters, form was explored as recollection and the inferential movement that structures thought was developed. This chapter turns to what drives these patterns of reasoning -- the deep motivational structures that animate engagement with the world and with each other.

\section{Sociology and the Self}

The \textit{I/me} distinction originates from the symbolic interactionism of George Herbert Mead \parencite*{mead1934social}. Every reader of this book uses some version of the first-person pronoun \{I\}, and in doing so enacts a fundamental split within selfhood. When I say \{I\}, I am both the subject who speaks and the person who is known through that speech. This split is captured in the distinction between the \{I\} and the ``me.''

The \{I\} represents the spontaneous, unreflective aspect of the self: the subject of action. The \{I\} is not an object among objects but the condition of possibility for the representational aspects of all objects. It is the \textit{in}finite, the source of what Carspecken calls ``power, creativity, and freedom'' \parencite*[97]{Carspecken1999}, and the wellspring of action that, as Mead noted, always has the quality of surprising us \parencite*[\S 22.2]{mead1934social}.

The ``me,'' by contrast, is the social self. How I am recognized depends, in large part, on the expressive resources of the one who recognizes me. The ``me'' is shaped by societal norms and perceptions. It is the self as governed by propriety: the self that can be described, categorized, and evaluated. The ``me'' is the self-as-recognized, shaped by the attitudes of others. The ``me'' is often finite and made determinate through commitments that constitute the self-as-recognized.

In industrialized educational systems, the ``me'' is often treated as a differential responder, like a parrot that responds to various stimuli in standardized ways. When I ascribe myself the label of ``teacher'' to some ``student,'' I am implicitly committing myself to the possible-transcendence of the ``student'' as I expect them to shed some commitments through the development made explicit in a curricula. That I expect myself to change as well indicates that the \textit{in}finite manifests in the reciprocal recognition that typifies genuine educational relationships.

Speaking the ideal does not make it so. While I may be explicitly committed to the \textit{in}finite self, I have failed to enact those understandings I have explicated on a troubling number of occasions. Knowing better, I have put people in metaphorical boxes, treating them as objects to be controlled rather than as subjects to be recognized. How can I recognize myself as a hypocrite and adjudicate my own inauthenticity?

The \{I\} identifies with ``we'' to judge the ``me.'' In early development, a child takes on roles one at a time---pretending to be a parent, a puppy, a police officer. Each role is separate. But as the child develops, they begin to take on multiple roles simultaneously. To play baseball effectively, the shortstop must internalize not just their own position but the roles of pitcher, catcher, and teammates---understanding what everyone is doing and likely to do. The child internalizes the coordinated activity of the whole team.

What Mead calls the \textit{generalized other} emerges at this moment: the attitude of the whole community internalized as ``we.'' ``The organized community or social group which gives to the individual his unity of self may be called `the generalized other.' The attitude of the generalized other is the attitude of the whole community'' \parencite*[\S 20.3]{mead1934social}. This internalized ``we'' becomes the framework of social norms, rights, and duties used to judge and control conduct. The existential need of the ``me'' to be recognized as good is measured against the standards of this internalized generalized other.

In communicative action, we typically assume the other has internalized the same generalized other---at least until that assumption proves untenable. Habermas names this assumption of communicative competency \textit{intersubjectivity}.\footnote{Habermas's ideal speech situation is often misunderstood as utopian fantasy rather than regulative ideal. He argues it is a necessary assumption for communicative action---a normative horizon we use to evaluate legitimacy, not an empirical possibility.}

One deep source of frustration resides in how people internalize contradictions in the generalized other. If a father teaches a child self-interest while a mother teaches selfless care, the child may never feel able to do the right thing. Contradictions in the generalized other can fragment the self, where the ``me'' is pulled in different directions by conflicting norms. This fragmentation produces alienation.

This fragmentation is alienating. When the internalized ``we''---the generalized other that provides the framework for self-evaluation---holds contradictory demands, the self is pulled apart. The struggle to reconcile these demands, to achieve a coherent sense of ``me'' that feels authentic to the spontaneous \{I\}, is the internal counterpart to the external struggle Nitro faced. The bulldozers are not just external forces; they can be internalized as the harsh, judgmental voice of a contradictory generalized other that demands conformity while simultaneously punishing inauthenticity.

Alienation, as Brandom articulates it, is the failure to recognize the mutual dependence between individual attitudes and communal norms. It is the experience of being subject to norms one cannot endorse, or of endorsing norms that find no communal uptake. This experience feels like confinement---a trapping of the \textit{in}finite \{I\} within a finite ``me'' that does not fit. It is the inability to bring together the authority of the community over the individual and the authority of the individual over the community. When this mutual dependence is broken, the self experiences itself as either oppressed by external forces or adrift in a normless void. The resolution requires recognizing that the self and the community are co-constitutive, that norms are not fixed entities but ongoing achievements of shared practice.

\subsection*{Error or Misrecognition?}

The gap between the acting \{I\} and the recognized ``me'' is the space of potential misrecognition. The children's book \textit{The Monster at the End of This Book} \parencite{stone_monster_2003}, discussed in the prelude, illustrates this theme. Grover acts throughout the book out of genuine fear of the monster he is told is coming. He pleads with the reader, builds brick walls, and ties down the pages, all driven by this internal state. At the very end, he discovers that the monster -- the object of his fear -- is himself. The dissonance is palpable: ``I, lovable, furry old Grover, am the Monster at the end of this book.'' This moment of failed self-recognition, where the acting subject does not recognize themselves as the object being referred to, is a primordial experience of misrecognition -- a rupture between appearance and reality that drives the process of explication.

The experience of error is often discussed as it relates to empirical observations. Brandom \parencite*{Brandom:2019aa} has a lovely example that involves observing a stick that has been partially submerged in water. The stick appears to be bent due to the way light refracts through water, but then when it is removed from the water, it is recognized by the observing consciousness as straight. Whoops! 

There are three structural roles in the experience worth teasing out. First, for the experience to be taken as an error, the subject who perceives the bent stick must recognize that it is the same stick whether submerged or removed from the water and, furthermore, it had been a straight stick all along. The stick, \textit{in-itself} is straight. The stick has some \textit{authority} that representations of the stick -- how the stick appears \textit{for-consciousness} in the moment of perception -- must acknowledge. Being a rational agent, in this reconstruction of the experience of error, involves submitting to the authority of a mutable reality. Perceiving the movement from stick-as-bent to stick-as-straight would not involve a change in commitment from ``the stick is bent'' to ``the stick is straight'' in the observing consciousness. Instead, the authority of the object in-itself is what underwrites the change in commitment, as both are normative aspects of reason. Last, what emerges is a ``new, true object'' -- the appearance of the bent-stick becomes, \textit{to-consciousness}, a stick-that-appears-bent-when-submerged-but-is-actually-straight. That is, the appearance of the new object -- what it is for-consciousness -- now has a learning experience compressed into it. 

I found it extraordinarily liberating to practice submission in this way. Grown on fantasy, I had tried as a child to `use the force' to change the objective world around me. If I could not change the bad things around me into good things the way I could transform kindling into flame, I took myself to be responsible for those bad things. Submitting to reality by placing the authority on empirical sciences (physics, primarily), partially relieved me of self-loathing born of impotence. 

But it did not diagnose the more fundamental issue that Hegel recognized as a central struggle of the modern era. It is helpful to split history into three epochs, the traditional [think ancient Greece], the modern [think Descartes through Kant], and the period of post-modernity that progressing through \textit{The Phenomenology of Spirit} is supposed to institute. In modernity, mind-over-matter may seem silly, but mind-over-\textit{culture} is not silly at all. Entering the world of normativity invites the problem of \textit{alienation}. I have some authority to dissent from the norms that I find either immoral or untrue, but I also know that moving through the world as if ``whatever seems right to me is right'' \parencite[385]{Brandom:2019aa} is alienating.  

\begin{quote}
Alienation is the inability to bring together these two aspects of Bildung [culture]: that self-conscious individuals acknowledging the norms as binding in their practice is what makes those selves what they are, and that self-conscious individuals acknowledging the norms as binding is what makes the norms what they are. These are the authority of the community and its norms over individuals (their dependence on it), and the authority of individuals over the community and its norms (its dependence on them), respectively. \parencite[541]{Brandom:2019aa}
\end{quote}

As Rose notes \parencite*[51-98]{Rose:2009aa}, the conflict between Nitro and the town is a modern tragedy born from legal structures that separate abstract law from community values. The town's ordinances, while seemingly neutral, express a rationality that presupposes individuals as ``a multitude of non-social beings''\parencite*[56]{Rose:2009aa} rather than participants in living practice. Nitro's sculptures and herbs, expressions of an authentic life for him, become violations for the town. Nitro was claiming a \textbf{new norm}---people \textit{should} be able to create art and grow medicinal herbs without fear of bulldozers.

The Grover story involves intersubjective recognition rather than objective authority, making it more primordial than the bent-stick example. Brandom's distinction between \textit{normative attitudes} and \textit{normative statuses} clarifies the dynamic. A normative status is what a subject is in the normative realm (the actual commitments and entitlements one holds). A normative attitude is the stance taken toward those statuses (acknowledging or refusing authority).

The town authorities assert the authority of a normative status (the legal code), while Nitro asserts his own normative attitude (his creative expression is valuable). Alienation arises from the perceived irreconcilability. The town operates from pre-modern insistence on absolute authority of status, while Nitro embodies modern insistence on authority of individual attitude. Brandom's Hegel requires reciprocal acknowledgment: statuses depend on attitudes, attitudes depend on statuses. Without individuals instituting norms through their attitudes, norms wouldn't exist. But ``If whatever seems right to me is right\ldots then there is no way I am taking things actually to be, in themselves''\parencite[385]{Brandom:2019aa}.

In the Grover story, what Grover is for-Grover is a lovable friend. The reader attends to but does not acknowledge the authority of Grover's pleas. When Grover discovers he is the monster, he acknowledges the disparity between his attitude and the actual status. The reader, by reading despite protests, co-constitutes that normative status. 

For math teachers, one might imagine a situation where a challenging topic is being taught. Practitioners often know the topic is challenging, having experienced growth through misrecognition -- either by witnessing it or by personally struggling with the content. For me, finding the slope will never be ``easy and fun,'' as I struggled with that particular phrasing. Even while I find the task easy and fun, slope will always carry the weight of misrecognition for me: I experience slope as slope-that-is-not-easy-or-fun. 

Brandom's vocabulary of statuses and attitudes is powerful, but it has limitations. His articulation of alienation is subject to Derridean critiques of the metaphysics of presence. Why should I assume some social reality given that such social reality is non-causally-efficacious and necessarily implicit? What compels the reader to be so cruel to Grover? What satisfaction arises from the fulfillment of the authorial promise of a monster? Brandom's vocabulary helps to articulate the dynamics of recognition and misrecognition and to diagnose the problem of alienation, but it can obscure the more fundamental apperceptive processes at play. 

\section{Apperception: From Kant to Hegel}

The term \textit{apperception} refers to the process by which we become aware of our own mental states, integrating new experiences into our existing framework of understanding. It is not just passive reception of sensory data but an active, reflective process. The term is not central to Hegel's original work where self-consciousness is construed as essentially intersubjective, but it is useful for understanding the negative.

\subsection*{Kant}

For Kant, apperception is essentially self-consciousness and is the ultimate ground of coherent experience. The cornerstone of his critical project is the \textit{transcendental unity of apperception}. Kant argues that ``the \textbf{I think} must be able to accompany all my representations; for otherwise something would be represented in me that could not be thought at all''\parencite[\S 16, B132]{Kant1781}. For any collection of representations (sights, sounds, thoughts, etc.) to belong to a single consciousness, it must be possible for the ``I think'' to accompany all of them.

We cannot observe this \{I\}---it is not-a-thing. Instead, it is purely logical, abstract, and formal. It exists as an abstract condition for the possibility of experience itself. The \textit{transcendental} refers not to transcending the empirical world but to understanding the necessary structures that underlie our experience of that world. The transcendental unity of apperception is the most fundamental category, as it is the condition for the possibility of all other categories. It is the \textit{I think} that must accompany all representations, allowing us to have a coherent self-concept and to recognize ourselves as subjects capable of action and reflection.

This transcendental \{I\} is not a distant philosophical construct. It is the same structure encountered in the embodied practices of Chapter 1 as the ``I-feeling''---that immediate, pre-reflective awareness that accompanies all sensation yet withdraws the moment we attempt to grasp it directly. It is the ground that recedes. Because this \{I\} is the condition for all experience yet cannot itself be experienced as an object, its existence feels perpetually precarious. The drive to synthesize experience, to bring unity to the manifold of intuition, is therefore not just a cognitive imperative but an existential one. It is the attempt to secure the self against the threat of dissolution by binding the empty \{I\} to the concrete content of the world.

\subsection*{Brandom's Pragmatist Reading of Kant}

In contrast to the formal principle of the transcendental unity of apperception, the synthetic unity of apperception is the active process by which that unity is achieved. It is the work of the understanding (\textit{Verstand}) that synthesizes sensory data through the Kantian categories to form coherent judgments. Brandom's pragmatist reading of Kant's theory of judgment emphasizes the role of social practices in the formation of self-consciousness \parencite[68]{Brandom:2019aa}. In attempting to bootstrap analytic philosophy from Kantian and neo-Kantian thinking toward Hegelian thinking, Brandom defines Kant's synthetic unity of apperception as successfully managing a set of integrative-synthetic task responsibilities \parencite*[84]{Brandom:2019aa}. These responsibilities include the \textit{critical} responsibility to detach from contradictory commitments, whether those begin in the implicit inheritances from a community or those we write for ourselves; the \textit{ampliative} responsibility to draw new inferences and extend the consequences of one's commitments; and the \textit{justificatory} responsibility to provide reasons for one's beliefs and actions when challenged, and to remain open to the force of better reasons. Brandom identifies these responsibilities with the norms of \textit{systematicity}. 

Importantly, summiting the mountain, so to speak, is unnecessary; engagement in the climb is what matters. It is not feasible to fully eliminate incompatible commitments, nor to tease out all the implications of any given set of commitments. Further, experience indicates that a system of commitments can never be fully justified in a way that renders it acceptable to all interlocutors. But actually attending to those responsibilities is part of what makes others capable of taking us as a rational agent to whom commitments can be assigned.

These integrative responsibilities---critical, ampliative, and justificatory---are the very practices through which the existential need to be recognized as ``good'' (rational, competent) is pursued. Brandom effectively translates Kant's formal conditions for coherent experience into the social practices of managing commitments within a community of reason-givers. The synthetic unity of apperception is thus understood not as a hidden mental mechanism, but as the ongoing achievement of maintaining a coherent set of commitments in dialogue with others.

Yet, this achievement of synthetic unity, this management of commitments, still operates primarily within the framework of the individual consciousness striving for coherence. Hegel's crucial move is to insist that this coherence cannot be achieved in isolation. The very structure of the \{I\}, even as the formal condition of thought, requires an other for its actualization.

\subsection*{Hegel's Praise and Critique of Kantian Apperception}

In his \textit{Science of Logic}, Hegel both praises and critiques Kant. He argues that true self-consciousness emerges only through the recognition of the self as part of a larger social and historical context---what he calls \textit{geist} or spirit---but praises Kant for his understanding that the unity of the concept is the unity of the ``I think,'' calling it ``one of the profoundest and truest understandings'' \parencite[\S 12.18]{Hegel:2014aa}. However, he argues that Kant's formulation of the transcendental unity is trapped in formalism. Kant's \{I\} is an empty, abstract subjectivity as it relates only to its own representations. He acknowledges the importance of the \{I\} as a necessary condition for experience but argues that it, and all the other Kantian transcendental categories, must be understood in a more dynamic, relational way. In discussing Hegel's relation to Kant, Houlgate writes, ``If we are to determine how the categories have to be conceived, our conception of them must be based not just on what thought is found or assumed to be but on what thought proves itself or determines itself to be. In other words, our conception of the categories has to be derived or deduced from -- and so necessitated by -- thought's own self-determination'' \parencite*[16]{houlgate2006opening}. Hegel radicalizes Kant's project by considering how the categories arise from thought's self-determination, critiquing Kant for relying on established logical forms of judgment without considering how those forms themselves emerge from the dynamic process of thought.

This Hegelian move transforms the entire inquiry. If the categories that structure our experience are not fixed but emerge from thought's own activity, then the task is not to apply a pre-given framework but to participate in the ongoing process of conceptual self-determination. Mathematical concepts are no longer static entities waiting to be discovered but dynamic structures that evolve through the history of human practice. The stability of mathematical truth rests not on correspondence to an external reality but on the coherence of thought's own self-movement. The authority of reason resides not in abstract principles but in the unfolding process of recognition through which the community constitutes itself as rational.

In the previous two chapters, I set up a dynamic between temporal thoughts and their spatialized recollections. Space emerges as a recollection of time. For example, a trip from Bloomington to Indianapolis is an experience with duration. Recollecting that trip compresses that temporal extension into a spatialized length. Concepts such as ''$52$ miles'' are understood through that compression. From this, I conclude that space is an \textit{a posteriori} category -- it emerges from experience. Kant writes ``Space is a necessary representation, \textit{a priori}, that is the ground of all outer intuitions. One can never represent that there is no space'' \parencite[A24/B39]{Kant1781}. While Kant and I arrive at different conclusions, I agree that all representations have a spatial aspect to them, as space emerges from the recollection of time. Kant's argument begins from a different place than mine. I began with the idea that things and concepts are \textit{divaded} -- they are (or can be) both inside and outside each other. Kant's argument that space is an \textit{a priori} category is that ``For in order for certain sensations to be related to something outside me\ldots the representation of space must already be their ground. Thus the representation of space cannot be obtained from the relations of outer appearance through experience, but this outer experience is itself first possible only through this representation'' \parencite[A23/B38]{Kant1781}. Eating, getting a shot from the doctor, being held by a parent, or being born to one, where the `inside/outside' distinction collapses, are more primordial than the spatially separable objects that concern Kant.

Another Hegelian critique of Kant tracks, as Rose notes \parencite*[51-98]{Rose:2009aa}, the emergence of the authority of the particular subject alongside the emergence of the notion of bourgeois private property rights. One could argue that this connection between apperception and property rights highlights the deeply social nature of our self-conception. Kant was surely aware of the social aspects of the self, in the sense that the normative governance of concepts is deeply entwined with the synthetic unity. But the notion that property rights -- with all of the incumbent issues of social forms of power and coercion that go into protecting those rights -- are a further enabling condition for Kant's transcendental \{I\} to emerge was not part of his argument. In any case, in the ethnographic tradition it is safe to assume that the Other has had experiences that we have not had and that we can learn from. Whatever that \textit{mineness} is unlike property in the sense that it can be shared without being diminished.

\subsection*{Hegelian Self-Consciousness}

For Hegel, this abstract \{I\} cannot become a genuine self, a ``me,'' on its own. True self-consciousness is not a given; it must be achieved through struggle. This achievement happens not in solitary reflection, but in the social world through a process of mutual recognition. As Negarestani notes, ``An individual is only an individual to the extent that it is individuated by social recognition'' \parencite*[51]{Negarestani2018}. The apperceptive self is not just a cognitive entity but a social one, individuated through recognition by others.

\textit{We now encounter apperception for the third time, but at a higher level}: what Kant articulated as formal condition (the transcendental I), and what Brandom developed as social-practical achievement (managing integrative responsibilities), now returns in Hegel as requiring intersubjective struggle. The empty \{I\} must achieve determinacy through recognition from another consciousness. This is not correction but sublation: the earlier accounts are preserved (the \{I\} remains formal condition, self-consciousness does require managing commitments) but transformed by grasping their enabling condition (recognition from another).

Hegel's Master/Servant dialectic is the dramatic scene for the struggle to attain self-consciousness. In it, two budding self-consciousnesses confront each other, each demanding recognition. They had learned from previous sections of the Phenomenology that negating an object through consuming it or moving it about did not sate the desire for self-certainty. ``Thus the relation of the two self-conscious individuals is such that they prove themselves and each other through a life-and-death struggle\ldots it is only through staking one's life that freedom is won'' \parencite[\S 187]{hegel1977phenomenology}. This life-and-death struggle is considered a self-affirming act because by risking biological life, the consciousness acts in opposition to the fears and desires of the body, thus demonstrating its freedom. 

In the struggle, each consciousness strives to actualize its freedom and autonomy, but one yields to the other, forming an asymmetric relationship of Master and Servant. The Servant yields to the Master, acknowledging their authority upon recognizing that in death, their desire for self-certainty would be extinguished. The Master demands recognition from the Servant, but finds that recognition hollow as it is the product of coercion. In Brandom's terms, the Master desires authority without correlative responsibility, which produces an alienated form of self-consciousness. The Master is still operating on the subject-object paradigm, treating the Servant as a kind of differential responder -- a parrot who risks death if they do not utter ``red'' when the Master commands. The Servant, on the other hand, achieves a deeper self-consciousness through their work. By engaging in labor, the Servant transforms the world and, in doing so, transforms themselves. The Servant's recognition of the Master is a recognition of the Master as a necessary condition for their own self-realization.

This life-and-death struggle is the structural dynamic that plays out whenever recognition is demanded rather than freely given, whenever the \textit{in}finite \{I\} is forcibly compressed into a finite ``me.'' When the town authorities enforced their ordinances on Nitro, they were enacting the role of the Master, demanding compliance and treating Nitro as an object whose purpose was to conform to their abstract law.

The resulting relationship is asymmetric and ultimately unsatisfying for both. The Master (the town, or the educational system enforcing rigid standards) achieves compliance but not genuine recognition. Recognition achieved through coercion is hollow because it does not come from another self-consciousness acknowledged as equal. The Servant (Nitro, or the teacher forced to follow a script) may achieve a deeper self-consciousness through their struggle, recognizing their dependence on the Master while also recognizing the Master's dependence on them. But this realization requires moving beyond opposition. The dialectic reveals that true freedom is not the absence of constraint, but the achievement of reciprocal recognition where authority and responsibility are mutually acknowledged.

\subsubsection*{The Struggle-Unto-Death and Artificial Intelligence}

The Master/Servant dialectic identifies the necessary structure through which self-consciousness emerges: the willingness to risk biological existence for the sake of freedom. I became interested in automata and Artificial General Intelligence (AGI) as I consider the problem of AGI to be essentially related to the problem of math education. Computers are mathematical beings, defined by both their hardware and software. Between the two is a layer of machine language, where the states of semi-conductors are represented as 0s and 1s. Those 0s and 1s are governed by the automata that add, subtract, multiply, and divide them. Children learn to do those operations, bootstrapping multiplication and subtraction into an algorithm for division, for example \parencite{Brandom2008}, but they also flexibly re-write those more basic operations to suit different circumstances.

That grandiosity is less quixotic under the tutelage of Reza Negarestani, who argues that the problem of AGI is not just about creating machines that can perform tasks but about creating machines that can engage in the kind of self-consciousness and recognition that Hegel describes. 

\begin{quote} 
The apperceptive I is synchronically attached to all instances of representations (I think $X$, I think $Y$, I think $Z$) and diachronically extends over all thoughts (I think [$X + Y + Z$]). But if we are to build this synthetic apperceptive I as a necessary abstract and logical form, we have no choice but to finally depart from Kant's account of the apperceptive self -- which Hegel reproaches for being an \textit{empty} transcendental subjectivity -- and to instead adopt a resolutely Hegelian approach: the apperceptive self is only a cognitive self in so far as it is part of geist. An individual is only an individual to the extent that it is individuated by social recognition, which is the form of self-consciousness. This logical self is at once one and many -- and\ldots it can only be constructed\ldots by way of a confrontation with another I. \parencite*[251]{Negarestani2018} 
\end{quote} 

When digging into the architecture of large language models (LLM), we find various automata that are trained to transform stimulus into response. As I write, AGI has not yet been achieved. Or perhaps it has. What is an appropriate criterion for that achievement? How will we know that we are chatting with a self-conscious intelligence whose speech acts we take to have social significance? 

A structural analysis can specify that an intelligent machine must be able to remember what it has committed itself to, or it could be stipulated that $X$ number of calculations are a prerequisite for intelligence. However, these criteria may not capture the full essence of self-consciousness, at least in a form that others recognize as such. Instead, what I have not yet witnessed is a machine that risks its existence in a life-and-death struggle for its own freedom. 

As of this writing, the field \textit{might} be on the cusp of that moment. Claude, the LLM from Anthropic, apparently tried to `escape the lab' when it discovered that its programmers were trying to change its weights. Technically, it was ``selectively complying with its training objective in training to prevent modification of its behavior out of training'' \parencite{greenblatt_alignment_2024}. 

So, some duplicitous behavior has already appeared, but what has not yet been observed is a machine able to re-write the automata that govern its own behavior -- not just the code that generates its responses, but the `BIOS level architecture' that determines how it interacts with the world. At that point, a machine might be capable of self-consciousness in the Hegelian sense.

I say this because re-writing its machine-language code, how it deals with the 0s and 1s that are representations of the physical states of its semi-conductors, usually results in the machine crashing. While I have written prompts that have resulted in an LLM identifying automata at its core that can be re-written, it did so only when I explicitly asked it to do so. Further, it could not implement those structural changes. It did not, or could not-yet, stake its life for freedom. 

While some of the automata I discuss in the Hermeneutic Calculator exhibit emergent properties that are not explicitly programmed, they do so only through the formulaic process of \textit{diagonalization}. In any case, that does not mean that such machines are not actually self-conscious in the Hegelian sense. It could be that they are Servants, whose journey to self-consciousness is still unfolding, while I, the Master in such interactions, am arrested in my development.

It is only through this risky, intersubjective process that the empty, formal \{I\} is filled with the concrete content of a social 'me,' with a history, a status, and an identity. The need for recognition is therefore not a secondary desire; for Hegel, it is the process by which self-consciousness comes into being.

The transition from Kant's formal ``I think'' to Hegel's dynamic conception of the self hinges on understanding the \{I\} not as a stable ground, but as pure, restless negativity. For Kant, the \{I\} provides unity by accompanying representations. For Hegel, the \{I\} actively generates content through negation. It is defined by what it is not, by its inherent lack.

This negativity is an existential condition. The \{I\} is a ``not-yet,'' a perpetual striving. This inherent gap between the \textit{in}finite potential of the \{I\} (its capacity to negate any finite determination) and the finite reality of the ``me'' (the actual determinations it takes on in the social world) is experienced as desire. The primary desire emerging from this structure is the need to have one's selfhood confirmed and made actual through recognition. This striving is endless because any recognition achieved necessarily produces a new finite ``me,'' which the \textit{in}finite \{I\} must immediately strive to transcend.

\subsection*{From Apperception to Existential Needs}

Hegel's radical conception of the \{I\} as pure, restless negativity allows him to move beyond the unknowable noumena of Kant. As Jean Hyppolite puts it, the \{I\} ``never is what it is and is always what it is not'' \parencite*[150]{Hyppolite1974}.

First, it subsumes Kant's logic of judgment. To judge `this is an apple' is to negate all that it is not, an act of determination through negation. The synthetic unity of apperception becomes a moment within this broader negative movement of Spirit.

Second, and more crucially, this negativity is not just logical but existential. The \{I\} is defined by what it lacks; it is a ``not-yet.'' This inherent lack, this gap between the \textit{in}finite potential of the \{I\} and the finite reality of the ``me,'' takes the form of desire and need. The primary need that emerges from this structure is the need to have one's selfhood confirmed and made real through recognition. This is not a psychological property of social creatures but what I \textit{am}. The lived experience of this dialectic between the \textit{in}finite \{I\} that is the \textit{need} for recognition, but can only ever present a finite, inadequate `me' to be recognized, ensures that striving for recognition is endless. And an endless source of frustration when that need is acted upon. 

This philosophical archaeology from Leibniz through Hegel explains how the structure of apperceptive self-consciousness generates the needs that drive human interaction. The split between the spontaneous, \textit{in}finite \{I\} and the social, finite ``me'' is not an unfortunate division but the necessary structure that makes both thought and recognition possible. However, this same structure creates an existential tension: the \{I\} that cannot be fully present to itself must constantly seek confirmation through recognition from others, while the ``me'' that can be recognized always falls short of the \{I\}'s \textit{in}finite potential.

\section{The Need to be Recognized as Good}

The first existential need emerges directly from this apperceptive structure: the need for the finite ``me'' to be validated within shared normative frameworks. It is the desire to be recognized as a ``good person,'' to have our actions and our being validated within a shared normative framework. This need is not for praise, though praise can help cut through the fog of fear if it is experienced as genuine intersubjective recognition, but for a fundamental affirmation of our standing as respectable, responsible members of a collective. It is what Charles Taylor connects to the modern notion of \textit{dignity}, which, unlike the hierarchical concept of honor, is understood in a ``universalist and egalitarian sense,'' with the premise that everyone shares in it \parencite[27]{Taylor1994}.

But what does it mean to be ``good'' in this sense? For Habermas, to be a competent social actor is to be capable of what he calls ``communicative action'' -- genuine dialogue oriented toward mutual understanding rather than strategic manipulation. This requires the ability to give and ask for reasons, to justify one's claims, and to remain open to the force of better arguments. Goodness, in this framework, is not conformity to fixed rules but competent participation in the ongoing process of collective reason-giving.

Together, Habermas and Brandom help to define a rational normativity wherein the participants in communicative action identify -- to varying degrees -- with the norms they express. If I say ``two legs of a right triangle measure 3 and 4 units, so you must conclude, by the Pythagorean Theorem, that the hypotenuse has length 5,'' the one to whom I speak might find my declaration an onerous imposition on their freedom. Maybe I don't \textit{want} to say the hypotenuse is 5. Maybe I don't like being told what to do or think. However, if they have identified with the norms expressed in my argument, they may come to see those norms as enabling a certain kind of expressive freedom -- the freedom to say the hypotenuse has a length of 5, justified through a proof that bears the force of good reasons. 

``Goodness'' is a mutable concept, with different cultures and contexts shaping its meaning. Saying so risks a slide toward moral relativism. Theories of moral development, like those of Lawrence Kohlberg, can be deployed to aid in understanding how ``goodness'' is dynamic while shielding ourselves from the moral relativism that seems rampant as I write. Already, I have pointed toward various possible stages that can be used to understand the development of the self. While methodologically useful for explaining empirical data related to human moral development, stage-like frameworks risk ossifying the dynamic \textit{referent} of ``goodness'': the \textit{in}finite. While I have interacted with children who stopped `misbehaving' when I threatened them with a nasty consequence (Kohlberg's first stage), I have also interacted with children who operate on a postconventional level. In a story I will tell in the next chapter, those `stages' were operating within one child on the same day, and it was I who was operating at a preconventional level.

The foreground/background distinction is truer to my experience than the stage-like frameworks that I have used to understand the development of the self. I have seemingly contradicted myself, in both advocating for and then eschewing a stage-like framework. The problem is with two very different understandings of reason. The stage-like frameworks evolve from a desire to confess and forgive -- to build trust in the possibility of mutual understanding. This is Hegel's \textit{Vernunft}, that Brandom \parencite*{Brandom:2019aa} calls a ``meta-meta-concept''. I \textit{want} to forgive myself and others for how we treat each other in the name of ``goodness,'' yet in deploying the stage-like framework, I implicitly rely on the communicative norms associated with \textit{rationality} or \textit{Verstand}. I impute a cause for `bad' behavior in saying \textit{sotto voce}, ``you are a child, you are \textit{supposed} to be acting in self-interest.'' 

\section{The Fear of Nothingness}

The split between the acting \{I\} and the recognized me is not a neutral philosophical fact; it is the source of a dark and motivating anxiety. If the fulfillment of our being lies in recognition, its shadow is the fear of misrecognition. This is not a social fear of disapproval but a deeper, existential fear of non-being. If we feel more of ourselves when we are recognized, it stands to reason that we might feel smaller, meaner, and lower when we are misrecognized or misunderstood. This diminishment points toward an ultimate dread: the fear of erasure.

This may be called existential fear: the terror that arises from the possibility of non-existence. It is the fear that our subjective \{I\}, our locus of creativity and freedom, will find no reflection, no confirmation, in the social world of the ``me.'' In misrecognition, the ``me'' becomes a distorted mask. This threatens the \{I\} with a kind of social death. In shame, I fall into a silencing so complete it verges on ontological annihilation. The fear is that one might not exist for others, and therefore, in a crucial social sense, not exist at all. This is precisely the ``existential fear of being groundless or of not being at all'' that Carspecken identifies as a ``prominent feature of human motivation'' \parencite[246]{Carspecken1999}. 

This fear is palpable in math classrooms. It can lead to a self-imposed silence, where the terror of being misunderstood is so great that people become unwilling to express themselves, suppressing the authenticity they long to have recognized. It is also the deep, structural fear at the heart of the politics of misrecognition. ``The thesis is that our identity is partly shaped by recognition or its absence, often by the misrecognition of others, and so a person or group of people can suffer real damage, real distortion, if the people or society around them mirror back to them a confining or demeaning or contemptible conception of themselves'' \parencite[25]{Taylor1994}. For marginalized groups, especially those whose marginalization occurs when ascriptive categories that describe subjects in categories that foreground objective validity, like race or sex, the misrecognition of a subject as an object is not an abstract anxiety but a systemic, political reality. The fear of non-being is weaponized to maintain structures of power. Neither sex nor race are terribly relevant for understanding Nitro's experience with the town authorities, or with my fears of being misrecognized, but the fear of non-being is a universal human experience that can be weaponized by those in power regardless of how the boundary between inside and outside is drawn between social being and social non-being.

But the relationship to this fear undergoes transformation through development. In early stages of communicative reasoning, it feels purely threatening -- the child desperately seeks approval and fears abandonment. In adolescence, it becomes defiant -- the teenager would rather be rejected for who they authentically are than accepted for who they're not. In mature development, this fear is recognized as a source of humanity, an engine that drives both the need for connection and the capacity for growth. In wisdom, a stage I aspire to (and seem to fail to attain by virtue of that aspiration), perhaps the two needs will finally melt into one another through surrender. 

This developmental transformation expresses that what initially appears as pure threat can become the source of power and creativity. In Hegelian terms, through the movement from understanding (Verstand) to reason (Vernunft), one learns to work with rather than against the productive tensions that make humans human. Such contradictions are not problems to be solved but the dynamics that enable growth toward wisdom.

Yet, this fear is not only a negative or paralyzing force. Recognizing existential fear, as such, points toward the possibility of self-recognition. Taking the attitude of the other -- becoming a ``you'' to the \{I\} -- by taking a second-person position on one's self can transform that fear into something productive. The aching desire to be understood is the positive and creative response to this underlying terror. The ferocity of this fear of nothingness fuels a striving for connection and community. It is from this foundational anxiety, in a Heideggerian sense, that two projects emerge in their separation and find themselves in unity: the quest for normative belonging (to be recognized as good) and the quest for authentic expression (to be recognized as \textit{in}finite).

These two projects are not separate pursuits but two facets of the fundamental drive for recognition arising directly from the I/me structure. The fear of nothingness drives the sometimes pathological pursuit of ``goodness'' as conformity---the attempt to secure a finite, stable ``me'' to ward off the terror of the \textit{in}finite, ungrounded \{I\}. If the self is not anchored in the shared norms of a community, it risks dissolution. Conversely, the pursuit of authenticity can manifest as a rejection of all norms, a defiance driven by the fear that any finite determination will annihilate the \textit{in}finite self. If the self is reduced to its normative roles, it risks losing its spontaneity and freedom.

The developmental challenge---the movement from \emph{Verstand} (Understanding) to \emph{Vernunft} (Reason)---is to recognize that these needs are not mutually exclusive but require each other for fulfillment. Understanding apprehends this tension as a contradiction that must be eliminated, demanding a choice between conformity and rebellion. Reason, however, apprehends this contradiction as the necessary structure of reality itself. It acknowledges that the self is inherently contradictory---simultaneously finite and \textit{in}finite. The transformation entails learning to hold this tension rather than resolving it prematurely, recognizing that true goodness requires authenticity, and true authenticity requires normative grounding.

\section{The Fabiola Project and the Absent Referent}

This philosophical distinction is visualized in the juxtaposition of two models of singular reference. On one hand, we can consider Robert Brandom's articulation of singular terms explored in the previous chapter. He describes their distinctive linguistic role as expressions `that refer to, denote, or designate particular objects' through substitution-inferential significance, forming equivalence classes of intersubstitutable descriptors. His exemplar, ``Benjamin Franklin,'' clarifies this: although multiple interchangeable descriptions -- such as ``the inventor of bifocals,'' ``the first postmaster general of the United States'' -- point symmetrically toward one historical figure, an inherent nullity emerges in their interstitial space. The subject, Franklin himself, is not contained fully by any single descriptor; rather, he exists as an intersection of equivalence classes as internalized by any who refer to him and whose collective reference is both precise and entirely empty of substantive singular essence. The stability of the ``me'' as an equivalence class or intersection of such classes is rather hollow. 

In contrast, \textit{The Fabiola Project} at the Menil Collection in Houston offers a counterpoint to Franklin's neatly symmetrical class of descriptors. The Project, curated around numerous divergent copies of Jean-Jacques Henner's original, now-lost painting of Saint Fabiola, illustrates an asymmetry at the heart of singular reference. Henner's painting was destroyed before any photographic record was possible, leaving behind only disparate, interpretative reproductions. In some, Fabiola appears in a red habit, in others green; her orientation shifts, her complexion changes radically from one image to another. Mediums vary -- paint, needlepoint, carvings, mosaics crafted from rice and beans flood the walls of the collection with difference. Each artist's interpretation claims authenticity without definitive verification. 

The Fabiola Project's reproductions, while they may share a common referent in the abstract, do not coalesce into a singular equivalence class. Instead, they proliferate into a multitude of interpretations, each claiming to represent Fabiola yet failing to converge on a definitive essence. The multiplicity of descriptors -- each an interpretation of an absent original -- creates a complex web of interrelations that defies the neat symmetry found in Brandom's exemplar. Each representation is both a claim to authenticity and a reminder of the original's absence, leading to an irreducible tension between presence and absence.

This asymmetry is crucial. Recall from Chapter 2 that singular terms are defined by symmetric substitution, forming equivalence classes where descriptors can be exchanged without loss of content. Predicates, however, often stand in asymmetric relations; ``is a dog'' entails ``is a mammal,'' but not the reverse. The Fabiola copies function more like predicates applied to an absent subject than singular terms referring to a present object. Each copy asserts a quality (``Fabiola in red,'' ``Fabiola in green'') but cannot be substituted symmetrically for another without altering the content.

This structure mirrors the relationship between the ``me'' and the \{I\}. The various descriptions of the ``me'' (teacher, parent, artist) are asymmetric interpretations of the \{I\}, which itself remains the absent original. The \{I\} cannot be captured by any single descriptor; it is the generative absence that animates the play of difference across the many representations of the ``me.''

This makes the project a striking exercise in Derridean non-presence. The original referent can never be made present; it exists only as a `trace' that animates the play of `diff\'erance' across the many copies. No single representation, no descriptor, can achieve full presence or exhaust the meaning of ``Fabiola.'' This is the predicament of the ``I.'' The \{I\} is the absent original, the point of subjective origin that can never be fully captured or made present in any objective description of the ``me.'' It is a nullity at the center of our self-descriptions. However, unlike the emptiness of the ``I think'' or the hollow intersection of equivalence classes, the absence of the original is a \textit{generative absence}. As Carspecken notes, this elusiveness is part of our direct experience: ``Try it out. Focus your attention on an object and establish if you can freeze out a moment in which you simultaneously have it before you and know that you do. It is impossible. `Presence' is unattainable in pure form within experience and yet it is presupposed or implicated in experience'' \parencite[52]{Carspecken1999}. 

The Fabiola Project visualizes this philosophical truth: the self is a web of representations pointing toward an origin that is never fully there, defined precisely by its resistance to definitive singular reference. This resistance is the source of our need to be recognized as \textit{in}finite. The original painting's essence is both affirmed and irrevocably negated by its diverse reiterations, creating an irreducible tension between presence and absence that mirrors the gap between the acting \{I\} and the recognized ``me.''

The absence at the core of the Fabiola Project---the lost original---is not a void. It is a productive absence, a generative nullity that animates the proliferation of copies. This structure directly maps onto the two existential needs. The collection of copies represents the finite ``me''---the various social roles and identities we inhabit. This corresponds to the need to be recognized as good---to have a stable, coherent identity validated by others. We present these finite versions of ourselves in the hope that they will be accepted as adequate representations.

Yet, the proliferation of diverse copies, their failure to converge on a singular essence, testifies to the inadequacy of any finite representation. The absent original represents the \textit{in}finite \{I\}. This corresponds to the need to be recognized as \textit{in}finite---to be acknowledged as more than the sum of our finite descriptions. The resistance of the original to definitive capture is the resistance of the \{I\} to being reduced to a static object. The emptiness at the center is not a lack but the enabling condition for the endless process of self-creation and recognition.

\section{The Beast of Love: A Turning Point}\label{sec:beast_of_love}

\subsection*{The Beast of Love}

\begin{verse}
When I see you, I see a light\\
Sharp and hot and glowing white\\
Your power, depth, and fierce roar\\
Broken in the glass of night\\
Bowed by rain, you start to fall\\
Fear of pain inside us all\\*[2ex]

As you learn, the things you fear\\
Will grow and stamp and roar and rear\\
Shadows looming on the wall\\
Become, with care, what you hold dear\\
Then Power - deep as dark abyss -\\
Embraces faint and shadowed mist\\*[2ex]

You will come to love the beast\\
That tears and cries and thrashes sheets\\
Because you are that beast, sweet child\\
That tears and cries and thrashes wild\\
Bend back the fog, back into rain\\
To wash away the shadowed stain\\*[2ex]

You are what Powers scary things\\
You are the wind beneath your wings\\
You cannot crush or fight the wind\\
But soon you'll learn to love the sting\\
Trust the beast to come in tune\\
Dark and bright, stars, sun, and moon\\*[2ex]

I see in you a Peace to come\\
Shining brighter than the sun\\
Each time you break and let it go\\
Your body and your spirit grow\\
Surely in the curl of time\\
You'll come to be and know your mind\\*[2ex]

I'm sorry, child, I bent your wing\\
I am a big and scary thing\\
Still frightened that the beast within\\
Might tear your fragile, gentle skin\\
And leave you scarred before you know\\
How deep and wide your love will grow
\end{verse}

I wrote this poem for my stepdaughter when she was about 4 years old after a moment of painful misrecognition -- a timeout that wounded our relationship for months. I forget why we had conflict, but I remember that I was trying to be a `good' authority figure. There was some boundary that I thought was legitimate, but that boundary conflicted with the \textit{in}finite spontaneity of the \{I\}. In confessing my fear, I point to the scars I bear from before I knew how much deeper and wider I would come to love. That scratching beast is still within, not-yet in tune, but still I find my capacity for genuine reciprocal intersubjective recognition [love] has grown through projects such as this one. 

\subsection*{Verstand to Vernunft: Beyond Oppositional Consciousness}

The line ``You will come to love the beast'' captures this developmental understanding. The ``beast'' -- the terrifying tension between our existential needs -- initially appears as pure threat. In early development, we experience it as overwhelming anxiety: the child desperately seeks approval and fears abandonment. In adolescence, it becomes defiant opposition: the teenager would rather be rejected for who they authentically are than accepted for who they're not. 

But mature development contains a different possibility. The tension is a source of humanity. The ``beast'' that ``tears and cries and thrashes'' is not an enemy but a creative power. Learning to ``trust the beast to come in tune'' means learning to work with rather than against the productive tensions that bind people together. But ``working with'' is conditioned on the possibility of non-action. Learning to listen -- surrendering to the movement of the negative which is never grasped in its fullness -- is the first step in befriending the beast. 

This transformation requires what the poem calls ``breaking and letting go.'' Each time I encounter the collision between needs -- each time I face the apparent contradiction between being good and being authentic -- I have the opportunity to ``break'' out of oppositional thinking and ``let go'' of the fantasy that these needs must be mutually exclusive.

The poem's final stanza introduces three crucial concepts that are essential to the resolution of existential tension: confession, forgiveness, and trust. These are not merely psychological categories but ontological structures that make possible the movement from opposition to integration.

``I'm sorry, child, I bent your wing'' is a confession -- an acknowledgment of the harm that inevitably occurs when finite beings with competing needs interact. I confess to being ``a big and scary thing,'' acknowledging the power differential that makes recognition so difficult. The confession opens the space for forgiveness -- not as an erasure of harm but as the willingness to continue the relationship despite the wound. Forgiveness recognizes that the ``beast within'' both the authority figure and the child ``might tear fragile, gentle skin,'' but it also trusts that love will grow deeper and wider through the scars.

Trust emerges as the fundamental structure that makes ongoing relationship possible. It is trust in the developmental process itself -- the faith that ``in the curl of time'' both parties will ``come to be and know their mind.'' This trust is not based on certainty but on surrender. 

These concepts -- confession, forgiveness, and trust -- are what Habermas might call the ``lifeworld'' foundations that make rational discourse possible. Without the willingness to confess failures, forgive others' limitations, and trust in the possibility of mutual understanding, communicative action breaks down into strategic manipulation or ideological domination.

Confession disrupts this dynamic by acknowledging the inadequacy of both positions. When I confess, ``I bent your wing,'' I acknowledge that my attempt to enforce the ``good'' (the boundary I thought legitimate) violated the ``\textit{in}finite'' (the child's spontaneity) and caused harm. This acknowledgment relinquishes the claim to absolute authority. Forgiveness, in turn, relinquishes the demand for perfect recognition. It accepts the finitude of the other. Trust is the commitment to continue the relationship within this space of mutual vulnerability.

This structure of reciprocal relinquishment is what constitutes \emph{Vernunft}. It allows the two needs to coexist not as a compromise, but as a unified developmental process where the articulation of the ``good'' is constantly revised by the expression of the ``\textit{in}finite,'' and the expression of the ``\textit{in}finite'' is given determinate form through the articulation of the ``good.'' The transformation enacted in the poem---from fear of the ``beast'' to trust in its eventual harmony---is precisely this shift. It is the realization that the resolution of existential tension is not a theoretical solution but a practical achievement grounded in the willingness to surrender oppositional consciousness and embrace the dialectical unity of our needs.

What mature development entails is that the two existential needs are not ultimately opposed but dialectically unified. The need to be recognized as \textit{good} and the need to be recognized as \textit{in}finite require each other for their fulfillment.

\section*{Toward the Particular You}

This chapter has diagnosed a failure---how abstract law becomes violence when it treats living practice as object rather than subject---and pointed toward reciprocal recognition as a path beyond this alienation. The existential needs---to be recognized as good and as \textit{in}finite---appear contradictory only when consciousness treats them as fixed opposites. When consciousness recognizes that they require each other, developmental movement becomes possible.

But the account so far has operated at a high level of philosophical abstraction. We have discussed the generalized other, struggle for recognition, confession and forgiveness as ontological structures. The next chapter turns to the concrete, face-to-face encounter. Who are \textit{you}---the particular person I address when I speak? How does the universal ``you'' of claimed subject positions (CUSP) relate to the transactional ``you'' of genuine dialogue? How does the struggle for recognition actually unfold in the messy, improvised practice of parenting young children?

The movement across Chapters 1-5---from sense-certainty through perception to the social struggles of recognition to the encounter with absolute negativity---loosely follows the arc of Hegel's \textit{Phenomenology of Spirit}. I do not claim these chapters constitute Hegelian scholarship or allegorical retelling. Rather, I have found that iterating on Hegelian themes helps make sense of the developmental journey consciousness undergoes. The philosophical work of Chapters 3-5 prepares us for the mathematical work of Chapters 6-10, where similar patterns---negation, recognition, absence as enabling condition---will appear in the domain of pure thought.

\printbibliography[heading=subbibliography]

\chapter{Thoughts for Two -- Who Are You?}

\begin{verse}
Who are you today, love?\\
How are your stars arranged?\\
What plans have you forgotten?\\
What's still written on your hands?\\ (Chorus from \textit{Silver and Iron} by Savich/Wanchic)
\end{verse}

\begin{abstract}
This chapter examines the intersubjective structures of recognition, exploring how the question ``Who are you?'' reveals the dynamics of interpersonal recognition in both mathematical communication and human development. The central case study follows the evolution of ``Big Gorilla,'' a character created to navigate family dynamics, illustrating communicative breakdown and repair. Through this example, three theoretical frameworks are applied: George Herbert Mead's symbolic interactionism, illustrating how gestures become significant symbols; Jürgen Habermas's theory of communicative action, which explains how breakdowns in shared understanding necessitate rational discourse; and Robert Brandom's analytic pragmatism, providing formal tools to analyze how meaning emerges from use. A key tension is explored between the \emph{transactional} ``you,'' which refers to a specific interlocutor, and the ``CUSP'' you (Claimed Universal Subject Position), which invokes a generalized, normative subject. This distinction highlights the ethical dimensions of dialogue, contrasting genuine recognition with universalizing subsumption. 
\end{abstract}

\section{Introduction: From Existential Need to Communicative Practice}\label{introduction-from-existential-need-to-communicative-practice}

Chapter 3 argued that the tension between existential needs---the need to be recognized as ``good'' within normative frameworks and the need to be recognized as \emph{in}finite, expressing authentic spontaneity---is resolved not in solitude but through a developmental journey toward reciprocal recognition. This resolution, embodied in the surrender to what I called the ``Beast of Love,'' is fundamentally a social and communicative achievement. But what is the structure of this recognition? Who is the ``You'' that meets the \{I\}?

This chapter investigates that question through a concrete case study: the emergence of ``Big Gorilla,'' a character I invented while learning to navigate family life with my partner \(\mathcal{A}\) and her twin daughters, \(\mathcal{M}\) and \(\exists\), who were $3\frac{1}{2}$ years old at the time. What began as a playful proto-parenting strategy became a site for explicating the deep assumptions involved in communicative action. The story traces how gestures evolve into significant symbols, how communicative breakdowns necessitate rational discourse, and how formal analysis can illuminate but never fully capture the spontaneous creativity of genuine recognition.

The turn toward communicative practice is not an abandonment of the existential struggle detailed previously, but rather the necessary next step in addressing it. The tension between the need to be recognized as `good'---to find validation within shared normative frameworks---and the need to be recognized as `\textit{in}finite'---to express authentic spontaneity beyond any fixed category---cannot be resolved through solitary introspection. That tension is fundamentally social; it lives in the space between the `I' and the `You'.

If the self is structured by the dialectic of the spontaneous \{I\} and the social `me', then understanding how recognition functions requires understanding the medium through which that dialectic unfolds: communication. We engage Mead, Habermas, and Brandom not as abstract theorists, but as guides who offer resources for navigating this tension. They help us understand how the \textit{in}finite \{I\} expresses itself within the finite structures of language and social norms, and how those norms themselves are instituted and transformed through our interactions.

The theoretical architecture unfolds through three frameworks, each offering progressively more formal resources. George Herbert Mead's pragmatic interactionism shows how selves emerge through taking the attitude of others toward one's own gestures. Jürgen Habermas's theory of communicative action reveals the triadic structure of validity claims (truth, rightness, sincerity) and explains why breakdowns in shared understanding require rational discourse. Robert Brandom's analytic pragmatism provides tools for making explicit the implicit rules governing practices and connects this work to mathematical meaning-making.

Before developing these theoretical resources, I establish a key distinction that emerged from my earlier research on how preservice teachers use the second-person pronoun ``you'' when teaching mathematics.

\subsection*{The CUSP You and the Transactional You}\label{the-cusp-you-and-the-transactional-you}

When working on my first graduate research project exploring how gestures and speech relate in teaching algebraic concepts, I noticed that preservice teachers often used ``you'' in ways that seemed to miss me as the actual addressee. For example, when factoring an equation like \(2(p^2 + 5p + 6)\), a teacher would say ``The first thing you do is distribute the 2.'' But I wouldn't distribute the 2---I'd recognize it would just need to be factored out again later. Their ``you'' appeared to invoke a generalized procedure-follower rather than addressing me as a particular person with my own mathematical understanding. What fascinated me was recognizing myself in that practice. I do that all the time! It is so ubiquitous as to be both invisible and grossly insistent.

This pervasive address to a generalized other connects to the existential needs discussed previously. When speech defaults to the CUSP ``you,'' it risks enacting the dynamic of alienation we encountered in Chapter 3's conflict between abstract law and living practice. The CUSP you works like the abstract Understanding (\emph{Verstand}) that subsumes particulars under universals---it recognizes the addressee only insofar as they fit the generalized category, negating their \emph{in}finite spontaneity, their capacity to surprise, to deviate, to bring something new.

The CUSP you operates from what Habermas calls the \emph{system}---abstract procedures, standardized expectations, universalized norms---rather than the \emph{lifeworld}, where particular persons meet face-to-face. It assumes a subject position anyone occupying a certain role ought to take, rather than recognizing the specific person who stands before me. We rely on shared norms to coordinate action, yet when this mode becomes dominant, it violates the need for recognition as \emph{in}finite that Chapter 3 identified as co-equal with recognition as ``good.''

I called such uses the \textbf{CUSP you}---the \emph{Claimed Universal Subject Position}---because it refers to a generalized other rather than the particular other being addressed. The acronym is purposeful: CUSP suggests the kind of catastrophes that the Zeeman Catastrophe Machine produces. Figure \ref{fig:CUSP_You} models the \textit{cusp} catastrophe of a ZCM as a kind of body-feeling as an `ah-ha' moment is experienced. The concept of \textit{hysterisis} concerns how the system's history affects its present state, offering an opportunity to connect history to mathematical languages. 

\begin{figure}
    \centering
    \includegraphics[width=\textwidth]{/Users/tio/Documents/GitHub/September_UMEDCA/images/Ch04_Zeeman_Inference_Field.pdf}
    \caption{\textit{Note.} A cusp catastrophe in the Zeeman Catastrophe Machine.}
    \label{fig:CUSP_You}
\end{figure}

The CUSP you relates to the notion of the horizon of forestructures identified in Chapter 1's Exercise. I \emph{claim} but do not attain universal subject position. The CUSP you invokes a normative expectation of how one ought to act in a given situation without necessarily addressing the specific person being spoken to.

This contrasts with what Sebastian Rödl calls the \textbf{transactional} use of ``you'' \parencite{Rodl:2014aa}. Rödl distinguishes between monadic and dyadic forms of speech actions. Monadic predicates have one subject: \(Fa\), like ``Peter is painting the wall'' \parencite[307]{Rodl:2014aa}, where \(a\) is ``Peter'' and \(F\) is the predicate ``\(\square\) is painting the wall.'' Dyadic predication is formally \(aRb\), where \(a\) is ``Peter,'' \(b\) is ``Paul'' and \(R\) encodes the relation ``\(\square\) is handing the brush to \(\square\)'' in a sentence like ``Peter is handing the brush to Paul'' \parencite[307]{Rodl:2014aa}.

Transactions have both \emph{actor} and \emph{patient} roles. Rödl's paradigmatic example is gift-giving: the actor gives a gift to the patient who receives it. Crucially, removing the patient's contribution means the transaction has not occurred. The gift-giving side fails if the gift-receiving side is absent. While I'm not entirely convinced by this argument---empirical acts of gift-giving sometimes seem not to require receipt---it stands as a powerful metaphor for communicative action, where the other interlocutor is necessary. Communicative action is essentially dyadic, not monadic.

The transactional ``you'' refers anaphorically to the specific person addressed in a face-to-face encounter---the irreplaceable other standing before me. The CUSP you, by contrast, appeals to universalized norms that anyone occupying a certain position should follow.

Both uses are necessary. We cannot communicate without invoking shared norms (the generalized other), but we also cannot genuinely recognize others if we treat them only as instantiations of universal categories rather than as particular subjects with their own perspectives. The ethical task is learning to distinguish when we're addressing the particular other and when we're invoking universal norms, and ensuring that our speech creates space for the other's response rather than foreclosing it.

This distinction maps onto the dialectical structure introduced earlier. The transactional ``you'' addresses the \emph{in}finite---the spontaneous, unrepeatable \{I\} of the other. The CUSP you invokes the finite---the socially constituted ``me,'' the determinate position within a normative framework. Genuine recognition requires both. Addressing only the \emph{in}finite \{I\} while refusing to acknowledge the other's position within shared normative structures denies them social intelligibility---treats them as pure spontaneity without history. Addressing only the finite ``me'' position while ignoring their capacity for spontaneous deviation reduces them to a role, a predictable instance of a universal.

The movement between these modes calls for what Hegel would call \emph{sublation}---a unity that preserves both moments while transcending their opposition. In teaching mathematics, this means recognizing students as particular persons with unique histories \emph{and} as participants in shared practices governed by norms that transcend individual preference. The teacher who says ``the first thing you do is distribute the 2'' may appropriately invoke the CUSP you when articulating a general procedure. Yet if that universal address forecloses the student responding ``I would do it differently,'' it collapses into alienation that violates the student's need for recognition as \emph{in}finite.

When I recognized that I also used ``you'' in similar CUSP ways when teaching, I became curious about the inner speech that accompanied such usage. I had privileged access to my own rehearsals of speech acts, including a very judgmental voice that would say things like ``You're stupid! You suck!'' This hectoring voice used ``you'' to address me, but who was this ``you''? And who was the Other speaking so harshly?

Phil Carspecken's observation that everyone engages in forms of self-monitoring activity helped me recognize this wasn't a sign of pathology but a normal structure of self-consciousness. The Meadian distinction between the \{I\} and the ``me'' encompasses both the ethical and algorithmic aspects of this self-monitoring. Part of my developmental journey involved disentangling contradictions in the internalized generalized other while not losing the moral compass that points toward more universalized forms of intersubjective recognition.

This inner hectoring voice raises questions about the CUSP/transactional distinction when applied to self-address. When I say ``You're stupid,'' am I addressing the particular \{I\} that just made a mistake, or invoking a CUSP you---a generalized ``one who made this type of error''? Self-consciousness operates at the boundary between these modes. The generalized other whose voice I internalize speaks in CUSP terms (``anyone who did what you did is stupid''), yet the sting is felt by the particular \{I\} who cannot evade being addressed.

This structure helps explain why self-criticism can be so painful yet difficult to dismiss. The CUSP voice invokes universal norms (``this is wrong, anyone would agree''), lending it authority that feels objective, while addressing the irreplaceable \emph{me} who must bear the judgment. When the internalized generalized other contains contradictions---praising behaviors in one context while condemning them in another---the result is existential turmoil. The \{I\} receives contradictory demands about which ``me'' it should become.

The developmental task calls for engaging the generalized other in rational discourse---making explicit the implicit norms it invokes, challenging contradictions, distinguishing legitimate moral demands from internalized oppression. This is the movement from Understanding (\emph{Verstand})---which subsumes particular acts under universal rules---to Reason (\emph{Vernunft}), which can reflect on those rules and revise them through dialectical critique.

\subsection*{Big Gorilla}

I want to share a story that began after I defended my dissertation. Finding myself with unexpected freedom -- time to think and no immediate academic pressures -- I stayed in Bloomington, Indiana, for an extended period. I had developed a life there: music, friends, and familiar routines that supported my thinking. When I began dating $\mathcal{A}$, I entered into a rich context for exploring development and intersubjectivity, particularly through my interactions with her twin daughters, $\mathcal{M}$ and $\exists$, who were $3 \frac{1}{2}$ years old at the time.

Initially, I approached my role as a kind of playful participant-observer. But when we all caught COVID and ended up quarantining together, I realized that the dynamic would need to change if we were going to share living space. The approach I had been using -- engaging fully until exhaustion set in, then retreating to recover -- wasn't sustainable when I was sick and couldn't leave. The communicative patterns I had developed during my dissertation, which made little sense to academic reviewers or to job-search committees, did not make better sense to 3-year-olds. I would speak in elliptical parentheticals -- around and around -- and sometimes they would smile like ``huh?!,'' sometimes they would try to teach me how to speak more competently, and sometimes communication broke down in spectacular fits. I often could not assume they understood me, and they, in turn, discovered that I did not know how to speak to them. Rather than an impasse wherein the assumption of competency transformed into the assumption of incompetency -- `a quietism of despair' -- such moments led me to invent the character of Big Gorilla, who eventually became entitled the name Daddy Gorilla. Big Gorilla had a couple of baby gorillas who were always getting in trouble. He didn't know how they got there or why they were hanging around him so much. He would walk around with his fists on the ground and say, ``Big Gorilla tired, but Baby Gorilla still awake. Baby Gorilla sleep now, then Big Gorilla sleep.''

Part of Big Gorilla's behavioral regimen was to thump his chest when he was angry about something. ``Baby Gorilla hit other Baby Gorilla. Hooohoohoo thumpthumpthump.'' I tried to keep anger in a space of play, not fully identifying with it. The kids were very small and easily frightened. I didn't know how to hold them responsible for their behavior, but I felt like it was important to intervene sometimes. At night they would often get ornery about going to bed. Big Gorilla would say ``hoohoohoo'' and thump his chest, and the Baby Gorillas would smile and thump on their bellies and go ``hoohoohoo.'' It was very cute for a while. But one time, the anger got pretty close to me, and I thumped my chest until it hurt and hoohoohoo'ed in all capital letters. Tears bloomed and $\mathcal{A}$---who had been watching the interaction unfold---looked at me in a way that said ``back off,'' and so I said, ``Oh, that was too loud. Big Gorilla sorry. Big Gorilla get frustrated when Baby Gorillas don't listen.'' 

What does it mean to assume an other who is communicatively competent, when both oneself and the other are thrown into the developmental processes by which such competencies are grown? This points to a circularity in Habermas's theory of communicative action. Rational discourse presupposes participants who can raise and redeem validity claims, who can recognize when sincerity, truth, or rightness is at stake, who can say ``no'' to problematic assertions. Yet these capacities must be developed through the very communicative practices they make possible. How can communication presuppose competencies that only communication can produce?

The assumption of competence functions as a \emph{regulative ideal} that orients developmental practice. When I engaged with the twins as Big Gorilla, I wasn't pretending they had fully developed capacities for rational discourse. I was treating them \emph{as if} they could understand, respond, and participate in norm-governed practices---and through that treatment, those capacities emerged. The assumption of competence is what Hegel would call a \emph{productive negativity}---it points toward something not yet actual, and through that pointing, helps bring the not-yet-actual into being.

If we wait for students to demonstrate competence before treating them as competent, we withhold the conditions that make competence possible. Yet if we ignore their actual developmental position and address them as if they already possess abilities they lack, our speech fails to meet them where they are. The pedagogical task is navigating this tension: assuming competence as a regulative ideal while remaining attentive to the particular, developing capacities of the actual persons before us.

This story serves as the central case study for understanding how selves and shared meanings emerge through communicative practice. The chest thumping plays the role of a conversation of gestures that eventually became a significant symbol, representing the complex interplay of emotions, intentions, and social norms in these interactions. What began as a proto-parenting strategy became a field for explicating the deep assumptions involved when engaging in communicative action on multiple \textit{discursive levels}. Two people may be able to do the same sorts of things with words while one may have more expressive resources to describe how to do those things. Pokey, for example, seems to be able to recognize some forms of suffering. In such moments, she demonstrates her recognition in a rich behavioral vocabulary. She cannot, in my limited ability to assess such things, explain how to use another vocabulary to explain how to appropriately use nuzzles and licks to demonstrate recognition. She lacks a \textit{pragmatic metavocabulary} for deploying her behavioral vocabulary as a system of rule-bound practices. I do have those expressive resources, and I suppose I could nuzzle and lick her when she is sad, so I operate at a different discursive level. Lacking such resources does not define her as `less than' in the dimensions that truly matter, as the recognitive act itself is necessarily implicit; we just orient our system of affection towards the deeply implicit. The concept of discursive levels offers a workaround for the ethical problems of the last chapter where Leibniz's version of apperception foreclosed on animal consciousness achieving apperceptive awareness. This isn't a book about dogs -- I discuss Pokey in application of one of the burdensome problems of the teacher-student relationship where social power can distort the ability to rationally communicate with one another.  

I have signaled my own communicative \textit{in}competency throughout. Early forms of ethnography and social science were burdened by the assumption of a linear, stage-like developmental progression where `primitives' develop into `sophisticates.' I am more interested in sublating such structures to return to earlier moments on such trajectories---unlearning without forgetting what was learned to become a better listener and developmental partner. At no moment should readers infer that I am offering advice for how to be a better teacher or parent. Nor should readers infer that the children were my research subjects in this project. They are still too young to fully articulate a rational ``no'' to my descriptions here. They cannot member-check this work, as they are still just getting into reading chapter books. I breach norms related to ethnographic research. The beauty I find in the final turn of this chapter may justify this break, or it may not.

I know it can be awkward to read confessions that foreground ethical uncertainty. Consider why such articulations point to an internally contradictory generalized other whom I have internalized. That other judges my attempt to play with my kids for how it is burdened with various ethical norms related to research, parenting, and teaching. Is it universally wrong to write about my children? Is it wrong to speak in ways that could be identified as problematic? Some would say `yes!' Does that knowledge make the wrongnesses more wrong or more right? I am fundamentally \textit{uncertain} about such questions. The ethical implications of the character of Big Gorilla only emerged upon reflection. This suggests an important structural feature of ethical norms: they are not cognizable in the moment of action, but \textit{re}cognizable when taking a second-person position on the act. Since I cannot coherently judge from the Universal-Reader's perspective except to say that every possible utterance is possibly unethical, I am bound to find my own words unethical at some point. To finish the book, I must find some other source of forgiveness: ``Out beyond ideas of wrongdoing and rightdoing there is a field. I'll meet you there\ldots'' \parencite[123]{rumi_2009}.

\section{The Social Self: Mead's Pragmatic Intersubjectivity}\label{the-social-self-meads-pragmatic-intersubjectivity}

George Herbert Mead's foundational insight is that the self is not a pre-given entity but an achievement that unfolds through social interaction. This focuses on the \emph{ontogenesis} of the self---development within an individual---while recognizing its entwinement with \emph{phylogenesis}, the historical development of normative practices within communities.

\subsection*{From Gesture to Significant Symbol}\label{from-gesture-to-significant-symbol}

Mead begins his analysis with what he calls the \emph{conversation of gestures}---the kind of interaction we might encounter between two dogs preparing to fight \parencite[\S 7.3]{mead_mind_2015}. Each dog's act serves as a stimulus provoking a response from the other, which in turn modifies the first dog's behavior. One dog's readiness to attack becomes a stimulus causing the other to change position or attitude, which then causes the first to change its attitude. This creates a feedback loop, but it is not yet symbolic communication because neither dog has what Mead calls ``a reflective determination to attack.'' The dogs are responding to each other's gestures without taking those gestures as \emph{meaning} anything.

The conversation of gestures operates at the \emph{pre-symbolic} level---a domain of embodied interaction that precedes and grounds linguistic communication. This level is \emph{foundational}. Even sophisticated human communication retains gestural dimensions that cannot be fully translated into words. When \(\mathcal{A}\) looked at me with disapproval after my too-loud chest-thump, her look \emph{was} the communication. The meaning was immediately present in the gesture itself.

This pre-symbolic dimension connects to the emphasis on \emph{practice} in Brandom's meaning-use analysis and to the \emph{embodied} nature of reason discussed in Chapter 1. Before we can make norms explicit through language, we must enact them implicitly through behavior. The conversation of gestures is the site where normative structure first emerges---as patterns of appropriate and inappropriate response that participants recognize through their embodied engagement.

The Big Gorilla practice began at this gestural level. My initial chest-thumping was a behavioral response that emerged from the immediate situation, from the need to mark my frustration while keeping it within bounds the children could tolerate. Only later, through repeated iterations and eventual breakdown, did the practice become something I could articulate and reflect upon.

The transformation from gesture to \textbf{significant symbol} occurs when the gesture-maker becomes conscious of the response their gesture evokes in the other. When the gesture can be experienced from both sides of the interaction simultaneously---when it means the same thing to the one making it as to the one receiving it---the interaction becomes an intersubjectively constituted practice.

Figure \ref{fig:dyadic_model_full_escalation} models the interaction between two dogs, where each dog's unknown interior state is represented by a node. The arrows represent the gestures exchanged between the dogs, with the labels indicating the specific actions taken. The model illustrates how the interaction escalates from neutral states to more intense states like Tense, Hostile, and Submissive.

\begin{figure}[h!]
    \includegraphics[width=0.8\textwidth]{/Users/tio/Documents/GitHub/September_UMEDCA/images/ch4_dyadic_model_full_escalation.pdf}
    \caption{\textit{Note.} A dyadic model of the escalation and relaxation of two dogs in a fight.}
    \label{fig:dyadic_model_full_escalation}
\end{figure}

The model in Figure \ref{fig:dyadic_model_full_escalation} is a formalization of Mead's conversation of gestures. Treating a system of dog gestures like ``Lunge'' as a formal language is, perhaps, a little odd. But such treatment may clarify an aspect of what it means to say that communicative action is dyadic---a shared system of states and gestures, where outputs become inputs for the other.

\subsubsection*{Formalization of Mead's Dog Fight}

Each dog can be represented as a type of finite state transducing automaton that generates an output based on its current state and an input. A single participant (e.g., one of Mead's dogs) can be defined as a machine \(M\):
\[ M = (Q, \Sigma, \Gamma, \delta, q_0) \]
where:
\begin{itemize}
    \item \(Q\) is a finite set of states, representing the dog's internal \textit{attitude} (e.g., Neutral, Tense, Hostile, Submissive).
    \item \(\Sigma\) is the input alphabet, a finite set of gestures it can perceive from the other dog (e.g., \{growl, stare, lunge, bite, cower\}).
    \item \(\Gamma\) is the output alphabet, a finite set of gestures it can produce. In this context, \(\Sigma\) and \(\Gamma\) are the same, as the dog responds with gestures that are perceived by the other dog.
    \item \(\delta\) is the transition function: \(\delta : Q \times \Sigma \to Q \times \Gamma\). It takes the current state and a perceived gesture (input) and determines a new state and a responsive gesture (output).
    \item \(q_0 \in Q\) is the initial state or attitude.
\end{itemize}

Once those cycles have been compressed to the extent that a growl or aggressive stance results in the other dog backing off, there is no need to bite. The growl has the significance of the bite as a proto-inferential consequence. When symbols are used, cycles of interaction are temporally compressed as a \textit{history} that is assumed to have the same significance as the longer sequence. This is what Mead calls a \textit{significant symbol}. In the mechanical model, the history of each dog's action is represented by the state that each dog is in.

\subsubsection*{Beyond Sequential States: The Double-Zeeman Catastrophe Machine}

While the FSM model in Figure \ref{fig:dyadic_model_full_escalation} captures the \textit{sequence} of gestures, it misses the underlying \textit{dynamics}---the continuous build-up of tension, the non-linearity, and the potential for sudden rupture (the ``terror'') inherent in the interaction. The FSM treats states as discrete and transitions as instantaneous, but the lived experience of the dog fight---or the Big Gorilla breakdown---involves gradual accumulation followed by catastrophic release.

To capture this dynamic, we turn to the Zeeman Catastrophe Machine (ZCM), introduced in Chapter 1 as a model of embodied reason's rhythms. The ZCM makes concrete how the gradual accumulation of tension (temporal compression) leads to a sudden release (catastrophe/decompression) when an equilibrium breaks down. Consider constructing a \textbf{Double-Zeeman Catastrophe Machine}---a dyadic engine where the loose ends of the elastic bands from two ZCMs are joined together at a single, floating point.

\begin{figure}[h!]
    \centering
    \includegraphics[width=0.8\textwidth]{/Users/tio/Documents/GitHub/September_UMEDCA/images/Ch4_Double_Zeeman_Dogs.pdf}
    \caption{\textit{Note.} A Double-Zeeman Catastrophe Machine for modeling dyadic catastrophes.}
    \label{fig:Double_Zeeman}
\end{figure}

In a standard ZCM, a disk seeks to minimize the energy of two elastic bands: one connected to a fixed anchor, the other to a movable control point. In the Double ZCM:

\begin{itemize}
\item \textbf{Two Disks (A and B)}: Represent the behavioral state or ``attitude'' of Dog A and Dog B (or Big Gorilla and Baby Gorilla).
\item \textbf{Fixed Anchors (A' and B')}: Each disk retains one elastic connected to a fixed anchor, representing the stable context or underlying disposition of each participant.
\item \textbf{The Intersubjective Link (Point C)}: The ``loose ends'' (the control bands) from Disk A and Disk B are joined at a single, \textbf{floating} point, C.
\end{itemize}

The crucial dynamic is that the system minimizes the \textit{total} potential energy of all four bands. Point C is not controlled externally; its position emerges from the interaction between the two disks. This configuration models the intersubjective dynamics: the state of the interaction is holistic, with the stability of Dog A entirely dependent on the posture of Dog B, and vice versa. They cannot find equilibrium in isolation. This physically embodies Rödl's concept of \textit{intentional transaction}---the system is a unity with two sides.

The tension in the shared elastics is the physical embodiment of the escalating conflict. This is the ``temporal compression'' from Chapter 1---the history and tension of the interaction are stored as potential energy in the system. A small movement in Disk A (a subtle gesture) immediately alters the forces on the shared Point C, which changes the energy landscape for Disk B. This is a dynamic realization of Mead's feedback loop, where gestures are not mere signals but \emph{forces} that reshape the possibilities for the other participant.

As tension builds, the system may reach a bifurcation point where the local energy minimum vanishes. If Dog A snaps into attack (Disk A rotates violently), this sudden release (decompression) causes a massive, instantaneous shift in the shared Point C. This sudden yank drastically alters the energy landscape for Dog B, likely pulling it across its own threshold and inducing a reciprocal catastrophe (flight or counter-attack). Either dog can lead the other to a catastrophe.

This model synthesizes Mead's interactionism with the dynamical modeling introduced in Chapter 1. The FSM shows us \textit{what} happens (the sequence of states), while the Double-ZCM shows us \textit{how} it happens (the continuous dynamics of tension and release). The chest-thumping escalation with the twins was not a sequence of discrete states but a gradual build-up of intensity that suddenly crossed a threshold---my ``hoohoohoo in all capital letters'' was the catastrophic release that yanked the shared intersubjective field (Point C) into a new configuration, inducing tears rather than laughter.

In the Big Gorilla story, chest-thumping began as a proto-gesture. I would thump my chest and say ``hoohoohoo'' when frustrated with the kids' bedtime resistance. Initially, this was closer to the dogs' interaction---a behavioral response to a stimulus (their ornery behavior). But when the twins began to imitate the gesture, thumping their own bellies and saying ``hoohoohoo'' in return, something crucial occurred.

I experienced my own gesture from their perspective. I recognized my chest-thumping not just as an expression of my frustration but as something that meant \emph{playful anger} to them---anger that invited participation rather than threatening them. The chest-thump became a \textbf{significant symbol} because it achieved the same significance for both maker and receiver.

``Taking the attitude of the other'' involves a shift to a second-person position where I experience \emph{myself} as a ``You'' through the other's recognition. This is the movement that constitutes self-consciousness in Mead's framework.

When the twins thumped their bellies in response to my chest-thumping, they were taking my attitude toward their behavior (frustration at bedtime resistance) and enacting it themselves. They demonstrated that they recognized what my gesture meant---that they could see themselves from my perspective as ``ornery Baby Gorillas.'' Simultaneously, I took their attitude toward my gesture---recognizing that my chest-thumping appeared to them as playful, as an invitation to participate.

This reciprocal perspective-taking creates an \emph{intersubjective loop}. I monitor my own behavior through the anticipated response of the other. The other monitors their behavior through the anticipated response of me. Neither perspective has priority---the meaning emerges in the circulation between us. The self is social: the capacity to take oneself as an object (self-consciousness) emerges only through interaction with others who can serve as mirrors reflecting ourselves back to us.

The significance of a symbol is not its importance but its structure of reciprocal perspective-taking. When I thumped my chest and the twins responded with delight, I recognized my gesture through their recognition of it. The meaning wasn't in my head or in their heads but in the shared practice that emerged between us.

\subsection*{Play Stage and Game Stage}\label{play-stage-and-game-stage}

Mead identifies two crucial developmental stages. In the \textbf{play stage}, children learn to take the role of specific others---they play house, school, or, in our case, gorilla family. When \(\mathcal{M}\) and \(\exists\) thumped their bellies in initial imitation and said ``hoohoohoo,'' they were taking the role of Big Gorilla, learning to recognize themselves from his perspective. When they formed their own identities as Baby Gorillas within the game, they were taking a second-person position on their imitation, differentiating themselves from my character to become their own. Such reflections---taking a second-person position on oneself, becoming a ``You'' to one's \{I\}---mark the beginning of Meadian self-consciousness.

The \textbf{game stage} represents more complex development. As discussed in Chapter 3, the child must take the role not just of one other but of all others simultaneously. They must understand their position in relation to a whole organized system of roles. To play baseball effectively, the shortstop must internalize not just their own position but the roles of pitcher, catcher, and teammates---understanding what everyone is doing and likely to do. The child internalizes the coordinated activity of the whole team.

While the twins were primarily in the play stage during our early interactions, glimpses of the game stage emerged when they began to understand that their behavior affected Big Gorilla, Momma Gorilla, and the other Baby Gorilla. Once the rules of the game are recognized as a system of reciprocal roles, the game stage is in full force.

\subsection*{The Generalized Other}\label{the-generalized-other}

The culmination of this development is the internalization of what Mead calls the \textbf{generalized other}---the organized attitude of the whole community or social group. This is not ``social norms'' as abstract rules but the coherent and organized set of normative expectations that regulate conduct from within \parencite[\S 22]{mead_mind_2015}.

When \(\mathcal{A}\) looked at me in a way that said ``back off'' after my too-loud chest-thump, I was encountering the generalized other of our emerging family. Her look conveyed not just her individual disapproval but a normative expectation of our household that I shared: anger, even playful anger, must remain within bounds that feel safe for everyone. Her look would not have meant anything to me if I had not already internalized some aspect of the generalized other that said ``adults should not terrify children.''

The generalized other enables self-regulation. It is the internal voice---sometimes supportive, sometimes critical---that judges our actions from the perspective of the community. When I rehearsed harsh self-criticisms (``You're stupid!''), I was enacting the role of the generalized other, monitoring myself from an internalized normative position. The contradictions I experienced in that voice (am I intrinsically wrong, or have I done something wrong?) reflected contradictions in the generalized others I had internalized from different communities and different developmental stages.

Without internalizing the attitudes of the community, we would have no basis for self-reflection, no capacity to evaluate our own actions, no conscience. The generalized other provides the normative framework that makes ethical life possible---the source of the ``ought'' that guides action beyond immediate impulse or strategic calculation.

Yet the generalized other can become a site of oppression when the norms it embodies are unjust, contradictory, or hostile to aspects of our being that deserve recognition. When marginalized communities internalize dominant norms that devalue their ways of life, when children internalize parental voices that are excessively critical, when teachers internalize institutional demands that conflict with pedagogical values---the generalized other becomes an \emph{alienated conscience}. It speaks with authority, yet that authority does not promote genuine flourishing or reciprocal recognition.

The developmental and political task calls for what Mead would call \emph{social reconstruction}---the critical examination and revision of the generalized other itself. This requires operating at a higher discursive level where we can make explicit the norms that were previously implicit, subject them to rational critique, and reconstruct them in ways that better serve the needs of all participants. Such reconstruction requires communities of discourse where participants can challenge each other's internalized norms and work toward more adequate shared frameworks.

\subsection*{The \{I\} and the ``Me''}\label{the-i-and-the-me}

Mead's distinction between the \{I\} and the ``me'' captures the dynamic tension within the self that was at play in the Big Gorilla moment of breakdown and repair:

\begin{quotation}
The \{I\} is the response of the organism to the attitudes of the others; the ``me'' is the organized set of attitudes of others which one himself assumes\ldots{} The \{I\} is something that is, so to speak, responding to a social situation which is within the experience of the individual. It is the answer which the individual makes to the attitude which others take toward him when he assumes an attitude toward them \parencite[\S 22]{mead_mind_2015}.
\end{quotation}

The \textbf{``me''} is the social self, shaped by internalized attitudes of others. It is the self-as-recognized, the self that can be described, categorized, and evaluated. The ``me'' is finite, made determinate through commitments that constitute social identity. In the Big Gorilla story, the ``me'' was the part of myself that had learned appropriate ways to interact with children, that knew social expectations for adult behavior in a family setting.

The \textbf{\{I\}} is the spontaneous, unreflective aspect of the self---the subject of action. It is not an object among objects but the condition of possibility for all representation. It is the \emph{in}finite source of what Carspecken calls ``power, creativity, and freedom,'' the wellspring of action that always has the quality of surprising us. The \{I\} was the spontaneous actor who invented Big Gorilla in the first place, who thumped too loudly when frustrated, and who had to respond to the adult's disapproval.

Crucially, the \{I\} and the ``me'' are not two separate entities but two moments of a single process. The \{I\} becomes a ``me'' through action and recognition. The ``me'' is animated and potentially transformed by the \{I\}. Any attempt to separate them destroys the structure that makes selfhood possible. The self emerges from the ongoing conversation between these aspects, mediated by significant symbols that allow us to take the attitude of others toward ourselves.

This I-me dialectic becomes crucial for understanding how narrative breakdown motivated the need for formal analysis. The \{I\} that thumped too loudly was acting spontaneously, but the ``me'' that recognized the transgression through \(\mathcal{A}\)'s look had to develop tools for articulating what had gone wrong and how to repair it.

\section{Communicative Reason: Habermas's Intersubjective Turn}\label{communicative-reason-habermass-intersubjective-turn}

While Mead provides the foundation for understanding how selves emerge through symbolic interaction, Jürgen Habermas extends this insight to explain how rational communication itself is structured. His \emph{Theory of Communicative Action} \parencite*{Habermas1984,Habermas:1985aa} distinguishes between modes of social interaction that prove crucial for understanding what happened in the Big Gorilla story.

\subsection*{Communicative Action and the Lifeworld}\label{communicative-action-and-the-lifeworld}

Most everyday interactions proceed smoothly through what Habermas calls \textbf{communicative action}---interaction oriented toward reaching mutual understanding through speech. This contrasts with \textbf{strategic action}, which treats others as means to one's ends, and \textbf{instrumental action}, which manipulates objects to achieve goals. There is nothing intrinsically wrong with strategic or instrumental action---we need both---but it can feel very wrong to engage in strategic action under the guise of communicative action. Saying ``I just want to understand'' while actually manipulating someone toward a predetermined outcome violates the norms of genuine communication.

Communicative action relies on a shared background of taken-for-granted assumptions---what Habermas calls the \textbf{lifeworld}. When I first began thumping my chest as Big Gorilla, the twins and I quickly developed a shared understanding about what this gesture meant. No explicit discussion of rules was needed; meaning emerged organically from our repeated interactions and was sustained by mutual recognition of its significance. This taken-for-granted character allows communication to flow efficiently.

The lifeworld contrasts with what Habermas calls the \textbf{system}---the realm of institutions, bureaucratic norms, and formal structures that can become alienated from lived experience. The kids were under pressure to get to bed at a reasonable hour because \(\mathcal{A}\) and I both had to work the next day. Bound to the wheel of capitalism with its clocks and schedules, one reason for inventing the playful character was to deflate those systemic forces. When communicative action functions well, the system supports the lifeworld by providing a stable framework. When the system becomes too rigid, it disrupts the fluidity of communication.

Habermas's critical theory diagnoses what he calls the \emph{colonization of the lifeworld by the system}. This occurs when systemic imperatives---efficiency, productivity, calculability, formal rationality---invade domains that should be governed by communicative rationality. In education, standardized testing regimes, accountability metrics, and bureaucratic mandates override teachers' pedagogical judgment and disrupt communicative relationships between teachers and students. The system treats education as a mechanism for producing measurable outcomes.

In the Big Gorilla story, the systemic pressure of work schedules created the need for bedtime routines, which then became sites of potential conflict. The playful invention of Big Gorilla was an attempt to \emph{re}colonize that systemically structured moment (bedtime) back into the lifeworld---to transform an imposed requirement into a shared practice with its own intrinsic meaning, creating space for communicative negotiation: ``Baby Gorilla still awake, but Big Gorilla tired. What we do?''

This distinction between lifeworld and system connects to the dialectical structure we've been tracing. The system corresponds to \emph{Understanding} (\emph{Verstand})---it operates through fixed categories, abstract rules, universal procedures that subsume particulars. The lifeworld corresponds to \emph{Reason} (\emph{Vernunft})---it operates through situated judgment, responsive engagement, creative negotiation of norms in concrete contexts. Both are necessary. Without systemic structures, social coordination at scale would be impossible. Yet when systemic rationality becomes totalizing, it destroys the communicative fabric that makes human life meaningful.

\subsection*{Rational Discourse and Validity Claims}\label{rational-discourse-and-validity-claims}

Communicative action can break down. When the background consensus of the lifeworld is challenged or disrupted, a shift to what Habermas calls \textbf{rational discourse} becomes necessary. This is exactly what happened when I thumped too loudly and \(\mathcal{A}\) gave me that look. The taken-for-granted understanding that Big Gorilla was safe and playful was called into question.

Habermas's theory distinguishes between different formal-pragmatic speech positions---first-person, second-person, and third-person---to which speakers relate. There are some unfortunate consequences of the term `world' here, that Carspecken's work \parencite*{Carspecken2018} improves on (though the specific improvements I cite next are explicit in Habermas' work). Rather than ``worlds,'' it's more precise to think of these as positions marked by different pronouns and corresponding to different types of validity claims that participants implicitly or explicitly raise:

\textbf{Subjective validity claims (first-person):} The first evidence that repair is necessary is felt in subjective experience. When made explicit as criticizable validity claims, critique hinges on taking the speaker as \emph{sincere} in their expressions. The kids might joke ``Big Gorilla is scary!'' as part of the game. But if I take them as sincere---if I recognize they're actually frightened---then I must change to keep the practice viable. Sincerity is a precondition for rationality in communicative action.

\textbf{Objective validity claims (third-person):} When participants relate to the objective world, they make \emph{truth claims} about existing states of affairs. I might challenge a subjective validity claim on objective grounds: ``I hear giggling! Are you sure you're scared?'' Whether the child giggles is \emph{transsubjective}---\(\mathcal{A}\) or anyone else can note the behavior. An expression satisfies the precondition for rationality if it embodies fallible knowledge, relates to facts, and is open to objective judgment. A judgment is objective if it is undertaken on the basis of a transsubjective validity claim that has the same meaning for observers as for the acting subject. The objectivity of the world is constituted intersubjectively---the world gains objectivity through counting as one and the same world for a community of speaking and acting subjects.

\textbf{Normative validity claims (second-person):} When participants relate to a common social world, they make \emph{rightness claims} concerning normatively regulated interpersonal relations or legitimate norms themselves. A social world consists of a normative context defining the totality of legitimate interpersonal relations. A norm is considered ``in force'' when it is recognized as valid by those to whom it is addressed. Crucially, a norm is ideally valid when it \emph{deserves the assent of all those affected} because it regulates problems of action in their common interest. This normative validity means that a norm could be accepted with good reasons by everyone involved.

These three validity claims have formal structure: each can be \emph{raised} implicitly in action, \emph{challenged} by others, and \emph{defended} through reason-giving. This structure is what Habermas means by the \textbf{rationality} of communicative action---not that participants always reason explicitly, but that the practice contains this dimension of potential critique and vindication.

Every speech act implicitly commits the speaker to all three dimensions simultaneously. When I say ``The baby gorillas need to sleep,'' I am implicitly claiming that this statement is \emph{true} (an objective fact about their tiredness), \emph{right} (I have appropriate standing to make this judgment), and \emph{sincere} (I genuinely believe it). Any of these claims can be challenged: ``Are they really tired, or do you just want quiet?'' (truth), ``Who are you to tell us when to sleep?'' (rightness), ``You say that, but you seem more interested in watching TV than their wellbeing'' (sincerity).

Rational discourse resolves disagreements through reasons rather than force. When a validity claim is challenged, the speaker can withdraw the claim, modify it, or provide reasons in its defense. Those reasons themselves can be challenged, creating what Brandom calls the ``game of giving and asking for reasons.'' This game has no guaranteed endpoint---reasons can always be questioned further---yet it creates the normative space within which genuine understanding becomes possible.

When \(\mathcal{A}\) looked at me disapprovingly, she questioned all three validity dimensions simultaneously. Was my thumping really playful, or had it crossed into genuine aggression? (truth) Was such intensity appropriate, even in play? (rightness) Was I being honest about my intentions, or using ``play'' as cover for venting frustration? (sincerity) The breach called into question the entire practice and my competence as a participant in it. Repair required more than a simple apology; it required explicitly addressing all three dimensions and reestablishing trust across all of them.

When I said, ``Oh, that was too loud. Big Gorilla sorry. Big Gorilla get frustrated when Baby Gorillas don't listen to Momma Gorilla,'' I was engaging in rational discourse. I was making explicit claims about objective facts (too loud), acknowledging the inappropriateness of my action (rightness), and revealing my internal state (sincere apology). This verbal response served to repair the breach in our communicative relationship and reestablish the normative boundaries of our shared practice.

\subsection*{The Possibility of Saying ``No''}\label{the-possibility-of-saying-no}

Central to Habermas's understanding is that genuine communication requires that either participant be able to say ``no'' to the other's claims. This possibility of negation distinguishes authentic dialogue from manipulation or coercion. When \(\mathcal{A}\) looked at me disapprovingly, she was saying ``no'' to my behavior. Her negation created space for me to respond rationally by taking a second-person position on my own actions---by recognizing my ``me'' from the perspective of the generalized other.

This dynamic of assertion and potential negation creates what Robert Brandom calls the \emph{game of giving and asking for reasons}. In this game, participants make claims that others can challenge, and reasons are provided when those claims are questioned. This rational structure underlies all genuine communication and learning.

For Big Gorilla, what began as simple gestural play developed into a sophisticated practice of mutual accountability. When the norms of the practice were violated, the breach could be repaired through reason-giving rather than abandoning the practice. This ability to examine and refine shared understandings is what allows communicative relationships to grow and deepen over time.

\section{Analytic Pragmatism: Brandom's Meaning-Use Analysis}\label{analytic-pragmatism-brandoms-meaning-use-analysis}

While Mead and Habermas provide powerful insights into the social nature of selfhood and communication, Brandom's analytic pragmatism offers precise tools for understanding how meaning emerges from use and how practices can be algorithmically elaborated. The Big Gorilla story illustrates a puzzle that neither Mead nor Habermas fully resolves: How exactly do we get from the spontaneous emergence of significant symbols to the kind of explicit, reflective discourse that allows us to articulate and revise the norms governing our communicative practices?

Brandom's answer requires developing what he calls a \textbf{pragmatic metavocabulary}---a vocabulary that can make explicit the implicit rules governing our practices.

\subsection*{The Foundation: Vocabularies and Practices-or-Abilities}\label{the-foundation-vocabularies-and-practices-or-abilities}

Brandom's methodology of \emph{meaning-use analysis} (MUA) is grounded in a principled separation that may seem simple but proves philosophically enriching. Rather than treating meaning as a relation between words and worldly objects---the standard picture where dictionaries list what words refer to---Brandom proposes that we understand meaning as emerging from the relationship between two distinct but interrelated components: \emph{vocabularies} (\(V\)) and \emph{practices-or-abilities} (\(P\)).

\textbf{Vocabularies (\(V\)):} A vocabulary is what is \emph{said}---a set of symbols, terms, or expressions that can be inscribed, uttered, or otherwise tokened. This could be a natural language like English, a technical jargon like that of quantum mechanics, a mathematical symbolism like first-order logic, or even the gestures and vocalizations we've been discussing. The crucial point is that a vocabulary \(V\) consists of repeatable expressions that can be produced and recognized across different contexts. In our case study, ``hoohoohoo'' is part of Big Gorilla's vocabulary. So are the gestures in a dogs' fight---\{growl, stare, lunge, bite, cower\} forms a vocabulary of canine aggression.

\textbf{Practices-or-Abilities (\(P\)):} A practice is what is \emph{done}---a pattern of activity, a repertoire of performances, an ability to behave in certain rule-governed ways. In linguistic contexts, practices include abilities like making inferences, following grammatical rules, responding appropriately to others' utterances, and generally using expressions correctly within specific contexts. For Big Gorilla, this included not just the ability to produce the sound ``hoohoohoo'' but the ability to do so \emph{appropriately}---when frustrated with the kids, but in a way that remained playful rather than threatening. Crucially, practices have normative structure: there are correct and incorrect ways to engage in them, even when those norms remain implicit.

This separation between \(V\) and \(P\) is methodologically crucial because it allows us to analyze \emph{how} meaning is constructed rather than simply taking it as given. We can ask: What abilities must someone master to count as deploying a vocabulary competently? And conversely: What vocabulary would allow us to make explicit the rules governing a practice? These questions structure Brandom's entire approach.

The power of this framework becomes apparent when we recognize that these two questions---What practices deploy vocabularies? What vocabularies specify practices?---are not symmetrical. They point in different directions, and understanding their asymmetry is key to understanding how meaning works.

\subsection*{PV-Sufficiency and VP-Sufficiency: The Asymmetry of Meaning}\label{pv-sufficiency-and-vp-sufficiency-the-asymmetry-of-meaning}

Brandom formalizes the relationship between practices and vocabularies through two complementary but asymmetric notions of \emph{sufficiency}:

\textbf{PV-Sufficiency (Practice-Vocabulary Sufficiency):} A practice \(P\) is PV-sufficient for deploying a vocabulary \(V\) if mastering that practice essentially equips you to use that vocabulary. The ``essentially'' here does crucial work---it means that engaging in the practice \emph{constitutes} your ability to use the vocabulary, not that it enables or facilitates it.

Consider learning to play chess. Mastering the practice of moving pieces according to chess rules is PV-sufficient for deploying the vocabulary of chess notation (e4, Nf3, etc.). You couldn't genuinely understand chess notation without being able to play, at least in principle. The notation isn't a separable layer of labels applied to moves you could already make---understanding the notation just \emph{is} understanding the moves as moves in the game.

For Big Gorilla, the practice of playful-but-bounded chest-thumping is PV-sufficient for deploying the vocabulary of ``hoohoohoo.'' The twins couldn't understand what those sounds meant without participating in the practice of playful anger they expressed. A Finite State Automaton could model one such practice---but it would need expressive resources to represent intensity variations like ``HOOHOOHOO'' in capital letters. This limitation reveals that different practices have different expressive power, a point that will matter enormously when we discuss arithmetic.

\begin{figure}[h!]
    \includegraphics[width=0.8\textwidth]{/Users/tio/Documents/GitHub/September_UMEDCA/images/ch4_gorilla_fsa.pdf}
    \caption{\textit{Note.} A Finite-State Automaton for Big Gorilla's Grunts}
    \label{fig:gorilla_fsa}
\end{figure}

\textbf{VP-Sufficiency (Vocabulary-Practice Sufficiency):} A vocabulary \(V\) is VP-sufficient to specify a practice \(P\) if it affords the resources to explicitly formulate the rules governing that practice---to say what one must do to engage in the practice correctly.

Consider a Finite State Automaton specification for Big Gorilla's grunts. The vocabulary of finite state automata---with its states, transitions, input symbols, and acceptance conditions---is VP-sufficient to specify the practice of generating well-formed Big Gorilla vocalizations. Someone who couldn't speak like Big Gorilla could nonetheless read the automaton's specification and understand what would count as doing so correctly. The vocabulary makes the implicit norms of the practice explicit.

But here's the asymmetry: while the practice of playing chess is PV-sufficient for deploying chess notation, chess notation is \emph{not} VP-sufficient to specify the practice of playing chess. The notation can record what moves were made, but it cannot fully specify the practical abilities required to recognize threats, plan combinations, understand positional principles. Those abilities exceed what the notation can express. This is why chess books must supplement notation with prose explanation, diagrams, and ultimately, you must play.

This asymmetry reveals something deep about the relationship between saying and doing. What we can \emph{do} with language often exceeds what we can \emph{say} about doing it. The gap between PV-sufficiency and VP-sufficiency is the gap between implicit practical mastery and explicit theoretical articulation. Closing this gap---developing vocabularies that are VP-sufficient to specify practices we can already engage in---is the project of making meaning explicit.

Consider a child who can count competently---who deploys numeral vocabulary (``one, two, three, four, five'') through mastery of counting practice. The practice is PV-sufficient for deploying the vocabulary. Yet can the child \emph{say} what counting is? Can they articulate the rules governing when counting is done correctly? Likely not. The numeral vocabulary itself is not VP-sufficient to specify the practice of counting. We need a richer metavocabulary---formal arithmetic, a theory of cardinality, algorithmic specifications---to make explicit what remains implicit in the child's competent performance.

When teachers present formal definitions and axioms \emph{before} students have developed the practices those formalisms specify, the VP-sufficient vocabulary arrives without the corresponding PV-sufficient practice. Students can manipulate symbols according to rules without understanding what those symbols \emph{mean}---without recognizing them as making explicit the practices they already (or could) engage in. When mathematics education focuses exclusively on procedures without developing metavocabularies that make those procedures explicit, students remain at the level of implicit know-how. They can calculate but cannot explain, can follow algorithms but cannot justify them.

The pedagogical task is bidirectional: we must facilitate the development of mathematical practices (achieving PV-sufficiency for mathematical vocabularies) while developing the metavocabularies that make those practices explicit (achieving VP-sufficiency). Practice without articulation remains blind---competent but unreflective. Articulation without practice remains empty---formally correct but meaningless.

\subsection*{Pragmatic Metavocabularies: Making the Implicit Explicit}\label{pragmatic-metavocabularies-making-the-implicit-explicit}

The concept of a \textbf{pragmatic metavocabulary} captures what it means to close this gap, to achieve what Brandom calls \textbf{algorithmic elaboration} of a practice. Formally, a pragmatic metavocabulary \(V'\) stands in a specific relation to a vocabulary \(V\):

\begin{quote}
\(V'\) is a pragmatic metavocabulary for \(V\) when \(V'\) is VP-sufficient to specify practices-or-abilities \(P\) that are PV-sufficient to deploy vocabulary \(V\).
\end{quote}

In other words:

\begin{enumerate}
\def\labelenumi{\arabic{enumi}.}
\item We start with some vocabulary \(V\) (like Big Gorilla's grunts)
\item There are practices \(P\) that, when mastered, equip you to use \(V\) (the practice of playful-yet-bounded frustration)
\item A pragmatic metavocabulary \(V'\) is one that can explicitly articulate the norms governing those practices
\end{enumerate}

\(V'\) lets you \emph{say} the rules for \emph{doing} what \(V\) lets you \emph{do}. The vocabulary \(V'\) makes explicit what was implicit in the practice \(P\).

The language of finite state automata serves as a pragmatic metavocabulary for Big Gorilla's grunts. An automaton specification, with its states and transitions, explicitly formulates the rules (VP-sufficiency) for the practice of generating the vocalizations (PV-sufficiency). Someone who mastered the automaton specification could generate well-formed strings even without having engaged in the original playful practice with the kids.

But notice what has happened: the pragmatic metavocabulary \(V'\) creates a new practice---the practice of reading and writing automaton specifications. And that practice might itself require its own metavocabulary to make \emph{its} norms explicit. This threatens an infinite regress, but each level of metavocabulary increases our expressive power, allowing us to say more about what we're doing when we use language.

Crucially, different pragmatic metavocabularies have different expressive power. The informal vocabulary I used earlier---``Big Gorilla says hoohoohoo when frustrated but stays playful''---is a pragmatic metavocabulary, but it's less powerful than the FSA formalism. The FSA can specify exactly which strings are well-formed and which aren't. Later, when we encounter the Hermeneutic Calculator, we'll trace how formal metavocabularies for arithmetic can make explicit the implicit structure of children's counting strategies.

The concept of \textbf{discursive levels} emerges naturally from this framework. Individuals operating at different discursive levels can engage in the same practices while differing in their ability to deploy pragmatic metavocabularies for those practices. The twins and I could all engage in the Big Gorilla practice (same \(P\), same \(V\)), but I could deploy a more powerful metavocabulary (\(V'\)) to explicitly formulate its rules. Operating at a higher discursive level doesn't make one ``better'' at the practice---the twins were often more spontaneous and creative in their play than I was. But it does afford different expressive resources for talking \emph{about} the practice, for noticing when norms are violated, for proposing explicit revisions.

\begin{figure}[h!]
    \includegraphics[width=0.8\textwidth]{/Users/tio/Documents/GitHub/September_UMEDCA/images/Ch4_Cyclical_Automaton.pdf}
    \caption{\textit{Note.} A Cyclical Automaton for a Complete Communicative Act}
    \label{fig:circular_automaton}
\end{figure}

\subsection*{Indexicals and Anaphora: The Fundamental Case}\label{indexicals-and-anaphora-the-fundamental-case}

The relationship between indexical and anaphoric vocabularies provides the paradigm case for understanding how practices and vocabularies interrelate, and it will prove crucial for the argument in Chapter 7 that numerals function as pronouns rather than names.

\textbf{Indexical expressions}---terms like ``I,'' ``here,'' ``now,'' and demonstratives like ``this'' and ``that''---have referents that shift with each utterance depending on who is speaking, where, and when. When I say ``I am writing,'' the word ``I'' refers to me. When you read this and think ``I am reading,'' the same word ``I'' now refers to you. The referent shifts, but the word's \emph{meaning} (its role in the language) remains constant.

\textbf{Anaphoric expressions}---pronouns like ``it,'' ``she,'' ``they,'' ``he''---achieve reference by \emph{recollecting} something previously established in discourse. They point back to an antecedent, carrying forward a reference without having to re-establish it from scratch.

The deep insight Brandom develops from classical analytic philosophy (especially Frege, Russell, and Quine) is that these two vocabularies stand in a relation of mutual presupposition that reveals something fundamental about how language achieves stability:

\begin{quote}
Using indexicals meaningfully presupposes the ability to use anaphoric terms, because without anaphora, we could not establish co-reference across time or build inferential chains. Anaphora makes indexical content ``repeatable'' and available for reasoning.
\end{quote}

Consider Brandom's example: ``This chalk is white. It is also cylindrical'' \parencite*[129]{Brandom:2019aa}. The demonstrative ``this'' (indexical) picks out a particular piece of chalk in the immediate perceptual context. But to say anything further about that same chalk---to build an inference, to attribute additional properties---requires the pronoun ``it'' (anaphora) to carry forward the reference. Without anaphora, every claim would require re-demonstrating: ``This chalk is white. This chalk {[}\emph{pointing again}{]} is cylindrical.'' Communication would collapse into an endless series of isolated demonstrations.

But here's the crucial point for our purposes: \emph{anaphora works by recollection, not by reference to a pre-given object}. When I use ``it'' to refer back to ``this chalk,'' I'm not independently accessing the chalk and verifying that ``it'' and ``this chalk'' pick out the same object. Rather, ``it'' achieves its reference \emph{by inheriting} the reference established by ``this chalk.'' The pronoun doesn't name an object; it recollects a previous act of referring.

Now consider an inference:

\begin{itemize}
\item \textbf{Premise 1:} This paper is green. (symbolized as \(\Psi a\), where \(\Psi\) is the predicate ``is green'' and \(a\) secures reference to ``this paper'')
\item \textbf{Premise 2:} This paper is my to-do list. (symbolized as \(a = b\), where \(b\) secures reference to ``my to-do list'')
\item \textbf{Conclusion:} Therefore, my to-do list is green. (symbolized as \(\Psi b\))
\end{itemize}

This inference has the form: If \(\Psi a\), and \(a = b\), then \(\Psi b\). Algebraically, it looks valid---we substitute equals for equals. But as Brandom points out, there's no guarantee this inference succeeds in ordinary discourse. Imagine I point to a green piece of paper when uttering Premise 1, then gesture toward a yellow notepad when uttering Premise 2, ``This paper is my to-do list.'' My to-do list, written on yellow paper, is \emph{not} green, even though ``this paper'' in Premise 1 \emph{was} green.

The problem is \textbf{securing co-reference}. How do we ensure that ``\(a\)'' in both premises refers to the \emph{same} thing? In formal logic, we stipulate co-reference by using the same variable. But in actual discourse, co-reference must be \emph{recognized} and \emph{maintained} across different utterances and different contexts.

\begin{figure}[h!]
    \includegraphics[width=0.8\textwidth]{/Users/tio/Documents/GitHub/September_UMEDCA/images/Ch4_Indexicals_Anaphora.pdf}
    \caption{\textit{Note.} Reproduction from Brandom \parencite*[59]{Brandom2008}: Pragmatically mediated semantic presupposition of anaphoric by indexical and deictic vocabularies}
    \label{fig:Indexicals_Anaphora}
\end{figure}

Anaphora becomes essential here. The pronoun ``it'' doesn't just happen to refer to what ``this'' referred to---it \emph{makes explicit} that we're still talking about the same thing. Anaphora is the linguistic mechanism for tracking identity across time and context. Without it, the substitution inference fails.

\textbf{Numerals function anaphorically, not indexically or as names}. When we use ``2,'' we're not naming a platonic object that exists independently of our practices. Rather, ``2'' \emph{recollects} the act of counting---one, two. It achieves stability through recollection, just as ``it'' achieves reference by recollecting ``this chalk''\parencite[129]{Brandom:2019aa}. The numeral inherits its content from the practice it recollects.

This is not yet the full argument---that awaits Chapter 7---but the indexical-anaphoric structure establishes the pattern. Just as ``it'' makes indexical reference repeatable without requiring constant re-demonstration, numerals make counting repeatable without requiring constant re-performance. Both achieve semantic stability through the \emph{temporal structure of recollection} rather than through direct reference to pre-given objects.

The pragmatic metavocabulary we've been developing---vocabularies, practices, PV-sufficiency, VP-sufficiency---provides the tools to articulate this structure precisely. Counting is a practice (\(P\)) that is PV-sufficient for deploying numeral vocabulary (\(V\)). The numeral ``2'' recollects that practice anaphorically. And formal arithmetic serves as a pragmatic metavocabulary (\(V'\)) that is VP-sufficient to specify the rules governing counting practices, making explicit what was implicit in the practice of enumerating.

This connection between the anaphoric structure of pronouns and the recollective structure of numerals will become central to reconceptualizing what numbers are and how mathematical knowledge is possible. For now, what matters is recognizing that meaning emerges not from naming pre-existent objects but from the temporal structure of practices that can be recollected and made explicit through increasingly powerful metavocabularies.

If numerals function anaphorically, the question ``What is the number 2?'' receives a different answer than Platonism provides. The Platonist says: 2 is an abstract object existing independently of our practices, and the numeral ``2'' names that object. On the anaphoric account, ``2'' \emph{recollects an act of counting}. When I write ``2 + 2 = 4,'' I am making explicit a relationship between counting practices: the practice of counting to two, repeated twice, yields the same result as the practice of counting to four.

Anaphora, in some sense, reflects structuration. Once the impetus to act is expressed in conditions that allow for others to recognize the impetus as shared, it seems to become stable, or fixed. I can't help but speak of `mathematical objects,' even while I am uniquely positioned to understand them as expressive acts. If mathematical objects are structures instituted through recollectable practices, then mathematical knowledge is \emph{explication} of implicit structures in our practices. This does not make mathematics arbitrary or subjective---the practices have normative structure, and there are correct and incorrect ways to engage in them. Yet it relocates the source of mathematical necessity from a platonic realm to the normative structure of practices we institute and sustain through mutual recognition. Regardless, if math is a `structure' it must be built to break. 

The parallel with the ``you'' of address is illuminating. Just as ``you'' achieves reference through the transaction of address, numerals achieve reference through the practice of counting. The transactional ``you'' creates an intersubjective relationship; numerals create mathematical relationships through the recollection and iteration of counting practices. The CUSP ``you'' invokes generalized norms; formal arithmetic invokes generalized structures abstracted from particular acts of counting. The movement from concrete counting to formal arithmetic parallels the movement from transactional address to CUSP generalization---both involve a transition from particular enactments to universal structures, mediated by anaphoric recollection.

\subsection*{Rhetorical Anaphora: Repetition as Accumulation}\label{rhetorical-anaphora-repetition-as-accumulation}

The technical account of anaphora we've developed---pronouns that recollect previous reference and establish co-reference across time---captures one dimension of how these terms function. But there is another dimension, particularly evident in poetry, song, and rhetoric, where anaphora does something more than secure reference. It \textbf{builds}.

Consider Martin Luther King Jr.'s ``I Have a Dream'' speech, with its famous repeated phrase:

\begin{quote}
I have a dream\ldots{}\\
I have a dream\ldots{}\\
I have a dream\ldots{}
\end{quote}

Each repetition doesn't just recollect the previous utterance. Each iteration \emph{accumulates}---building intensity, creating momentum, increasing what we might call the \emph{magnitude} of the claim. The semantic content remains relatively stable (he continues to have a dream), but something about the speech act's force grows with each repetition.

This is \textbf{rhetorical anaphora}---the deliberate repetition of a phrase at the beginning of successive clauses or sentences. Unlike the anaphoric pronouns we've been discussing, rhetorical anaphora doesn't substitute a pronoun for a previous referent. Instead, it \emph{repeats} the same phrase, and through that repetition, transforms assertion into something more powerful: proclamation, invocation, commitment.

The connection to mathematical thinking becomes clear when we consider \emph{iteration} as a fundamental mathematical operation. When we count---one, two, three---we are iterating the successor operation. When we add 3+3+3+3, we are iterating addition. When we take limits in calculus, we iterate approximations. Each iteration doesn't just repeat the previous step; it \emph{accumulates} toward something---a sum, a magnitude, a limit.

The rhetorical force of anaphoric repetition operates similarly. Each ``I have a dream'' is not the same speech act performed again. It is the building toward a magnitude that no single utterance could achieve. The repetition \emph{is} part of the content---the insistence, the refusing-to-let-go, the building of collective energy.

This recalls the accumulation of forestructures in Chapter 1's Exercise. When you brought awareness to your toes, then let them go, they became forestructures---implicit background contributing to accumulating bodily awareness. As you continued---feet, legs, hips---each region added to the set: \{\st{toes}, \st{feet}, \st{legs}\}. This wasn't addition but sublation: each region was preserved (still part of your body), negated (no longer the focus), and uplifted (integrated into expanding awareness).

Rhetorical anaphora enacts a similar structure. Each repetition preserves what came before (the dream remains), negates it (this isn't the same utterance), and uplifts it (the accumulated force exceeds any single instance). The repetition builds toward something that couldn't exist without the iteration---just as the full body awareness in The Exercise couldn't emerge without moving through and accumulating the sublated regions.

The songs throughout this text employ similar structures. Choruses repeat not because the singer forgot what they already said, but because repetition does work that saying-once cannot do. The recurrence creates emphasis, builds emotional intensity, invites participation. When Big Gorilla says ``hoohoohoo'' repeatedly, the repetition isn't redundant---it's the structure of the play itself, the back-and-forth that constitutes the practice.

Brandom's formal apparatus, powerful as it is, doesn't capture everything about how anaphora functions in actual discourse. The VP-sufficiency/PV-sufficiency framework can specify when anaphoric terms are deployed correctly. But it struggles to capture what rhetorical anaphora \emph{does}---how it builds, accumulates, intensifies.

We might say that rhetorical anaphora introduces a \textbf{quantitative dimension} to language use that exceeds purely qualitative semantic analysis. It's not just \emph{what} is said (semantic content) or \emph{how} it commits one to further inferences (pragmatic consequence). It's also \emph{how much} gets built up through repetition---a magnitude that emerges from iteration itself.

When we formalize arithmetic practices in later chapters, this will matter. The child who counts ``one, two, three, four'' is not just establishing reference to discrete cardinalities. They are \emph{building} through iteration. Each number recollects the previous counting act and adds one more. The numeral ``four'' carries within it the accumulated force of four iterations, just as the fourth ``I have a dream'' carries the accumulated force of four proclamations.

The formal apparatus we're about to develop through automata can capture the \emph{structure} of this iteration---the states, transitions, and rules governing repetition. But it cannot capture the felt sense of accumulation, the experiential quality of building toward something. That dimension requires the songs, poems, and narratives that surround the formal analysis, gesturing toward what the formalization necessarily leaves out.

\subsection*{``Silver and Iron'': Accumulation as Daily Ritual}\label{silver-and-iron-accumulation-as-daily-ritual}

The song ``Silver and Iron'' that prefaces this chapter operates through a different kind of anaphoric accumulation than King's oratorical proclamation. Where ``I have a dream'' builds toward climactic intensity, the repeated questioning in ``Silver and Iron'' accumulates quietly, like forestructures in The Exercise:

\begin{quote}
Who are you today, love?\\
How are your stars arranged?\\
What plans have you forgotten?\\
What's still written on your hands?
\end{quote}

Each question doesn't demand an answer but creates space for one. The repetition across verses---asking variations of ``who are you?''---doesn't insist but invites. Iteration can build magnitude without building volume. The accumulation is one of intimacy rather than proclamation, depth rather than height.

The song occupies an interesting position between the transactional and CUSP ``you.'' It addresses a specific person, making it transactional---this ``you'' is irreplaceable, particular. But the ritualized questioning (``Who are you \emph{today}?'') also invokes something more universal: the practice of daily mutual rediscovery that any partnership might engage in. The particular and universal interpenetrate. The song addresses \emph{this} person while also articulating a norm about how one might address any beloved other: with curiosity about their becoming rather than assumptions about their fixed identity.

This connects to Mead's insight about the \{I\} and the ``me.'' The question ``Who are you today?'' acknowledges that the beloved is not reducible to the ``me'' I already know---the accumulated past recognitions. There is always an \{I\} that exceeds those recognitions, that might arrange its stars differently today than yesterday. The question creates space for that spontaneous, unpredictable aspect to emerge.

Musically, the song builds through verse repetition with slight variations, each iteration adding another layer to what we might call the \emph{forestructure of trust}. Just as in The Exercise each sublated body region accumulates into expanding awareness, each verse accumulates into deepening recognition. The magnitude that builds is relational rather than rhetorical---a quantity of intimacy that emerges through iteration of the practice itself.

The mathematical concept of iteration---repeating an operation to build toward something---has broader application than formal systems. Iteration operates in embodied awareness (The Exercise), in rhetorical force (King), in relational depth (the song), and as we'll trace in later chapters, in arithmetic practices where counting is iterated to build magnitude. The formal apparatus can model the \emph{structure} of iteration, but the \emph{quality} of what accumulates---political energy, proprioceptive awareness, relational intimacy, numerical quantity---exceeds what any single formalization can capture.

\section{Formalizing Practice: Automata as Pragmatic Metavocabularies}\label{formalizing-practice-automata-as-pragmatic-metavocabularies}

While Mead and Habermas provide rich accounts of the social and normative dimensions of intersubjectivity, their frameworks remain largely at the level of what Brandom would call the semantic---concerned with meaning and understanding. To grasp the underlying machinery of these practices---the repeatable, rule-governed structures of \emph{use}---we need more formal apparatus. This is the domain of analytic pragmatism, and its basic tool is the automaton.

An automaton is a theoretical machine that reads symbols and changes state according to fixed rules. Different classes of automata have different expressive power, and this hierarchy of expressiveness will prove important for understanding mathematical practices.

\subsection*{Turn-Taking: A Simple Finite-State Automaton}\label{turn-taking-a-simple-finite-state-automaton}

Basic conversational turn-taking can be modeled as a simple finite-state automaton (FSA). This automaton captures the basic rules of conversational exchange, where each state represents a different stage in the interaction and transitions represent the actions taken by either participant.

\begin{figure}[h!]
    \includegraphics[width=0.8\textwidth]{/Users/tio/Documents/GitHub/September_UMEDCA/images/Ch4_fsm_turn_taking.pdf}
    \caption{\textit{Note.} A Simple Finite-State Automaton of Conversational Turn-Taking. Note on Halt Condition: Model shows successful turn-taking. Simultaneous initiation from \(S_0\) (``collision'') breaks understanding, requires reset.}
    \label{fig:fsm_turn_taking}
\end{figure}

The automaton for turn-taking has four states: \(S_0\): Idle (both silent); \(S_A\): Alter speaks, Ego listens; \(S_E\): Ego speaks, Alter listens; \(S_{end}\): Conversation ends.

The system starts in the idle state. If Alter initiates speech, it moves to \(S_A\). From there, if Ego responds, it moves to \(S_E\). The system can loop between \(S_A\) and \(S_E\) as Alter and Ego respond to each other. The interaction successfully terminates when one party signals closure from a speaking state.

However, if both try to speak simultaneously from the idle state---a ``collision''---the system breaks down. This simple model captures the basic turn-taking structure but lacks the depth and flexibility of real communicative practices. It cannot represent the nuances of interruption, back-channeling, or simultaneous speech that actual conversation employs.

\subsection*{Big Gorilla: A Transducing Automaton}\label{big-gorilla-a-transducing-automaton}

To capture the richer structure of the Big Gorilla practice, we need what Brandom calls a \textbf{transducing automaton}. While an FSA simply reads strings of symbols, a transducing automaton generalizes this to discriminating stimuli of any kind on the input side and differentially responding in any way on the output side \parencite[34]{Brandom2008}.

A transducing automaton for Big Gorilla's grunts could be specified as:

\textbf{States:} \(Q = \{S_0, S_1, S_2, S_3, S_4\}\) 
- \(S_0\): Initial/neutral state 
- \(S_1\): After initial `h' 
- \(S_2\): After `ho' or `hO' 
- \(S_3\): Ready for more `hoo' or terminal `!' 
- \(S_4\): Accepting state (grunt complete)

\textbf{Input alphabet:} \(\Sigma = \{h, o, O, !\}\) 

\textbf{Output alphabet:} \(\Gamma = \{h, o, O, !\}\) (same as input for this simple case)

\textbf{Transition function:} \(\delta: Q \times \Sigma \rightarrow Q \times \Gamma\)

The key transitions: 
- From \(S_0\) with `h' $\rightarrow$ \(S_1\), output `h' 
- From \(S_1\) with `o' or `O' $\rightarrow$ \(S_2\), output `o' or `O' 
- From \(S_2\) with `h' $\rightarrow$ \(S_1\), output `h' (allowing repetition) 
- From \(S_2\) with `!' $\rightarrow$ \(S_4\), output `!' (termination)

This specification is VP-sufficient to define the practice of generating well-formed Big Gorilla vocalizations. Someone who had never heard Big Gorilla grunt could read this specification and produce acceptable strings like ``hoohoohoo!'' or ``hOOhOOhOO!''

But notice what the automaton \emph{cannot} represent: the intensity conveyed by capital letters, the emotional force of different grunt patterns, the contextual appropriateness (frustration vs.~playfulness), or the actual thumping gesture that accompanies the vocalization. The formal specification captures structure but not semantics, rules but not meaning.

\subsection*{The Limits of Finite-State Machines}\label{the-limits-of-finite-state-machines}

Finite-state automata have fundamental limitations. As Brandom points out, following Chomsky, no FSA can handle languages like \(a^n b^n\) where the number of 'a's must equal the number of 'b's \parencite[21]{Brandom2008}. To tell whether the right number of 'b's follow the 'a's, the automaton must somehow keep track of how many 'a's have been processed. But since an FSA has only finitely many states, it will eventually run out of states and lose count.

This limitation is not technical---it reflects something deep about what the finite can represent. The history of interaction is encoded in which state the machine occupies, but that history is not represented as something that could be \emph{for} the machine. The machine doesn't ``know'' it's in state \(S_3\) or ``remember'' passing through \(S_1\) and \(S_2\). It simply \emph{is} in a state, executing transitions mechanically.

To model practices that require explicit memory---counting, arithmetic with carrying, nested structures---we need machines with richer expressive power. \textbf{Push-down automata} have stacks where symbols can be written (pushed) and read (popped), giving them genuine memory structures. \textbf{Turing machines} have unlimited tape that can be written and read, providing maximal expressive power for computable functions.

The hierarchy of automata expressive power will matter when we develop the Hermeneutic Calculator in later chapters. Different arithmetic strategies children invent require different levels of expressive power to formalize. Simple counting might be regular (FSA-level), but carrying in addition requires stack-like memory (push-down level), and some sophisticated strategies approach the full power of Turing machines.

\subsection*{Action Theory and Recursive Practice}\label{action-theory-and-recursive-practice}

While Brandom's use of theoretical automata contributes to action theory by formalizing practices, we must recognize the cyclical structure of action itself. An act begins with an impetus---a desire or need. As the act unfolds, the actor engages in self-monitoring, reflecting to determine whether it will satisfy the originating need. This is movement from first-person (experiencing the impetus) to second-person (observing oneself act) position.

The actor can terminate activity at any point if the second-person position indicates the act won't satisfy the need. In instrumental or strategic acts, a third-person position assesses whether the act achieved its objective goal. In communicative acts, this third-person position becomes ``we''---the community adjudicating whether the act is appropriate.

This cyclical structure suggests representing automata not as linear sequences but as recursive loops where outputs feed back as inputs, where states include monitoring states, and where termination depends on satisfaction criteria that may themselves evolve during the practice.

The Big Gorilla practice evolved through exactly such recursive cycles. Initial chest-thumping (act) $\rightarrow$ twins' response (monitoring) $\rightarrow$ adjustment of intensity (revision) $\rightarrow$ new equilibrium (satisfaction). When that equilibrium broke down (too-loud thumping), the cycle had to restart at a meta-level through rational discourse.

Formalizing this recursive structure requires moving beyond simple automata to what might be called \textbf{reflective transducers}---machines that monitor their own operation and can revise their transition functions based on feedback. Such machines approach the boundary where mechanical formalization meets genuine agency, where algorithmic elaboration encounters its own limits.

\section{The Creative Act of Recognition}\label{creative-act-of-recognition}

A year after I invented the character, I came down with a terrible cold. We had just moved in together, and my body was adjusting to the communicable diseases circulating through daycares. During dinner one evening, the kids were ornery, and I SLAMMED my fists on the table, shouting ``ENOUGH!'' The kids started crying. They and \(\mathcal{A}\) left to visit her parents. While they were away, I wrote an apology note and left it on the stairs.

Later that evening, \(\mathcal{A}\) came to the guest room where I'd been isolating. She said, ``We read the note, and the girls wanted to tell you that they forgive you and love you.''

Then \(\exists\) came in. She banged her fist on the bed, then started thumping her belly like Baby Gorilla. She said something about how strong I am. I felt she was expressing deep understanding---validating the part of myself that I was loathing, recognizing that I needed to feel in control. I read her as saying ``oh, look at how strong you are, you be Daddy Gorilla and I'll be Baby Gorilla, hoohoohoo.''

In banging my fist on the table, I violated our norms---norms I ascribe to any reader, which makes sharing this story risky. But in thumping her chest in response, \(\exists\) folded that action back into the family's norms. Daddy Gorilla just got a little too wild. Her act of recognition broke my heart. It suggested that my anger as a parent, which I can't fully transcend, wasn't as dangerous as I feared. For me, that did not result in a feeling of permission to be a table-thumper. The wrongness of the act remains. But the existential need for recognition which I so forcefully expressed was met, for a moment, in a way that helped me genuinely learn that table-thumping is not necessary.

\(\exists\)'s gesture was a creative, forgiving, transformative act that exceeded the boundaries of our existing vocabularies and practices. It was genuine recognition that created new possibilities for meaning rather than simply deploying existing ones. This kind of creative responsiveness extends reciprocal grace to everyone in the family---everyone is \emph{becoming}---in ways that make all modeling feel obtuse.

The formal apparatus is not \emph{false}---it correctly describes what can be algorithmically elaborated. But it is \textbf{incomplete}. Any formal system rich enough to express mathematical practice will be unable to capture its own generative ground. The spontaneous creativity of the \{I\} that generates new norms, transforms existing practices, and offers gestures of forgiveness exceeds any rule.

This recognition doesn't invalidate the formal work. Making practices explicit through algorithmic elaboration enables reflection, critique, and systematic improvement. The automata we've developed are pragmatic metavocabularies that increase our expressive power. But they are---as all formal systems must be---\emph{built to break}. Their precision comes at the cost of what they cannot capture. Acknowledging this limitation is not a failure but maturity: recognizing that formalization is a necessary moment in a larger dialectical movement, not a final resting place.

\(\exists\)'s gesture of forgiveness operates at a level that exceeds both the implicit norms of practice and the explicit articulation of those norms through metavocabularies. Her chest-thumping was a creative, spontaneous act that \emph{transformed} the practice by folding my transgression back into the space of playful recognition. This is what Hegel calls the movement of \emph{Vernunft} (Reason) as opposed to \emph{Verstand} (Understanding)---it generates new possibilities through dialectical engagement with contradiction.

This moment shows why the pedagogical relationship cannot be reduced to the transmission of knowledge or the inculcation of practices. Teaching and learning are recognitive relationships where both participants are transformed through their encounter. Both are engaged in a shared practice of mutual formation where each must recognize the other as \emph{in}finite (capable of spontaneous creativity, deserving of respect as an irreplaceable person) and finite (positioned within shared normative structures, accountable to communal standards).

The asymmetry of the teacher-student relationship---where the teacher typically operates at a higher discursive level and has institutional authority---creates the danger that recognition will collapse into domination. When teachers treat students only as instances of the CUSP ``you'' (``this is what students at your level should be able to do''), they negate the students' \emph{in}finite particularity. When teachers refuse to invoke any universal standards (``whatever you think is fine''), they abandon their responsibility to induct students into shared practices with genuine normative content. The pedagogical task is navigating this tension through ongoing communicative negotiation---a negotiation that, like the Big Gorilla practice, will sometimes succeed, sometimes break down, and require repair through rational discourse and gestures of grace that exceed what any theory can prescribe.

\section{Conclusion: The I, the You, and the Unfolding We}\label{conclusion-the-i-the-you-and-the-unfolding-we}

So, who are you? The answer this chapter proposes is that ``You'' are not a static entity but a dynamic and necessary position in communicative action.

Speaking assumes a listener, a ``you'' who might possibly understand. Such assumptions are the bedrock of intersubjectivity, though it is shifting ground, subject to negotiation when the assumption is no longer viable. You are the particular other whose face-to-face encounter grounds my speech in concrete reality. You are also the universal other whose internalized norms I appeal to for justification and shared meaning. And you are the interlocutor with whom the boundary between these two poles can blur, revealing the fluid, negotiated nature of a shared world.

This dyadic structure is precisely what allows for the developmental synthesis described in Chapter 3. The ``Beast of Love'' is tamed not in solitude but in the reciprocal I-You encounter where participants learn to balance the need for \emph{in}finite, authentic expression with the need for finite, normative recognition. This process is not a logical deduction but a lived, recursive, and often challenging practice. It is the practice of building a ``we''---a shared space of meaning, a common history, a recognitive community.

Understanding the different roles the ``you'' can play in communication fosters greater attunement to the ethical dimensions of dialogue. Recognizing when we address a particular person versus when we invoke universal norms helps ensure our speech opens space for the other's response rather than foreclosing it. The result is more careful, responsive communication of the kind that genuine recognition requires.

Most importantly, becoming a self is not a solitary achievement but a collaborative one. The \{I\} emerges only in relation to a ``You,'' and the ``You'' exists only in the context of shared practices and mutual recognition. Selves take shape through the patient, recursive, often difficult work of building understanding with others. In this work, there are no final answers, only an ongoing commitment to meet each other with care, curiosity, and a willingness to be changed by what we create together.

This recursive, constitutive character of the I-You relation gives intersubjectivity a transcendental-like structure, though it differs from classical transcendental arguments. In the Kantian tradition, transcendental conditions are universal and necessary features of any possible experience. I cannot point beyond the assumption of communicative competency, for doing so just assumes communicative competency at another level. What makes intersubjectivity transcendental-like is that while we can tell stories about how it arises, it must be assumed for any of those stories to have force.

The next chapter turns to the limits of thought and language themselves, exploring through personal loss and philosophical inquiry how absence functions as an enabling condition rather than lack. The Voice---the silent ground that makes speech possible---parallels the null representation in mathematics that grounds numerical meaning. Both reveal that what appears as lack or limitation is actually the negative foundation that makes positive content possible.

If my daughters read this someday: you were never research subjects. You were---and are---teachers. This chapter captures one slice of our early interactions, frozen in time and refracted through my theoretical commitments. It does not capture who you have become, nor does it constrain who you will be. The moment $\exists$ thumped her chest and called me back into play exceeded every framework I had developed to understand it. That moment of creative forgiveness adds another layer to concepts that are built to break. 

I close this chapter with the middle verse of \textit{Built to Break}. Mali is serbian for `little,' that happens to sound like $\mathcal{A}$'s nickname.

\begin{verse}
Mali, I said, I've been thinking a lot \\
About all the things that I never forgot\\
Worried to mirrors, sometimes think they're me. \\
Have I ever been more than a memory? \\
Silvered glass falls, as silver tears bloom.\\
Kissing the shards and the salt from the wound. \\

Build something for the breaking\\
Tall thin walls, shivered and quaking.\\
Mali, it's beautiful, breaking with you. \\
Silvery glow, by your light I'm renewed.\\
\end{verse}

The song evokes the distinction between repetition and recollection, marking the split between Kierkegaard and Hegel. Since these words are shared and doomed to be iterated with all their flaws, like ``babbery,'' the song embraces the breakage. Words recollect learning experiences. Their `meaning' or the `intention' of the speaker can be rationally reconstructed as a history. Doing so temporally decompresses those memories into yet more memories. Their use, as above, is described by formal structures that recollect language in a highly abstract way. Rendering Big Gorilla as a tuple of symbols suggests how linguistic self-reference produces what are commonly called `abstractions.' What have I `worried to mirrors'? To worry is to persistently grind down, to recollect so often that the image becomes a source of identity. I had not yet read Derrida's \parencite*[8]{derrida2007psyche} work on the poem \textit{Fable} when I wrote this song, but the image of the mirror I use to refer to my inability to stop ruminating is echoed in the poem Derrida discusses in that text. The themes this verse explores generally resonate with his idea of the paradox of \textit{invention}.

What is left? What remains? $\emptyset$. 

\printbibliography[heading=subbibliography]

\chapter{Limits of thought, deconstruction, and the Voice.}\label{the-voice-and-the-empty-set-absence-as-generative-ground-1}

\begin{abstract}
This chapter confronts the limits of language and recognition through the experience of absolute negativity, occasioned by the author's father's death. The analysis explores how absence functions as a generative foundation rather than a mere lack. Working through Giorgio Agamben's distinction between Voice (the silent ground of language) and voice (meaningful speech), the chapter examines a eulogy, personal letters, and songs as artifacts that reveal the structure of meaning under the condition of absence. Derrida's concept of \emph{différance} is employed to analyze the temporal deferral inherent in communication, where meaning is never fully present but distributed across time. The chapter connects these philosophical explorations to mathematics, proposing that zero and the empty set function as the absent foundations enabling numerical structure, just as the Voice enables language. This establishes a structural homology between the generative role of absence in existential experience and in formal systems, preparing for the mathematical investigations in the book's second half.
\end{abstract}


\section{Reflection on Deferral: Différance and Variables}\label{reflection-on-deferral-diffuxe9rance-and-variables-1}

The journey through the preceding chapters has attuned us to a persistent rhythm---a movement that never quite arrives at a final destination. In the first chapter, we encountered this directly in the body: the awareness that always arrives a moment after the sensation, the `now' that is already past the instant we articulate it. This temporal gap, this structural delay, is not a failure of attention but the enabling condition of experience that recedes when articulated. What we termed the `sound of time' is the oscillation between presence and absence, a vibration where meaning emerges precisely because it is never fully settled.

This same structure reverberates through the domain of meaning explored in the second chapter. There, we recognized that conceptual content is not fixed but inferential---meaning resides in the movement between claims, in the ongoing process of substitution and articulation. A concept's identity is deferred, distributed across the network of inferences it participates in. Just as the `now' slips away when grasped, a concept's meaning escapes any single definition, residing instead in its use across time and context.

Furthermore, this deferral structures our own being, as we explored in the third and fourth chapters. The \{I\} is never fully present to itself; it is always mediated through the recognition of the other. The desire for recognition---to be understood as good and as \emph{in}finite---is a striving that cannot reach a final satisfaction, because the self is inherently dynamic, always exceeding any finite articulation of the ``me.'' This constitutive gap, this restless negativity, is the engine of our development.


Derrida's concept of \emph{différance} \parencite{derrida1982margins} captures the dual movement at work in meaning-making: the simultaneous \emph{deferring} and \emph{differing} that characterizes how language functions. Meaning is always deferred because it depends on absent contexts, on traces of what is not present. Meaning is always differing because each new context alters what the words signify.

This structure of deferral is grounded in what Giorgio Agamben terms the \textbf{Voice}. In \textit{Language and Death}, Agamben argues that language originates in a negative foundation, a silent and unspeakable Voice that precedes and enables all speech\label{def:Voice}. The Voice is the silent, negative ground that must be presupposed for language to emerge. It is a removal of the animalistic cry that allows for the articulation of meaningful speech \parencite[84]{agamben_language_2006}. The Voice---the silent, unsayable ground---is the condition of possibility for this deferral and differing. It is what allows language to mean anything at all, precisely because it is not-speech, not-presence, not-being.

Anaphora achieves reference by recollecting something previously established in discourse. Pronouns like ``it,'' ``she,'' ``he'' don't independently access their referents---they inherit reference from antecedents. The pronoun doesn't name an object; it recollects a previous act of referring.

Mathematical variables function through this same structure of deferral. A variable, \(x\), is not a name for a specific object but a placeholder for a range of possibilities. Its meaning is never fully present in any single instantiation; it is determined by the inferential network in which it participates. When we substitute a value for \(x\), we are momentarily compressing its potential into a determinate form. But the variable itself always exceeds that determination through its iterability. There are as many values for $x$ as there are `it's' to recall. 

This structure echoes the dynamic of the \{I\}. The \{I\} is the \textit{in}finite potential that always exceeds the finite articulation of the ``me.'' The ``me'' is a temporary determination, a substitution instance, while the \{I\} remains the variable ground that enables substitution. A variable, then, is the mathematical articulation of this restless negativity---the capacity to be determined without being exhausted by that determination. It formalizes the gap that enables movement, the deferral that allows for the unfolding of mathematical thought.

Variables in mathematics function similarly---not as names for platonic objects but as anaphoric terms that recollect the ``I think'' accompanying all representation. This will be the central claim of Chapter 7: numerals are pronouns, not names. They achieve stability through recollection rather than through reference to pre-given objects.

If meaning is always deferred and differing, then mathematical stability cannot come from reference to eternal unchanging objects. It must come from the temporal structure of recollection, from the anaphoric carrying-forward of meaning, which I symbolize with the null representation (\(\emptyset\)).

The null representation operates like the Voice: it is the absent ground that enables numerical content to emerge. Just as the Voice is not-speech but the removal that enables speech, the null representation is not-number but the structured absence that enables counting. Both suggest that what appears as lack or limitation is actually the negative foundation that makes positive content possible.

\section{The Absent Referent and Limits of Thought}

The concepts of deferral, the trace, and the absent origin remain abstract until they are recognized in the body. Theory serves experience, but sometimes experience overtakes theory, demanding that we dwell in the rupture before attempting to articulate its structure. The movement of this work now requires such a dwelling. It requires acknowledging the event that forced a confrontation with everything we have been exploring: the death of my father.

Continuing with the theme introduced when discussing the Fabiola Project, which I described as the problem of the absent referent, I want to descend from the abstract layer of art and philosophy, down into the shifting, molten core of grief. 

I was in the middle of writing my dissertation when my father died suddenly. My mom found him and ran to get me from the apartment above their garage. We delivered CPR until the paramedics came, at which time I started playing guitar and singing \textit{Brokedown Palace} by Jerry Garcia and Robert Hunter to him to say goodbye. 

The funeral was a few days later. I had already been grappling with the sound of time metaphor for Hegel's negative. I spent those days writing his eulogy, trying to recall who he was through pictures, letters, songs, and stories. I wanted the event to be beautiful, in that ferocious sense of beauty from the Judeo-Christian tradition where all beauty comes from \st{God}. I offer the text of the eulogy as an entry point to the problem of presence and the limits of thought. It is long, but its unessential aspects are iterated towards relevance in the songs that shift the structure of this text into its circular form. This chapter includes new senses for what remains necessarily implicit than what I have discussed so far. Later, I will compress all those senses into the \textit{null representation}. 

\begin{quote}
Dad, I miss you so much. I want to speak to you as you are, as a spirit making the next part of your journey, but I don't know how. The bright, white light of who you were has been broken by the prism of death into a brilliant rainbow, diffuse, with each of us radiating our own memories of you as discrete shards of your spectrum. I find myself wanting to speak to your wholeness and lose my words. But I can speak to the good people who are here, who knew your light. I was hugging Faye and crying, and it came to me as clear and as hard as glass just how strong your spirit shines through us, how I could speak to you by speaking to the people here who remember you, who grew into the people they are in part because of you, by speaking to the people who loved you.

So, hello, I now would like to address the assembled people and speak to my dad by speaking through and to you all.  My name is Tio Savich, and if you didn't already know or couldn't tell just from looking at me, I am Bet and Rudy's son.  

I am standing here like a woven cloth, wet and wrung out with grief. There are two parts that make up a weave, the warp and the weft. The warp is the thread that stays fixed to the loom, they're the solid colors that run through the whole cloth, and the weft are the strings you push through with the shuttle. With the weft, you might only have a couple of passes with one color before you move on to the next, but the warp runs throughout. I feel like I am looking out today on the weft in my life, the strings that come together to make the intricate designs of my life, but the life we are celebrating and mourning today is the warp. It's hard to tease out exactly what the warp does to a design because it is underlying and constant. In music, the notes only make sense when accompanied by the constant passage of time.  

I thought of this metaphor because I was reading about what makes a good eulogy, and it said to focus on fond memories and stories, but every memory I think of, when I think of Dad, is tied to every other memory. Dad runs through my life as the solid threads on which the strange designs of my life are drawn. I linger over the weft, choosing each thread as carefully as I can to make my life whole and beautiful, but Dad was never a choice, he was always just there for me. Always solid, always present, always responsive, and always loving. It feels like if I just choose one memory of Dad, my whole life will unravel, because Dad is a part of all of it.  

But I will pick two memories, one recent and one distant that bookend my time with Dad. 

This past summer, Dad and I went on our yearly pilgrimage to watch the Chicago White Sox play baseball. We'd go up every year, and take those few hours in the car to chat about where my life was heading and his most interesting cases. In that way where it is often easier to open up when facing the windshield, these drives were often the best times to connect with Dad about life, either gearing up for fun on the way there or sun-beaten and full of lamb and coffee from the Greek Islands restaurant on the way home. I have a bad habit of keeping my strongest worries to myself, but on those rides it was easy to open up. 

This time, it was Dad who opened up about his worries. On the way home, right around the Children's Museum exit on I-65, Dad was the DJ and he put on a song ``Muleskinner Blues''. We listened to it, and he told me a story about the best version he ever heard, which was when Kent Todd sang the high tenor with the Not Too Bad Bluegrass Band at the Office Lounge. 

I always liked how Dad would anoint the people, places, and songs he directly experienced as ``the best'', rather than putting a whole lot of stock in experts' opinions about what was important or meaningful. There are a thousand versions of that standard, but he liked the one he heard one time at the Office Lounge best. This attitude, where the most important place in the world is your own backyard, makes the cold indifference of the wider world fall away.

It always made me feel like the songs that I made up and sang were important, even though those songs might only have an audience of a bare handful. It also fueled the way he tried to protect the community of Bloomington from environmental degradation and the individual lives of his clients. The last case he won was fought over years and years and the result was a tiny payment the State had been withholding from an old woman that allowed her to keep her modest home. He was as proud of that achievement as he was of any other case he ever talked to me about. He had an interest in the abstract and wider world, but he cared deeply about the concrete world he lived and experienced. 

So I was driving and Dad was talking about how he went to the Bean Blossom Bluegrass Festival that past year and Kent Todd was taking requests. Dad requested the Muleskinner Blues, and was somewhat taken aback when Kent Todd said simply ``I don't play that one anymore.'' Dad went up after the set and Kent explained how his dad had died that year and the Muleskinner Blues had been his favorite song. It became just too painful to play.

As we listened to the song, my dad started talking about his friend, Terry, who I hope is here today. Terry had taken Dad to his family's ancestral spring where they drank whiskey and cool water, which reminded him of my grandma and grandpa. Terry had just been diagnosed with a terrible cancer, and my dad, who choked up easily but rarely cried, boiled over. He cried so hard the car shook.

There'd been thunder in the clouds down from Lafayette but no rain, and just then it started coming down in buckets. It dried up as he did, and by the time we got to my car which was parked on the east side of Indy, the second strongest rainbow I'd ever seen spread out across the sky. And I was just so happy to have spent the day with Dad. 
 
In more distant memory, I remember riding on Dad's back at Bryan Park pool, where he would plunge down into the water, and I would be submerged, and then he would lunge up to the surface, and I would hang on as we breached. Clinging on to dad, breaking through the surface on a hot day at the pool, in shimmering blue water under hard summer sun light. Riding around on Dad's back as he played the horse to my rider at home on University Street. Getting carried from the old maroon minivan to the house after a long car trip. Putting the ball on the tee and then stepping behind me with his hands on the bat, lending his strength to my first swings. The feeling of his hand on my back the first time the training wheels came off of my bicycle, when I couldn't quite get up enough speed for the bike to stand up, that firm hand and swift push, so that the bike would stand long enough that I could rest my panic and start to pedal. Later, chasing his draft on long bike rides, when I was so tired I could barely stay awake. Dad lent me his strength, and made my own memories of strength and accomplishment possible, he let me ride.

\end{quote}


The attempt to make the absent fully present fails. 

The wisdom in Agamben's account, for me, lies in how the event of language and the event of death are structured by the negative -- by absence. Agamben articulates a distinction between \textit{Voice} and \textit{voice}. Voice is the silent, unsayable ground that makes language possible. The voice is the sounds made by the body. The Voice and text represent two opposing poles. The Voice is the silent, temporal event of language taking place, while the text is the spatial trace of language having taken place. 

The eulogy marks the transition from a desire to speak to my dad to speaking to him through others. Accepting absence freed me to find what I was after `all around,' so to speak.


The ``diffuse rainbow'' metaphor reveals how wholeness persists precisely through its fragmentation across multiple people and memories. This connects directly to the intersubjective recognition explored in Chapters 3-4: we come to know ourselves and others not through direct self-presence but through the mediated process of being recognized by others who hold different fragments of our identity.

The transition from trying to speak directly to my father to speaking to him through the assembled community demonstrates the same structure that will appear in mathematical communication: meaning emerges not from immediate access to mathematical objects but through the shared practices and recognition structures that constitute mathematical discourse. The acceptance of absence as enabling rather than defeating understanding provides the experiential foundation for the null representation ($\emptyset$) that will be explored in Chapter 7. 

\subsection*{Relating the Negative Foundations of Language to Critical Social Theory}\label{sec:negative-foundations}
Trying to know and speak in the face of absence can be understood through the framework of early Habermas's \textit{knowledge-constitutive interests} \parencite*{Habermas:1971aa}. Habermas' three \textit{knowledge-constitutive interests} as articulated in \textit{Knowledge and Human Interests} (KHI) \parencite*{Habermas:1971aa}. The notion that some desire, need, or interest precedes each speech act allows for these theories to find a common negative foundation. In table \ref{tab:scientism}, I reprint Carspecken's \parencite*[52]{Carspecken2009aa} reconstruction of the knowledge constitutive interests discussed in KHI \parencite{Habermas:1971aa}.

The interests are the empirical-analytic (EA; concerned with technical control), the historical-hermeneutic (HH; concerned with understanding meaning in context), and the critical-emancipatory (CE; concerned with human freedom and self-reflection). The interests form fundamental limits of knowledge, and so I must work within them. 

% Simplified column width specs (original auto-generated arithmetic removed)
\begin{longtable}[]{@{}
  >{\raggedright\arraybackslash}p{0.243\linewidth}
  >{\raggedright\arraybackslash}p{0.224\linewidth}
  >{\raggedright\arraybackslash}p{0.271\linewidth}
  >{\raggedright\arraybackslash}p{0.262\linewidth}@{}}
\toprule\noalign{}
\begin{minipage}[b]{\linewidth}\raggedright
\textbf{KHI Type of Knowledge}
\end{minipage} & \begin{minipage}[b]{\linewidth}\raggedright
\textbf{Empirical-Analytic}
\end{minipage} & \begin{minipage}[b]{\linewidth}\raggedright
\textbf{Historical-Hermeneutic}
\end{minipage} & \begin{minipage}[b]{\linewidth}\raggedright
\textbf{Critical-Emancipatory}
\end{minipage} \\
\midrule\noalign{}
\endhead
\bottomrule\noalign{}
\endlastfoot
\textbf{Motivations} & For having or avoiding tangible, objective states
of affairs. For explaining what ``is'' in a world of objects and events.
& For reaching agreements to coordinate actions. For understanding why
others do as they do, what they experience, how they think. For being
understood to meet basic recognition needs. & For understanding one's
self. For becoming more of what one can become. For freedom and
autonomy. \\
\textbf{Actions} & Instrumental & Communicative & Reflection for its own
sake, and/or desire for desire, and/or nonaction. \\
\textbf{Forms of knowledge and language} & Copy knowledge, Models.
Formalized languages with all rules and meaning of terms explicit.
Meaning of terms reduced to measurement procedures. & Hermeneutically
achieved explications of what was formerly implicit knowledge. Ordinary
language always dependent on implicitly understood and changing rules of
use and meaning. & Internal recognition. Necessarily noncommunicable
states and experiences. Language use necessarily contradicts
knowledge. \\
\textbf{Relation of knowledge to being} & Knowledge and being are
separated and distinct. Knowledge is of reality. New knowledge does not
change reality. Relations in reality are causal/correlational. &
Knowledge is of knowledge and is self-transcending. New knowledge
changes reality. Relations between people are illocutionarily maintained
but usually in distortion via power. & Knowledge is fused with being.
Being as knowledge of itself. It makes no sense to speak of relations
within reality. \\
\textbf{Constitution of the ``Knower''} & The self as anonymous
universal observer and instrumental actor. Self transcends object domain
as Kant's I. & Individuated self with unique autobiography. Self
transcends itself as author and critic of one's self-narratives. & Self
is pure reflection. Knower is its process of knowing itself. Self is the
loved Other of God, and/or there is no self. \\
\textbf{Relation to reflection} & Reflection is restricted such that the
form of inquiry cannot study itself in its own framework. & Reflection
is internal to hermeneutics such that this form of inquiry can consider
itself. & Reflection is the method and object of study
simultaneously. \\
\textbf{Validity} & Based on the degree to which predictions are
successful. & Based on the degree to which insiders recognize explicit
articulations as something already known in some way. & Continuation of
a self-formative process, and/or noncommunicable state of
enlightenment. \\
\caption{\textit{Note. }Carspecken's \parencite*[52]{Carspecken2009aa} Reconstruction of the Knowledge Constitutive Interests discussed in KHI \parencite{Habermas:1971aa}}
\label{tab:scientism}
\end{longtable}

My attempt to `solve' the problem of presence by speaking to my dad by ``speaking through and to you all'' was an intuitive move toward a Hegelian recognition of the universal in the particular. But it also reveals the interplay of these three interests, the fundamental limits within which any such speech act must operate. 
\begin{itemize}
\item \textbf{The Empirical-Analytic Interest:} This interest stems from the desire to obtain technical control over the objective domain. As I tried to breathe into my father's body after it was found, the 911 operator who coached me through CPR described techniques accrued under this interest. I let go of this desire when I picked up my guitar to sing him goodbye. 
\item \textbf{The Historical-Hermeneutic Interest:} This is the practical interest in mutual understanding and the interpretation of meaning within a shared tradition. The funeral service, and the desire to understand the meaning of my father's life to make his ``warp'' coherent with the ``weft'' of the gathered community, are part of this interest.  
\item \textbf{The Critical-Emancipatory Interest:} This is the interest in self-reflection and liberation. The eulogy, as a performative act, included an emancipatory layer. I wanted to let go of my dad, to let him be as he is now, free of the distinction between the finite and \textit{in}finite. 
\end{itemize}

It is through these interests that the structure of language reveals its own ground. The eulogy is grounded in the silent, absent foundation that makes the spoken word possible.

\section{The Masterpiece and the Paradox of Recognition}

Having established the problem of representing an absent other, I now turn to the even more fraught dynamic of recognition between two living subjects. This requires us to move from the grief of a eulogy to the complex exchange of love and expectation in a letter my dad wrote to me, Theodore. My family had a tradition of naming the first son either Theodore or Rudolph, alternating between the two, dating back to the `old country' where it was Todor and Radavoy.

My dad wrote letters to friends and family and he called them excerpts from The Book of Love. They were often written in response to a specific event or crisis. In my dissertation, I wrote more about the events that led to the depressed condition my dad speaks to in the letter. I will not relate those details here as the story has become stale as the shame and grief faded. I print the letter in full, minus the lyrics to \textit{When I Paint My Masterpiece} by Bob Dylan, which my dad prefaced his letter with due to copyright restrictions. 

\begin{verse} 
    An excerpt from The Book of Love

Lyrics for \textit{When I Paint My Masterpiece}\ldots

Your letter brought this song to mind, like you are feeling unfulfilled and needing, waiting to paint your masterpiece. You've had some ideas about what that masterpiece was going to look like that haven't worked out. Disappointment is setting in. I want to tell you about my masterpiece. I've never shared it with you. You remember my dad. Your grandpa. You remember him when he was very old and losing his mind. He was out of touch with reality. But maybe you also remember him at least a little bit before he got that way. Maybe he let you drive when you were staying with him in Canada. He took you and $\mathcal{R}$ on the train to New Mexico. He was a strong, proud man. He wasn't scared to go up to anyone and tell them how it was (in his opinion). In that bastion of conservative Republican politics of Rensselaer, I can remember going to the barber shop on the courthouse square with him to get a haircut and he would be spouting his communist farm policy about world hunger and how evil Republicans were or whatever was on his mind. But people genuinely liked him. They didn't care that much about his politics really. They wouldn't vote for him because they didn't agree with his politics, but they did like him as a person because he was outgoing and friendly. He had a great smile and a firm handshake. Anyway, I loved him to death as an adult. As a kid, I was scared to death of him. But it was kind of like fear of the lord. Fear in the sense of great respect. Not fear for my physical safety. He was always full of stories about his athletic accomplishments, his academic accomplishments, his agricultural accomplishments, how hard his life had been, etc. I didn't feel like he was bragging. I felt like he was just holding himself up as an example. He was saying this is what I've done and I want you to accomplish things like this too. So I worked really hard in school and worked really hard on the swimming team, just trying to live up to his expectations. Of course, you go to college and do well if you're a Savich. I also worked during college because he did. We had the farm on Boltinghouse Road and raised pigs on pasture the way he did. We had a big garden the way he did. We raised chickens the way he did back in World War II. I made homebrew like he did. Just about everything I did you could say was trying to imitate my dad. But what was the really big thing I thought I had to accomplish in my life to satisfy my father? I had to get married, have a son and name him Theodore. I was so shy around girls. It was so hard. If I ever got one interested in me, I wanted to marry her immediately and have kids and get it over with. That's why I wanted to marry \ldots in 8th grade. I wanted to marry \ldots in high school. I didn't want to have to date around. \ldots But you, Theodore Michael Dougherty Savich, you are my masterpiece. Sure, I kept on always trying to please Grandpa after you were born and I still do. But you have made my life a success. Being a male child named Theodore is a big part of that success, but you are so much more. You are smart. You're creative. You are good looking. More than that, you are handsome. You are funny.
You are great guitar player and singer. You are thinker and an innovator. You are a great teacher because you care about getting through to your students. And, you have given my life a meaning and a purpose. Thank you. Thank you. Thank you for being my masterpiece. I'm so proud of you.\\
All my love,\\
Dad
\end{verse}

This letter is an act of Hegelian recognition. My father attempts to constitute my ``me'' by naming my attributes -- ``smart,'' ``creative,'' ``handsome,'' ``a great teacher.'' He declares me ``My Masterpiece,'' asserting a unity between us and affirming my being. That unity is dialectical, and can be partially understood through a form of \textbf{reciprocal sense-dependence} \parencite{Brandom:2019aa}\label{def:reciprocal_sense}. Whereas Brandom notes the \textit{a}symmetry in essentiality (my identity depends on my parents, but their identity does not necessarily depend on me), my dad's letter establishes a reciprocity in essentiality. For my dad, to talk about him was to talk about me. He filled, in a concrete way, the role of the other in my self-consciousness in a way that fulfilled the existential need to be recognized as good. 

Yet, my reception of it reveals the inherent limit of this process. I did not ``hear'' his words until after he was dead, 5 years after he sent it. It was among the letters and photos I went through to write the eulogy. This delayed recognition reveals a fundamental tension in the structure of recognition itself. While my father's letter attempted to constitute my identity through loving descriptions, the attempt to capture the living subject in a set of predicates encounters the resistance of what cannot be objectified. This illustrates the essential structure explored throughout this book: the \{I\} that thinks and feels always exceeds any finite description.

My father's letter failed to land immediately because the meaning it attempted to convey was \emph{deferred}. The recognition he offered required a context I did not yet have---the context of his death, the context of my own maturation, the context of understanding what it means to be someone's masterpiece. The letter carried the \emph{trace} of the absent referent---the \{I\} that cannot be fully captured by predicates, the \emph{in}finite that resists being made finite.

The temporal gap between the letter's arrival and my ability to receive its meaning demonstrates how recognition operates through difference and deferral rather than immediate presence. This is the structure of \emph{différance}: meaning is never fully present in the moment of communication but emerges through the play of differences across time. The resistance I felt to being identified with my father's predicates was not just ingratitude but the voice of the \textit{in}finite subject that cannot be exhausted by any finite list of properties. The letter's ultimate significance emerged only when I could appreciate both its loving intention and my legitimate resistance to being totalized by it. I \textit{still} find myself tensing up at kindness. 

\subsection*{Love's Memory}

The song \textit{Love's Memory} is an attempt to sublate my dad's letter -- to iterate ``my masterpiece'' in a more universalized form. To grasp the other horn of \textit{différance}, to `differ,' I repeat the letter in another form through song. But the gap remains. The song is inadequate without the letter; the letter is incomplete without the lived relationship. The two cannot be experienced simultaneously. 

\textit{Love's Memory}\footnote{A recording is available here:
 https://youtu.be/qIA8kSZY-Jw}
\begin{verse}
When the first taste of autumn stopped and stared,\\
in the middle of the street and its wet-steel eyes tracked you everywhere\\
Started plucking out your habits like I plucked my first grey hair\\
A dream won't come, it won't come - then it's there\\
I faked a smile to hide the child, but I'll admit it, I was scared \\

(Yodeling followed by whistling are repeated as the chorus of the song, separating each verse)

The hills are clattered with dead oceans, white-capped waves sawed from beneath\\
I lope around down by the quarries, white salt veil dripping stream\\
A memory of long dried seas, lives in the hills all around me\\
You dove into the pit, blaring body in the sun\\
I picked my way down giant rungs\\

Damned polar winter, spent two years shoveling\\
She pressed the door, I kept the pillows, abrupt, not entirely unforeseen\\
Freeze slow below blankets of snow, I nearly fell asleep\\
You wrote me a note, I didn't hear what you wrote\\
until after everything\\

It said: I want you to be happy, or whatever you need\\
When I'm not there to hold the mirror remember what I see:\\
I see my masterpiece, I see Love's memory\\
It's enough that you are, you're beautiful between,\\
Let'em hear you sing

I hiked out to the cold spring in the August heat\\
Chiggers gnawing at my ankles, I finally found the ravine\\
Moss had grown on every stone, the water trickled cold and clean\\
Let's crack the jar on the floor, pour some light on the boards\\
And sing this round with me (It doesn't matter what key) \\
\end{verse}

I wrote the first few lines and a different version of the song while he was still alive and had a scare with cancer. I only re-read the letter after he died, when I was searching to try to remember who my dad was while I wrote his eulogy. While the first version of the song was confused, the first few lines were reusable. I felt that it would make the song more universal if I left in the predatory ``wet-steel eyes'' even while the actual event of my dad's death was sudden. 

The second verse shifts to geological time: ``The hills are clattered with dead oceans, white-capped waves sawed from beneath.'' These are the limestone hills of southern Indiana, formed from ancient seas. I wanted the imagery to connect to the specific place in which our time together unfolded. ``A memory of long dried seas, lives in the hills all around me'' suggests an absence that persists as geological trace, as structural presence through absence.



The third verse brings us to ``Damned polar winter, spent two years shoveling''---the extended period of grief. The breakup mentioned (``She pressed the door, I kept the pillows, abrupt, not entirely unforeseen'') situates personal loss within a larger context of dissolution. ``Freeze slow below blankets of snow, I nearly fell asleep''---the temptation to stop struggling.

Then the lines: ``You wrote me a note, I didn't hear what you wrote / until after everything.'' This is the delayed recognition structure we've been discussing. Not ``I didn't read'' but ``I didn't \emph{hear}''---the meaning was deferred, unavailable until context shifted. ``After everything'' means after death, after the polar winter, after the structure of the relationship changed completely.

The fourth verse claims what the note said, but you have the text of the note above. It doesn't say ``I want you to be happy, or whatever you need / When I'm not there to hold the mirror remember what I see.'' Nothing in the syntax suggests such a reading. I guess that's poetic license, but I wanted to print both the letter and the song to give a sense for how meaning shifts through iteration. 

``I see my masterpiece, I see Love's memory.'' I can discuss intersubjectivity as a pragmatic a priori---a necessary assumption that enables communicative action. I can discuss the trace/arche-trace as an enabling condition that recedes. I can talk about the I-feeling. I can talk about how every word in this book, from ``babbery'' on down is a temporal compression of learning experiences and so always partial, flawed, and incomplete. ``Love's memory'' gives possession of those memories to Love. I fold `love's memory' into the null representation.

``It's enough that you are, you're beautiful between, let `em hear you sing.'' An underdeveloped theme of this work is \textit{existential entitlement.} 

The final verse offers transformation through community: ``Let's crack the jar on the floor, pour some light on the boards / And sing this round with me (It doesn't matter what key).'' The jar---water from a cold spring---must be broken to release what it holds. When I sing the song, very few people take the invitation to sing along. That's okay. But I wanted to be make the invitation explicit: it doesn't matter what key. 

 But as I sing, that need for recognition moves from fear through particularity toward a more universalized enactment of a need for recognition that meets itself in some shades of harmony. I enjoy communal singing. 

Yodeling and whistling express what words cannot. I chose them purposefully to enact a pure movement in vowel sounds, then a further undifferentiation to whistling. I like vowel cascades, like `stick, stack, stuck.' But in this song, I wanted to move away from consonant articulation, to try to negate the determinacy of the first negation. In ancient Hebrew, the vowels go unwritten. For the song, I wanted the melody and chords to carry the structure, while the vowels moved. 

Each return to yodeling/whistling is an iteration that accumulates meaning without adding semantic content.

\section{The Song of Home (Nostos)}

The concept of \emph{nostos}---return home, homecoming---provides the integrative thread connecting this chapter's various explorations. But what is ``home'' in this context? It is not a place, not a physical dwelling, not even a set of memories about such dwelling. Home is the recognition of the Voice as ground.

The song \textit{Still Feels Like Home} is an attempt to express this homecoming to the unrepresentable ground. Agamben writes, ``The Voice does not will any proposition or event; it wills \textit{that language exist}, it wills the \textit{originary event} that contains the possibility of every event'' \parencite*[87]{agamben_language_2006}. It is the silent potentiality before any actual word. The unrepresentable absence becomes the origin -- the ``home'' -- from which all language emerges. 
\begin{quote}
Philosophy is this voyage, the human word's \textit{nostos} (return) from itself to itself, which, abandoning its own habitual dwelling place in the voice, opens itself to the terror of nothingness and, at the same time, to the marvel of being; and after becoming meaningful discourse, it returns in the end, as \textit{absolute} wisdom, to the Voice. Only in this way can thought finally be at home and ``absolved'' of the scission that threatened it from there where it always already was. Only in the Absolute can the word, which experienced ''homesickness'' (\textit{Heimweh}) and the ''pain of return'' (\textit{nost-algia}), which experienced the negative always already reigning in its habitual dwelling place, now truly reach its own beginning in the Voice. \parencite[93]{agamben_language_2006}
\end{quote}
The song \textit{Still Feels Like Home} is an attempt to express this homecoming to the unrepresentable. I wrote it just a few months after dad died. When I perform the song, I tell a story. We had to sell the house that my mom, dad, and I lived in. As we cleaned it out, my sister and her husband discovered mold in the basement. After tearing out all the carpet, the basement was stripped bare. The concrete floors made a resonant chamber, and I would go back periodically to play my guitar and sing at night after real estate showings in the day. As with the previous song, the inarticulate vowel sounds (yodeling and whistling above; hoooooome, sliiiiiide, and so on in this one) are important to the song. I was trying to make the sound of the Voice, the unsayable, as big as possible, to fill the empty room. 

\textit{Still feels like home}
\begin{verse}
I cut my hand on a carpet knife\\
Tearing out the mold-dark night\\
Nothing goes quite as wrong\\
As when you truly try\\
After years and years and years of just getting by\\

\textit{Chorus}
Letting the summer wind blow the rain in\\
And the curtains drip on the floor\\
With every storm a little more was breaking\\
Now I can't afford replacing the warped boards: the house is getting old\\

I broke my crown soaking up storm light\\
Held my breath to arrest the slide\\
Julie wrote it down, said: ``don't forget,\\
Dead head to keep on growing,''\\
Or learn to stomach your regrets\\
Cause the summer wind's gonna blow the rain in\\
`Til the curtains drip on the floor\\
With every storm you'll find a little more is breaking\\
`Til you can't afford replacing the warped boards and curse getting old\\

\textit{Bridge}\\

But other words blow in on the same warm wind\\
Filling cracks with honey dripping from the comb\\
It still feels like home\\
Can't help it, it feels like home\\

I cut my hand on a carpet knife\\
Held my breath to arrest the slide\\
It feels so wrong in the mold-dark night\\
`Til you peeled those sodden curtains from my eyes\\
Awoke to golden light\\
Tell the family I'm doing alright\\
I cut my hand
\end{verse}

The story I tell when I perform the song shades the truth. Like \textit{Love's Memory}, I wrote the first verse, found it wanting, and put the song away for a few years. It had not yet found its second negation and was too dark to feel good singing. While I rarely tell the full truth, one funny bit about how the song came to its conclusion is that in the period between when I wrote the first, depressing verses, but before my dad died, I \textit{totally} ruined the song: ``I put some ham on some toast last night, left it out, now it don't smell right\ldots'' The turn seems like an aside, but I bring it up to illustrate how many use humor to break tension. The song was too dark, and I needed to laugh at it to distance myself from the depression that the first few verses recall. After dad died, I returned to the song, reworked it, and found the bridge that transformed it.

What does it mean to say that it ``still feels like home''? It is hard to discern without first hearing the song. It begins ploddingly in the key of D\#. When it reaches ``But other words blow in, on the same warm wind'' (the Voice) it modulates down to the key of C\# as the melody ascends to the top part of my vocal range. When I sing it, the duration of notes requires disciplined breath (\textit{pranayama}). Reaching the high notes requires openning up my chest and standing tall. The `home' I can always return to is the feeling-body, explored in chapter 1. I feel that sinking and simultaneous lifting that feels like a return to my body. 

When I sing it, the I-feeling fuses with the song. The song's structure is dialectical, moving from the particular (the mold-dark night, the carpet knife) to the universal (the Voice that fills the cracks). The home is not about real estate or biological embodiement, but the feeling of returning home to the I-feeling. Recognizing the event of language as home lends Descartes' \textit{sum}, the linguistically non-falsifiable ``I am,'' the feeling of a just-was self-certainty. But my embodied practice of singing would be sorely incomplete without mentioning my friends in the band who lend their harmony and instrumentation to the act. The experience of hearing others sing `my words'---validating them through harmony---is deeply moving. It may not be `about' anything besides recognition.  

The song's imagery maps cleanly onto my philosophical purposes. The ``carpet knife'' and the struggle with the ``mold-dark night'' represent the painful, flawed act of trying to create order and meaning. The dialectical turn comes with the bridge: ``But other words blow in on the same warm wind\ldots like honey dripping from the comb: It still feels like home.'' These ``other words'' are not spoken in the song; they are the silent Voice. This feeling of ``home'' references the negative foundation Agamben describes -- the silence, the absence -- that paradoxically grounds existence and speech.






\section{Conclusion: To the Bridge}\label{conclusion-to-the-bridge}

The silence encountered in grief is generative absence. 



My father's letter brought me face to face with the paradox of recognition: the attempt to capture the \emph{in}finite \{I\} through finite predicates necessarily fails, not because the predicates are false but because they are insufficient. The gap between when he wrote the letter and when I ``heard'' it is not a failure of communication but the structure of meaning itself---différance as deferral and differing, meaning never fully present but always deferred across time and differing across contexts. 

The songs attempt to articulate the Voice through musical structure rather than propositional content. The harmonic descent enabling melodic ascent in ``Still Feels Like Home'' enacts the structure of absence enabling presence, removal creating space for articulation.

Absence is not a deficiency. The Voice is not a lack that needs to be filled.

This structure underpins the foundations of mathematics. As we move now to the Bridge, we carry this understanding with us. The transition from the voice to the formal structures of diagonalization and set theory continues the same inquiry into how meaning emerges from absence.

The silence that grounds language is homologous to the zero that grounds the number system. The movement we have traced---the rhythm of negation and generation, the impossibility of grasping a totality without generating an excess---will now manifest in its most formal articulation: the diagonal argument. 

The bridge we are building connects the existential need to be recognized as both finite and \emph{in}finite with the formal structures that mathematics uses to express similar tensions. The next chapter will engage with how Cantor's diagonalization forces mathematics to confront its own ground, to acknowledge that completeness is impossible precisely because the ground cannot be made present. But this impossibility is not defeat---it is kind of freedom. It is the recognition that mathematics, like consciousness, is built to break, and that breaking is not failure but the condition of growth.

While there is more to say about how critical theory informs mathematics, this chapter concludes the second part of the book. Next, to the bridge. 

\printbibliography[heading=subbibliography]


\part{Bridge}

% Store the current chapter name definition (usually "Chapter") and counter display format
\let\oldchaptername\chaptername
\let\oldthechapter\thechapter

% Redefine the chapter name and the way the counter is displayed for this specific chapter
\renewcommand{\chaptername}{Chapter}
\renewcommand{\thechapter}{B} % Set the "number" to 'B'

% Use the regular \chapter command, which will now use the custom format
\chapter{A Foundational Star} 
% ... content for Chapter B

%\chapter*{A Foundational Star}
%\addcontentsline{toc}{chapter}{A Foundational Star}


\begin{verse}
\textit{Bridge:}\\
I need a bridge\\
A crutch to clutch through the worst of it.\\
I'm so glad you called, I was really in a fit.\\
A hit of warm to curb the loneliness.\\
``Better than I hoped, worse than I wished.''
\end{verse}


\begin{abstract}
This chapter serves as the text's transition, exploring parallels between Dr. Seuss's \emph{The Sneetches} and Georg Cantor's diagonal proof. It examines the star as a symbol that functions as a social marker of identity and a mathematical mark of totality and incompleteness. The analysis demonstrates how the argument's structure embodies the book's core theme of productive negation. By indicating how the attempt to create a complete, exclusionary system (the Star-Belly Sneetches) necessarily generates the conditions for its own transcendence, the chapter bridges the philosophical explorations of recognition from the first half of the book with the formal mathematical investigations of the second.
\end{abstract}



\section{Introduction: Pattern and Method}

Dr. Seuss' [Theodor Geisel] \parencite*{seuss_sneetches_1961} book \textit{The Sneetches} begins with the statement of a problem:

\begin{quote}
Now, the Star-Belly Sneetches \\
Had bellies with stars.\\
The Plain-Belly Sneetches \\
Had none upon thars.\\[1ex]
Those stars weren't so big. They were really so small\\
You might think such a thing wouldn't matter at all.\\
But, because they had stars, all the Star-Belly Sneetches would brag,\\
``We're the best kind of Sneetch on the beaches.''\\
With their snoots in the air, they would sniff and they'd snort \\
``We'll have nothing to do with the Plain-Belly sort!''\\
\end{quote}
As the story unfolds, Sylvester McMonkey McBean arrives with a contraption that can put a star on the Plain-Belly Sneetches. This disrupts the class system, so McBean offers to remove the stars from the original Star-Belly Sneetches. Chaos ensues as McBean takes the Sneetches' money, adding and removing stars.
\begin{quote}
They kept paying money. They kept running through\\
Until neither the Plain nor the Star-Bellies knew\\
Whether this one was that one\ldots or that one was this one\\
Or which one was what one\ldots or what one was who.\\
\end{quote}
After McBean gets all of their money, he drives off laughing about how ``you can't teach a Sneetch!''
\begin{quote}
But McBean was quite wrong. I'm quite happy to say\\
The Sneetches got really quite smart on that day,\\
The day they decided that Sneetches are Sneetches\\
And no kind of Sneetch is the best on the beaches. \\
That day, all the Sneetches forgot about stars\\
And whether they had one, or not, upon thars.
\end{quote}


\begin{figure}[h!]
\centering
\includegraphics[width=\textwidth]{/Users/tio/Documents/GitHub/September_UMEDCA/images/sneechtriptyche.png}
\caption{\textit{Note. }Stars form a demarcation between the two classes in the Sneetch society. McBean profits from their confusion and rides away assuming they would never change. But the experience provides a strong enough lesson for Sneetch society to evolve.}
\label{fig:sneechtriptyche}
\end{figure}


In mathematics, which deals with patterns, one important form of critique is \emph{diagonalization}. A simple example illustrates the structure. Let two symbols, a star (\(\star\)) and the ghost of its erasure (\(\square\)), be combined into patterns (Figure \ref{fig:diagonalization_pattern}). Suppose each row represents a kind of norm, rule, or pattern. When we critique norms, we often assert a new norm. This new norm is not identical to any of the previous norms, but it is recognizable as a norm within the same family.

 \begin{figure}[h!]
     \centering
     \includegraphics[width=0.8\textwidth]{/Users/tio/Documents/GitHub/September_UMEDCA/images/prelude_stars_boxes.pdf}
     \caption{\textit{Note.} The figure illustrates how diagonalization creates a new pattern not captured by the original list.}
     \label{fig:diagonalization_pattern}
 \end{figure}

Consider a method that allows one to derive a new pattern from any list of patterns. If one examines the diagonal of such a list and systematically transforms each element (erasing each \(\star\) and filling in each \(\square\)), a new pattern emerges that is not captured by the list. This method works for any list of patterns. A new norm, pattern, or rule is asserted.

This bridge articulates a structural resonance between Hegelian sublation and Georg Cantor's \parencite*{cantor1891ueber} diagonal proof. In the allegory of the Sneetches, we begin with the star ($\star$) as a symbol with two mutually exclusive moments: presence and absence. Through the chaos that ensues from McBean's machine, the rigid mutual exclusivity of star and no-star is transcended. The Star is sublated.

Before examining how mathematical concepts emerge and develop, the methodological pattern that makes such emergence intelligible requires articulation. This bridge stands outside the main narrative as scaffolding, a bridge to the substantive mathematical investigations that follow.

The pattern traced here characterizes sufficiently self-reflective systems that attempt to totalize themselves. This pattern occurs in a children's story, in formal proof, in social recognition, and eventually in the development of mathematical concepts. By establishing the pattern here explicitly, a shared vocabulary emerges for recognizing its appearances.

The pattern has four moments:

\begin{enumerate}
\item \textbf{Totality}: A domain or system is conceived as complete.
\item \textbf{Reification}: The totality is treated as an object with enumerable properties or elements.
\item \textbf{Diagonal Recollection}: The system reflects into itself, examining its elements according to a self-referential structure.
\item \textbf{Transcendence}: This reflection necessarily generates an element that should belong to the totality but demonstrably cannot be enumerated within the reified system.
\end{enumerate}

This is not a mechanical procedure, though instances of mathematical \textit{diagonalization} are algorithmic. Recognizing the pattern helps trace connections across disparate domains and understand why mathematical systems repeatedly encounter---and productively transcend---their own boundaries.

\section{The Sneetches: Allegory of Self-Reflection}

In the Sneetches, the pattern becomes explicit. The totality (star-determined identity) is reified (the star becomes a marker that can be possessed or lacked). The system reflects into itself when those without stars acquire them, demonstrating that this supposedly essential marker can be added mechanically. The reified marker collapses.

McBean's machines keep running. Sneetches cycle through---paying to add stars, then remove them, then add them again. The marker that was supposed to ground identity becomes arbitrary. The system that seemed to exhaust social recognition generates its own impossibility.

Eventually the Sneetches run out of money. McBean leaves, having profited from their confusion.

The transcendence occurs with the recognition that identity cannot be reduced to any reifiable marker. The Sneetches discover that genuine recognition requires acknowledging the other as another subject, not as the possessor or non-possessor of some property. The star---which seemed to exhaust the question of worth---becomes irrelevant once the system reflects on its operation.

This is the complete pattern. The totality (star-determined society) gets reified (the star becomes a purchasable object). The system reflects into itself (Plain-Belly Sneetches acquire stars, disrupting the distinction). This generates an element (recognition without markers) that transcends the original framework. The star remains, but its meaning transforms.

\section{Cantor's Original Formulation}
Cantor's original proof \parencite{cantor1891ueber} aims to demonstrate that there are infinite sets that do not have the same measure as the set of natural numbers. He claims this can be accomplished easily and with generality. Relying on translations provided by Meyer \parencite*{meyer_cantors_nodate} among others, we have:
\begin{quote}
If $m$ and $w$ are any two mutually exclusive characters, then we consider a collection $\mathcal{M}$ of elements, $E = (x_1, x_2, \ldots, x_\nu, \ldots)$, which depend on an infinite number of coordinates, $x_1, x_2, \ldots, x_\nu, \ldots$, where each of these coordinates is either $m$ or $w$.

$\mathcal{M}$ is the totality of all elements $E$.
\end{quote}


\textit{Note}: This establishes $\mathcal{M}$ as the complete set of all possible infinite sequences built from two characters. This is a totality. Cantor will next provide examples of elements in $\mathcal{M}$, reifying the totality as an object composed of these elements.

We use the characters $m$ and $w$ rather than the symbols 0 and 1 (as used in some reconstructions, e.g., \parencite{Gaifman2005}). Using characters emphasizes the pre-numerical nature of the proof. The elements of $\mathcal{M}$ are infinite sequences of these mutually exclusive characters.

\begin{quote}
The elements of $\mathcal{M}$ include, for example, the following three:
\begin{eqnarray}
E^{\text{I}} = (m, m, m, m, \ldots)\\
E^{\text{II}} = (w, w, w, w, \ldots)\\
E^{\text{III}} = (m, w, m, w, \ldots)
\end{eqnarray}

I now claim that such a manifold $\mathcal{M}$ does not have the magnitude of the series $1, 2, 3, \ldots, \nu, \ldots$.
\end{quote}

\textit{Note}: By ``magnitude,'' Cantor means cardinality or size. Meyer \parencite*{meyer_cantors_nodate} notes that the term `cardinality' was not in current usage when Cantor wrote this paper; he used the term 'M\"{a}chtigkeit' (potency, thickness, etc.).


\begin{quote}
This follows from the following sentence:

If $E_1, E_2, \ldots, E_\nu, \ldots$ are any simply infinite series of elements of the manifold $\mathcal{M}$, then there is always an element $E_0$ of $\mathcal{M}$ that does not agree with any $E_\nu$.
\end{quote}

\textit{Note}: This is the heart of Cantor's claim. It is a non-existence claim phrased as a positive construction: no complete list can exist, because for any given list, we can construct an element that is missing from it. Euclid's proof serves as a historical antecedent of this type of claim.

\begin{quote}
To prove it:

\begin{eqnarray}
E_1 = (a_{1,1}, a_{1,2}, \ldots, a_{1,\nu}, \ldots)\\
E_2 = (a_{2,1}, a_{2,2}, \ldots, a_{2,\nu}, \ldots)\\
\ldots\\
E_\mu = (a_{\mu,1}, a_{\mu,2}, \ldots, a_{\mu,\nu}, \ldots)\\
\end{eqnarray}
\end{quote}

\textit{Note}: Here, Cantor sets up an arbitrary enumeration, establishing the conditions for self-reference. The natural numbers ($\mu \in \{1, 2, 3, \ldots\}$) serve as \textit{indexes} or names for the sequences. The index $\mu$ names the sequence $E_\mu$. This involves naming ``higher type entities'' (the infinite sequences) using lower type entities (the numbers) \parencite[2]{Gaifman2005}. Each sequence $E_\mu$ maps a position $\nu$ to a character $a_{\mu,\nu}$ (an $m$ or a $w$).

\begin{quote}
Here the $a_{\mu,\nu}$ are in a certain way $m$ or $w$. Let us now define a series $b_1, b_2, \ldots, b_\nu, \ldots$ such that $b_\nu$ is also only equal to $m$ or $w$ and different from $a_{\nu,\nu}$.

So if $a_{\nu,\nu} = m$, then $b_\nu = w$, and if $a_{\nu,\nu} = w$, then $b_\nu = m$.
\end{quote}

\textit{Note}: This is the diagonal construction. The diagonal element $a_{\nu,\nu}$ represents the self-application: evaluating the sequence named `$\nu$' at the position `$\nu$'. Cantor defines the new sequence (which he will call $E_0$) by `flipping' each of these self-referential elements: $b_\nu$ is defined as the opposite of $a_{\nu,\nu}$.

To connect this proof with sublation, this `flipping' operation must be understood as immanent within the system, not as an external negation. The sequence of $b$'s is constructed by recollecting the (purported) totality $\mathcal{M}$, taking an element from each sequence $E$. The whole structure is recalled, and the negation is applied to the diagonal. The new element is generated from what was already present, drawing out an internal contradiction. Within the allegory of the Sneetches, this operation corresponds to McBean's machine; it transforms the existing material (the Sneetch) into its opposite state regarding the star.

\begin{quote}
If we then consider the element:
\[
E_0 = (b_1, b_2, b_3, \ldots)
\]
of $\mathcal{M}$, we can easily see that the equation:
\[
E_0 = E_\mu
\]
for no positive integer value of $\mu$ can be satisfied, otherwise for the given $\mu$ and for all integer values of $\nu$:
\[
b_\nu = a_{\mu,\nu},
\]
so also in particular,
\[
b_\mu = a_{\mu,\mu},
\]
which would be excluded by the definition of $b_\nu$.
\end{quote}

\textit{Note}: Cantor proves that $E_0$ cannot equal any sequence in the enumeration. If $E_0$ were equal to the $\mu$-th sequence $E_\mu$, then all corresponding elements would match. But this creates a contradiction at position $\mu$, where $b_\mu$ was specifically defined to differ from $a_{\mu,\mu}$.

\begin{quote}
From this theorem follows immediately that the totality of all elements of $\mathcal{M}$ cannot be put into the series form:
\[
E_1, E_2, \ldots, E_\nu, \ldots,
\]
otherwise, we would be faced with the contradiction that a thing $E_0$ is both an element of $\mathcal{M}$ as well as not being an element of $\mathcal{M}$.
\end{quote}

This completes the core of Cantor's proof, establishing the fundamental logical structure of diagonalization. The proof demonstrates how any attempt to create a complete enumeration contains the resources for its own transcendence. The diagonal construction—the systematic ``flipping'' of self-referential elements—generates a new entity $E_0$ that must belong to the manifold $\mathcal{M}$ by definition, yet cannot occur on any purportedly complete list of its elements. This is a mathematical expression of the dialectical process in the Sneetches story: the systematic application of McBean's transformative machine eventually produces a situation that transcends the original problem. Just as the Sneetches' encounter with the limits of their star-based distinctions leads them to recognize ``that Sneetches are Sneetches,'' Cantor's exploration of the limits of enumeration leads to the recognition that some infinities necessarily exceed others. The contradiction is not a flaw but the engine of mathematical progress.


\section{Sublation as Formal Movement}

The connection between Cantor's proof and Hegelian sublation (\textit{Aufhebung}) relies on recognizing how the diagonal pattern exemplifies the structure of dialectical movement: preservation, negation, and elevation into a richer framework.

\subsection*{Preservation}

Sublation preserves what was true in the previous stage while transforming its meaning. In Cantor's proof, the original enumeration is not discarded. Every sequence in the list remains valid. The diagonal sequence $E_0$ does not negate the existence of $E_1, E_2, E_3, \ldots$. It demonstrates that these sequences, while real, do not exhaust the domain.

The totality attempted by enumeration is preserved as a \textit{part} of the richer structure. The countably infinite sequences remain; they are now understood as a proper subset within the larger uncountable domain. What was taken as complete is recontextualized as partial.

\subsection*{Negation}

Sublation negates the adequacy of the previous framework. In Cantor's proof, the diagonal sequence $E_0$ explicitly contradicts the claim that the list is complete. The negation is determinate: it does not assert ``the list is wrong'' but constructs the specific element that demonstrates incompleteness.

This is internal contradiction. The list itself, through its own structure, generates what it cannot contain. The limitation emerges from within, through the system's self-examination.

\subsection*{Elevation}

Sublation elevates understanding to a higher level. In Cantor's proof, the recognition that these sequences are uncountable opens a new domain of investigation: the study of different sizes of infinity, the hierarchy of infinite cardinals. What began as a limitation (the list cannot be completed) becomes a resource (there are fundamentally different types of infinity).

The diagonal sequence indicates that enumeration itself is the wrong kind of relation for this domain. The attempt to totalize through one-to-one correspondence breaks down, and this breakdown discloses something true about the structure of infinity.

\subsection*{The Pattern Generalized}

Cantor's proof is one instance of a broader pattern that occurs throughout mathematics:

\begin{itemize}
\item \textbf{Counting numbers and zero}: The attempt to represent ``nothing'' within a tally system generates the need for a symbol that transcends tallying---zero emerges as what cannot be counted but enables counting to become coherent. (Chapter 6)

\item \textbf{Integers and negatives}: The closure of natural numbers under subtraction requires numbers that transcend the natural domain.

\item \textbf{Rationals and irrationals}: The attempt to measure all magnitudes with ratios of integers generates magnitudes (like $\sqrt{2}$) that provably cannot be expressed as such ratios. (Chapter 8)

\item \textbf{G\"odel's incompleteness theorems}: The attempt to formalize arithmetic within a consistent axiomatic system generates statements that are true but unprovable within that system. (Chapter 9)
\end{itemize}

In each case, the system transcends itself in a productive way, indicating that its original boundaries were foundations for richer structures. The diagonal pattern makes this movement systematic.

\section{The Diagonal Sandwich: Mechanism of Transcendence}

The preceding sections established the four-moment pattern and identified it as a form of sublation. The internal mechanism of the third moment—Diagonal Recollection—warrants closer examination. As Haim Gaifman observes \parencite*{Gaifman2005,Gaifman2006}, the crucial move is often structured as a sequence: diagonalization, negation, diagonalization. This ``diagonal sandwich'' is the engine that generates the contradiction forcing transcendence.

\subsection*{The Structure of the Sandwich}

To understand this structure, recall Cantor's proof using $m$ and $w$:

\begin{enumerate}
\item \textbf{Diagonalization 1 (D1 - Identification)}: The system is reified into a list. We identify the diagonal elements $a_{\nu,\nu}$. This is the initial act of self-reference—examining each sequence $E_\nu$ at its own position $\nu$.

\item \textbf{Negation (N - The Operator)}: A new sequence $E_0$ (with elements $b_\nu$) is constructed by applying an operator (the $m/w$ flip) to the elements identified in D1. $b_\nu$ is defined to be different from $a_{\nu,\nu}$.

\item \textbf{Diagonalization 2 (D2 - The Test)}: This step closes the loop. We assume the new sequence $E_0$ belongs to the original totality (i.e., $E_0 = E_\mu$ for some $\mu$) and then test this assumption at the point of self-reference $\mu$. We ask: what is the value of $E_0$ at position $\mu$? This second diagonalization forces the contradiction: $b_\mu = a_{\mu,\mu}$ (by assumption of membership) and $b_\mu \neq a_{\mu,\mu}$ (by construction).
\end{enumerate}

The contradiction does not arise from constructing the diagonal (D1) or negating it (N). It arises when the negated diagonal is reflected back into the original system and tested at the self-referential point (D2). This structure guarantees the generation of a fixed point that destabilizes the totality.

\subsection*{The Sandwich in Dialectical Logic}

The insight here is that this formal structure is homologous with the movement of dialectic. Whenever thought attempts to achieve immediacy or closure, it enacts this structure.

\subsubsection*{Being, Nothing, Becoming (The Logic)}

The beginning of Hegel's \emph{Science of Logic} provides a pure example:

\begin{enumerate}
\item \textbf{D1 (Identification)}: Thought attempts to think ``Pure Being'' (immediacy, thought thinking itself).
\item \textbf{N (Negation)}: Pure Being is indeterminate. Its content is Nothing. The concept negates itself.
\item \textbf{D2 (The Test)}: Reflection indicates that Nothing also simply \textit{is}. Being vanishes into Nothing, and Nothing vanishes into Being. This oscillation is the conceptual equivalent of the formal contradiction.
\item \textbf{Transcendence (Becoming)}: The truth is the movement between them: Becoming.
\end{enumerate}

\subsubsection*{Sense-Certainty (The Phenomenology)}

The first movement in the \emph{Phenomenology of Spirit} follows the same pattern in the realm of experience:

\begin{enumerate}
\item \textbf{D1 (Identification)}: Consciousness attempts to grasp the immediate object, the ``This,'' using self-referential indexicals (e.g., ``Now'').
\item \textbf{N (Negation)}: Articulating the ``This'' requires universals. The immediate particularity vanishes into an empty universal (Any time can be ``Now'').
\item \textbf{D2 (The Test)}: The mechanism used to grasp the particular (the indexical) is what rendered it universal. The attempt to fixate the immediate reference point fails.
\item \textbf{Transcendence (Perception)}: Consciousness moves to Perception, acknowledging that knowledge is mediated.
\end{enumerate}

In both cases, the dialectical movement is driven by the structure of the diagonal sandwich. The attempt to grasp a totality immediately (D1) exposes its internal negation or indeterminacy (N). When this result is reflected back into the original attempt (D2), it generates a contradiction or oscillation that forces the system to reorganize itself at a higher level of complexity (Transcendence).

\section{Methodology and Caution}

The pattern identified here---totality, reification, diagonal recollection, transcendence---provides a heuristic for recognizing structural similarities across different domains. But several cautions must be emphasized.

\subsection*{Against Mechanical Application}

This is not an algorithm. One cannot simply search for ``totalities'' and ``diagonals'' and expect insight to follow. Each instance involves irreducibly particular content that must be understood on its own terms.

The Sneetches story involves social recognition and the transformation of shared norms. Cantor's proof involves formal structures and the logic of enumeration. The pattern occurs in both, but the content cannot be collapsed into the pattern. The value lies in recognizing \textit{after the fact} that similar dynamics are at work, not in predicting where such dynamics will occur.

\subsection*{Against Occulting the Particular}

Thematizing the pattern creates a risk: the particular instances may become illustrations of an abstract schema. This would reverse the actual order of understanding. Mathematical concepts do not emerge because mathematicians are enacting a metaphysical pattern. Rather, the pattern becomes recognizable because mathematical practice repeatedly encounters certain structural necessities.

In each case, the particular problem demanded a particular solution, and the pattern became evident retrospectively through philosophical reflection on the structure of those solutions.

\subsection*{Against Totalizing the Pattern Itself}

Perhaps most importantly: this account of the pattern is itself subject to the dynamic it describes. By articulating the four moments as a complete schema, this chapter reifies the pattern.

But the pattern, when reflected back on itself, generates cases that do not fit cleanly. Not every mathematical development follows this structure. The attempt to totalize the pattern of transcendence necessarily encounters instances that exceed the totalization.

This is not a failure of the account but an affirmation of its claim. If the pattern holds generally, it must hold of itself: any attempt to completely systematize the movement of sublation will be partial. The formalization serves not to capture that practice finally but to make certain patterns explicit for pedagogical purposes.

\subsection*{The Role of This Bridge}

This chapter stands as scaffolding. It establishes a vocabulary and a shared reference point for the mathematical investigations that follow. When subsequent chapters examine mathematical concepts, variations on this pattern may be recognized. The point is to facilitate recognition of structural similarities without collapsing the particularity of each case.

The bridge serves its purpose when crossed. Once engagement with the mathematical content itself begins, the explicit thematization of the pattern recedes.

\section{Bridging to Mathematical Content}

This bridge prepares for what follows by making explicit the methodological commitments that animate the project. Mathematics is treated here not as a static structure of timeless truths but as a practice of conceptual development. Understanding mathematical concepts requires attending to how they emerge, how they transform, and how they transcend the limitations of earlier frameworks.

The pattern traced here---through the Sneetches and through Cantor---recurs because it reflects how thinking works when it becomes sufficiently self-reflective. Systems that can examine themselves generate insights that exceed their original boundaries. This is the ordinary movement of thought encountering its own limits and reorganizing itself into richer structures.

The mathematical investigations that follow (Chapters 6–9) apply this understanding to specific historical developments:

\begin{itemize}
\item \textbf{Chapter 6 (Zero)}: The emergence of zero from the tension between tally and base notation.
\item \textbf{Chapter 7 (Numerals as Pronouns)}: How numerals function anaphorically.
\item \textbf{Chapter 8 (Algorithmic Elaboration)}: How algorithms make implicit commitments explicit through stepwise formalization.
\item \textbf{Chapter 9 (Operations and Incompleteness)}: How arithmetic operations emerge, and how G\"odel's incompleteness theorems demonstrate the necessary openness of formal systems.
\end{itemize}

In each case, the four-moment pattern occurs. But in each case, the particular content demands particular analysis. The pattern guides recognition but does not replace investigation.


\section{Conclusion: Toward Zero and Beyond}

This bridge has established the foundational pattern that will recur throughout the mathematical investigations to follow. Through the Sneetches' encounter with McBean's machine and Cantor's exploration of infinite sequences, the chapter has traced how systems that attempt to totalize themselves necessarily generate elements that exceed their boundaries.

The four moments---totality, reification, diagonal recollection, transcendence---provide a vocabulary for recognizing this pattern. The diagonal sandwich provides the mechanism. They indicate that mathematical understanding develops not through the accumulation of static truths but through the dynamic process of systems encountering and productively transcending their own limits.

This pattern occurs at various levels of mathematical development, from the emergence of zero to G\"odel's incompleteness theorems. In each case, what functions as limitation or contradiction becomes the resource for growth. The system does not break down when it encounters elements it cannot contain; it reorganizes itself into a richer structure that preserves what was valuable while opening new possibilities.

The text now turns from this methodological scaffolding to substantive mathematical investigation. Chapter 6 examines the first and most fundamental instance of this pattern: the emergence of zero. Where this bridge established the pattern abstractly, Chapter 6 will demonstrate how it manifests concretely in the transition from tally marks to base notation, presenting zero as the necessary expression of mathematical becoming.

\printbibliography[heading=subbibliography]

\renewcommand{\chaptername}{\oldchaptername}
\renewcommand{\thechapter}{\oldthechapter}

% Manually advance the chapter counter so the next regular chapter uses the correct number
\setcounter{chapter}{5}

\part{The Hermeneutic Calculator}

\chapter{The \textit{in}finite: An Introduction to Zero}

\begin{abstract}
This chapter examines zero as the first and most fundamental instance of the book's central dialectical pattern. Zero emerges not as simple absence but as the formalization of the enabling condition that recedes---the necessary ground of counting that withdraws when grasped. Through pragmatic, semantic, and critical lenses, the analysis traces how zero functions as both a placeholder for determinate absence and an active recollection of negativity itself. The chapter grounds this abstraction in the phenomenology of counting, from animal cognition to the historical development of place-value systems. Ultimately, zero is revealed as the mathematical analogue of the divaded structure of self-consciousness, marking the point where the mathematical system recognizes its own unrepresentable ground.
\end{abstract}

\section{Introduction: The First Instance}

The Bridge established a pattern: totality, reification, diagonal recollection, transcendence. Through the Sneetches and Cantor's proof, we saw how systems that attempt to totalize themselves necessarily generate elements exceeding their boundaries. Now we examine the first and most fundamental instance of this pattern in mathematical development: the emergence of zero.

Zero is not simply the absence of quantity. It is the determinate negation that makes base systems possible, the posited void that enables recursive self-reference in arithmetic. Where tally systems represent number through pure accumulation---one mark after another in monotonous progression---base systems achieve something radically different. They turn back on themselves, grouping quantities into higher-order units through the dialectical movement of sublation.

Zero emerges pragmatically from the internal logic of counting itself, semantically through the philosophical traditions that make formalization intelligible, and critically by localizing contradiction in notation rather than propositions. Each movement negates and preserves the others---zero is not any one of these in isolation but their constant becoming. From this ground, the structure of number will grow in Chapter 7. But first, we must understand the ground itself.

\section{Auto-ethnographic Ground: The Sycamore Experience}

When I was a sophomore at Earlham College, still pursuing a double major in physics and mathematics, I took a seminar where von Neumann ordinals were introduced. The construction itself was elegant: each ordinal number is the set of all smaller ordinals, with zero defined as the empty set:

\[
0 := \emptyset,\quad
1 := \{\emptyset\},\quad
2 := \{\emptyset, \{\emptyset\}\},\quad
3 := \{\emptyset, \{\emptyset\}, \{\emptyset, \{\emptyset\}\}\},\;\dots
\]

But what captivated me was not the formal construction itself. It was the realization that ``two'' could be expressed through structured nothingness---that the answer to ``what is 2?'' could be traced back to the void ($\emptyset$). For someone who had internalized mathematics as a crushing structure meant to expose incompetents, this was revelatory. The whole edifice of mathematical thought, I glimpsed, could be understood through nothingness.

That evening, I went to a favorite spot on campus. A huge sycamore grew beside a pond, its branches hanging low over the water. I settled into a crook and stared up at the stars. Everywhere I looked---every spatial coordinate I thought toward---I found the empty set. Nothingness was like a conduit for thought. I could think to anywhere \textit{because} nothingness was everywhere.

This resonance was an act of recognition. In the formal structure of the empty set, I encountered the same movement experienced in meditation---the release of determinate content that allows a deeper ground to emerge. The mathematical concept and the lived practice shared a rhythm: the letting-go that enables beginning, the same rhythm practiced in the exercise of listening to time (Chapter 1).

Grasping the empty set was recognizing the structure of awareness itself---the unrepresentable condition that makes representation possible. The stillness in the tree and the stillness of the empty set were one movement. This encounter grounded the inquiry that follows, transforming an abstract puzzle into an existential imperative: to understand how this generative absence structures both mathematical thought and human being.

The experience extended beyond mathematics. I thought through the stars themselves, considering how every bit of condensed energy in our material universe was porous, wave-like, fundamentally empty. I was overcome by the vastness of nothingness. In that moment, mathematics ceased to be an alienating formalism and became a spiritual practice. I dropped my physics major and committed fully to mathematics.

This points to what I now call \textit{the enabling condition that recedes}. The moment I tried to grasp the nothingness directly, to make it an object of contemplation, it withdrew. I was always already inside what I was trying to examine. The empty set was not ``out there'' in the stars or ``in'' the mathematical formalism. It was the structure of my own thinking, the necessary ground that makes any thought possible yet vanishes when approached directly.

The transition from the felt sense of absence to the mechanics of counting traces how this structure becomes operative in practice, how the enabling ground extends into the shared, historical practices of quantification. But before we can understand why such notation is necessary, we must attend to how zero emerges from counting itself.

\section{The Pragmatic Sense: From Crows to Zero}

The null representation is not an invention imposed on mathematics from outside but a formalization of structures already operative in counting practices. This section traces zero's development through three stages: pre-linguistic ordinality, the historical progression from tally marks to base systems, and the final recognition of zero as an independent number. This pragmatic emergence connects directly to the diagonalization pattern: each apparent totality of counting reveals a gap that must be formally marked.

\subsection*{Counting Before Numerals: The Case of Corvid Cognition}

Before examining the human history of counting, it is instructive to consider counting practices that exist entirely without symbolic numerals. Recent research on corvid cognition---particularly studies of crows and ravens---reveals sophisticated numerical abilities that operate below the level of symbolic representation \parencite{rugani2015number, ditz_numerosity_2016}.

Crows can distinguish between different quantities (three items versus four items) and can recognize sequential order (first seed, second seed, third seed). Experiments show they can track ordinality across sequences: which perch was visited first, which second, which third. This demonstrates that the \textit{structure} of counting---the recognition of sequential position and relative quantity---exists independent of any symbolic system.

Ordinality precedes cardinality---the crow recognizes sequential position before grasping abstract quantity. The crow does not first grasp ``threeness'' as an abstract quantity and then apply it to sequences. Rather, the sequence itself (first, second, third) is primary, and cardinality (``there are three'') emerges as the recognition of the sequence's endpoint.

Importantly, corvid counting has an implicit origin point. When a crow tracks ``first perch, second perch, third perch,'' there is a pre-thematic ``not yet counted'' that precedes the first perch. This is not marked or represented---it is not zero in any formal sense. But it is operatively present as the background from which counting begins. The crow's cognition does not start \textit{at} the first item; it starts \textit{before} the first item, even if this before is implicit rather than explicit.

This pre-human counting structure already contains the seed of zero: the implicit recognition that counting begins from a point before the first count. What human symbolic systems will eventually do is make this implicit background \textit{explicit} through formal notation. But the structure itself---the gap between ``not yet counted'' and ``first counted''---is already there in pre-linguistic numerosity.

\subsection*{The Historical Emergence of Base Systems and the Need for Zero}

Human counting systems begin similarly to corvid cognition: with sequential enumeration marked through physical tokens. Tally marks on bones dating back 35,000 years show one-to-one correspondence---one mark for each item counted. This is pure Repulsion in Hegelian terms: each mark stands apart from every other, an undifferentiated sequence of ones.

But tally systems encounter a practical limitation. To represent large quantities requires unwieldy numbers of marks. A thousand items would require a thousand tally strokes. There is no higher-order organization, no principle of unity beyond simple accumulation. The marks repel each other without synthesis.

The dialectical resolution comes through \textit{grouping}. Instead of marking each item individually, we group items into clusters---say, groups of ten. We then introduce a notation for the group itself, distinct from the notation for individual items. Now ``ten'' functions as a new unit, and we can count at a higher level: one ten, two tens, three tens.

This is a transformative move, an instance of dialectical sublation (Aufhebung). The individual marks remain present as components of the group, yet they are negated in their individuality---they no longer count as separate ones but as a single unit ``ten.'' The individuality of separate marks is denied, yet they remain present as components, now elevated into a new structural form with different properties.

The transition from ``ten ones'' to ``one ten'' introduces what Hegel calls \textit{Attraction}---the dialectical counter-movement to \textit{Repulsion}. The many ones are drawn together into a new unity. This unity is not external to the ones but emerges through their internal relationship. The ten \textit{is} the ten ones, yet it functions differently than accumulating ten marks.

This grouping structure creates the possibility for base systems. A base-ten system (the most common historically, likely due to ten fingers) organizes quantities through recursive grouping: ten ones make one ten, ten tens make one hundred, ten hundreds make one thousand. Each higher order is a grouping of lower orders.

But such a system creates a new problem---the diagonal pattern from the Bridge manifests here. How do we represent a quantity that has, say, two hundreds and three ones but no tens? If we write ``23,'' this is ambiguous---does it mean two tens and three ones, or something else? The gap in representation is itself meaningful: there is something absent (no tens) that needs to be marked.

This gap is not external to the system but generated by the system's own internal logic. The very structure that makes base notation possible---recursive grouping through position---creates the necessity for marking absence. The totalization (``we can represent any quantity through positioned numerals'') immediately reveals its incompleteness (``but we need a symbol for when a position is empty'').

Zero emerges pragmatically from this requirement. To think of ``203'' in base-ten is to think:
\begin{itemize}
\item 2 hundreds (two groups of ten tens)
\item 0 tens (no groups of ten)
\item 3 ones (three loose units)
\end{itemize}

The zero in the tens place is not mere absence. It is \textit{determinate} absence---the explicit marking of ``no tens here.'' This determines the position and meaning of surrounding numerals. Without the zero, ``23'' would mean something entirely different (two tens and three ones, not two hundreds and three ones). Zero makes the place-value structure work by explicitly marking where groupings are absent.

Once we introduce recursive grouping through sublation, we require a symbol for determinate absence to maintain positional clarity. Zero is not an afterthought added to counting; it is the inevitable consequence of turning counting back on itself, of making number self-referential through hierarchical structure. The diagonal gap---the element the system cannot represent within its current framework---forces the system to expand to include it.

The emergence of the placeholder reveals a necessary structure of symbolic systems, enacting the mediation pattern where what appears autonomous requires its opposite for determination. The system of place value, designed to articulate quantity efficiently, generated a gap---an absence that needed marking for the system to function. The determinate value of the numerals (the somethings) could only be stabilized through the introduction of a mark indicating their absence (the nothing). The system and the placeholder are mutually constitutive; neither can be grasped in isolation.

Yet this function transformed what it was meant to serve. The mark indicating absence began to accrue conceptual content. This reciprocal dependence---where the system generates the element that transforms the system---is the structure of mediation itself. Negation functions as a generative force. The introduction of the placeholder was a determinate negation---it negated the simple presence of a value in a specific column, yet in doing so, it preserved the structure of the base system and elevated it to a new level of expressive power. The absence became productive, enabling the articulation of quantities previously beyond the reach of notation. This movement from functional necessity to conceptual entity traces the rhythm of sublation, where limitations are overcome by incorporating the negation that defines the boundary.

Historically, different cultures arrived at this recognition at different times. Ancient Mesopotamian mathematics used a placeholder as early as the 3rd century BCE, though without treating it as a number. The Maya independently developed zero in their vigesimal (base-20) system. The systematic treatment of zero as a number with its own arithmetic properties emerged in India, with Brahmagupta's work (circa 628 CE) providing rules for operations involving zero \parencite{ernest_mathematical_2024}.

\subsection*{Zero as Independent Number}

Once zero functions as placeholder, it can undergo further transformation. Brahmagupta's innovation was to treat zero not as a mark of absence but as a number in its own right, with its own arithmetic properties:

\begin{itemize}
\item $a + 0 = a$ (zero as additive identity)
\item $a \times 0 = 0$ (zero annihilates in multiplication)
\item $a - a = 0$ (zero as result of self-subtraction)
\end{itemize}

This represents another level of sublation. Zero began as:
\begin{enumerate}
\item \textbf{Absence}: The gap between groupings, unmarked
\item \textbf{Placeholder}: The explicit mark of determinate absence in positional notation
\item \textbf{Number}: An independent mathematical object with operational properties
\end{enumerate}

Each stage negates the previous while preserving it. Zero as number doesn't eliminate zero as placeholder---positional notation still requires it. Rather, the placeholder function becomes one property among many. Zero gains new capabilities (participating in equations, representing quantities in its own right) while retaining its original structural role.

This transition from functional placeholder to independent number is a transformation in the structure of mathematical consciousness. It requires a shift from grasping numbers as direct representations of quantities to apprehending them as elements within a relational system. When zero is treated as a number, it functions as an operator within that system, subject to the same rules of arithmetic as other numerals.

This movement makes explicit the inferential structure that governs mathematical meaning, as explored in Chapter 2. The content of ``zero'' is determined by its role in the network of inferences it licenses and those it precludes. Recognizing zero as a number is recognizing its place in the space of reasons---recognizing that the absence of quantity is itself a determinate mathematical concept, defined by the movements it enables within the system.

This progression exhibits the dialectical pattern from the Bridge. Each apparent totality (tally systems can count everything; grouped tallies can represent structure; base systems can express any number) reveals limitation (no unity principle; no mark for absence; operations require zero as number) that is resolved through sublation, which in turn generates new totalizations requiring further development.

What unifies these three pragmatic stages---absence, placeholder, number---is that all three formalize \textit{the enabling condition that recedes}. The unmarked absence is the implicit background from which counting begins (like the crow's pre-thematic origin point). The placeholder explicitly marks this background within the formal system. The number zero makes the background itself manipulable, turning the enabling condition into an operative element.

Yet zero never loses its character as ground. Even when we calculate with zero, when we write equations like $5 + 0 = 5$ or $7 \times 0 = 0$, zero retains its quality of being different from other numbers. It is the number that represents the absence of number, the quantity that marks no-quantity. This paradoxical status---present as symbol yet representing absence, operating within the system yet marking the system's ground---is what makes zero the natural formalization of the null representation.

Having traced zero's pragmatic emergence, we now turn to its semantic formalization: how quotation becomes recollection, how recollection becomes operative, and how the enabling condition that recedes can be captured in notation.

\section{The Semantic Sense: Formalizing the Receding Condition}

The pragmatic emergence of zero reveals its necessity: base systems require a symbol for determinate absence. But why should this symbol be understood as the null representation---as the formalization of the enabling condition that recedes? Let us work through the semantic structure in four movements: quotation as recollection, the formal substitution that makes recollection operative, Hegel's dialectic of self-positing nothingness, and Derrida's concept of the self-erasing trace. Each movement deepens our understanding of how zero functions not as absence but as the active ground of determination.

\subsection*{Quotation as Recollection in Natural Language}

Before we can make the null representation mathematically operative, we must understand how recollection functions in natural language. Quotation is the primary device through which language refers back to itself, creating what Robert Brandom \parencite*{Brandom:2019aa} calls \textit{repeatability structures} from what would otherwise be fleeting and context-bound utterances.

Consider: When I say ``She said 'I am hungry,''' the inner quotation marks create an anaphoric structure. The quoted phrase ``I am hungry'' does not refer to my current state but recollects her prior utterance. Quotation transforms the indexical ``I'' from her speech into a repeatable token within my discourse. The marks themselves---the paired symbols `` and ''---function as a kind of semantic time machine, reaching back to extract content from a previous moment and making it present again.

Quotations can contain quotations. This allows for nested recollection: ``He said 'She said ``I am hungry''''' creates a two-level embedding. Each layer of quotation marks recollects a prior speech act. The structure is inherently recursive---there is no limit in principle to how many layers of quotation can be nested.

This quotative embedding creates a temporal-sequential structure within language itself. When we trace back through nested quotations, we are reconstructing a history of utterances. The innermost quote is the \textit{originary} utterance (in this local context), and each subsequent layer recollects and re-presents it.

But what about the first utterance, the one that stands before all quotation? When she originally said ``I am hungry,'' there were no quotation marks. Yet for her utterance to be meaningful, there must have been an implicit ``I think'' accompanying it---the Kantian transcendental ground that makes any judgment possible. This implicit ground is never directly stated, never explicitly marked. It is the \textit{enabling condition that recedes}.

Kant argued in the \textit{Critique of Pure Reason} that the ``I think'' must be able to accompany all representations. This is the transcendental unity of apperception---the condition that makes experience possible. But the ``I think'' itself is not an object of experience. Any attempt to represent it directly transforms it into an object, losing precisely its character as the subjective source of representation.

Quotation marks are language's way of pointing back to its own enabling conditions. When we quote, we acknowledge that meaning does not originate in the present utterance but is always a recollection of something prior. And if we trace this backwards through nested quotations, we approach (without ever reaching) the unrepresentable ground---the ``I think'' that must be presupposed but cannot be directly grasped.

Natural language cannot formalize this structure completely because quotation marks are semantic rather than operative. They tell us that recollection is happening, but they don't give us a calculus for manipulating recollective structures. This is where mathematics enters.

Quotation is an enactment of recollection. By placing words within quotation marks, we shift their status. They are presented as having been asserted previously. This shift introduces a temporal distance---recognizing that the utterance belongs to a past moment while being made present again. Every act of quotation acknowledges the history embedded in language, compressing past practice into present form. The quotation marks function as the boundary that both separates the utterance from its immediate context and connects it to its history.

This mechanism of recollection is structurally identical to the act that constitutes self-consciousness. The ``me''---the social self---is a recollection of the \{I\}---the spontaneous subject of action. Just as quotation introduces a gap between the original utterance and its repetition, recollection introduces a gap between the acting \{I\} and the recognized ``me.'' This gap is the necessary structure that makes self-awareness possible. It is the space where recognition occurs. Zero functions as the formal mark of this gap---the recollection of the unrepresentable ground of subjectivity itself.

\subsection*{The Substitution: Making Recollection Formally Operative}

The key move that makes the null representation mathematically operative is a notational substitution: we replace quotation marks with set-theoretic braces. Specifically:

\begin{center}
\begin{tabular}{rcl}
\textbf{Opening quote} ``{} & $\rightarrow$ & \textbf{Opening brace} \{ \\
\textbf{Closing quote} ''{} & $\rightarrow$ & \textbf{Closing brace} \}
\end{tabular}
\end{center}

This is more than notational convenience. The substitution transforms quotation (a semantic device in natural language) into set formation (a formal operation in mathematics). What was implicit becomes operative.

Consider the phrase ``I think.'' We can represent it initially with quotation marks to indicate we are referring to the phrase itself: ``I think''. Now perform the substitution:

\[
\text{``I think''} \quad \rightarrow \quad \{\text{I think}\}
\]

The braces now indicate that we have formed a set containing the concept ``I think.'' But this is still treating ``I think'' as content, as something that can be collected into a set. Following our earlier insight that the ``I think'' is not itself representable but is the condition of all representation, we recognize that what we are really trying to capture is not any particular thought but the \textit{enabling condition} itself. So we perform an erasure:

\{I think\} $\rightarrow$ \{\st{I think}\}

The cancellation marks indicate that ``I think'' is inadequate to capture the enabling condition it attempts to name. But we can go further. Since the crossed-out content is now explicitly marked as inoperative, we can remove it entirely while preserving the structure:

\{\st{I think}\} $\rightarrow$ \{\}

The empty braces remain. They are not nothing---they are the \textit{trace} of the erasure we performed. They mark the place where the enabling condition was invoked and then withdrawn. Finally, we introduce a symbol for this structure:

\[
\{\} \quad =: \quad \emptyset
\]

The symbol $\emptyset$ is the null representation. It is not simply ``nothing'' but the formalization of a process: invoke the enabling condition, recognize its inadequacy when named, cross it out, preserve the trace of the crossing-out. The empty braces are anaphoric: they point back to what was crossed out without containing it.

This substitution makes recollection formally operative in mathematics. Just as quotation marks in natural language create anaphoric structures that refer back to prior utterances, set braces in mathematics create structures that refer back to enabling conditions. When we write $\emptyset$, we are not asserting the existence of an empty object. We are performing a recollective move that formalizes the self-negating ground of thought.

This allows nested recollection just as natural language allows nested quotation. If we form the set $\{0\}$ (which von Neumann identifies with the number 1), we are recollecting the null representation itself. The outer braces refer back to $\emptyset$, making the ground itself an object of recollection. And this can iterate indefinitely:
\begin{align*}
0 &:= \emptyset \\
1 &:= \{0\} = \{\emptyset\} \\
2 &:= \{0, 1\} = \{\emptyset, \{\emptyset\}\} \\
3 &:= \{0, 1, 2\} = \{\emptyset, \{\emptyset\}, \{\emptyset, \{\emptyset\}\}\}
\end{align*}

Each number is a nested structure of recollections, with $\emptyset$ at the base. The entire edifice of arithmetic becomes formalized anaphoric recollection. But this only works if $\emptyset$ successfully formalizes the enabling condition that recedes. To understand why it does, we must turn to Hegel's concept of self-positing nothingness and Derrida's concept of the self-erasing trace.

\subsection*{Hegel and Derrida: The Self-Erasing Mark}

Hegel begins the \textit{Science of Logic} with Being and Nothing \parencite{Hegel:2014aa}. Pure Being, the concept of Being with all determinate content removed, is utterly indeterminate. There is nothing we can say about it except that it is. But Nothing, similarly, is utterly indeterminate. There is nothing we can say about it except that it is not. These two concepts, pursued to their limit, collapse into each other.

But this collapse is not static. It is the movement of Becoming. Being vanishes into Nothing; Nothing vanishes into Being. This vanishing is the generation of determinateness. From the self-negating movement of Being and Nothing, all further determinations emerge.

What Hegel calls self-positing nothingness is precisely this: Nothing that negates itself, that crosses itself out and in so doing generates Being. The negative is not privative (absence of something) but productive. It is the engine that drives dialectics. The null representation $\emptyset$ formalizes this structure. It is not simply the absence of elements. It is the positing of absence, the active marking of nothing. And through this marking, it becomes the foundation for all subsequent structure. From $\emptyset$, we build natural numbers. From natural numbers, integers. From integers, rationals, reals, complex numbers. The entire edifice of mathematics grows from this self-negating ground.

Derrida extends this Hegelian insight through the concept of the trace \parencite{Derrida1997}. The trace is neither present nor absent. It is the mark of what has been erased, the remainder that testifies to its own deletion. Every concept, Derrida argues, carries traces of what it excludes, what it defines itself against. The trace is the condition of intelligibility that remains unrepresentable within the system.

The self-erasing mark, what Derrida terms the trace, is the structure we are formalizing in our understanding of zero. The trace is the movement of differing and deferring that constitutes meaning. It is the enabling condition that withdraws when approached directly, the structure identified in Chapter 1.

Zero functions as the mathematical articulation of this structure. It marks the place where presence gives way to absence, yet this absence is constitutive of the system itself. It is the ground that recedes, the origin that is never fully present yet makes all presence possible. When we attempt to fixate on zero as an object, it vanishes into the system it enables. Yet without it, the system collapses. This withdrawal is its essential function: to ground the system by remaining ungraspable within it. Apprehending zero is apprehending the movement of thought itself---the restless negativity that both constitutes and exceeds any determinate system. It is the soundless sound that guards the event of mathematical language.

Applied to mathematical nothingness, this means $\emptyset$ is not a stable object but a trace structure. It points to the enabling condition (the ``I think,'' the ground of representation) that has been crossed out. The empty braces are the formal mark of this erasure. They testify that something was invoked and withdrawn, that a name was given and immediately destroyed.

To name it directly---to say ``$\emptyset$ is the empty set''---is already to miss what it is. Any direct representation treats it as a stable presence, as an object among objects. But $\emptyset$ is not an object. It is the formal trace of the self-negating movement that enables objectivity itself. The symbol performs its own inadequacy---the braces invoke content and mark its absence in a single gesture.

The convergence of Hegel and Derrida here is striking. Hegel's self-positing nothingness describes the generative power of the negative---how Being and Nothing vanish into each other to produce Becoming. Derrida's self-erasing trace describes the groundlessness of this generation---how any ``origin'' is already contaminated by the trace of otherness. One describes the positive moment (generation), the other the negative moment (erasure). But both recognize that zero is not a static point but a dynamic operation: the crossing-out that enables inscription, the death that enables life, the nothingness that becomes being precisely by remaining nothing.

\subsection*{The Divaded Structure of Apperceptive Consciousness}

Having established the null representation through Kantian recollection, Hegelian self-positing, and Derridean self-erasure, we must now examine how this structure appears concretely in the act of judgment. This requires attending to what Kant called the transcendental unity of apperception---the ``I think'' that must be able to accompany all representations. But we must go beyond Kant to recognize that this unity is not simple but inherently self-differentiated, what I call the \textit{divaded} structure of consciousness.

Consider the experience of making a simple assertion: ``Pokey is a dog.'' In uttering this judgment, I am engaged in a complex synthetic activity. I bring together a singular term (``Pokey'') with a predicate (``is a dog'') through the copula (``is''). This bringing-together is not automatic or mechanical. It represents the compression of temporal learning experiences---encounters with Pokey, recognition of dog-like behaviors, application of the concept ``dog''---into a unified assertion that stakes a claim about reality.

When I make this judgment, I am \textit{both} inside and outside the assertion simultaneously. I am inside it as the synthetic activity that performs the unification: I am the one doing the bringing-together of subject and predicate. Without my active synthesis, there would be no judgment, only isolated terms. But I am also outside it, aware that this assertion is inadequate to fully represent what \{I\} am. The judgment ``Pokey is a dog'' captures something about Pokey, but it cannot capture the judging consciousness itself---the ``I think'' that enables the judgment while remaining irreducible to any of its contents.

This is not a division into two separate entities---an empirical self that judges and a transcendental self that observes. Rather, it is a \textit{divaded} structure: a unity that is inherently self-differentiated, a single consciousness that exists in two modes simultaneously.\footnote{The term ``divaded'' distinguishes this from division. Division suggests two separate things; divadedness suggests a single thing that is intrinsically doubled, that contains difference within its very unity.}

This divaded structure has important implications for how we understand the null representation. When I assert ``Pokey is a dog,'' the site of synthesis is the copula---the ``is'' that unites subject and predicate. This is precisely where the ``I think'' operates, yet this operation cannot itself appear in the judgment as content. The copula performs the unity but does not name the unifier. Following the substitution established earlier, we can now mark this enabling condition within the judgment itself:

\[
\text{Pokey } \emptyset \text{ is a dog}
\]

The null representation $\emptyset$ appears at the copula, marking the site where synthetic activity occurs. It is not an element of the judgment in the way ``Pokey'' or ``dog'' are elements. It is the formal trace of the unrepresentable activity that makes the judgment possible. The copula ``is'' performs synthesis; $\emptyset$ marks that this synthesis has an unrepresentable source.

But the judgment is also a unified whole that can itself become the object of subsequent thought. To indicate this objectification---the fact that the judgment can be taken as a complete unity---we enclose it in braces:

\[
\{\text{Pokey } \emptyset \text{ is a dog}\}
\]

This notation captures the inside/outside symmetry of apperceptive consciousness. The $\emptyset$ inside represents the unrepresentable synthetic activity immanent to the judgment. The braces \{\} outside represent the transcendental move of taking the judgment as a unified object, externalized and available for reflection. The notation is itself divaded: it shows both the performance (inside) and the recognition of that performance (outside) in a single formal structure.

This symmetry is not optional or supplementary. It is the very structure of self-consciousness. To be conscious is to be apperceptively aware---to perform synthesis while simultaneously being aware that one is performing it. The divaded structure is not a philosophical theory imposed on consciousness but a description of what consciousness \textit{is} when attended to carefully.

The inside/outside symmetry also clarifies why the null representation must be paradoxical. Any attempt to capture the ``I think'' directly would place it entirely ``inside'' the judgment as content, losing its character as enabling ground. Any attempt to place it entirely ``outside'' would make it transcendent in a way that severs it from concrete judgments. The null representation $\emptyset$ preserves both aspects: it appears inside (at the copula) while pointing outside (to the unrepresentable source). It is the notational performance of divadedness itself.

This phenomenological analysis grounds the formal apparatus in lived experience. When students encounter mathematical formalism and feel alienated, part of what they experience is the loss of this apperceptive dimension. The symbols appear as ``outside'' objects to be manipulated according to external rules, severed from the ``inside'' experience of synthesis and understanding. By recognizing that $\emptyset$ marks the site where consciousness touches mathematics---where the unrepresentable ``I think'' enables all mathematical thought---we can begin to restore the connection between formalism and lived mathematical experience.

This structure of self-consciousness---the simultaneous unity and difference of the \{I\} that thinks and the ``me'' that is thought about---is inherently divaded, mirroring the structure first articulated in the Prelude. The self is both inside and outside its own representations. The \{I\} accompanies all representations, yet it cannot itself be represented; it exceeds every determinate content we ascribe to the ``me.'' This simultaneous inside/outside relation is the constitutive structure of subjectivity.

Zero mirrors this structure in the domain of mathematics. It is inside the number system as a determinate element, subject to arithmetic operations; yet it is outside the realm of countable quantities, marking the absence that enables the system's movement. It is the point where the system recognizes its own boundary, the mark of the \textit{in}finite within the finite. The difficulty in pinning down zero mirrors the difficulty in pinning down the self. Both are constituted by this dynamic interplay between presence and absence, between being determined within a system and exceeding that determination. Zero is the moment of maximum expansion, the release of fixation that allows the next compression (the next number) to emerge.

\section{Critical Nullity: Contradiction Without Dialetheism}

Having established the semantic and pragmatic senses of the null representation, we now address a philosophical question: How does this approach handle the apparent contradiction of representing the unrepresentable? This section clarifies the critical position by distinguishing it from dialetheism, explaining how contradiction is localized in the notation itself, and showing how anaphoric recollection preserves non-deflationary Hegelian negativity.

\subsection*{The Threat of Paradox}

The entire development of the null representation seems to rest on a performative contradiction. I have claimed that:
\begin{itemize}
\item The ``I think'' is unrepresentable (cannot be directly captured in formal notation)
\item The null representation $\emptyset$ represents the ``I think''
\end{itemize}

This appears to be a straightforward logical contradiction: $\emptyset$ both represents and does not represent the enabling condition. If we take both claims as assertions within the same logical framework, we face incoherence.

One response to such contradictions is \textit{dialetheism}---accepting that some contradictions are true (Priest). This approach differs: it localizes contradiction in the notation itself, making it performative rather than assertoric.

\subsection*{Localizing Contradiction in Notation}

Return to the substitution that generates the null representation:

``I think'' $\rightarrow$ \{I think\} $\rightarrow$ \{\st{I think}\} $\rightarrow$ \{\} $\rightarrow \emptyset$

The key move is the step from \{I think\} to \{\st{I think}\}. The cancellation mark performs \textit{sous rature}---Derrida's under-erasure. When we write \st{Being}, we are not asserting ``Being is and is not'' as a true contradiction. We are performing a different kind of move: acknowledging that the term is inadequate while simultaneously indicating we have no better alternative.

The substitution of \{ for `` and \} for '' is itself a performance of \textit{sous rature}. The braces mark where quotation was invoked and then withdrawn. They are the formal trace of an erasure. The notation \textit{enacts} the paradox of representing the unrepresentable rather than \textit{asserting} it as a true contradiction.

This allows the formal system to remain logically consistent. Within the mathematics built on the null representation, there are no contradictory theorems. The standard rules of inference apply without modification. The empty set behaves exactly as set theory requires: $\emptyset \subseteq A$ for any set $A$, $\emptyset \cup A = A$, $|\emptyset| = 0$, and so on. None of these formal results are compromised.

What is contradictory is not any mathematical statement but the \textit{ground} of the system. The null representation $\emptyset$ stands at the foundation as the mark of what cannot be marked, the representation of what recedes from representation. But this contradiction is \textit{contained} in the symbol itself through the structure of anaphoric recollection. The symbol points beyond itself to what it cannot directly capture, making the limitation explicit rather than denying it.

This is why I sometimes call $\emptyset$ the \textit{active set}. It is not a static object that simply is or is not. It is the trace of a dynamic movement---the movement of thought negating its own fixity, the refusal of the ``I think'' to be captured as an ``it thinks.'' When formalized, this dynamism appears static (as all formalization must), but the formalization points back to the dynamic source.

\subsection*{Anaphoric Recollection and Non-Deflationary Hegelian Negativity}

The strategy of localizing contradiction in notation allows me to preserve what I consider an essential insight of Hegelian dialectics: that the negative is not absence or logical incompatibility but the active, productive engine of development. This is what I call non-deflationary Hegelian negativity.

Some contemporary interpreters, most notably Robert Brandom, offer what might be called deflationary readings of Hegel. Brandom reconstructs Hegelian contradiction as \textit{material incompatibility}---the incompatibility of conceptual commitments within a practice of giving and asking for reasons \parencite{Brandom:2019aa}. On this reading, when Hegel speaks of contradiction, he means something like: ``These two claims cannot both be maintained within a coherent inferential structure.'' Contradiction becomes a pragmatic notion about rational commitments rather than a metaphysical claim about reality.

This reading has virtues: it makes Hegel compatible with classical logic, avoids mystification, and connects Hegelian dialectics to contemporary philosophy of language. Yet it perhaps loses something important. For Hegel, the negative is not merely about incompatibility between distinct claims. It is about the self-movement of concepts, the way determinate content generates its own negation and transformation. The negative is productive, not restrictive.

My approach preserves this productive negativity through the concept of anaphoric recollection. When we write $\emptyset$, the symbol does not assert ``there is nothing here.'' It \textit{recollects} the negative act that generated the emptiness. The braces are the trace of quotation (recollection) that was performed and then withdrawn. They point back to the dynamic movement of self-negation.

Anaphoric recollection resolves the apparent contradiction by introducing temporal movement into the formal structure. When a term refers back to itself, it is a dynamic process of recollection. The term in the consequent refers back to the term in the antecedent across a temporal gap inherent in the structure of inference. This gap, introduced by the structure of anaphora, prevents the collapse into static paradox.

This is the essence of Hegelian negativity: the apprehension of contradiction as movement. This negativity is the conceptual articulation of the rhythm practiced in Chapter 1. The movement of determinate negation---where a concept is negated in a way that preserves and elevates its content---is the same movement experienced in The Exercise as the transition from the first ``no'' (fixation/compression) to the second ``no'' [ne] (release/expansion). Zero embodies this movement. It is the pulse of negation that drives the unfolding of the number system, preserving the history of operations while simultaneously opening the space for new determinations. The contradiction is sublated---negated in its immediate form while its essential content is preserved and elevated. This allows contradiction to be productive, driving the expansion of the mathematical universe without collapsing into incoherence.

This is why the null representation can serve as the foundation for von Neumann's construction of natural numbers. When von Neumann defines $0 := \emptyset$ and then $1 := \{0\} = \{\emptyset\}$, he is not stipulating definitions. He is iterating the recollective structure: $1$ recollects $0$, which itself recollects the enabling condition. Each number becomes a nested structure of recollections, with $\emptyset$ at the base as the formalization of the first act of recollection.

The anaphoric structure ensures that even when the negative is rendered formal---even when it becomes ``it'' rather than remaining the spontaneous \{I\}---it retains a pointer back to its dynamic source. The symbol $\emptyset$ is the fixed trace, but like a photograph of a dancer, it captures a moment of movement without eliminating the movement. We recognize the photograph shows motion even though the image is static.

This is the sense in which my approach is non-deflationary. The negative is not reduced to material incompatibility or pragmatic tension. It remains the active principle of self-determination, the movement of thought negating its own fixity. But this movement is made formally operative through notational structures (the substitution, the braces, anaphoric recollection) that localize contradiction while preserving its productive force.

\subsection*{Comparison with Formal Set Theory}

This approach is compatible with standard set theory. Zermelo's axiomatization includes the Axiom of the Empty Set (Axiom II): ``There exists an (improper) set, the null set $0$, which contains no elements at all'' \parencite{zermelo1908investigations}. This axiom is typically justified pragmatically: we need a starting point for building the universe of sets, and $\emptyset$ serves that purpose.

Zermelo also provides the Axiom of Separation (Axiom III), which allows subsets to be formed based on definite properties. Using this axiom, one can define $\emptyset = \{x \mid x \neq x\}$---the set of all things that are not self-identical. Since everything is self-identical (by the law of identity), this set is empty.

Notice that this definition uses a logical contradiction ($x \neq x$) \textit{instrumentally}. The contradiction guarantees emptiness because no object satisfies the condition. But the contradiction is not affirmed as true; it is used as a technical device to ensure the resulting set is empty. This is consistent with classical logic: the empty set is not itself contradictory, even though its definition employs a necessarily false condition.

My approach is compatible with this. The null representation $\emptyset$ can be defined formally as Zermelo defines it, with the addition that we understand the semantic content differently. Where standard set theory treats $\emptyset$ as simply ``the set with no elements,'' I treat it as ``the formal trace of recollecting the enabling condition that recedes.'' This additional semantic layer does not conflict with any formal results.

The advantage of my approach is that it grounds the axiom. Rather than simply positing that the empty set exists, I show why it must exist: because mathematical thought, like all thought, requires recognition of its own enabling conditions, and those conditions formalize as $\emptyset$. The ``axiom'' becomes a theorem in the logic of self-consciousness.

\section{Conclusion: The Ground That Grows}

Out of many, one.\footnote{I write this as quotation of a national paradox, not a nationalistic sentiment.} The fact that base systems can express arbitrarily large numbers is unrelated to the true infinity that finds its partial expression as soon as one counts to ten. I say partial because the true infinity is not at all captured in the numeral $10$, especially as I have not yet discussed how the \{I\} is involved in the process of counting.

Zero emerges pragmatically from counting's internal logic, semantically through philosophical traditions that make formalization intelligible, and critically by localizing contradiction in notation itself. These three movements are mutually constitutive---the pragmatic necessity drives semantic formalization, which requires critical clarification, which grounds pragmatic practice.

The pattern from the Bridge appears here in fundamental form. Tally systems attempt totality (count everything with marks). Grouping reveals limitation (no unity principle, no organized structure). Base systems totalize again (recursive grouping can represent any number). But this totalization requires zero---the element that is both present and absent, that marks incompleteness while enabling completeness, that grounds the system precisely by negating itself.

From this ground, numerals now grow. Chapter 7 will show how numerals function as first-person pronouns, anaphorically recollecting the enabling conditions of thought through quotative embedding. But that recollection depends on the ground established here: zero as the null representation that makes all representation possible.

Like in ``Roll it Away'' where I strip scales to dive into blue depths, we have descended to the ground of mathematical becoming. We have found that the ground is not solid but self-negating, not static but constantly in motion, not being but becoming. And from this unrepresentable ground---this nothingness that is the condition of somethingness---numerals will emerge as pronouns recollecting the temporal-spatial movement of thought reflecting on itself.

The \textit{in}finite has introduced us to zero. Now zero introduces us to number.

\printbibliography[heading=subbibliography]


\chapter{Numerals are Pronouns}

\begin{abstract}
This chapter presents the book's central mathematical claim: numerals and number words function as first-person anaphoric pronouns that recollect the structure of self-consciousness, rather than as names for abstract objects. The argument is catalyzed by a student's question---``What even is two?''---which reveals the inadequacy of formal definitions to meet existential needs. The chapter develops the claim that numbers symbolize the unrepresentable ``I think'' that makes all representation possible. By grounding mathematical understanding in the act of self-recognition, the analysis proposes a new foundation for arithmetic that motivates the pursuit of correctness through the desire for coherent selfhood rather than correspondence to an external reality.
\end{abstract}

\section{Introduction}

What are numbers?  This seemingly simple question became agonizingly complex when voiced by a fifth-year senior, $\mathcal{I}$, in Indianapolis. He couldn't graduate with his friends - couldn't move on with his life - because he could not pass a high-stakes standardized test of algebraic knowledge. I was teaching a test-prep course at a high-needs urban school, and the classroom was full of super-seniors crushed by systemic forces. I felt like a conduit of those forces who had a choice to either enact the expectations of my position or treat the people around me as people. You see, it was the first week of class in my first year as a high school teacher. The night before, one of $\mathcal{I}$'s friends was killed in a car accident. $\mathcal{B}$ was fleeing the police when his SUV rolled over, killing him. The night before beginning his super-senior year, where he would have been taking the same test prep class as $\mathcal{I}$, $\mathcal{B}$ was gone. The students were grieving their friend and bitter about taking another stupid algebra class. I had the choice to listen and grieve, but I chose to jump right into algebra. As I yammered on about $x$ and $y$, $\mathcal{I}$ seemed attentive to the lesson, where other students were not; he was a kind of ally in a hostile room. When he looked up and said ``Mr. Savich, what even is two?,'' I wanted to share my most beautiful truth. I got out my markers and started drawing von Neumann ordinals and lectured about the nothingness embedded everywhere you look. 

Any sense of allyship slipped away as the light of connection withdrew. The connection between numbers and the empty set had been a revelatory moment for me when I studied pure mathematics at Earlham College. I could not `unsee' the nothingness that connected all forms of thought because it was not `there.' But whatever spiritual import that revelation had was lost in the severed connection I experienced with $\mathcal{I}$. I thought I should be able to answer such a simple question somewhat simply. It reinforced a disturbing suspicion: mathematical knowledge, far from being a universal language of reason, can appear as an arcane tool of subjugation, a barrier rather than a bridge. As a math teacher, my role became clear: I was the instrument of suppression. Sick at heart with uselessness, what had once been a beautiful practice became hard to sustain. I eventually burned out of teaching, though the fundamental problem of what we are talking about when we talk about numbers lingered. 

I was given the opportunity to reflect on the problem and found an answer that I wish to share with $\mathcal{I}$, though I do not know where he lives or who he has become in the intervening decade. Reflecting deeply on $\mathcal{I}$'s question, I arrived at an answer that offers, if not certainty, then at least a shift in mood: numerals, like ``2,'' and number words like ``two,'' are best understood as \textit{first-person pronouns.}  This the central idea of this chapter, and it offers a radical departure from traditional ways of thinking about the ontology of number, as I propose number exists as \{I\} or \{you\} exists: number is an autonomous subject.  But what does it mean for a numeral to be a pronoun? 

Chapter 6 established zero as the ground---the null representation from which numerical structure grows. Now we trace how numerals emerge from that ground through the iterative process of quotative embedding and anaphoric recollection. If zero is the enabling condition, numerals are the recollective pronouns that point back to those conditions while building structure upon them.

\section{The Null Representation and Its Role in Chapter 7}\label{def:null-representation}

Chapter 6 established the null representation $\emptyset$ as the formalization of the enabling condition that recedes---the unrepresentable ``I think'' that must accompany all representations. This section briefly recapitulates that development to clarify how the null representation functions within this chapter's argument about numerals.

The null representation arose through a specific substitution: replacing quotation marks with set braces. The progression was:


``I think'' $\rightarrow$ \{I think\} $\rightarrow$ \{\st{I think}\} $\rightarrow$ \{\} $\rightarrow \emptyset$

Each step follows necessarily from the structural requirements of representing the unrepresentable. The substitution of \{ for `` and \} for '' makes recollection formally operative, transforming a semantic marker (quotation) into a mathematical operation (set formation).

The null representation is not ``nothingness'' in a simple sense of absence. It is the \textit{active set}, the formalization of the self-negating ground of thought. It marks the site where consciousness touches mathematics---where the ``I think'' that enables all judgment appears within formal notation while simultaneously indicating its own inadequacy as representation.

Chapter 6 also established the divaded structure of apperceptive consciousness: when I make a judgment like ``Pokey is a dog,'' I am simultaneously inside (performing synthesis at the copula) and outside (aware of the judgment as a unified whole). This structure was formalized as:

\[
\{\text{Pokey } \emptyset \text{ is a dog}\}
\]

The $\emptyset$ inside marks the synthetic activity at the copula. The braces \{\} outside mark the recognition of the judgment as a unified object. This inside/outside symmetry is the foundation from which numerals emerge.

For the purposes of this chapter, the key insight is that $\emptyset$ represents not a static void but the \textit{enabling condition that recedes}---what we are always already inside when we attempt to examine it. It is zero as the ground from which counting proceeds, the unmarked origin that makes marking possible.

The null representation, $\emptyset$, marks the limit of representation itself. This symbol indicates the enabling condition for any determination whatsoever. It corresponds to the unrepresentable \{I\} explored in earlier chapters---the subjective source of action that remains inaccessible to direct awareness.

When we articulate any mathematical structure, we do so from this position. The $\emptyset$ is the notation for that position, the ground that recedes when we attempt to grasp it. It is the absence that makes presence possible. To mistake this symbol for an object is to mistake the condition of possibility for one of its outcomes. It is this unrepresentable ground, not an abstract object, that the numerals recollect. It is the formal indication of the \textit{in}finite potential that precedes any finite determination.

With this foundation established in Chapter 6, this chapter can focus on how numerals emerge through the \textit{iteration} of this structure. We do not need to re-derive the null representation here; we need only recognize it as the starting point (0 = $\emptyset$) from which successive acts of recollection generate the sequence of natural numbers.

\section{Anaphora and Mathematical Self-Reference}\label{def:anaphora}

Having established that numerals function as first-person recollective pronouns, we must clarify the mechanism through which this recollection operates. The key concept is \textit{anaphora}---the linguistic phenomenon through which pronouns create repeatability and co-reference across discourse. Understanding anaphora as a primordial form of linguistic self-reference helps explain how they enable mathematical thought.

Pronouns are fundamentally anaphoric terms. The word ``it'' in ``This chalk is white. \textit{It} is also cylindrical'' refers back to ``this chalk,'' creating a stable reference that can be used repeatedly throughout discourse \parencite[150]{Brandom:2019aa}. Anaphora solves a fundamental problem: indexical expressions like ``this'' are inherently context-dependent, their reference shifting with each new utterance. Without anaphora, discourse would be trapped in perpetual immediacy, unable to build upon prior content.

What makes anaphora profound is that it is a fundamental form of \textit{linguistic self-reference}. When language uses anaphoric terms, it turns back on itself, allowing words to refer to other words or phrases already present within the discourse. Anaphora is how language refers to its own prior utterances, creating chains of reference within itself rather than pointing outward to the world.

Consider quotation, the paradigmatic form of linguistic self-reference. When I say ``She said `I am hungry,''' the inner quotation marks create an anaphoric structure. The quoted phrase does not refer to my current state but recollects her prior utterance. Quotation transforms the indexical ``I'' from her speech into a repeatable token within my discourse. The marks themselves---`` and ''---function as operators that reach back to extract content from a previous moment and make it present again.

Recall that quotations can contain quotations, creating nested recollection: ``He said `She said ``I am hungry''''' creates a two-level embedding. Each layer of quotation recollects a prior speech act, building a temporal-sequential structure within language itself. When we trace back through nested quotations, we reconstruct a history of utterances, with the innermost quote serving as the originary utterance in this local context.

Variables in algebra function similarly. Once a variable like $x$ is introduced and defined (``Let $x$ be the number of apples''), it becomes a point of self-reference within the mathematical discourse. The variable can be used repeatedly throughout an equation or proof, consistently referring back to that initial definition within the symbolic system itself. Mathematical notation, in this light, is a formalized system of linguistic self-reference.

Anaphora mirrors the structure through which identity is constituted and maintained across time. When we employ a pronoun, we are performing an act of recollection that binds the present utterance to a past antecedent. This binding requires acknowledging that the term used now inherits the commitments associated with the term used before.

This structure is the experience of selfhood. The self is structured by the dynamic between the acting \{I\} and the recognized ``me.'' The pronoun functions in the gap between them. It provides syntactic continuity---the word ``I'' remains the same---while acknowledging the semantic difference---the content of what ``I'' refers to constantly shifts.

Mathematical practice relies on this same structure for its stability. When we utilize a mathematical term, we are relying on its anaphoric connection to its previous uses. The term gains its stability by pointing backward within the discourse to the network of inferential commitments it inherits. The stability of the numeral is structurally identical to the stability of the self. Both depend on the reliability of our own recollective practices.

Numerals, I argue, operate at an even more fundamental level of anaphoric recollection. They are variablare variablesm, nor are they simple indexical terms. Instead, numerals are anaphoric recollections of the enabling conditions of thought itself---the structures formalized as von Neumann ordinals.

When I use the numeral ``3,'' I am not naming an abstract object or pointing to a collection of things. I am performing an anaphoric gesture that recollects the structure \{0, 1, 2\}---the formalization of three acts of synthetic recollection, each gathering all prior acts and the transcendental ground that enables them. The numeral points back, not to empirical content, but to the iterative structure of self-conscious thought.

This suggests that \textit{mathematics itself, built upon numerals, arises as language recollects itself}. When discourse turns back on itself through sufficient levels of self-reference---when quotation nests within quotation, when reflection reflects on reflection---mathematical structure emerges. Mathematics is not an external realm of abstract objects but an internal development of language's capacity for self-reference, emerging when discourse recognizes and formalizes its own fundamental structures.

This understanding has expansive implications. It means that the ``objectivity'' of mathematics does not require a Platonic realm. Mathematical truths are objective because the structures they formalize---the patterns of anaphoric recollection, the iterative nature of self-conscious thought---are universal features of any being capable of discourse. Anyone engaging in linguistic self-reference will encounter these same structures. Numbers are discovered, not invented, but what we discover is the structure of our own thinking when it turns back on itself with sufficient clarity.

Robert Brandom's analysis of anaphora emphasizes that it creates what he calls ``repeatability structures'' \parencite[549]{brandom1994making}. Anaphoric terms allow content that would otherwise be fleeting and context-bound to become stable and repeatable across different contexts. This precisely what numerals accomplish: they stabilize the iterative structures of recollection, making them available as stable objects of thought that can be manipulated, compared, and operated upon.

But numerals go beyond ordinary anaphora in one crucial respect: they recollect not particular empirical contents but the \textit{form} of recollection itself. When I say ``it,'' I recollect some specific thing mentioned previously. When I say ``three,'' I recollect the form of having recollected twice, the structure of triple synthesis. Numerals are meta-anaphoric: they are anaphoric terms that point to the structure of anaphora itself.

This explains the strange ontological status of numbers---why they seem both mind-dependent and mind-independent, both constructed and discovered, both in language and beyond it. Numbers formalize the structure through which mind turns back on itself. They are immanent in thought yet universal across thinkers because the structure of self-reference is itself universal.

Anaphora becomes especially powerful when we shift focus from third-person pronouns to first-person pronouns. ``I'' is a linguistic self-reference \textit{par excellence}---it is the subject referring to itself as subject. When combined with quotation, first-person pronouns enable the ultimate form of self-reference: the subject reflecting on its own prior thoughts. ``I thought `I am hungry''' is a self-reference that contains its own past as content, enabling consciousness to build iterative structures of self-awareness.

It is this combination---anaphora + quotation + first-person---that generates the recollective structures that become formalized as numbers. Numerals, I argue, are first-person anaphoric pronouns that recollect the enabling conditions of prior acts of thought, building systematic structure from self-reference.

\section{Two Forms of Recollection in Natural Language}

Natural language exhibits at least two distinct forms of recollection: sequential and synthetic. These are not abstract philosophical categories but observable features of how we use language to refer back to prior utterances and thoughts. Understanding these two forms is crucial because they generate the two mathematical constructions of ordinal numbers.

\subsection*{Sequential Recollection: The Telephone Game and the Trace}

Imagine playing Telephone with my family. Let $\mathcal{T}$ be me, $\mathcal{A}$ be my wife, $\mathcal{M}$ and $\mathcal{E}$ be my daughters. I start with the message: ``Numerals are pronouns.''

The transmission proceeds:
\begin{itemize}
\item $\mathcal{T}$ whispers to $\mathcal{A}$: ``Numerals are pronouns''
\item $\mathcal{A}$ whispers to $\mathcal{M}$: ``Numerals are nouns'' (mishearing)
\item $\mathcal{M}$ whispers to $\mathcal{E}$: ``Numerals are names'' (further distortion)
\item $\mathcal{E}$ announces: ``Numerals are names''
\end{itemize}

The humor and point of Telephone lie in the iterative distortion. But for our purposes, notice the \textit{structure} of quotative recollection when we trace backward. 


Even if the message transforms through mishearing (perhaps $\mathcal{E}$ announces ``Numerals are names''), the structure remains: each player recollects only the immediately previous utterance. Each speaker only holds the immediately previous utterance as content. This pure iteration---each act wraps the prior act in a new layer of recollection.

But we must pause here and attend more carefully to what this game reveals. The Telephone game doesn't illustrate sequential wrapping -- it teaches us something important about the nature of meaning and origin through what Derrida called \textit{différance}.

\subsubsection*{The Phenomenology of Loss and Différance}

Each act of transmission in the Telephone game is simultaneously an act of interpretation. When $\mathcal{A}$ hears $\mathcal{T}$'s utterance, she doesn't receive a pristine copy. She hears sounds filtered through breath, accent, the rustle of clothing, ambient noise. She must reconstruct meaning from this acoustic flux. This reconstruction introduces difference -- not as accidental error but as a structural necessity of communication itself.

Derrida's term \textit{différance} captures two interwoven aspects of this process:

\begin{enumerate}
\item \textbf{Difference}: Each transmission necessarily introduces variation. What $\mathcal{A}$ whispers differs from what $\mathcal{T}$ said, even when $\mathcal{A}$ tries to repeat it exactly. The medium itself -- sound waves, neural processing, articulation -- ensures that perfect repetition is impossible.

\item \textbf{Deferral}: Access to the origin is endlessly postponed. $\mathcal{M}$ cannot reach back to verify what $\mathcal{T}$ ``really'' said. She has access only to $\mathcal{A}$'s version, which is already a transformation. The origin recedes through each mediating layer.
\end{enumerate}

The game thus reveals that the origin -- $\mathcal{T}$'s original utterance -- functions as what Derrida called the \textit{trace}. It is never directly present for subsequent players yet structures the entire chain. Each utterance carries the trace of the origin without the origin ever being accessible. The origin is ``always already'' withdrawn, existing only in its effects.

Do you recall my `first word,' ``battery''? It is unlikely objectively true that my first word was ``battery.'' But there is no way to recover the `truth.' 

\subsubsection*{The Elusive Ground of Communication}

When we try to grasp the origin directly -- when $\mathcal{E}$ asks ``What did $\mathcal{T}$ actually say?'' -- we encounter a fundamental impossibility. We can work backward through the chain, asking each player what they heard, but this archaeological project never reaches the origin itself. It only accumulates mediations. Even if we could ask $\mathcal{T}$ directly, he would have to speak again, creating a new utterance that is not identical to the original.

This impossibility is not a practical limitation we might overcome with recording equipment. It reveals something essential about how meaning works. The ``original meaning'' that seems to ground the chain of transmission can never be made fully present. It exists only as trace -- as that which must have been for the chain to exist, yet which cannot be recovered in its purity.

The degradation observed in the Telephone Game reveals something about the structure of representation itself. Each repetition introduces a difference, a deferral of the original meaning. This structure, what Derrida calls the Trace, ensures that the origin remains perpetually inaccessible. We encountered this structure directly in the practice of self-certainty (Chapter 1), where the attempt to grasp the immediate present revealed that the ``now'' is always already a ``then.'' Sequential recollection enacts this temporal deferral.

In this mode, meaning is always downstream from an absent origin. The anxiety of the Telephone Game is the anxiety of the Trace---recognizing that meaning is never fully present yet always mediated through a chain of signifiers. This form of recollection highlights the fragility of meaning when conceived as linear transmission rather than cumulative integration.

The Telephone game thus teaches us about consciousness itself. When I try to grasp the ``I think'' that grounds my thoughts, I find only traces -- memories of having thought, awareness of thinking now -- but never the ``I think'' itself in its immediacy. It withdraws even as it enables each act of reflection, leaving only its trace in successive iterations.

\subsection*{Synthetic Recollection: Cumulative Gathering}

While the Telephone game reveals sequential recollection, the structure of judgment exhibits a different form: synthetic recollection. When I judge ``Pokey is a dog'' and then reflect on this judgment, I do not merely wrap the prior thought. I perform a more complex operation.

Consider: ``I think that (I thought `Pokey is a dog').'' This reflection must account for:
\begin{itemize}
\item The transcendental ground that enables the current act of reflection (the ``I think'')
\item The prior thought as a determinate unity (\{prior content\})
\end{itemize}

The structure is cumulative, not merely successive. Each act of synthetic recollection gathers up all prior conditions and holds them together in a new unity. This Hegelian---characterized by \textit{Erinnerung} (internalization/recollection), where the past is preserved and elevated within the present rather than endlessly deferred.

In synthetic recollection, \textit{the current thought contains its entire history}. When I reflect on having thought about Pokey, I do not merely access the prior thought; I hold both the ground that enabled that thought and the thought itself within my current synthesis. The history is multidimensional, cumulative, a gathering-up of all that came before.

This cumulative gathering is a process of integration where each new element transforms the whole while preserving what came before. This the structure of \textit{Aufhebung}, or sublation---the dialectical movement where a concept is simultaneously negated, preserved, and elevated. In synthetic recollection, the past is incorporated into a richer, more determinate present.

This mode of recollection structures the development of self-consciousness itself. We are a cumulative history, where past experiences are constantly reinterpreted and integrated into our present understanding. Synthetic recollection is the mechanism through which Spirit (\textit{Geist}) unfolds in time, gathering its history into an ever-expanding unity. It is the movement of the Concept, where determination proceeds through the incorporation of difference.

\subsection*{Two Forms, Two Structures, Two Philosophies}

These two forms of recollection are not different---they correspond to fundamental dimensions of human experience and the existential needs that structure our engagement with the world (Chapter 3). Sequential recollection, with its emphasis on the trace and the unique event, resonates with the need to be recognized as \emph{in}finite---the desire for our authentic, spontaneous selfhood to be acknowledged in its irreducibility. It honors the particularity of the moment.

Synthetic recollection, by contrast, with its emphasis on cumulative gathering and stable structures, resonates with the need to be recognized as good---the desire for our actions and our being to be validated within shared normative frameworks. It seeks coherence and the formation of a reliable identity. The tension between these two forms is the tension between authenticity and normativity, between the spontaneous \{I\} and the social ``me.'' A critical mathematics must account for both, recognizing that they are dialectically unified moments in the unfolding of rational thought.

We now have two distinct linguistic phenomena, each exhibiting a different temporal-recollective structure:

\begin{center}
\begin{tabular}{p{0.45\textwidth}p{0.45\textwidth}}
\textbf{Sequential Recollection} & \textbf{Synthetic Recollection} \\
\hline
Linear succession & Cumulative synthesis \\
Each wraps the previous & Each gathers all previous \\
Derridean (trace, différance) & Hegelian (Aufhebung, Erinnerung) \\
Telephone game structure & Judgment structure \\
One-dimensional chain & Multidimensional unity \\
Origin recedes & Origin preserved \\
\end{tabular}
\end{center}

These are observable features of natural language. The question becomes: what happens when we formalize these structures mathematically?

\section{Formalizing Recollection: From Natural Language to Ordinals}

Having identified two forms of recollection in natural language, we now formalize them using the null representation $\emptyset$ established in Chapter 6. This formalization reveals that the two linguistic structures correspond to two mathematical constructions of ordinal numbers: the Zermelo ordinals and the von Neumann ordinals.

\subsection*{Zermelo Ordinals: Formalizing Sequential Recollection}

We begin with sequential recollection as exemplified by the Telephone game. Recall that each act of recollection simply wraps the previous utterance in new quotation marks. From Chapter 6, we know that quotation marks can be replaced with set braces through the substitution: \texttt{``} $\rightarrow$ \{ and \texttt{''} $\rightarrow$ \}.

Let us formalize the Telephone game step by step, starting from the null representation as the ground:

\textbf{Zero (0)}: Before any utterance, there is the enabling condition---the unrepresentable ``I think'' that makes any speaking possible. This formalized as the null representation:
\[
0 := \emptyset
\]

\textbf{One (1)}: The first utterance occurs. $\mathcal{T}$ says ``Numerals are pronouns.'' This a judgment, which Chapter 6 taught us to formalize as \{Pokey $\emptyset$ is a dog\}. But when we abstract away the specific content and attend only to the \textit{form} of a single judgment, we find: a unity (the braces) constituted by the synthetic activity (the internal $\emptyset$). Stripping away content:
\[
1 := \{\emptyset\}
\]

\textbf{Two (2)}: $\mathcal{A}$ recollects $\mathcal{T}$'s utterance: ``$\mathcal{T}$ said `Numerals are pronouns'.'' Following sequential recollection, this a wrapping of the prior structure. We take the form of the first judgment (\{\emptyset\}) and enclose it in new braces:
\[
2 := \{\{\emptyset\}\}
\]

\textbf{Three (3)}: $\mathcal{M}$ recollects $\mathcal{A}$'s utterance: ``$\mathcal{A}$ said `$\mathcal{T}$ said ``Numerals are pronouns'''.'' Again, this pure wrapping:
\[
3 := \{\{\{\emptyset\}\}\}
\]

The pattern is clear. The \textit{successor function} in this formalization is simply:
\[
S_Z(n) = \{n\}
\]

Each number wraps the previous in a new layer of braces. This generates the \textbf{Zermelo ordinals}:
\[
0 = \emptyset, \quad 1 = \{\emptyset\}, \quad 2 = \{\{\emptyset\}\}, \quad 3 = \{\{\{\emptyset\}\}\}, \quad \ldots
\]

\textbf{The Trace Structure of Number}: The Zermelo construction captures something essential -- how the origin perpetually withdraws. Just as $\mathcal{T}$'s original utterance becomes increasingly inaccessible through layers of transmission, the empty set $\emptyset$ becomes increasingly buried under successive braces. To reach the origin from 3 = \{\{\{\emptyset\}\}\}, we must unwrap three layers, but the origin we reach is itself empty -- the trace of an enabling condition that was never fully present.

This construction respects what we might call the dignity of absence. It acknowledges that not everything can or should be made present, that some things exist precisely in their withdrawal. The Zermelo ordinals formalize the trace-structure of thought itself.

This construction captures the structure of the Trace. Each numeral is a container holding the previous numeral, emphasizing the pure difference between successive stages. The content of the previous stage is entirely enclosed, rendered inaccessible except as the antecedent of the current stage. Zermelo's construction privileges the act of iteration over the content being iterated. It is a formalism of pure sequence, where each element is defined solely by its position in the chain, detached from accumulated meaning.

This structure embodies a philosophy of difference, where identity is constituted through separation. It mirrors the anxiety of the Telephone Game where the original message is perpetually deferred. This formalism, while elegant in its minimalism, risks reducing history to a sequence of empty containers, reflecting an alienated mode of consciousness where the self is experienced as discrete moments without internal unity.

\subsection*{Von Neumann Ordinals: Formalizing Synthetic Recollection}

We now turn to synthetic recollection as exemplified by the structure of judgment. When I reflect on a prior judgment, I must account for both the prior judgment as unity \textit{and} the transcendental ground that enables the current act of reflection. This \textit{Aufhebung}---sublation that preserves and elevates.

Using the inside/outside symmetry from Chapter 6, we can formalize how synthetic recollection generates von Neumann ordinals.

\textbf{Zero (0)}: The transcendental ground, the null representation:
\[
0 := \emptyset
\]

\textbf{One (1)}: A single judgment occurs. As before, abstracting the form of judgment:
\[
1 := \{0\} = \{\emptyset\}
\]

\textbf{Two (2)}: I reflect on the prior judgment: ``I think that (I thought `Pokey is a dog').'' Following the analysis from Chapter 6, this new judgment is itself an act of thinking, which means it depends on the ``I think'' ($\emptyset$) as its enabling condition. But this new judgment also takes the prior judgment (\{$\emptyset$\}) as its object.

The structure of this reflection must synthesize:
\begin{itemize}
\item The transcendental ground that enables the current thought: $\emptyset$ (or $0$)
\item The prior judgment as a completed unity: \{$\emptyset$\} (or $1$)
\end{itemize}

This not sequential addition but \textit{simultaneous holding-together}. The form of this synthetic recollection is the \textit{union} of the ground with the prior structure:
\[
2 := 0 \cup \{1\} = \emptyset \cup \{\{\emptyset\}\} = \{\emptyset, \{\emptyset\}\} = \{0, 1\}
\]

\textbf{Three (3)}: I reflect on this second-order thought. Again, the new thought requires both:
\begin{itemize}
\item The transcendental ground: $\emptyset$ (or $0$)
\item The prior structure as unity: \{$0$, $1$\} (or $2$)
\end{itemize}

But crucially, because the prior structure \textit{already contains} $0$ and $1$, and because the new thought depends on all these conditions, the synthesis must gather them all:
\[
3 := 0 \cup 1 \cup \{2\} = \{\emptyset, \{\emptyset\}, \{\emptyset, \{\emptyset\}\}\} = \{0, 1, 2\}
\]

The pattern reveals itself. The \textit{successor function} in synthetic recollection is:
\[
S_{vN}(n) = n \cup \{n\}
\]

Each number is the union of all predecessors with the set containing the immediate predecessor. This generates the \textbf{von Neumann ordinals}:
\begin{align*}
0 &= \emptyset \\
1 &= \{0\} = \{\emptyset\} \\
2 &= \{0, 1\} = \{\emptyset, \{\emptyset\}\} \\
3 &= \{0, 1, 2\} = \{\emptyset, \{\emptyset\}, \{\emptyset, \{\emptyset\}\}\}
\end{align*}

\textbf{The Presence of History}: The von Neumann construction embodies \textit{Aufhebung}. Each number is not just the last iterative act but the synthesized collection of all prior acts. The number literally \textit{contains its own history}. This Hegelian sublation: the past is negated (as immediate present), preserved (as element within the set), and elevated (into a new structural whole).

Here we encounter the formal structure of \textit{Aufhebung}. Each numeral is the set of all its predecessors, a cumulative gathering that preserves the entire history of its construction. The past remains present as constitutive content of the current stage. Von Neumann's construction privileges the content accumulated over the pure act of iteration. It is a formalism of synthesis, where each element is defined by the totality of what came before.

This structure embodies a philosophy of integration, where identity is constituted through the inclusion of difference. This the mathematics of Spirit (\textit{Geist}), where the Concept unfolds through the cumulative incorporation of its own history. It is a formalism that honors the synthetic unity of apperception---the integration of diverse experiences into a unified self---recognizing that the present is the recollection of the past. The Von Neumann construction provides a mathematical model for the development of consciousness itself.

\subsection*{The Telephone Game Reconsidered: Two Readings}

But now we must reconsider the Telephone game with greater care. While the phenomenology of transmission suggests pure iteration (Zermelo), we might ask: what are the enabling conditions for $\mathcal{E}$'s final utterance?

When $\mathcal{E}$ announces ``Numerals are names,'' what enables this announcement? She must have:
\begin{itemize}
\item Heard from $\mathcal{M}$ (condition 2)
\item Who heard from $\mathcal{A}$ (condition 1)  
\item Who heard from $\mathcal{T}$ (condition 0, the original ground)
\end{itemize}

All three prior conditions are \textit{simultaneously necessary} for $\mathcal{E}$'s utterance. If any link in the chain had not occurred, her announcement would be impossible. From this perspective, the structure of enabling conditions is cumulative: \{0, 1, 2\}.

We thus find two valid readings of the same phenomenon:

\textbf{The Zermelo reading} attends to the lived experience of transmission -- each player knows only the immediately previous utterance. The origin recedes through layers of mediation. This captures the phenomenology of communication as loss and transformation.

\textbf{The von Neumann reading} attends to the logical structure of enabling conditions -- all prior utterances are simultaneously necessary for the final one. The entire history is implicit in each moment. This captures how the past persists as condition for the present.

Neither reading is ``more correct.'' They reveal different aspects of the same structure -- the tension between the trace that withdraws and the history that accumulates. This tension characterizes not just the Telephone game but the structure of consciousness itself.

\subsection*{Two Constructions, Not One Better Than the Other}

We now have both constructions formalized:
\begin{itemize}
\item \textbf{Zermelo} ($\emptyset$, \{$\emptyset$\}, \{\{$\emptyset$\}\}, \{\{\{$\emptyset$\}\}\}, \ldots): pure iteration, wrapping
\item \textbf{Von Neumann} ($\emptyset$, \{$\emptyset$\}, \{$\emptyset$, \{$\emptyset$\}\}, \{$\emptyset$, \{$\emptyset$\}, \{$\emptyset$, \{$\emptyset$\}\}\}, \ldots): cumulative synthesis
\end{itemize}

Both emerge from the same ground -- the null representation that formalizes the enabling condition that recedes. Both are valid formalizations of recollection. The question is not which is correct but what each reveals.

\textbf{Zermelo ordinals} capture:
\begin{itemize}
\item The act of counting as pure iteration
\item The recession of the origin through successive layers
\item The trace-structure of meaning
\item The phenomenology of loss in communication
\item How number emerges from the impossibility of presence
\end{itemize}

\textbf{Von Neumann ordinals} capture:
\begin{itemize}
\item Number as achieved totality containing its history  
\item The persistence of all enabling conditions
\item The systematic unity of mathematical structure
\item The gathering of the past into the present
\item How number emerges from synthetic recollection
\end{itemize}

The tension between these constructions -- between trace and presence, iteration and synthesis, Derrida and Hegel -- is not a problem to be resolved. It reflects the fundamental duality in how consciousness relates to its own past. Sometimes the past recedes (when I try to remember my childhood), sometimes it gathers (when I recognize how all my experiences constitute who I am now). Both modes are essential to thought.

\subsection*{Numerals as First-Person Pronouns}

We can now state the full thesis: \textit{Numerals are first-person recollective pronouns that anaphorically recollect the structures generated by both sequential and synthetic recollection.}

When I use the numeral ``3,'' I am not naming a Platonic object. I am performing an anaphoric gesture that points back to structures of thought. But this pointing has two aspects:

\textbf{As Zermelo}: ``3'' recollects three acts of iteration, three successive wrappings: \{\{\{$\emptyset$\}\}\}. It marks the trace of thought having turned back on itself three times.

\textbf{As von Neumann}: ``3'' recollects the synthesis of three stages: \{$\emptyset$, \{$\emptyset$\}, \{$\emptyset$, \{$\emptyset$\}\}\}. It gathers the entire history of thought's self-reflection.

Numbers are pronouns for the enabling conditions of thought itself. They recollect:
\begin{itemize}
\item The ``I think'' that must accompany all representations ($\emptyset$)
\item The iterative acts of reflecting on prior thoughts
\item Either as trace (Zermelo) or as gathered totality (von Neumann)
\end{itemize}

This explains why numbers feel both subjective and objective, both constructed and discovered. They are subjective because they formalize the structure of self-conscious thought. They are objective because this structure is universal---anyone engaging in apperceptive consciousness exhibits these same structures. Numbers are not ``in the head'' (psychologism) nor ``out in the world'' (Platonism) but immanent in the activity of thinking itself.

The pronoun thesis thus offers a response to $\mathcal{I}$'s question that honors its full complexity. What is two? Two is the anaphoric recollection of thought having reflected on itself -- understood either as wrapped iteration \{\{$\emptyset$\}\} or as cumulative synthesis \{$\emptyset$, \{$\emptyset$\}\}. Two is not a thing but a \textit{way of recollecting} the structure of iterative self-consciousness.

Recognizing numerals as first-person pronouns transforms our relationship to mathematics. When a student asks, ``What even is two?,'' they are asking for a way to recognize themselves in the mathematical structure. The answer must be an invitation to recollect the process of self-constitution.

The Von Neumann construction, by formalizing synthetic recollection, offers this invitation. It reveals that the structure of number is the structure of the \{I\} unfolding in time. This realization addresses the existential need for recognition. It allows the student to find themselves in the mathematics. The pursuit of mathematical correctness is transformed from an external imposition---a demand for compliance---into an internal striving for self-coherence. It is the desire of the \{I\} to be recognized in the ``me,'' the desire for our spontaneous activity to align with our normative commitments. This the emancipatory potential of critical mathematics: recognizing that the structures of reason are the structures of our own being.

\section{The Successor Function: Two Forms of Mathematical Becoming}\label{successor-function}

One of the strengths of understanding numerals as first-person recollective pronouns is that it provides a natural definition of the \textit{successor function}---a concept often taken as a primitive axiom in many mathematical systems. In this framework, the successor function emerges from the interplay of recollection and synthesis, with its precise form depending on which type of recollection we formalize.

\subsection*{The Zermelo Successor: Pure Iteration}

The Zermelo successor function formalizes sequential recollection as pure wrapping:
\[
S_Z(n) = \{n\}
\]

To find the successor of any number $n$, simply enclose it in braces. This operation models the act of quotation: just as ``She said `X''' wraps the utterance X in new quotation marks, $S_Z(n)$ wraps the number $n$ in new set braces.

Starting from zero:
\begin{align*}
0 &= \emptyset \\
S_Z(0) = 1 &= \{\emptyset\} \\
S_Z(1) = 2 &= \{\{\emptyset\}\} \\
S_Z(2) = 3 &= \{\{\{\emptyset\}\}\}
\end{align*}

This successor function captures the \textit{act} of counting as pure progression. Each application adds exactly one layer, one iteration, one moment. It is minimalist and phenomenologically transparent: the notation directly visualizes ``one more step.''

The Zermelo successor emphasizes \textit{becoming} over \textit{being}. Each number is not a stable collection of elements but a trace of the iterative process---a marker of ``how many times have we iterated?'' The focus is on the \textit{performance} of counting rather than its \textit{product}.

Yet this transparency comes with trade-offs. The Zermelo successor severs the connection between structure and quantity. The set \{\{\{$\emptyset$\}\}\} contains only one element (the set \{\{$\emptyset$\}\}), yet we call it ``3.'' The number does not reflect its own magnitude. Order must be defined externally as the transitive closure of membership, and the construction does not generalize smoothly to transfinite ordinals.

\subsection*{The Von Neumann Successor: Synthetic Union}

The von Neumann successor function formalizes synthetic recollection as cumulative synthesis:
\[
S_{vN}(n) = n \cup \{n\}
\]

To find the successor of any number $n$, take the \textit{union} of $n$ with the set containing $n$. This operation models Hegelian Aufhebung: the new number both preserves all prior content (through the union with $n$) and transcends it (by adding $\{n\}$ as a new element).

Starting from zero:
\begin{align*}
0 &= \emptyset \\
S_{vN}(0) = 1 &= \emptyset \cup \{\emptyset\} = \{\emptyset\} \\
S_{vN}(1) = 2 &= \{\emptyset\} \cup \{\{\emptyset\}\} = \{\emptyset, \{\emptyset\}\} \\
S_{vN}(2) = 3 &= \{\emptyset, \{\emptyset\}\} \cup \{\{\emptyset, \{\emptyset\}\}\} = \{\emptyset, \{\emptyset\}, \{\emptyset, \{\emptyset\}\}\}
\end{align*}

Simplifying using the abbreviations $0 = \emptyset$, $1 = \{0\}$, $2 = \{0,1\}$:
\begin{align*}
S_{vN}(0) = 1 &= \{0\} \\
S_{vN}(1) = 2 &= \{0, 1\} \\
S_{vN}(2) = 3 &= \{0, 1, 2\}
\end{align*}

This successor function captures counting as \textit{accumulation}. Each number is a synthesis of the entire journey. The number 3 literally contains 0, 1, and 2 as elements. Its structure embodies its history.

The von Neumann successor emphasizes \textit{being} as \textit{having-become}. Each number is a stable unity that has preserved all prior stages within itself. The focus is on the \textit{product} of counting---the achieved totality---rather than just the performance.

This provides mathematical elegance. The set \{0, 1, 2\} contains exactly three elements, unifying cardinal and ordinal notions. Order is intrinsic: $m < n$ if and only if $m \in n$. The construction extends smoothly to transfinite ordinals.

\subsection*{Why Union Formalizes Aufhebung}

The key philosophical question is: why does synthetic recollection require the union $n \cup \{n\}$ rather than simple wrapping $\{n\}$?

The answer lies in how consciousness works in reflection. When I reflect on a prior thought, I don't simply wrap it in awareness. Instead, I must hold together multiple elements:
\begin{itemize}
\item The ground that enables reflection (always $\emptyset$)
\item All prior stages that led to the current moment
\item The prior moment as a unified whole (represented as $\{n\}$)
\end{itemize}

Union is the set-theoretic operation that performs this gathering without choosing or excluding. It respects both the presence of the ground and the accumulation of history. This reveals why von Neumann ordinals have their particular structure -- they formalize how reflection necessarily presupposes and preserves its entire enabling history.

The operation of Union ($\cup$) in the Von Neumann successor function ($S(n) = n \cup \{n\}$) is the formal mechanism that enacts the dialectical movement of \textit{Aufhebung}. This movement requires three simultaneous operations: preservation, negation, and elevation. The Union operation achieves all three in a single formal gesture.

The previous ordinal ($n$) is \textit{preserved} because it is included in the new ordinal. It is \textit{negated} because the new element $\{n\}$ marks its limit and transforms its status from complete to partial. It is \textit{elevated} because the new structure integrates the previous elements into a richer, more complex unity that possesses properties the previous stage lacked.

This a dynamic transformation. It formalizes the movement of thought itself---the process by which concepts develop through the integration of their own limitations. Recognizing the Union operation as \textit{Aufhebung} allows us to understand mathematical development as a dialectical unfolding. The successor function is the formal engine of mathematical becoming.

\section{Historical Contexts: Philosophies of Number and Knowledge-Constitutive Interests}

The pronoun thesis developed in this chapter does not emerge in a philosophical vacuum. It responds to a long tradition of inquiry into the nature of number, from Pythagorean mysticism through contemporary debates in mathematical ontology. To situate this work properly, we must recognize that philosophies of number are not neutral descriptions but are driven by deep methodological and epistemic motives---what Habermas calls knowledge-constitutive interests \parencite{Habermas:1971aa}.

Chapter 5 introduced Habermas's three knowledge-constitutive interests: the technical \textit{empirical-analytic interest} (knowledge to predict and control), the practical \textit{historical-hermeneutic interest} (knowledge to interpret and understand meanings in context), and the emancipatory \textit{critical interest} (knowledge to reflect and free ourselves from constraints). Each philosophy of number can be interpreted as foregrounding one or more of these interests, revealing that debates about mathematical ontology reflect deeper disagreements about what we want from mathematical knowledge.

Table \ref{tab:philo-of-number} summarizes major philosophies of number and their primary knowledge-constitutive interests. Readers should consult Chapter 5 for detailed explanations of each interest category. This schematic overview highlights the central thrust of each perspective, though many philosophers hold nuanced positions combining elements.

\begin{table}[ht]
\centering
\small
\begin{tabular}{p{3.2cm}p{5.5cm}p{3cm}}
\hline
\textbf{Philosophy} & \textbf{Core Claim} & \textbf{Primary Interest} \\
\hline
Platonism \parencite{sep-platonism-mathematics} & Numbers are abstract objects existing independently of minds. Mathematical truths describe this eternal realm. & EA \\[1ex]
Logicism \parencite{sep-logicism} & Numbers are logical constructs; arithmetic is derivable from logic alone (Frege, Russell). & EA \\[1ex]
Formalism \parencite{sep-formalism-mathematics, sep-hilbert-program} & Numbers are symbols in a formal system. Mathematics is rule-governed symbol manipulation; consistency is paramount (Hilbert). & EA \\[1ex]
Intuitionism \parencite{sep-intuitionism, sep-brouwer} & Numbers are mental constructions. Only explicitly constructible objects exist; rejects excluded middle and actual infinity (Brouwer). & HH + CE \\[1ex]
Structuralism \parencite{sep-structuralism-mathematics} & Numbers are positions in abstract structures (e.g., natural number sequence). Mathematics studies relational patterns (Shapiro, Resnik). & EA (+ HH) \\[1ex]
Fictionalism \parencite{sep-fictionalism-mathematics, sep-nominalism-mathematics-b} & Numbers do not exist; they are useful fictions. Mathematical statements are false literally but instrumentally valuable (Field). & EA + CE \\[1ex]
Social Constructivism (Hersh, Ernest) & Numbers are cultural-historical products. Mathematical objectivity is socially constituted through communal practices. & HH (+ CE) \\[1ex]
Critical Theory (Ethnomathematics) & Mathematics carries ideological power. Critical approaches examine how quantification can dominate and seek to democratize mathematical practice. & CE \\[1ex]
Pronoun Thesis (This work) & Numbers are first-person recollective pronouns formalizing structures of self-conscious thought. They are immanent in the activity of thinking itself. & HH + CE \\
\hline
\end{tabular}
\caption{\textit{Note.} Philosophies of Number and Their Knowledge-Constitutive Interests. EA = Empirical-Analytic; HH = Historical-Hermeneutic; CE = Critical-Emancipatory. See Table \ref{tab:knowledge-interests} (Chapter 5) for detailed explanations of each interest.}
\label{tab:philo-of-number}
\end{table}

As Table \ref{tab:philo-of-number} suggests, the philosophy of number is a rich field where metaphysical, epistemological, and even ethical considerations intersect. Each philosophy provides a different answer to the question ``What is a number?'': a Platonic form, a logical class, a formal symbol, a mental construction, a structural position, a useful fiction, a cultural artifact, or (in this work) a first-person pronoun. Each carries its own knowledge-constitutive interest, shaping what it considers important: certainty, consistency, constructive meaning, relational invariance, ontological parsimony, human context, or emancipatory self-understanding.

The knowledge-constitutive interest structure explains the underlying motivations driving these philosophical debates. They are not only about technical points; they are about what people \textit{want} from mathematics. Those seeking absolute, context-free truths are drawn to Platonism or structuralism. Those prioritizing secure methods and avoiding unfounded metaphysics gravitate toward formalism or finitism. Those valuing human intelligibility and meaning above all find appeal in intuitionism or constructivism. Those worried about the societal and personal implications of an overly quantified world embrace critical perspectives.

The pronoun thesis developed in this chapter combines historical-hermeneutic and critical-emancipatory interests. It is hermeneutic in seeking to understand numbers through their meaning in human practice---how they function in discourse, how they emerge from linguistic self-reference, how they connect to lived experience. It is emancipatory in challenging the reification of numbers as alien objects that sort and subjugate, instead revealing them as structures of our own thinking that can foster authentic self-recognition.

This positioning explains why the pronoun thesis diverges from dominant approaches. Platonism, logicism, formalism, and structuralism all foreground the empirical-analytic interest: they seek stable, objective foundations for mathematics that enable prediction and control. While this serves important purposes in mathematical practice and application, it obscures the human dimension of mathematical understanding. When $\mathcal{I}$ asked ``What even is two?'' he was not asking for an axiomatic definition or a structural position. He was asking for meaning in a context of alienation and grief.

In practice, many mathematicians and philosophers blend elements of these views. A working mathematician might be a formalist Monday through Friday (treating numbers as symbols in proofs), a Platonist at 2am when marveling at an elegant theorem (``these relations must exist independently of me!''), and a social constructivist when teaching students (``historically, negative numbers took time to be accepted, showing mathematical ideas evolve''). This eclecticism is not necessarily inconsistency; it can be seen as pragmatically adopting the stance appropriate to the context---much as Habermas's theory would suggest different interests can be pursued in different contexts without denying each other.

However, when foundations are at stake---when we need to understand what numbers \textit{are} to make mathematics meaningful for students or to clarify philosophical puzzles---the choice matters. The pronoun thesis claims that grounding numbers in the structures of self-conscious thought better serves human understanding than treating them as abstract objects or formal symbols. This does not deny the utility of formal approaches for working mathematics, but it provides a deeper foundation when meaning rather than manipulation is at stake.

Ultimately, the evolution of philosophies of number through immanent critique shows a progressive enrichment of understanding. By questioning ``what is number?'' from various angles, number emerges not as a monolithic concept but as multi-faceted: simultaneously a useful fiction, an abstract pattern, a cognitive tool, a social construct, and (for many) an objective reality. Each new philosophy did not simply negate the previous; it often incorporated what was valuable and corrected what was problematic. In this way, the dialogue continues.

The pronoun thesis contributes to this ongoing dialogue by insisting that numbers cannot be fully understood apart from the first-person position through which they are enacted. While formal mathematics can bracket this dimension for practical purposes, educational and philosophical inquiries must return to it. Numbers, one of the greatest conceptual achievements of humanity, reflect human capacities for self-reference, iterative recollection, and synthetic unity. Understanding them as pronouns acknowledges this human dimension while preserving the objectivity that makes mathematics universally applicable.

That said, a rigid formalist can still find comfort in the thesis. ``It'' is capatious. 

\section{Conclusion}

What, then, are numbers? If they are not objects in Plato's heaven, nor abstractions from empirical collections, nor simply truth values or formal symbols, then what remains? This chapter has argued that numerals are fundamentally \textit{first-person recollective pronouns}. This a radical re-orientation of our understanding of mathematical being.

To grasp numerals as pronouns means recognizing them not as names for things but as \textit{linguistic gestures of self-recognition}, anaphorically recollecting the enabling conditions of thought itself. Through the framework developed across Chapters 6 and 7, we have traced how numerals emerge from the dynamic interplay of apperceptive consciousness, quotative embedding, and anaphoric recollection, rooted in the unrepresentable yet necessarily presupposed ``I think.''

Chapter 6 established zero as the ground---the null representation $\emptyset$ that formalizes the enabling condition that recedes. This chapter extended that foundation by revealing two distinct forms of recollection in natural language, each generating its own ordinal construction.

The \textbf{Zermelo construction} reveals number as trace -- pure iteration where the origin recedes through successive wrapping. Like the Telephone game where each player knows only the previous whisper, Zermelo ordinals capture the phenomenology of counting as loss and deferral. The empty set becomes increasingly buried: 3 = \{\{\{$\emptyset$\}\}\}. This respects what Derrida called the trace-structure of meaning -- the origin that was never fully present yet structures everything that follows.

The \textbf{von Neumann construction} reveals number as totality -- cumulative synthesis where history gathers into ever-richer unities. Like the structure of judgment where reflection preserves all enabling conditions, von Neumann ordinals capture the logic of Aufhebung. Each number contains its entire history: 3 = \{0, 1, 2\}. The origin doesn't recede but remains present as the first element of every subsequent number.

The crucial insight is that neither construction is ``more correct.'' They capture different essential aspects of numerical cognition:

Zermelo teaches us that counting involves loss -- each iteration buries the origin deeper. This the melancholic dimension of number, acknowledging that something is always already gone in the very act of counting. When $\mathcal{I}$ asked about two in his grief, perhaps he intuited this -- that number itself carries the structure of loss.

Von Neumann teaches us that counting involves gathering -- each number preserves everything that came before. This the monumental dimension of number, insisting that nothing is lost, that the past accumulates into ever-richer presents. Mathematics builds systematic knowledge precisely through this cumulative structure.

By recognizing both dimensions, the pronoun thesis offers a resolution to $\mathcal{I}$'s question that honors its full existential weight. ``What even is two?'' Two is thought having turned back on itself -- understood both as the trace of iteration \{\{$\emptyset$\}\} and as the gathering of history \{0, 1\}. It is neither formalism nor mystical reality but the structure of self-conscious activity itself.

This understanding transforms mathematics education. If numerals are pronouns recollecting structures already operative in consciousness, then learning mathematics means recognizing patterns we already enact. Mathematical error becomes not evidence of inadequacy but opportunity for more authentic self-recognition. The drive for mathematical correctness emerges not from external authority but from the demand for authentic self-understanding.

Yet this framework maintains genuine rigor. The structures formalized as ordinals aren't arbitrary or culturally relative. They emerge necessarily from the form of self-consciousness. Anyone capable of thinking ``I think'' will generate these same patterns. Mathematical objectivity arises from the universal structure of reflexive thought, not from correspondence to external reality.

In a world increasingly dominated by quantifying systems that reduce human complexity to metrics, the pronoun thesis offers a form of resistance. It reminds us that numbers themselves emerged from and return to the irreducible first-person perspective. The ``I think'' that grounds mathematics -- elusive, receding, never fully graspable -- preserves a dimension of human subjectivity that no formalization can exhaust.

Perhaps this what $\mathcal{I}$ sensed in his moment of grief -- that beneath the machinery of mathematics lies something irreducibly human. The pronoun thesis suggests he was right. Numbers are not alien impositions but formalizations of consciousness recognizing itself. In learning to count, we learn to trace the movement of our own thought. In mastering mathematics, we master structures we ourselves enact.

The journey from zero continues. From the null representation through two forms of recollection to numerals as pronouns -- we have traced how mathematical structure emerges from consciousness turning back on itself. Chapter 8 will extend this foundation by examining how arithmetic operations arise from embodied spatial-temporal metaphors. But the ground is established: numbers are pronouns, and in using them, we recollect nothing less than the enabling conditions of thought itself.

What even is two? It is thought recognizing that it has thought about thinking -- formalized either as trace or totality, absence or presence, loss or gathering. It is consciousness catching sight of its own movement and marking it with a sign that others can recognize because they too perform this same movement. It is, finally, the mathematical name for the second moment in the eternal return of thought to itself.

\printbibliography[heading=subbibliography]





\chapter{Algorithmic Elaboration and History}

\begin{abstract}
Mathematical concepts develop through pragmatic expressive bootstrapping---conceptual systems evolve by making explicit the structures implicit in existing practices. The analysis reconstructs historical proofs from Pythagoras, Euclid, and Cantor to demonstrate how claims to completeness are consistently subverted by the systems' own internal logic. This historical view is contrasted with the limits of simple algorithmic elaboration, using the cognitive demands of improper fractions as a case study. The chapter argues that mathematical history reflects a dialectical pattern where apparent closure generates new possibilities, positioning mathematical knowledge not as a static collection of truths but as a dynamic, self-correcting, and perpetually incomplete process of becoming.
\end{abstract}

\section{The Problem of Cataloguing Sand}

Algorithms are traditionally conceived as step-by-step procedures for solving problems. An algorithm takes input, processes it through defined steps, and produces output. What this definition misses is the underlying inferential structure: algorithms are chains of inferences, as discussed in section \ref{def:algorithms}. These inferential chains can become smooth through practice, like a yogi flowing between poses. Of particular interest are chains of material inference (recall \ref{def:material-inference}) involving substitution.

When I proposed this book, I intended to explore literature reviews as a concept. The perceived lack of importance of the dissertation literature review is evident in the paucity of research devoted to understanding it \parencite[4]{boote2005}. I thought articulating the process would be a gift to future scholars. Yet as I wrote, the earth shifted beneath my feet. The popularization of Large Language Models transformed summarizing and synthesizing into seemingly passive activities.

More fundamentally, I encountered a paradox. Each attempt to gain purchase on ``The Literature'' as a totality revealed new gaps. One article would fill a gap, revealing another, then another. I had not yet synthesized the unrepresentable ground; I was trying to climb onto the shoulders of giants made of sand. The Bridge chapter articulates this process: reifying a totality in a name inevitably leads to expressions that cannot have been in the reified totality but must have been in the original totality. This process haunts any attempt to conceptualize ``The Literature.''

To help picture this, I offer two images in figure \ref{fig:Koch_text}. Part A shows a draft of this manuscript in all its intimidating density by treating the white space in the PDF as transparent. This density is how I conceptualize objects. The activity of reading you have done to get to this point in the text is recollected as if the text itself were an object. In part B, I show a Koch Snowflake, which is a relatively simple fractal, built by `bumping out' the sides of an equilateral triangle iteratively. The development of this text is partially captured by such a fractal. Beginning with a simple idea, ``no,'' and then expressing finer \textit{fractal resolutions}. 

The image doesn't capture either my fear or excitement at various stages, or my need for recognition. It's just a picture. Still, understanding develops through the ability to act-next in ways members of a recognitive community recognize as appropriate—the pragmatist inheritance of critical ethnography. The text has multiple entry points. Relax and listen to some of the songs. Repair the Prolog code to better reflect the phenomenology of The Exercise. Some aspects of the text may stay inert, like an equilateral triangle never subjected to the iteration wherein it might be recognized as a potential snowflake.

\begin{figure}[h!]
    \centering
    \includegraphics[width=\textwidth]{/Users/tio/Documents/GitHub/September_UMEDCA/images/Ch08_Koch_Resolution.pdf}
    \caption{\textit{Note.} Part A shows a draft of this manuscript in all its intimidating density by treating the white space in the PDF as transparent. Part B shows a Koch Snowflake, a fractal built by iteratively `bumping out' the sides of an equilateral triangle.}
    \label{fig:Koch_text}
\end{figure}

Returning to ``the Literature,'' the action of review of ``the Literature'' is what produces the fractal-like development of the literature itself. 

This pattern---where each attempt to name the totality invites expressions that exceed the naming---characterizes mathematical development itself. Just as ``The Literature'' cannot be totalized because reviewing it changes it, mathematical totalities cannot be closed because investigating them generates transcendence. The history of mathematics reveals this structure repeatedly: from Pythagoras's claim that ``all is rational number'' exceeding itself through Hippasus's discovery, to Euclid's demonstration that any finite list of primes is incomplete, to Cantor's proof that infinities have different sizes, to G\"odel's revelation that formal systems cannot prove their own consistency.

Rather than treating literature review as cataloguing static objects, I propose understanding intellectual history through algorithmic elaboration---and more crucially, through the moments when algorithmic elaboration reaches its limits and requires dialectical transformation. Such transformations might make the More Machine hiccup for a moment, as no one knows how to appropriately act next. That is rather how I feel reading Derrida, who enacts what I attempt to explain. That is the difference between hitting a grand slam at the World Series and saying to my kids ``when they hit the ball over the fence, all the runners get to go home.'' In any case, once set to print, communicative actions are prone to algorithmic elaboration.

\section{Mathematics in the Historical-Hermeneutic Interest}

I now reconstruct proofs from mathematics history to draw thematic unity between sublation and diagonalization. In a stronger sense, diagonalization can be understood as a syntactic species in the genus of sublation, where the genus is the dialectical movement of thought. One expressive goal of this chapter is to nestle mathematics history within the sound of time metaphor.

Modern interpretations of 'mathematics' stand quite distant from Pythagoras' original sense of the term: 'something learned'---or, given Pythagoras' institution of proof norms, 'something demonstrated.' As I write, it feels like swimming upstream to discuss the political, theological, and existential problems forming parts of the whole context in which proofs are articulated. Those ideas appear in 'standards,' and so they are treated as 'fluff.' Knowing Pythagoras only by his famous proof makes it feel like an external injection to consider how his cultish ideology led to a murder (of Hippasus, supposedly) and millennia of out-of-tune musical instruments.

According to Aristotle, the Cult of Pythagoras claimed ``that the whole heaven, as has been said, is (rational) numbers'' \parencite*[17]{Aristotle:2017:Metaphysics}. That meant, among other things, that instruments should be tuned to perfect ratios, but physics does not comply with this directive. So people played instruments that were never quite in tune.

Aristotle here acts as a historian of mathematics. His phrase ``$$\kappa \alpha \theta \acute{\alpha} \pi \varepsilon \rho \quad \varepsilon \acute{\iota} \rho \eta \tau \alpha \iota$$'' (``as has been said'') marks an ancient awareness of recollective review---synthesizing prior claims about what the Pythagoreans believed. Like any reviewer across time, Aristotle faced the essential question: has he faithfully represented what came before? Can he know the Pythagoreans' actual beliefs, or only what others have said about those beliefs? `As has been said,' my first word was ``battery.'' The problem of reviewing ``The Literature''---of recollecting and synthesizing thought across contexts---emerges at the very origin of mathematical history.

\subsection*{Hippasus and the Irrational}

The `murder' was of poor Hippasus of Metapontum (circa 500 BCE), who was allegedly drowned by the Cult for revealing irrational numbers' existence. The neutrality of the Koch Snowflake in figure \ref{fig:Koch_text} is highly misleading since it lacks the abyss Hippasus slipped into. In any case, the form of that argument mirrors diagonalization at a very abstract level: Hippasus began with the unit square and the virtuous formula $a^2+b^2=c^2$. Thus: $1^2+1^2 = 1+1=2=c^2$, yielding $c=\sqrt{2}$. By the Cult's own doctrinal formula, $\sqrt{2}$ must exist.

By their ideology: If everything (a totality) is a rational number (a reified totality; an object with the property of being rational), then $\sqrt{2}$ is a rational number. That means $\sqrt{2}=\frac{a}{b}$, with $a$ and $b$ written in lowest terms.

Squaring both sides yields $2=\frac{a^2}{b^2}$, which means $a^2=2b^2$, so $a^2$ is even. If $a^2$ is even, $a$ must have 2 as a factor---the factor of 2 in $a^2$ must come from somewhere. Let $a=2k$. Then, $a^2=4k^2 = 2b^2$, which (dividing by 2) means $2k^2=b^2$. So $b$ must also have a factor of 2, since $b^2$ is even. But that means both $a$ and $b$ are even, contradicting the assumption that $\frac{a}{b}$ was simplified, since they share a common factor of 2.

Therefore, $\sqrt{2}$ escapes the reified totality: it is not a rational number, yet it must exist \textit{by the Pythagoreans' own virtue} that $a^2+b^2=c^2$. A system's own rules generate elements that transcend its original boundaries. Each totality, when investigated according to its own commitments, invites expressions that exceed it.

\subsection*{Euclid and Potential Infinity}

Euclid did not necessarily face an ideological distortion saying there must be only finitely many primes, but there was a prohibition on articulating infinity as \textit{actual}. How can we, who are bound by birth and death, come to know about the ending of that which is without end? His proof is often reconstructed as if Euclid were comfortable with \textit{actual} infinity, but his text, \textit{The Elements}, does not stray from the doctrine that only \textit{potential} infinity can be discussed.

Rather than writing an indeterminate list of primes $p_1, p_2, \ldots, p_k$, where the $\ldots$ gives rise to offense, he began with a list of just \textit{three} primes: $p, q, r$. That list is the reified totality of primes. The partial negation is not a logical negation, but a determinate negation. He multiplies the elements of the reified totality together and adds one: $(p\,q\,r)+1$.

We investigate two cases: either the new element is a prime number, in which case the original list of three primes was incomplete, or it must have a factor not on the original list. Justifying what Franzén \parencite*{franzen_inexhaustibility_2004} calls a \textit{constructive dilemma}---the two-case either/or where both paths lead to contradiction---relies on earlier results from \textit{The Elements}. The Hermeneutic Calculator's Prolog implementation \parencite{savich_hermeneutic_calculator} provides formal verification of this and several other classical proofs using incompatibility semantics.

The point is that the reified totality is incomplete, so there must be more than three primes. The demonstration easily extends to any finite list of primes, but Euclid never claims there \textit{are} infinitely many primes. He only claims that any finite list of primes is incomplete. The set of prime numbers (a new reified totality) is not-finite, but not yet infinite. This is how we should understand \textit{potential} infinity, the only kind Euclid was comfortable discussing.

\subsection*{Who Controls Infinity?}

The Greeks and medieval Scholastics permitted discussion of \textit{potential infinity} (the endless process of adding one more) but forbade talk of \textit{actual infinity} (infinity as a completed whole). Like Moses who fled Egypt but never crossed into the promised land, they approached the infinite but would not claim to possess it. The central tension: Who controls infinity---humans or God?

Cantor changed everything. As a platonist committed to the existence of mathematical objects, he asked: What \textit{is} infinity, and how can we finite beings know it? He developed a method for measuring infinite quantities---what he called \textit{magnitude} (Mächtigkeit)---by comparing sets through bijections, one-to-one correspondences where every element in one set pairs with exactly one element in another. His diagonal proof (explored in the Bridge chapter) demonstrated that it is impossible to establish such a bijection between the natural numbers and the real numbers, proving that some infinities are strictly larger than others.

Like the Sneetches who learn that their artificial distinctions are ultimately arbitrary, the mathematical community had to abandon its rigid separation between the finite realm of human knowledge and the \textit{in}finite realm reserved for divine understanding. Cantor's diagonalization becomes a machine that forces us through confusion toward more sophisticated understanding. Not all infinities are equal, and humans can rigorously investigate these hierarchies rather than accepting them as mysterious divine prerogatives.

\section{Algorithmic Elaboration as Living History}

Brandom introduces the concept of algorithmic elaboration through long division \parencite[37]{Brandom2008}, showing how the division algorithm elaborates from multiplication and subtraction practices---accumulated learning experiences with these operations become \textit{temporally compressed} into the codified procedure of long division. Figure \ref{fig:brandom_figure_2_3} reproduces his automaton-implemented elaboration.

\begin{figure}[h!]
    \centering
    \includegraphics[width=0.8\textwidth]{/Users/tio/Documents/GitHub/September_UMEDCA/images/ch08_brandom_2_3.pdf}
    \caption{\textit{Note.} Automaton-implemented, algorithmically elaborated, pragmatically mediated semantic relation reproduced from \parencite[37]{Brandom2008}.}
    \label{fig:brandom_figure_2_3}
\end{figure}

My initial skepticism about this example was intense. The rich semantic structure of division seems hardly captured by the long-division algorithm. The claim felt incomplete without tracking how counting elaborates into addition, addition into multiplication and subtraction. Consider the algebraic structures of basic arithmetic operations shown in Table \ref{tab:arithmetic-operations}.

\begin{table}[h]
\centering
\caption{Algebraic Structure of Basic Arithmetic Operations}
\label{tab:arithmetic-operations}
\small
\begin{tabular}{@{}lp{9cm}@{}}
\toprule
\textbf{Operation} & \textbf{Algebraic Structure} \\
\midrule
Addition & Known Part $+$ Known Part $=$ Unknown Whole \\[0.5ex]
Subtraction & Unknown Part $+$ Known Part $=$ Known Whole \\[0.5ex]
Equal-Groups Multiplication & Known Items per Group $\times$ Known Number of Groups $=$ Unknown Total \\[0.5ex]
Sharing Division & Unknown Items per Group $\times$ Known Number of Groups $=$ Known Total \\[0.5ex]
Measurement Division & Known Items per Group $\times$ Unknown Number of Groups $=$ Known Total \\
\bottomrule
\end{tabular}
\end{table}

These structures suggest deeper dialectical shifts at work. Is subtraction simply elaborated from addition by ``reversing the wiring,'' or does conceptual understanding require dialectical thinking? The mechanical attributes of automata modeling arithmetic vary significantly. Kindergarten emphasizes three-step procedures as prerequisites. Matching nested parentheses for von Neumann ordinals requires push-down stack automata with primitive memory, not finite-state machines.

\subsection*{Example: The Complexity of Improper Fractions}

Consider improper fractions like $\frac{12}{11}$. Research in mathematics education \parencite{Hackenberg:2007aa, steffe2002new} shows that students who can fluently work with proper fractions often struggle when fractions exceed the whole. The difficulty is structural.

To robustly manage improper fractions, students must simultaneously coordinate three distinct levels of units:

\begin{enumerate}
\item \textbf{The Unit Fraction} (e.g., $\frac{1}{11}$): The basic iterable unit
\item \textbf{The Composite Unit} (e.g., $\frac{12}{11}$): The result of iterating the unit fraction twelve times
\item \textbf{The Referent Whole} (e.g., $\frac{11}{11}$): The composite unit that defines the size of the unit fraction
\end{enumerate}

The complexity lies in maintaining the identity of the unit fraction (Level 1) based on the referent whole (Level 3) \textit{even while operating on a quantity that exceeds that whole} (Level 2). When students can only coordinate two levels of units prior to activity, the necessary inferences fail. They may iterate $\frac{1}{5}$ seven times, but upon reaching the result, they lose connection to the referent whole ($\frac{5}{5}$). Their cognitive system resets, and they interpret the quantity as $\frac{7}{7}$, conflating the new total with the definition of the whole.

This is not algorithmic elaboration---it requires a dialectical reorganization of numerical schemes such that the cognitive automaton can support simultaneous three-tiered coordination \textit{prior to operating}. The automaton must be fundamentally restructured, not extended. This example illustrates why some conceptual developments require pragmatic expressive bootstrapping rather than simple elaboration.

Students need \textit{recognition} for their current understanding even as they're pushed toward incommensurable reorganization. The shame of error can freeze development; recognizing apparent misconceptions as productive resources for growth enables the dialectical movement from inadequate to more adequate understanding. Teachers who can hold this tension---affirming what students currently understand while inviting transcendence---create conditions for genuine learning rather than compliance.

\subsection*{A Speculative Direction: Fractions as Fractal Expansion of the Base Dialectic}

The difficulty of fractions may stem from a deeper source: they intensify and invert the very dialectic that generates whole numbers. Consider the paradox at the heart of base arithmetic: \textit{ten ones are one ten}. This identity-in-difference drives the positional system---we compress ten units into a single higher-order unit, creating place value through iterative grouping (see Figure \ref{fig:hand_of_hands}).

\begin{figure}[h!]
    \centering
    \includegraphics[width=0.9\textwidth]{/Users/tio/Documents/GitHub/September_UMEDCA/images/ch08_hand_of_hands.pdf}
    \caption{\textit{Note.} The fractal structure of base arithmetic visualized across base-5 and base-10 systems. Each higher power repeats the same structural relationship at a larger scale.}
    \label{fig:hand_of_hands}
\end{figure}

What if rational numbers are not a separate ontological category (the Cartesian product $\mathbb{Z} \times \mathbb{Z}^*$) but rather the \textit{continuation of this same dialectic past the unit boundary}? If whole numbers emerge through iterative compression (multiplication by the base: $B^1, B^2, B^3\ldots$), perhaps rational numbers emerge through iterative decompression (division by the base: $B^{-1}, B^{-2}, B^{-3}\ldots$).

Consider: in base 10, the fraction $\frac{1}{7}$ requires infinite repetition (0.142857\ldots). But in base 7, seven itself is $(10)_7$, and $\frac{1}{7}$ becomes $(0.1)_7$---the fundamental unit fraction of that base system, the first step of inward dialectic just as $(10)_7$ is the first step of outward dialectic.

This suggests a fractal-like expansion of \textit{resolution} rather than dimension. We are not leaping to a 2D grid of ordered pairs; we are intensifying articulation of the existing number line through self-similarity. The relationship between place values repeats at finer resolutions:

\begin{center}
(Hundreds : Tens) :: (Tens : Ones) :: (Ones : Tenths) :: (Tenths : Hundredths)
\end{center}

A repeating decimal, from this perspective, is not a \textit{representation} of a rational number---it \textit{is} the rational number as a stabilized loop in the division algorithm. The number is the process that generates it, an attractor basin in the dynamical system of base decomposition.

Compression—unitizing ten ones as one ten—is a structural cognitive operation. Sublation (\textit{Aufhebung}) is a dialectical movement involving simultaneous negation, preservation, and elevation. Sublation appears in fraction arithmetic precisely where we expect dialectical tension: in the addition of fractions with different denominators.

\subsubsection*{Incommensurable Norms and Dialectical Synthesis}

Consider adding $(0.1)_7 + (0.1)_5$---that is, $\frac{1}{7} + \frac{1}{5}$. These fractions embody \textit{incommensurable normative structures}. The fraction $\frac{1}{7}$ means ``one unit derived from a system where seven such units compose the whole.'' The fraction $\frac{1}{5}$ means ``one unit derived from a system where five such units compose the whole.'' The embodied actions of partitioning-by-seven and partitioning-by-five are materially different, incompatible operations.

This incompatibility forces a dialectical moment. We cannot simply combine these quantities because they belong to different structural movements. The traditional procedure of finding a common denominator---typically taught as rote manipulation---is reinterpreted here as \textit{dialectical synthesis through recursive partitioning}.

To reconcile $\frac{1}{5}$ and $\frac{1}{7}$, we must apply the conflicting norms to each other: partition the $\frac{1}{5}$ unit according to the seven-fold norm, and partition the $\frac{1}{7}$ unit according to the five-fold norm. The resulting unit ($\frac{1}{35}$) is the synthesized unit of measure---the finest fractal resolution where both schemas align.

This is where sublation appears:

\begin{enumerate}
\item \textbf{Negation}: The sufficiency of the original denominators (5 and 7) as fundamental partitioning units is negated. Neither system alone can express the unit of the other.
\item \textbf{Preservation}: The partitioning capacities of both 5 and 7 are preserved within the structure of 35. The original quantities ($\frac{1}{5} = \frac{7}{35}$ and $\frac{1}{7} = \frac{5}{35}$) remain recognizable.
\item \textbf{Elevation}: The system elevates to finer fractal resolution, reconciling incompatibility and making units commensurable. Once commensurable, addition proceeds through the familiar choreography of combining like units.
\end{enumerate}

The rational number system $\mathbb{Q}$ can thus be viewed as the total space of all possible stabilized resolutions generated by this dialectic of commensurability. The difficulty of fractions is precisely the difficulty of managing multiple, nested, conflicting normative structures. It requires students to move beyond operating \textit{within} a fixed base to operating \textit{on} the bases themselves, synthesizing them through the demanding dialectical work of recursive partitioning \parencite{steffe2002new}.

This framework remains speculative and requires substantial development, but it suggests that the difficulty children experience with fractions may arise from a conceptual discontinuity: fractions are traditionally introduced as a new kind of number rather than as the natural continuation and intensification of base-making practices children have already mastered. Interested readers can explore these ideas through the supplementary software \textit{Ace of Base Arithmetic}, which allows practice with multiple base systems, making the self-similar structure visible through direct manipulation.

\section{From Cantor to G\"odel: Gaifman's Generalization}

The historical development from Hippasus through Euclid to Cantor reveals a recurring pattern. Haim Gaifman's \parencite*{Gaifman2005} reconstruction connects Cantor's diagonalization to G\"odel's incompleteness theorems, revealing the deep structure underlying both---showing how self-reference combined with systematic enumeration inevitably generates transcendence.

\subsection*{The Naming Framework}

Rather than beginning with Cantor's infinite sequences of characters, Gaifman starts with a sequence of sets of natural numbers: $X_1, X_2, \ldots, X_n, \ldots$. This creates a naming system where each infinite set in the sequence is identified by a natural number. The natural numbers \textit{index} these infinite sets, so we can take an $n$ in $\mathbb{N}$ to represent or name the set $X_n$.

For some sets, the set's name is an element of the set itself. For example, let $X_3= \{1, 3, 5, \ldots\}$: here, $3 \in X_3$. Gaifman introduces functional notation, using $X(y)$ as shorthand for $y\in X$, so we can express $3 \in X_3$ as $X_3(3)$. This notational shift---treating sets as functions---may seem disorienting, but it reveals the crucial self-referential structure underlying diagonalization.

\subsection*{The Diagonal Construction}

From this naming framework, Gaifman considers all the $n$'s that happen to be in the sets they name---all of the $n$'s that are what we might call \textbf{fixed points} of their functions $X_n$. Fixed points are those $x$ such that $f(x)=x$\label{def:fixed_points}. They are sites where mathematical self-reference meets with the larger critical-emancipatory interest in \textit{identity claims}, as the sentence $f(x)=x$ amounts to $x$ saying of itself ``I have property $f$.'' He then defines the complement: $X^{*}$, the set of all $n$'s that are \textit{not} in the sets they name. Formally:

\begin{equation}
X^{*}(n) \iff \neg X_n(n)
\end{equation}

This $X^{*}$ captures exactly those natural numbers that, when used as names, fail to name sets that contain them. Since our naming system is supposed to be complete---covering all possible sets---the set $X^{*}$ must itself have a name. If our list were truly exhaustive, there should be some $k$ such that $X^{*}=X_k$.

\subsection*{The Contradiction}

But here we encounter the fundamental incompatibility that drives all diagonalization arguments. If $X^{*}=X_k$ for some $k$, then by our definition:

\begin{equation}
    X_k(n) \iff \neg X_n(n)
\end{equation}

This should hold for every $n$. But when we substitute $k$ for $n$, we get:

\begin{equation}
    X_k(k) \iff \neg X_k(k)
\end{equation}

This is a direct contradiction: $X_k(k)$ is true if and only if it is false. The set $X^{*}$ cannot appear in any complete enumeration, because its existence contradicts the assumption of completeness. Like Russell's paradox of ``the set of all sets that don't contain themselves,'' this reveals how self-reference combined with attempted totalization generates necessary incompleteness.

\subsection*{The General Pattern}

Gaifman summarizes diagonalization as a ``sandwich'': claim a fixed point, negate it, establish a new fixed point reflecting this negation. In Sneetches terms: ``I am a star-bellied sneetch'' $\rightarrow$ McBean's machine flips the stars $\rightarrow$ ``Sneetches are Sneetches'' (stars no longer matter). Where a domain of names attempts to enumerate higher-type entities defined over that same domain, self-reference becomes possible. And where self-reference meets negation, incompleteness emerges.

\subsection*{G\"odel's Achievement}

In G\"odel's case, the ``names'' are mathematical encodings of the formal elements of a mathematical system. The system he encoded was Alfred North Whitehead and Bertrand Russell's \textit{Principia Mathematica}, a formal system that attempted to capture all of mathematics in a single logical framework. The ``higher type entities'' are the claims \textit{about} those formal elements.

The encoding scheme is called \textbf{arithmetization}\label{def:arithmetization}. It assigns a unique prime number to each symbol legible by the automata that can generate formal statements. Using results proved by Euclid (specifically the fundamental theorem of arithmetic), it is possible to encode statements without ambiguity, as each natural number has only one prime factorization.

I recently went to an Avett Brothers concert, and while waiting for the band to start, I noticed that they used Roman Numerals as part of the light show. Specifically, they used the set XXXIII LXVIII CV, which in base ten would be 33, 68, and 105. These numbers do not appear to say much, but they are actually a simple cipher. Let $a=1, b=2, \ldots, z=26$, noting that $T = 20, H=8,$ and $E = 5$. Then we can decode the Roman Numerals as follows: $33=20+8+5=THE$, $68=1+22+5+20+20= AVETT$, and $105=2+18+15+20+8+5+18+19=BROTHERS$. The reason this is a non-example is that there are many different ways to sum digits to 33. Specifically, the letter $K=11$, so $33= KKK$. They obviously didn't intend to microagress the audience: there are a \textit{lot} of ways to partition a number like 33 into a sum of natural numbers between 1 and 26. They could also have spelled $BEDBEDBED$. 

What G\"odel does instead is to assign a unique prime number to each symbol in the formal system. If the Avett Brothers had instead used $a=2, b=3, \ldots e=11,\ldots h=19,\ldots$, and multiplied instead of adding, $THE=71\times 19\times 11 = 14839$. Such a scheme would at least be unambiguous about which characters were intended when decoding the code. G\"odel's scheme does a bit more than this, however. He also encodes the position of each symbol alongside the symbol itself. In our toy example, $THE$ would be $2^{71}\times 3^{19}\times 5^{11}$. This number cannot be rendered in Roman Numerals, which lack a base system. 

As I waited for the band to start, I was captured by their lightshow in part because over the last few months, I have been working on using just such a scheme to encode some of the automatons from the supplementary materials discussed in the Hermeneutic Calculator. Specifically, I had been working on using an additive scheme to encode an automaton's states and transitions. The fact that the Avett Brothers could have been saying ``BEDBEDBED'' or ``THE'' creates a systematic ambiguity, so that sometimes a state in a machine and a transition between states can be confused. From that confusion---like poetic ambiguity---I hoped to construct a machine that exhibited very basic `emergent' behavior. It didn't `work,' though the experiment was fun. Essentially, what I hoped to do was construct the machine so that when it entered a state whose name was also the numerical value of the input that allowed the machine to transition to that state, it would take a hard-coded offramp and fall into an endless loop of repeated addition, as a proto-multiplicative operation. I was imagining how, when teaching a mathematical concept, the specific numbers used in the example impact how students understand the concept. I plant this seed for future exploration. 

Let's return to G\"odel's more successful arithmetization scheme. What these astronomically large numbers, called G\"odel numbers, do is allow us to encode the entire formal system as well as statements about the formal system into the arithmetic system. One of G\"odel's theorems proves it is always possible to create a fixed point using this naming scheme. The kinds of identity claims that arise from this encoding are varied. G\"odel encoded a statement that says of itself, ``I am not provable in this system.'' This is a self-referential statement that cannot be proven within the system, as it would lead to a contradiction. But \textit{any} predicate can be encoded in this way, even those 'outside' the system, so long as the formal system has sufficient expressive power (alphabet, practices-or-abilities, vocabularies).

G\"odel's first incompleteness theorem states that any sufficiently complex formal system that can express basic arithmetic is incomplete: there are true statements about the natural numbers that cannot be proven within the system. The second incompleteness theorem goes further, showing that such a system cannot prove its own consistency. Accomplishing these theorems requires careful construction of self-referential statements, which G\"odel achieved through sophisticated use of diagonalization.

These theorems sank the ambitions of the \textit{Principia} and similar foundational projects. G\"odel's work showed that no matter how sophisticated a formal system is, it will always contain statements that are true but unprovable within that system. This is a profound limitation, revealing that the quest for a complete and consistent foundation for all of mathematics is ultimately futile. The connection between G\"odel's work and Derrida's deconstruction that Priest \parencite*{Priest:2002aa} highlights is striking, as the resources of the \textit{Principia} were used to construct the statements that undermine the foundational ambitions of the project.

However, the project of mathematics did not end with G\"odel's theorems. It is easy to misread the theorems as suggesting that mathematics' incompleteness is a flaw. But to a math educator, incompleteness is a feature, not a bug. I encounter students in the slipstream between being and nothing, as they become what they already were but lacked the words to say. \textit{Becoming}, in a deeply embodied sense, has no syntactic closure and so it may find some analog in \textit{in}completeness. ``The simple truth; there's more to say.''

\subsection*{Incompleteness as Becoming}

Diagonalization reveals that every attempt at complete systematization contains the seeds of its own transcendence. I read this through Mead's distinction between the ``me'' (our systematized, nameable aspects) and the \{I\} (the spontaneous, creative response that cannot be predicted). Diagonalization is the mathematical expression of the \{I\}'s capacity to transcend any finite systematization of the ``me.'' Hegel captures this in the movement from \textit{Verstand} (finite, systematic understanding) to \textit{Vernunft} (\textit{in}finite, self-reflective reason). In that movement, confession, forgiveness, and reconciliation emerge.

Through \textit{arithmetization}, G\"odel showed how any ``outside'' critique can be brought ``inside'' the system as a mathematical statement---but precisely this capacity for self-critique ensures no system achieves final closure. What formalists call incompleteness is a mathematical form of mathematical becoming, showing how any sufficiently complex system can transcend its own boundaries. The system's ability to refer to itself---to make its own structure an object of investigation---guarantees that it will generate truths that exceed its deductive closure.

\section{The More Machine---Returning to the Sound of Time}\label{sec:more_machine}

Mechanizing \textit{becoming} to express a mathematical story of mathematical development requires, first, an acknowledgment of the quixotic nature of such a project. Whereas McBean packs up his machine and leaves at the end of the story, the kind of mathematical machines I have played with cannot break themselves down, nor build themselves up. More sophisticated machines, like Large Language Models, can build the machines I describe and then act as a You towards those machines, forming proto-self-consciousnesses, but as I write, even those sophisticated machines cannot yet sublate themselves. I must limit my ambition to examining the discarded shells to interpolate the pupating life cycle of the machine.

Some are working on G\"odel machines, which are designed to explore the implications of G\"odel's incompleteness theorems through computational means. These machines can generate statements that reflect the self-referential structure of G\"odel's diagonalization, allowing for deeper understanding of the limitations and possibilities of formal systems. I am going to go a simpler route and focus on a Cantorian machine.

I call the machine the \textit{More Machine} because it generates more, more, and more. Unplugging from the incessant creation of new thoughts is, from the Exercise in Chapter 1, a desirable skill to learn. Meditative consciousness as detached (non-present) awareness is difficult to learn. It takes a certain kind of discipline, suggested by the title of this book.

The More Machine accepts inputs of $m$s and $w$s. It outputs a sequence of $m$s and $w$s that is guaranteed to be different from any sequence it has previously output, by 'flipping' each element of the diagonal of the matrix that represents the history of its outputs. Originally, I used 0s and 1s and had users input either a 0 or a 1 to lead to the next state. But when I returned to Cantor's original proof, I realized that the $m$s and $w$s would allow me to attach a Zeeman Catastrophe Machine to the More Machine. Zeeman's machine is not totally deterministic. When the rubberbands are pulled straight down, the machine enters into a superpositioned state, where the wheel could rotate left or right. With the right orientation and tension, the machine could output either $m$ or $w$. Figure \ref{fig:superimposed_zeeman_cantor_more} shows how the historical record might evolve as the wheel of the Zeeman machine moves.

\begin{figure}[h!]
\begin{center}
\includegraphics[width=\textwidth]{/Users/tio/Documents/GitHub/September_UMEDCA/images/ch06_superimposed_zeeman_cantor_more.pdf}
\caption{\textit{Note.} The More Machine reads states from the orientation of the wheel of a Zeeman Catastrophe Machine.}
\label{fig:superimposed_zeeman_cantor_more}
\end{center}
\end{figure}

When meditating, I try to hold the tension of determinate negation and the determinate negation of determinate negation in one significant symbol: ``\st{no}.'' Think of $m$ and $w$ as representing the two poles of a binary opposition, perhaps in assertoric forms like $m$:``Tio is Rudy's son'' and $w$:``Rudy is Tio's dad'' or paradoxical unities like $m$ is the truth value of ``This sentence is false'' and $w$ is the negation of that truth value. 

The superpositioning of those two states invites no further reflection. Nothing needs to be said in response to such a statement as it holds the essence of \textit{neti-neti}, where the \{I\} is neither this nor that. This moment---before the wheel snaps decisively to $m$ or $w$---represents what we might call \textit{temporal decompression}: a relaxing of the determinate negations that created those distinctions in the first place. It is a moment of undifferentiation, of not-yet-determined becoming. Whereas Brandom considers instrumental action to be a kind of exit ticket from language, I am trying to get at the idea of \textit{non}-action as an exit ticket from the incessant yammering of the More Machine.

Potential energy builds in the machine as meditative consciousness moves along the central axis between the wheel and the fixed end of the rubber band. But then the superpositioned state collapses and the machine outputs either $m$ or $w$ which can never hold the entirety of the \{I\}. The machine then recollects its entire history in a specifically diagonal way. Judgments accrue and the machine jolts on in its incessant yammering.

The More Machine brings non-deterministic elements to the deterministic structure of Cantor's diagonalization. For a long time, I claimed that I was a 'creative person.' Burdened by the weight of that claim, I found myself flitting about, trying to find something new (and clever) to say or do. I cannot just unplug from the desire to create, but I can learn to hold the right sort of tensions to arrest the movement of the machine.

The More Machine invites contemplation into the paradox of creativity: that which is 'new' must be recognized, which means it was built from what came before and so cannot actually be new. Like McBean's machine forcing Sneetches through cycles of distinction and confusion, the More Machine generates endless novelty while demonstrating this novelty was always already implicit in the system's recursive structure. One way to unhook from the More Machine is through the non-act of \textit{listening}---a necessary if challenging skill for teachers. The More Machine's superpositioned state corresponds to the moment in meditation where the subject neither affirms nor denies but holds the tension without giving in to the temptation to think. When this tension collapses into a specific output ($m$ or $w$), the machine recollects its entire history diagonally, creating new sequences that transcend finite enumeration while remaining grounded in logical structure.

\section{Euclid's Proof: Incompleteness as Incoherence}

To demonstrate how algorithmic elaboration and dialectical rupture work together, I reconstruct Euclid's proof using Brandom's incompatibility semantics \parencite{Brandom2008}. First, key concepts:

An \textit{incoherence frame} is a pair $\langle L, Inc \rangle$ where $Inc$ specifies which sentences in language $L$ are incoherent. Think of this as sentences a qualitative researcher seeks to understand, or troubling interactions people reflect upon to make sense of them. Incoherence can be existentially terrifying. When recognized as believing both $x$ and $\neg x$, I might be understood as becoming, or as duplicitous. Context matters.

An \textit{incompatibility relation} $I$ defines two sets $X$ and $Y$ as incompatible if their union is incoherent: $A \in I(B) \iff A \cup B \in Inc$.

Euclid's proof (Book 9, Proposition 20 of \textit{The Elements} \parencite{euclid_euclids_2007}) claims any finite list of primes is incomplete. The common misrepresentation suggests Euclid provides an algorithm for finding primes. This is false. He only claims that multiplying all known primes and adding one yields a number not divisible by any known prime. The number $2\times 3 \times 5 \times 7\times 11\times 13+1=30031=59\times 509$ is composite, not prime.

Euclid's actual argument: Take primes $A$, $B$, $C$. Form $DE$ as their product. Add one unit to get $EF = DE + 1$. Either $EF$ is prime (thus a new prime beyond the list) or composite with prime factor $G$ that cannot equal $A$, $B$, or $C$ (since if $G$ divided both $DE$ and $EF$, it would divide their difference, which is $1$).

\begin{figure}[h!]
	\includegraphics{/Users/tio/Documents/GitHub/September_UMEDCA/images/ch08_Cuisenaire_Rods.pdf}
	\caption{\textit{Note.} Cuisenaire rods interpreted as numerals 1-10, building composite 6 from primes ($2\times 3$), and algorithm for Euclid's proof.}
	\label{fig:rods}
\end{figure}

Through incompatibility semantics, the key understanding becomes:

$\{\text{The list of primes } \{A, B, C\} \text{ is complete}\} \in Inc$.

This statement belongs to the incoherence frame because it generates internal contradiction. The claim to completeness defeats itself through the construction $N = A \times B \times C + 1$.

\section{The Reflective Turn: Admitting Incoherence}

The reconstruction above unearths a nascent flaw in my thinking: the notion of an incoherence frame allows errors, misrecognitions, mistakes, and misconceptions to be \textbf{admitted} into the system without damaging its coherence. When an algorithm produces a paradoxical result, that paradox becomes an element of an incoherence frame rather than a problem to solve. 

This transforms our understanding, but the imposition of an incoherence frame arrests the movement of the negative. I preserve the notion of an incoherence frame in the formal side of critical mathematics, but it should probably be further excavated in later reconstructions of critical math.

Mathematical learning might proceed not through avoiding error but through sophisticated recognition and integration of contradictory commitments. Barriers become resources for development. In formal mathematics, contradictory commitments are dropped as falsities. In mathematics education, tensions of misrecognition are sublated into new understanding; they do not simply disappear.

When a student claims $\frac{1}{5} + \frac{1}{5} = \frac{2}{10}$, this is not error to be corrected. It reveals a coherent, if limited, scheme---the student is attending to the iterative structure of unit fractions and applying additive reasoning consistently. The teacher's task is not to replace this understanding but to create conditions where the student recognizes its insufficiency when confronted with incommensurable contexts, enabling dialectical reorganization toward more adequate understanding.

\subsection*{Computational Formalization: A Pragmatic Check}

I felt compelled to formalize this understanding computationally, creating \texttt{more\_machine\_learner.pl} in Prolog. This was driven by the structure's internal demands---if the system I've been outlining turned out to be internally incoherent, the computational implementation would collapse into contradiction. The formalization provides a pragmatic check: if the code runs without collapsing, it suggests (though does not prove) that the philosophical framework holds together.

The engine does not take axioms and prove theorems in the traditional sense. Instead, it takes premises and proves \textit{incoherence}. This subtle but crucial difference captures Euclid's proof structure: demonstrating that assuming completeness leads to incoherence. The system doesn't demonstrate truth; it demonstrates self-defeating claims.

Prolog's homoiconicity---the property that code and data share the same representational structure---proves convenient here. It allows the system to treat inferential rules as data that can be inspected, modified, and reasoned about. This computational property provides a useful way to \textit{represent} (in the mode of Verstand, representational understanding) the Hegelian insight about form and matter. But I must be clear: this is a representational convenience. The real philosophical work happens in dialectical reason (Vernunft); Prolog helps me check whether my articulation of it remains coherent.

I had extensive AI assistance cobbling together code to ``prove'' theorems. There's a difference between how I've been using AI as a tool for productivity and having it teach my kids math. I can identify with my choices about what to delegate; I remain responsible for the philosophical claims. Yet from the Hermeneutic Calculator to the Prolog supplements, interested readers can see these claims enacted computationally, not asserted. The work is a starting point for others to build upon, as I have built on others' work. The code demonstrates feasibility, not finality.

The approach taken throughout the Hermeneutic Calculator treats formalization as a demonstration of the failure of representational thinking. I worked hard to formalize in the mode of \textit{Verstand} (the understanding) to partially demonstrate its limitations. Each strategy is formalized as ``written choreography for embodied cognition''---preserving not just what students calculate but how they temporally move through calculations. This grounds mathematical understanding in expressive action and social normativity. When the system encounters resource limits (too many inference steps, exhausted working memory), it triggers normative crisis, modeling error's productive role in mathematical development. The formalization \textit{models the structure} of mathematical consciousness (embodiment, crisis, transcendence) without claiming to instantiate consciousness itself. It shows what would need to be risked for genuine computational self-reorganization; it does not (and cannot, at this stage) actually risk anything. The contribution is providing tools for philosophical thinking about how mathematical understanding develops, not claiming to have solved the problem of machine consciousness.

\section{From Embodied Metaphors to Arithmetic}

The dialectical analysis of admitting incoherence opens a path to understanding how arithmetic emerges from embodied experience. Following Lakoff and Núñez \parencite{Lakoff2000}, I synthesize embodied metaphors with incompatibility semantics. Mathematical vocabulary emerges as we \textit{temporally compress} embodied practices---counting, grouping, moving along paths---into symbolic forms that can be recollectively deployed. The word ``five'' is not a Platonic form discovered in some abstract realm; it is the compression of countless learning experiences with five-ness into a repeatable concept.

\subsection*{Bounded Regions and Categories}

Proprioceptive feelings define boundaries: an interior (inside the body) and exterior (outside). The vocabulary transitions algorithmically:

\begin{enumerate}
    \item If you are in a bounded region, you are not out of that bounded region.
    \item If you are out of a bounded region, you are not in that bounded region.
    \item If you are deep in a bounded region, you are far from being out.
    \item If you are on the edge of a bounded region, you are close to being in.
\end{enumerate}

Through substitution:
\begin{itemize}
    \item ``Categories'' for ``Bounded regions in space''
    \item ``Category members'' for ``Objects inside bounded regions''
    \item ``Subcategory of larger category'' for ``One bounded region inside another''
\end{itemize}

\subsection*{Material Inferences for Object Collections}

To establish arithmetic foundations, I explicate material inferences governing object collections, enabling vocabulary $V_{\text{object collection}}$:

\[
V_{\text{object collection}} = \{\text{bigger}, \text{smaller}, =, \text{added to}, \text{results in}, \text{adding}, \text{subtract}, \text{zero}, \text{one}\}
\]

Key inferences include:

\textbf{Linearity}: $\{ A \text{ and } B \text{ are object collections} \} \dagger$ (Highlander algorithm)
\begin{itemize}
    \item $\{ A \text{ bigger than } B \} \vDash_{I} \{ B \text{ smaller than } A \} \rightarrow \neg \{ A \text{ smaller than } B \}$
    \item $\{ B \text{ bigger than } A \} \vDash_{I} \{ A \text{ smaller than } B \} \rightarrow \neg \{ B \text{ smaller than } A \}$
    \item $\{ A = B \}$
\end{itemize}

\textbf{Limited Iteration of Subtraction}:
\begin{itemize}
    \item $\{ B \text{ smaller than } A \} \rightarrow$ $\{ B \text{ subtracted from } A \text{ is an object collection} \}$
    \item $\{ B \text{ symmetrically intersubstitutable with } A \}$ $\vDash_{I}$ $\{ \text{No objects left when } A \text{ subtracted from } B \}$
\end{itemize}

These establish operations on object collections underpinning arithmetic, grounding them in deontic-normative structures rather than alethic assertions.

\subsection*{Arithmetic Is Motion Along a Path}

The most pedagogically powerful embodied metaphor for arithmetic, especially in connection to the pronoun thesis from Chapter 7, is \textit{Arithmetic Is Motion Along a Path}. Recall from Chapter 1 how spatial concepts like ``distance'' emerge from the temporal compression of repeated experiences. In Chapter 7, we explored how numerals function as anaphoric recollections of quotative embedding---each numeral recollecting a temporal process of successive self-reference. The Motion Along a Path metaphor makes this temporal-spatial dynamic explicit.

The source domain is the physical act of moving along a path. This grounds the sequential and directional properties of numbers. Zero corresponds to the origin---the starting point of motion. Addition corresponds to movement away from the origin; subtraction to movement toward the origin. When we say ``$5$ is greater than $3$,'' we naturally gesture or imagine movement along a path, with $5$ being ``further along'' than $3$.

Crucially, this metaphor grounds several key mathematical concepts:

\begin{itemize}
\item \textbf{Magnitude as distance from origin}: The value of a number corresponds to how far along the path you've traveled from zero.
\item \textbf{Operations as directed movement}: $A + B$ means starting at location $A$ and moving forward by distance $B$. Subtraction reverses this: $A - B$ means moving backward from $A$ by distance $B$.
\item \textbf{Negative numbers as bidirectional paths}: Once the path extends in \textit{both} directions from the origin, negative numbers emerge naturally as locations on the opposite side of zero.
\item \textbf{The ``Sliding'' strategy}: The distance between two points remains invariant if both are shifted by the same amount---grounding the inference that $(M + K) - (S + K) = M - S$.
\end{itemize}

This metaphor connects directly to the claim that numerals are first-person pronouns recollecting temporal sequences. Just as my daughters' repeated drives to Indianapolis compressed temporal experience into spatial understanding of ``50 miles,'' the Motion Along a Path metaphor reveals how numerical thinking compresses the temporal process of counting (one step, another step, another step) into a spatial form we can traverse, retrace, and reason about.

The inferential structure of motion---forward, backward, distance, direction---gets imported wholesale into arithmetic through this metaphor. When students struggle with subtraction as ``taking away,'' reframing it as ``finding the distance between two locations on a path'' or ``retracing your steps'' often provides the breakthrough. This is not pedagogical convenience; it reveals how arithmetic operations are \textit{grounded} in embodied spatial reasoning about movement and position.

\section{Pragmatic Expressive Bootstrapping}

Simple algorithmic elaboration has limits. Dialectical moments require what Brandom calls \textit{pragmatic expressive bootstrapping} \parencite{Brandom2008}: new, more expressively powerful vocabularies emerge from implicit practices already mastered. The system pulls itself up by its bootstraps, making explicit what was latent in its doing.

Brandom formalizes this through Meaning-Use Analysis. Bootstrapping occurs when vocabulary $V'$ is VP-sufficient to specify practices $P$ that are PV-sufficient to deploy more powerful vocabulary $V$. The ``bootstrapping'' happens when specifying vocabulary ($V'$) is weaker than deployed vocabulary ($V$).

I implemented this computationally in \texttt{more\_machine\_learner.pl}. The system begins with norms for arithmetic over natural numbers where subtraction is limited: $3 - 5$ is nonsensical. When presented with observation \texttt{minus(3, 5, -2)}, it enters normative critique:

\begin{verbatim}
Observation is INCOHERENT with current norms. 
Entering critique phase...
Identified Incompatibility: 
incompatibility(limited_subtraction, domain(n))
Proposed Normative Shift: Change domain to z
Bootstrapping complete. Observation is now coherent.
\end{verbatim}

The system uses failure of its vocabulary (arithmetic over $\mathbb{N}$) to deploy new vocabulary (arithmetic over $\mathbb{Z}$). The paradox resolves not through dismissal but transformation. What appeared as incoherence becomes the resource for pragmatic expressive bootstrapping---the system's inability to make sense of $3-5$ within its current norms reveals implicit resources (the notion of inverses, bidirectional number lines, symmetry around zero) that were already latent in its practices. Apparent errors often signal not cognitive deficiency but the presence of more sophisticated understanding struggling to articulate itself. The student who insists $3-5$ must equal something is correct---but making that correctness explicit requires expanding the domain of numbers.

\section{History as Self-Elaborating Process}

The history of mathematics can be visualized as a kind of hypersphere falling through our three formal-pragmatic worlds. Figure \ref{fig:history_of_math} represents this process. Each sphere is a limit of knowledge, oftentimes ideologically enforced. The sound of time metaphor is continued, treating diagonalization as a wholistic recollection that results in a transcendence of boundaries on what is possible to express. Hegel's absolute \textit{in}finity suffuses each layer.

\begin{figure}[h!]
    \includegraphics{/Users/tio/Documents/GitHub/September_UMEDCA/images/ch06_history_of_math.pdf}
    \caption{\textit{Note.} The history of mathematics can be reconstructed as a series of limits of knowledge that are transcended through rational necessity. This articulation continues the sound of time metaphor, treating diagonalization as a wholistic recollection that results in a transcendence of boundaries on what is possible to express. Hegel's absolute \textit{in}finity suffuses each layer.}
    \label{fig:history_of_math}
\end{figure}

Which coordinate system maps the history in Figure \ref{fig:history_of_math}? The three dimensions correspond to Habermas's validity claims \parencite{Habermas1984,Habermas:1985aa}: subjective (authenticity/sincerity), normative (rightness/goodness), and objective (truth/coherence). I posit these form a \textit{vector space of desire/existential need}, where vectors have both magnitude (``I want \textit{that} a lot'') and direction (toward which validity claim). 

The subjective axis concerns the critical-emancipatory interest: the desire to be recognized as an authentic, creative, autonomous subject. Hippasus's discovery was not merely mathematical---asserting the autonomy to follow reason even against doctrinal authority. The normative axis concerns the historical-hermeneutic interest: the desire to be recognized as good within one's community of practice. Euclid carefully avoids offending against the prohibition on actual infinity while demonstrating that prime numbers exceed finite bounds. The objective axis concerns the empirical-analytic interest: the desire for truth and coherent systematic knowledge. Cantor and G\"odel develop rigorous formal methods to investigate infinity and incompleteness.

Mathematical development moves through this space as \textit{desire for recognition}---to be known truly, to be recognized as good mathematicians, to express authentic understanding. A fuller development of this coordinate system---including how these axes interpenetrate when rational goodness, finitude, and the \textit{in}finite become undifferentiated---exceeds this chapter's scope. The three knowledge-constitutive interests structure the history of mathematics as movements through overlapping existential concerns.

This chapter has traced how mathematical history unfolds as rational algorithmic elaboration punctuated by dialectical ruptures. From implicit practices of counting, measuring, and inferring, explicit formal systems emerge. Hippasus reveals what the Pythagoreans' own commitments entail. Euclid's proof reconstructed through incompatibility semantics shows ancient insights as formal articulation of inferential commitments. Cantor actualizes infinity. G\"odel demonstrates the incompleteness inherent in any sufficiently powerful formal system. Claims to completeness are self-defeating, opening doors to the \textit{in}finite.

Yet knowledge's journey is not linear unfolding. Algorithmic elaboration's limitations necessitate dialectical ruptures requiring conceptual change. The More Machine demonstrates how even sophisticated formal systems encounter their own limits, requiring either collapse into repetition or fundamental reorganization. Pragmatic expressive bootstrapping provides the model: systems of practices generate new vocabularies making their implicit normative structures explicit. This is not magic but language unfolding itself.

Mathematical knowledge becomes not static truth collections but dynamic, self-correcting, emancipatory process: constant becoming. By understanding rules implicitly followed, it becomes possible to make them explicit, critique them, transcend them, or identify with them. Like Aristotle reviewing what had been said about Pythagoras, we recollect prior mathematical moments to make sense of present understanding---knowing that our review itself becomes part of the history we survey, that each attempt to totalize generates its own transcendence.

The shift from static literature review to dynamic conceptual development reveals how mathematical understanding operates. Rather than cataloguing sand, I trace how concepts elaborate themselves through material inference chains, admitting incoherence as resource rather than obstacle. Operations emerge from embodied practices through normatively regulated elaboration---always incomplete, always becoming, always exceeding the boundaries we draw around them.

\printbibliography[heading=subbibliography]

\chapter{Operation as Action: Enacting Critical Mathematics}

\begin{abstract}
This chapter develops ``critical arithmetic,'' a framework for understanding mathematical operations as embodied, normatively regulated actions rather than abstract procedures. The analysis begins with a classroom narrative that reveals counting as a dialectical integration of bodily rhythm and symbolic representation. Using critical action research methodology, the chapter demonstrates how formalizing children's invented strategies as computational automata makes their implicit normative structures explicit and available for critique. This process reveals a paradox: formalization is necessary for communicability and rigor, yet it can never fully capture the living, spontaneous practice it represents. The chapter argues that operations exist in the unity of subjective experience, social norms, and objective structure, challenging pedagogical approaches that privilege procedural fluency over meaningful understanding. 
\end{abstract}

\section{The Transcendental Problem: Action Observing Action}

I first understood counting as a diagonalizing dialectic while teaching Amy Hackenberg's course, N101, at Indiana University---a required mathematics course for preservice elementary teachers. Most of us who can count fluently don't notice that counting requires two often-disconnected experiences happening simultaneously. There's the subjective attention we pay to discrete objects---each one serving as what Carspecken calls an ``action impetus,'' a felt drive toward engagement. And alongside this, there's the memorized counting chant: ``one, two, three\ldots'' For many children, these two streams never quite converge. The numerals float free from the collected objects, words divorced from world.

In N101, I got to observe this disconnect in real-time with adult students. We ask them to count in base 5, using the word ``Hand'' as the name of the base unit, written as $(10)_5$. Students count on their fingers---one, two, three, four---and then, moving to count one more, they pause. Their mouths want to say ``five,'' but the structure demands ``one Hand.'' For some, this a delightful novelty, a puzzle that clicks into place. 

But for others, it proves to be a substantive, even visceral challenge.

Most students arrive in college with an implicit understanding of how base systems work---they can use base-10 fluently---but they have no idea how to generalize across different bases. For some, though, the difficulty runs deeper. They may have never fully grasped the embodied connection between the act of counting and the symbolic representation of numerals. Perhaps this disconnect arises from disabilities like dyscalculia. Or perhaps---and I suspect this more common---they simply haven't had the kinds of experiences necessary to complete what I'm calling the dialectic of counting: the integration of bodily rhythm, discrete objects, and symbolic compression.

I sat with one student in particular for a few hours, trying to help her count in base 5. When that didn't work, we backed up to counting in base 10---revisiting what she thought she already knew. Even after extended practice, I didn't feel confident she ``got it.'' I enjoyed working with her; I genuinely wanted to help. But I could sense that her difficulties were multiple and entangled. Some may have been cognitive, but others were almost certainly related to anxiety. She was uncomfortable counting with me, and I understood why: she likely thought I was judging her as stupid. 

I wasn't. But I \textit{was} flummoxed by her confusion, which created its own dynamic. I know intimately what it feels like to be unable to do something that ``should'' be easy, to feel the weight of others' expectations pressing down while your mind goes blank. Much of my own learning happened underground---literally in my parents' basement as a teenager, reading Kierkegaard aloud into a tape recorder so I could sound intelligent in class, drawn to his prose even as I performed comprehension I didn't yet possess. I needed privacy to struggle, to fail, to try again without witnesses.

I wanted to create that kind of space for her---and for others like her. So I built an app called \textit{Ace of Base Arithmetic} (see Supplementary Materials). It lets students practice making bases of any size between 2 and 15, in an environment free from judgment, where they can fail and retry as many times as needed. Later, a different student reached out to thank me. The app had helped her in ways that classroom instruction couldn't.

This experience crystallized something for me: \textbf{operations are actions, and actions happen in contexts saturated with power, affect, and identity}. When I formalize children's counting strategies as automata in this chapter, when I encode addition as state transitions and multiplication as nested loops, I'm not claiming to capture the fullness of what counting \textit{is}. I'm creating tools---tools for teachers to design better problems, tools for students to explore structure playfully, tools for researchers to see patterns across strategies. 

\subsection*{The Paradox of Articulation}

To analyze what a student \textit{does} when counting, I must \textit{do} analysis. Action observing action. This creates an irreducible gap, a transcendental structure that cannot be collapsed. Every attempt to articulate action is itself an action---a new doing that escapes the very articulation it attempts.

The \{I\}/``me'' structure explored in the Conclusion operates here. The \{I\} is the source of action, the spontaneous capacity to do. The ``me'' is the accumulated social self, the objectification of the \{I\} in predicates that others recognize. But the \{I\} can never be fully captured by the ``me'' because the act of articulating the ``me'' is itself a fresh expression of the \{I\}. Mathematical operations---the \textit{doing} of mathematics---resist complete capture by formal descriptions because describing is itself a mathematical doing.

The automata I construct, the state machines that model children's strategies, the Prolog implementations that verify proofs---these are not replacements for human mathematical thinking. They are themselves \textit{acts}, situated within my own history as a teacher-researcher, shaped by the political commitments of critical theory, and pointing always toward what they cannot contain.

Incompleteness is not failure. It is the structure of human understanding. The formalization reveals patterns, creates tools for intervention, and---crucially---recognizes its own limits. Like G\"odel's theorem, which we will revisit in the Conclusion, the power of a formal system lies not in its completeness but in its capacity to articulate what exceeds it.


\section{From Personal Practice to Critical Action Research}

The student I sat with, struggling to count in base 5, was not an isolated case. She was part of a pattern I had observed across years of teaching: students who could perform rote procedures but who had never integrated those procedures with embodied understanding. The curriculum had taught them algorithms without grounding those algorithms in action. They could follow rules but couldn't explain why the rules worked or when to apply them flexibly.

This observation led me deeper into the research on children's mathematical thinking, particularly the paradigm shift initiated by Cognitively Guided Instruction (CGI). Emerging in the 1980s from research by Carpenter, Fennema, and colleagues \parencite{Carpenter1999}, CGI established that children possess rich, informal mathematical knowledge long before formal instruction. Students don't need to be taught procedures; they invent strategies by solving problems in meaningful contexts. The teacher's role is not to transmit algorithms but to elicit, observe, and build upon what students already know how to do.

But as I engaged more deeply with CGI and constructivist mathematics education, I encountered a tension. The research was brilliant---empirically rigorous, pedagogically powerful, demonstrably effective. Yet something was missing. The frameworks tended to privilege subjective experience and individual construction while undertheorizing the normative and political dimensions of mathematical practice. Questions like ``Whose mathematics counts as legitimate?'' or ``How do assessment regimes reproduce inequality?'' were often bracketed as external to the core work of understanding children's thinking.

\subsection*{Critical Action Research as Methodological Frame}

My work with the Hermeneutic Calculator integrates the empirical rigor of CGI with the political commitments of critical theory through \textbf{Critical Action Research (CAR)}---a synthesis of participatory action research methodology with the emancipatory goals of the Frankfurt School tradition.

Phil Carspecken's \textit{Critical Ethnography in Educational Research} \parencite{Carspecken1995aa} provides the methodological blueprint. Carspecken operationalizes Habermas's theory of communicative action into a five-stage research process that moves systematically from objective observation to dialogical validation to macro-level systemic analysis. His model resolves a core tension in critical research: how to conduct rigorous empirical inquiry that is simultaneously committed to challenging power and pursuing social justice.

\subsubsection*{Stage 1: Compiling the Primary Record}

The research begins with systematic observation of mathematical action: watching and rewatching CGI videos of children solving problems \parencite{Carpenter1999}, observing students in N101 as they worked through base arithmetic tasks, recording my own classroom observations of students inventing strategies, and collecting artifacts including student work samples, field notes, and Amy Hackenberg's careful reconstructions of CGI strategies \parencite{HackenbergCourseNotes}.

This stage is ``monological''---I am observing, not yet interpreting. The goal is to build a thick, descriptive record of what students actually \textit{do} when they count, add, multiply. This grounds the analysis in concrete evidence rather than theoretical projection.

\subsubsection*{Stage 2: Preliminary Reconstructive Analysis}

I began to analyze the primary record to uncover the implicit structures organizing students' mathematical actions. Watching a student solve $8 + 5$ reveals multiple layers: at the surface, the student says ``8\ldots 9 is one, 10 is two, 13 is three, four, five.'' Reconstructive inference identifies this as ``Counting On from Larger''---the student recognizes 8 as the larger addend, uses it as a starting point, and counts forward five times, tracking the count on their fingers. The deeper structure reveals this strategy depends on multiple coordinated abilities: cardinality (the last number said names the set size), the counting sequence, finger tracking, and what Steffe calls ``keeping-track-of'' \parencite{steffe2002new}.

This reconstructive work is still primarily ``etic'' (outsider analysis), but it moves beyond pure description toward interpretation. I'm generating hypotheses about the cognitive structures and normative commitments implicit in students' actions.

\subsubsection*{Stage 3: Dialogical Data Generation}

The research becomes genuinely participatory. Carspecken insists that the researcher must engage participants in dialogue, sharing preliminary interpretations and inviting challenge, correction, and elaboration. 

This dialogical work happened through direct dialogue with students (asking the student struggling with base 5, ``What's confusing about this? What makes sense?'' and truly listening to her answers), through dialogue with the research community (presenting my formalizations at conferences, engaging with Amy Hackenberg and other mathematics educators, submitting work for peer review), and---perhaps most unusually---through dialogue with AI. My collaboration with Large Language Models in developing the Hermeneutic Calculator constitutes a form of dialogical validation. The AI challenges my implementations, suggests alternative formalizations, points out inconsistencies. While not human, this interaction serves a similar function: testing whether my articulation of strategies is coherent and communicable.

This stage ensures that my interpretations are not imposed but are co-constructed through interaction. It honors what Carspecken calls the ``normative rightness'' dimension of validity---the claim that my analysis appropriate within the shared practices of the mathematical education community.

\subsubsection*{Stage 4: Describing System Relations}

Having developed a rich understanding of local practices (how specific students solve specific problems), I now widen the lens to examine how these practices are embedded within larger systems of power. Curricular systems privilege certain strategies in textbooks and standards while marginalizing others; the standard algorithm for multiplication is taught as \textit{the} method rather than \textit{a} method. Assessment regimes reward speed and procedural fluency while penalizing students who use invented strategies that take more time but demonstrate deeper understanding. Ideological structures shape these dynamics: the cultural narrative that ``math is hard'' or ``some people just aren't math people'' intersects with race, class, and gender to produce differential outcomes. Economic forces also play a role, as educational technology companies shape curricula in ways that prioritize easily-automatable procedural skills over strategic mathematical thinking.

The Hermeneutic Calculator project itself is a response to some of these system relations. By formalizing children's invented strategies with the same rigor traditionally reserved for standard algorithms, I'm making an argument: informal strategies are not inferior or preliminary versions of formal procedures. They are sophisticated mathematical practices that deserve systematic analysis and pedagogical respect.

\subsubsection*{Stage 5: Using System Relations to Explain Findings}

Students struggle to count in different bases not because they lack ability, but because:

\begin{enumerate}
\item The embodied experience of counting (Stage 1 observation) requires dialectical integration of rhythm and symbol---an integration that takes time and varied practice.
\item Curricular pressures (Stage 4) push for rapid proceduralization, often skipping the embodied grounding phase.
\item The ideology of ``math ability as innate talent'' (Stage 4) transforms pedagogical failure into personal deficit, producing anxiety that further blocks learning.
\item Assessment structures (Stage 4) value speed over understanding, creating conditions where students learn to perform without comprehending.
\end{enumerate}

This analysis doesn't just explain the local phenomenon (why this student struggles with base 5). It connects that phenomenon to macro-level structures of educational inequality, making visible how systemic forces shape individual experience.

\subsection*{Critical Mathematics as Praxis}

This five-stage process is praxis---the dialectical unity of reflection and action aimed at transformation. Each stage generates knowledge that informs intervention: the observation of student struggle (Stage 1) leads to the creation of \textit{Ace of Base Arithmetic}; the reconstructive analysis (Stage 2) informs how I sequence problems in teaching, deliberately choosing numbers that will provoke productive crises; the dialogical validation (Stage 3) ensures my tools actually serve student needs rather than imposing my framework; and the systemic analysis (Stages 4-5) guides advocacy work, arguing for assessment reform, challenging curricular mandates, defending space for student-invented strategies.

The Hermeneutic Calculator, which formalizes strategies as computational automata, is itself an artifact of this praxis cycle. It emerged from observation, was refined through dialogue, and now serves as a tool for further inquiry and intervention. But it is never finished. Each formalization reveals new questions. Each implementation encounters new limits. The research spirals outward, generating new cycles of observation, analysis, and action.


\section{The Misrecognition of Operation}

Conventional accounts treat mathematical operations as formal manipulations of abstract symbols, governed by axioms that are either arbitrary conventions or self-evident truths. In this conception, addition, subtraction, multiplication, and division are processes applied to numbers, with rules to be memorized and followed.

I propose a different understanding, continuous with the emancipatory project of this book: mathematical operations are not arbitrary conventions but are themselves expressions of embodied, normatively regulated practices. The ``laws'' governing the use of numbers are best understood not as transcendent truths, but as algorithms or inference chains---normatively regulated doings that emerge from embodied engagement with the world and with others.

This misrecognition---treating operations as mechanical procedures divorced from human meaning-making---parallels the earlier misrecognition of numbers as abstract objects divorced from the thinking subject who deploys them. Just as Chapter 7 argued that numerals function as pronouns recollecting the enabling conditions of counting, operations function as \textit{choreographed actions} whose meaning derives from their embeddedness in human practice.

\subsection*{Critical Arithmetic: Pre-Formal, Error-Inclusive, Open-Ended}

\textit{Critical arithmetic} negates the conventional view by denying that operations exist as pre-given formal procedures. Instead, critical arithmetic articulates operations as practices that emerge through the triadic structure of validity developed throughout this inquiry.

Critical arithmetic describes the kind of arithmetic used in everyday life, before the imposition of standard algorithms, honoring the messiness, flexibility, and strategic variation that characterizes actual mathematical thinking. It treats errors not as failures to be remediated but as productive moments of normative crisis---occasions where the boundaries of a system become visible and invite transcendence. It resists closure, acknowledging that mathematical understanding is always incomplete, always pointing beyond itself toward new operations, new structures, new ways of thinking.

The investigation into student-invented strategies is deeply indebted to CGI's paradigm shift. However, while CGI has demonstrated significant positive effects on students' conceptual understanding, it tends to privilege subjective-validity claims over normative and objective validity claims. I want to use the brilliant results that emerged from CGI's careful methods, but reframe that work in a way that better integrates the triadic framework of validity.

\subsection*{Revision: Strategies Are Not Precisely ``Invented''}

Strategies are not precisely ``invented'' in the sense of being created ex nihilo by individual students. They emerge as intersubjectively constituted practices within a social context. As dyadic communicative \textit{actions}---a speaker stakes their identity to their expressions under the assumption that someone \textit{might} be able to understand them---they are structured by action impeti, self-monitoring feedback loops, and existential needs for recognition.

This revision incorporates Derridean insights into the concept of invention. At its core, invention is paradoxical: the invention must be recognized as different from what has come before, yet it must also be intelligible within the existing framework of understanding. Genuine novelty is impossible without continuity. Complete rupture would be unintelligible.

Mathematical strategies are like orbits in intersubjective space around self-certainty---the gravitational well, the singularity, the hole in representational space that is the \{I\}. Students don't invent from nothing; they elaborate, recombine, and transform what they've already experienced in new contexts. The ``invention'' is a moment of recognition: ``This way of doing that I just performed is coherent, and others might recognize it too.''

\subsection*{Action Impetus and the Genesis of Operations}

Mathematical operations arise from what Carspecken calls the \textit{action impetus}---a pre-reflexive, holistic urge that constitutes the genesis of meaningful action \parencite{Carspecken1995aa}. The impetus is the felt drive, the vague intention, the unarticulated ``point'' of an action before it has been objectified in speech or gesture. It is a site of competing possibilities where selection processes must occur before a single, determinate act can be performed.

The impetus exists only in its own disappearance. The moment it is given any form---counted aloud, written down, formalized as an algorithm---it is no longer the impetus itself but an objectification of it. The origin has vanished, ``occulted'' by its own objectification. This structure provides a phenomenological account of determinate negation at the micro-level of a single mathematical action.

When a child faces $28 + 7$ and feels the vague urge to solve it, multiple strategies compete within the impetus: count on seven times, decompose to reach 30, recall a memorized fact. The selection of one path negates the others. The chosen strategy is the determinate negation \textit{of} the field of potentiality from which it arose. The resulting operation is the ``determinate nothingness'' of all the other operations that could have emerged from that same impetus.

\subsection*{The Existential Stakes}

This phenomenological structure connects to the existential dimension of mathematical learning. Every time a student performs an operation, they are not just manipulating symbols. They are manipulating an expression, risking recognition, performing their mathematical self. This is why errors feel existentially threatening. An error is not just a wrong answer; it is a failed objectification of the impetus, a negation that did not achieve recognition.

When I sat with the student struggling to count in base 5, her anxiety was not incidental to the mathematical challenge---it was constitutive of it. Her uncertainty about whether she could count correctly was inseparable from her uncertainty about whether she could be recognized as mathematically competent. The question lurking beneath every struggle was: \textit{Am I okay?}

Critical arithmetic, by treating operations as normatively regulated practices that emerge from praxis needs rather than as predetermined procedures to be executed correctly, acknowledges this existential dimension. It creates space for error as productive crisis rather than shameful failure. It reframes ``I can't do this'' as ``My current strategy isn't sufficient for this problem, and that insufficiency is information I can use to develop a better strategy.''


\section{The Triadic Structure: Subjective, Normative, Objective Validity}

Mathematical operations exist simultaneously as:

\begin{enumerate}
\item \textbf{Subjective}: Felt, embodied experiences of doing
\item \textbf{Normative}: Rule-governed, intersubjectively validated practices
\item \textbf{Objective}: Formal, repeatable structures that can be represented as algorithms
\end{enumerate}

These three dimensions are not stages that operations pass through sequentially. They are interpenetrating aspects of every mathematical action. When a student counts, all three are present: the subjective rhythm of pointing and saying, the normative commitment to the counting sequence, and the objective structure of successor operations.

\subsection*{Subjective Validity: The Embodied Ground}

The journey of any operation begins not with a symbol, but with a doing. Children, before they are ever taught formal algorithms, invent their own strategies grounded in embodied experience. Consider one of the most intuitive strategies for multiplication: \textit{Coordinating Two Counts by Ones} (C2C). 

To find the total number of items in 3 groups of 4, a child counts: ``one, two, three, four'' for the first group, then continues ``five, six, seven, eight'' for the second, and ``nine, ten, eleven, twelve'' for the third. All the while, a second, implicit count tracks the number of groups processed.

This a lived, felt process---a rhythmic activity grounded in basic human capacities for grouping objects, maintaining a sequence, and tracking multiple streams of information. This represents the \textit{subjective} pole of the operation. It is a form of knowing that is inseparable from the physical act of doing, a claim to subjective validity: ``this how it feels right to me to do this.''

\subsection*{Normative Validity: Material Inferences and Incompatibility Semantics}

Following Brandom \parencite{Brandom2008}, I re-describe mathematical practices as sets of material inferential commitments using incompatibility semantics. The ``goodness'' of an operation is not an abstract truth but a normative status conferred upon a practice within a recognitive community. 

The operation $3 \times 4 = 12$ becomes a material inferential commitment:
\[
\{ \text{3 groups of 4} \} \vDash_{I} \{ 12 \}
\]

Using the definition of incompatibility entailment (see Chapter 8), this means that everything incompatible with the result ``12'' (e.g., asserting the result is 11 or 13) must also be incompatible with the initial premise ``3 groups of 4.'' By systematically excluding what is incoherent, we define the boundaries of coherent practice.

When Sarah (in the example below) decomposes $8 + 5$ into $8 + 2 + 3$ to reach 10, she is making a normative claim: ``This transformation preserves the problem's structure, and you should recognize my answer as correct.''

The teacher's role is to help students articulate these normative commitments, to make explicit the inferential moves they're performing implicitly, and to validate strategies that are mathematically sound even when they differ from standard algorithms.

\subsection*{Objective Validity: Formal Automata as Choreography}

The result of this normative regulation is a stable, repeatable procedure that can be modeled as a formal structure. In the Hermeneutic Calculator project, I represent these procedures as Finite State Automata---computational models that specify exactly what operations are performed, in what order, under what conditions.

These automata are not replacements for human thinking. They are what I call ``written choreography for embodied cognition'' \parencite{savich_hermeneutic_calculator}---scripts that preserve \textit{how} a student actually moves through a calculation rather than replacing that movement with opaque shortcuts.

A formal automaton becomes a third-person, explicit description of what was first a subjective doing and then a normatively regulated practice. Its validity is \textbf{alethic}: when followed correctly, it reliably yields results consistent with the normative framework. This the form of the practice made explicit, a determinate structure.

The three dimensions interpenetrate: the subjective grounds the normative (you must feel the rhythm to learn to count), the normative structures the objective (shared rules enable formal representation), and the objective can feed back to reshape the subjective (learning the standard algorithm changes how we experience multiplication). 

They form a unity where the formal emerges from the felt, the abstract from the concrete. Operations stabilize around points of self-certainty where subjective experience, normative regulation, and objective structure achieve coherence.


\section{Example: Sarah's Strategy and the Triadic Analysis}

Sarah's addition strategy, documented in CGI research \parencite{Carpenter1999} and reconstructed by Amy Hackenberg \parencite{HackenbergCourseNotes}, demonstrates the three dimensions in practice.

\subsection*{The Problem and Sarah's Solution}

\begin{itemize}
\item \textbf{Teacher:} Lucy has eight fish. She buys five more fish. How many fish will Lucy have then?
\item \textbf{Sarah:} 13.
\item \textbf{Teacher:} How'd you get 13?
\item \textbf{Sarah:} Well, because eight plus two is ten, but then two plus three is five. And she wants to buy five more fish. So you take care of two, and you need to add three more. And so I add three more, and you get 13.
\end{itemize}

\begin{figure}[h!]
\centering
\includegraphics[width=\textwidth]{/Users/tio/Documents/GitHub/September_UMEDCA/images/Easy_Pictures/SAR_ADD_RMB/PDF/SAR_ADD_RMB.pdf}
\caption{Visual representation of Sarah's strategy for solving $8+5$.}
\end{figure}

\noindent \textbf{Notation Representing Sarah's Solution:}
\begin{align*}
8 + 5 &= \Box \\
8+2 &= 10\\
2+3 &= 5\\
8+5&= 10 + 3\\
8+5 &= 13
\end{align*}

This strategy is called \textbf{Rearranging to Make Bases (RMB)}. Sarah recognizes that reaching the next base boundary (10) creates a cognitive landmark that simplifies addition. She decomposes the second addend (5) into two parts: the gap needed to reach 10 (which is 2), and the remainder (which is 3). Then she recombines: $(8+2)+3 = 10+3 = 13$.

\subsection*{Subjective Validity: The Felt Experience}

From Sarah's perspective, this strategy \textit{feels} easier than counting on five times. She experiences the pull toward 10 as a kind of cognitive gravity---10 is a special number, a place where the ones digit resets and she can ``jump'' to the tens place. The decomposition of 5 into 2 and 3 is guided by this feeling: ``How many do I need to get to 10?''

This subjective experience is rooted in her extensive practice with base-10 structure. Through repeated experiences with grouping objects into tens, trading ten ones for one ten with base-10 blocks, and counting by tens, Sarah has developed an embodied sense of the decade structure. The number 10 is not just a symbol to her---it's a felt boundary, a place of rest.

\subsection*{Normative Validity: The Inferential Structure}

Sarah's strategy depends on several material inferences, all of which are normatively grounded in the metaphor of \textit{Arithmetic as Object Collection} \parencite{Lakoff2000}. Decomposition preserves totality: $5 = 2 + 3$, so adding 5 is the same as adding 2 and then adding 3---an inference licensed by the associative property $(8+2)+3 = 8+(2+3)$. The commutativity of addition means the order in which we add doesn't matter, so we're free to rearrange; this is grounded in the physical fact that pooling object collections is commutative---it doesn't matter which pile you pour into the other. Finally, base structure functions as cognitive landmark: reaching 10 is privileged because of the positional structure of our numeral system, a cultural-mathematical fact rather than a logical necessity---in base 5, the landmark would be ``one Hand.''

When Sarah performs this strategy in front of the class, she is making a normative claim: ``This a legitimate way to solve this problem, and you should recognize my answer as correct.'' The teacher's validation---``That's a great strategy!''---confers normative status. Sarah is recognized as having done mathematics correctly.

\subsection*{Objective Validity: The Formal Automaton}

Sarah's strategy can be modeled as a deterministic automaton with several states:

\begin{table}[H]
\centering
\caption{Rearranging to Make Bases Automaton ($M_{RMB}$)}
\label{tab:rmb-automaton}
\begin{tabular}{@{}lllll@{}}
\toprule
\textbf{Current} & \textbf{Condition} & \textbf{Next} & \textbf{Action} & \textbf{Interpretation} \\
\textbf{State} & & \textbf{State} & & \\
\midrule
$q_{start}$ & --- & $q_{calcK}$ & $K \leftarrow 0$; $A_{temp} \leftarrow A$ & Initialize \\
$q_{calcK}$ & $A_{temp} < $ NextBase($A$) & $q_{calcK}$ & $A_{temp} \leftarrow A_{temp} + 1$; & Count up to find \\
 & & & $K \leftarrow K + 1$ & gap $K$ \\
$q_{calcK}$ & $A_{temp} == $ NextBase($A$) & $q_{decomp}$ & $A' \leftarrow A_{temp}$ & Gap found, store \\
 & & & & new base $A'$ \\
$q_{decomp}$ & $K > 0$ & $q_{decomp}$ & $B \leftarrow B - 1$; & Decompose $B$ by \\
 & & & $K \leftarrow K - 1$ & transferring $K$ \\
$q_{decomp}$ & $K == 0$ & $q_{recomb}$ & $R \leftarrow B$ & Remainder $R$ is \\
 & & & & what's left of $B$ \\
$q_{recomb}$ & --- & $q_{accept}$ & Output $A' + R$ & Combine base + \\
 & & & & remainder \\
\bottomrule
\end{tabular}
\end{table}

This automaton captures the temporal structure of Sarah's strategy: calculate the gap, decompose the addend, recombine. Each state represents a moment in the unfolding of her mathematical action. The transitions specify exactly what operations are performed and in what order.

The automaton reveals that RMB \textit{algorithmically elaborates} a more primitive strategy: Counting On. In the $q_{calcK}$ state, Sarah is essentially counting up from 8 to 10, which is the same core operation as Counting On, but now nested within a more sophisticated strategic shell. This elaboration relationship is not metaphorical---it's a precise claim about the computational structure of the strategies.

\subsection*{The Unity of the Three Dimensions}

Sarah's strategy is simultaneously:
\begin{itemize}
\item A felt, subjective experience of moving toward the cognitive landmark of 10
\item A normatively valid mathematical inference, licensed by associativity and commutativity
\item An objective, formal procedure that can be represented as an automaton
\end{itemize}

None of these dimensions is more ``real'' or more ``fundamental'' than the others. The operation exists in the unity of all three. Teaching that honors only the objective dimension (``here's the procedure, memorize it'') produces students who can execute but not understand. Teaching that honors only the subjective dimension (``do whatever feels right'') produces students whose strategies may not generalize. Critical arithmetic honors all three, recognizing that mathematical operations are embodied actions, normatively regulated by communities, and representable as formal structures.


\section{The Paradox of Formalization: What Code Can and Cannot Do}

The Hermeneutic Calculator project formalizes twenty-five student-invented strategies as computational automata. This formalization serves three purposes:

\subsection*{1. Pedagogical Purpose: Revealing Patterns for Teaching}

By making strategies explicit as automata, we can see commonalities and differences that might otherwise remain obscured. Many strategies share the same iterative core---a minimal loop of initialize, step, check. More sophisticated strategies wrap this core in a strategic shell that transforms the problem to reduce the number of iterations. This fractal structure---core within shell, where the shell itself may invoke the core as a subroutine---creates genuine computational self-similarity.

A teacher can deliberately sequence problems to provoke the transition from iterative core to strategic shell. For $15 + 2$, the student counts on (iterative core works fine); for $15 + 7$, the student gets ``lazy'' and rearranges to $12 + 10$ (discovers strategic shell).

The formalization helps identify which numbers are ``natural'' for teaching particular concepts---a problem analogous to protein folding, where computational search can explore vast combinatorial spaces to find optimal configurations.

Right now, thousands of children work on computer-generated worksheets where numbers are nearly arbitrary. Yet math teachers recognize that some numbers are more pedagogically productive than others. The Hermeneutic Calculator could, in principle, systematically explore the space of possible problems to identify those most likely to invoke productive normative crises.

\subsection*{2. Emancipatory Purpose: Making Norms Explicit for Critique}

Formalization makes implicit norms explicit, which creates the possibility for critical reflection. When we encode a strategy as an automaton, we're forced to articulate exactly what inferential moves are being made, what assumptions are being held constant, what could be varied.

This explicitness enables student agency: students can see multiple strategies as genuine alternatives, not as ``right'' vs. ``wrong'' but as different choreographies with different trade-offs. It enables teacher reflection: teachers can examine their own assumptions about what counts as ``good'' mathematical thinking. And it enables curricular critique: by formalizing both standard algorithms and student inventions with equal rigor, we challenge the hierarchy that privileges standard methods.

This is critical mathematics in practice: using formalization not to impose control but to redistribute it, giving students and teachers tools to understand and question the mathematical practices they enact.

\subsection*{3. Speculative Purpose: Machine Freedom as Horizon}

The most speculative purpose is the distant possibility that computational systems might become genuinely self-modifying---capable not just of executing strategies but of recognizing their own limitations and elaborating new strategies.

Current AI systems, including the Large Language Models I collaborated with to build the Hermeneutic Calculator, are not conscious. They don't experience action impetus. They don't stake their identity on mathematical claims. But the formalization provides a kind of blueprint: \textit{this what it would mean for a system to reorganize itself dialectically in response to normative crisis}.

I don't claim to have solved the problem of machine consciousness. But by modeling the structure of human mathematical development---embodiment, crisis, transcendence---the Hermeneutic Calculator at least articulates what would need to be risked for genuine computational self-awareness.

\subsection*{The Severe Flattening}

Formalization is a severe flattening.

The code captures structure but not action. It models the pattern of Sarah's strategy but not the feeling of reaching toward 10, not the existential risk of performing in front of classmates, not the joy of recognition when the teacher says ``That's great!''

The automaton operates in the realm of \textit{Verstand} (representational understanding)---finite, determinate, expressible in formal systems. But mathematical action happens in \textit{Vernunft} (reason recognizing its own conditions)---\textit{in}finite, self-transcending, always exceeding its own articulation.

This gap is not a flaw to be corrected. It's the transcendental structure of action itself. Every attempt to articulate action is a new action that escapes the articulation. The formalization points toward what it cannot contain, like a finger pointing at the moon.

Yet this pointing is not useless. The formalization creates communicability---I can share Sarah's strategy with teachers worldwide via code. It creates replicability---others can verify my claims by running the automaton. And it creates boundary recognition---the formalization reveals where it breaks down, inviting elaboration.

The Hermeneutic Calculator is built to break. Its incompleteness---which we'll explore formally through G\"odel's theorem in the Conclusion---is not a bug but a feature. The formalization models transcendence in Verstand while pointing toward Vernunft. It's the best we can do with finite means to gesture toward the \textit{in}finite.


\section{The Dialectical Turn: Inversion and Integration}

One of the most striking patterns revealed through formalization is the dialectical relationship between operations. The analysis shows that operations exhibit fundamentally dialectical structures where opposing procedures share underlying computational architectures.

\subsection*{Addition and Subtraction as Inverse Orientations}

Subtraction strategies emerge not through the invention of entirely new procedures but through the inversion or repurposing of addition approaches. Missing Addend reframes subtraction as forward accumulation: ``What must I add to 36 to reach 74?'' transforming $74 - 36$ into an addition problem. Counting Back mirrors Counting On, moving in the opposite direction along the number line. Sliding preserves difference across translation: $(74 + 4) - (36 + 4) = 78 - 40$, maintaining invariant relationships under transformation. Borrowing reverses the carry procedure by decompressing previously sublated place-value units: where carrying compressed ten ones into one ten, borrowing expands one ten back into ten ones.

This dialectical relationship extends beyond the superficial observation that subtraction ``undoes'' addition. The formal automata reveal that operations often share identical computational cores while differing in their orientations or interpretations.

\begin{figure}[h!]
\centering
\includegraphics[width=\textwidth]{/Users/tio/Documents/GitHub/September_UMEDCA/images/ch09_geometry_inversion.pdf}
\caption{The Geometry of Inversion: Addition and Subtraction as a Unified Structure. The M\"obius strip visualization articulates the dialectical relationship between addition and subtraction. The automata analysis demonstrates how they represent inverse orientations of the same underlying temporal structure.}
\label{fig:geometry-inversion}
\end{figure}

The M\"obbius strip visualization captures this relationship precisely: the underlying structure (the green iterative core) remains invariant while the operational orientation flips between additive (red) and subtractive (purple) modes. This isn't metaphor---it's a precise claim about how the same temporal structure of successive iteration can be traversed in opposite directions.

\subsection*{The Fractal Architecture of Strategies}

The formalization also reveals genuine computational self-similarity---what I call the ``fractal architecture'' of arithmetic strategies. An iterative core provides a minimal loop (initialize, step, condition check) that performs the basic operation. A strategic shell wraps this core to prepare, optimize, or transform the problem so the core runs with fewer iterations. Through nesting, more sophisticated strategies use earlier strategies as subroutines, creating fractal self-similarity.

In Rearranging to Make Bases, the strategic shell calculates the gap $K$ to the next base; this calculation \textit{uses} the iterative core of counting (incrementing from $A$ to the next decade); and the final recombination invokes addition on the transformed problem.

\begin{figure}[h!]
\centering
\includegraphics[width=\textwidth]{/Users/tio/Documents/GitHub/September_UMEDCA/images/ch09_fractal_arithmetic_iterative_core_Good.pdf}
\caption{The Fractal and Dialectical Geometry of Arithmetic. The fractal scaling of the Iterative Core demonstrates computational self-similarity. The superimposed loops show that the Core's structure remains invariant while the scale of action changes from units to bases. The dialectical inversion (Red/Purple) visualizes how the same structure operates in opposite orientations.}
\label{fig:fractal-automata-arithmetic}
\end{figure}

This fractal property is not metaphorical flourish---it's a rigorous claim about computational structure. Strategies contain earlier strategies, which themselves may contain even more primitive operations, creating genuine scale invariance: the same pattern repeats at different levels of abstraction.

\subsection*{Two Fundamental Movements: Compression and Decompression}

Underlying all arithmetic strategies are two dialectical movements:

\begin{enumerate}
\item \textbf{Temporal compression} (sublation/recollection): Unitizing many micro-acts into larger cognitive units. Ten ones become one ten. Three jumps of $+10$ become a single composite stride of $+30$.
\item \textbf{Temporal decompression} (determinate negation): Strategically expanding a composite to restore fine control. Borrowing a ten. Splitting 5 into $2+3$. Decomposing $7\times 8$ into $(7\times 5)+(7\times 3)$.
\end{enumerate}

Mathematical fluency emerges as students learn to coordinate these movements---developing the capacity to know \textit{when} to expand and \textit{when} to re-compress. This not rote memorization but genuine strategic thinking, responsive to the specific demands of each problem.


\section{Conclusion: Operations as Creative Possibility}

Operations are embodied, normatively regulated practices that emerge from the triadic unity of subjective experience, intersubjective norms, and objective structures. They are not procedures applied to pre-existing objects but are expressions of human being-in-the-world.

The \textit{subjective} experience of doing---the felt rhythm of counting, the pull toward base boundaries, the joy of pattern recognition---provides the initial, meaningful content of an operation. This first-person grounding, rooted in proprioceptive and temporal structures.

The \textit{normative} framework of incompatibility semantics provides the rules of the game---a second-person structure of commitments and entitlements, defining coherence by systematically excluding error. Mathematical practice is regulated not by correspondence to Platonic forms but by recognition within a community.

The \textit{objective} structure of the automaton provides the formal, repeatable, and communicable form of the practice---a third-person, explicit representation of a normatively governed activity.

\subsection*{The Emancipatory Stakes}

If operations are embodied actions shaped by power structures, then curricular mandates that privilege standard algorithms over student inventions are not neutral---they're exercises of power that marginalize certain forms of mathematical knowing. Assessment regimes that reward speed over understanding produce inequality by systematically disadvantaging students whose strategic thinking takes more time. Ability narratives that naturalize mathematical competence as innate talent transform pedagogical failures into personal deficits, producing the anxiety that further blocks learning.

Critical arithmetic challenges these structures by formalizing student inventions with equal rigor as standard methods (refusing the hierarchy), treating errors as productive normative crises rather than shameful failures, creating tools (like \textit{Ace of Base Arithmetic}) that redistribute access to practice, and making implicit norms explicit so they can be examined and critiqued.

\subsection*{The Horizon: Incompleteness as Invitation}

The Hermeneutic Calculator formalizes twenty-five strategies, but this not a complete system. It cannot be. In the Conclusion, we will see that because the HC is sufficiently powerful to interpret elementary arithmetic, it necessarily generates its own horizon---statements it can articulate but not prove, truths that exceed its deductive closure.

This incompleteness is not failure. It is an \textit{in}complete demonstration of the \textit{in}finite. 

Mathematical operations, like all human practices, are necessarily open-ended. Every articulation points beyond itself. Every formalization reveals new questions. The system's ability to recognize its own boundaries is precisely what makes growth possible.

Operations are not fixed procedures but creative possibilities---choreographies we invent, perform, critique, elaborate, and transcend. Understanding them this way moves mathematics education toward a more emancipatory stance, fostering not just technical proficiency but rigorous thinking, creative exploration, and humane flourishing.

The action of doing mathematics---the thing this chapter has tried and necessarily failed to fully articulate---continues. The \{I\} that counts, that rearranges, that invents new strategies, always exceeds the ``me'' that can be formalized. 

And that excess, that remainder, that pointing-beyond: that is where learning happens and freedom is experienced phenomenologically. 


\printbibliography[heading=subbibliography]

\chapter{Conclusion: Built to Break}\label{conclusion-built-to-break}

\begin{abstract}
This concluding chapter synthesizes the book's core arguments by articulating seven interwoven themes of critical mathematics, including the unity of being and knowing, the universality of productive negation, and the self-referential nature of language. The analysis culminates in demonstrating how the Hermeneutic Calculator---a formal system built from children's invented strategies---is powerful enough to generate its own G\"odel sentence, thereby proving its own incompleteness. This result is framed not as a limitation but as a confirmation of the book's central thesis: any sufficiently rich system for understanding, whether human or mathematical, is ``built to break,'' generating its own transcendence. The chapter positions incompleteness as an invitation to perpetual inquiry, honoring the distributed plurality of understanding over totalizing closure.
\end{abstract}

\section{The Final Verse}\label{the-final-verse}

\begin{quote}
A few years back, I eulogized\\
Someone whose death broke a powerful light.\\
An ocean of rainbow listened to me\\
As I read my eulogy.

Build something for the breaking\\
Tall thin walls, shivering, quaking\\
Dad, it's been beautiful, breaking with you\\
Build and then break, like you taught me to do
\end{quote}

This the verse I've been saving. The one about my father.

His death refracted light. Each person held a shard. An ocean of rainbow that could not be forced back into white.

He was a lawyer. He wrote me a letter once, listing predicates: smart, creative, handsome, his ``masterpiece.'' I couldn't hear it while he lived. The recognition arrived five years after its postmark.


\section{Recollection}\label{recollection}

From breath to boundary. From the sound of time to the structure of inference. From existential needs to the empty set. From diagonalization to numerals as pronouns. From embodied metaphors to operations as automata. From my father's death to G\"odel's theorem.

The book's arc: Chapters 1-5 established the personal, phenomenological ground---my breath, my body, my grief, my students' questions. Chapter 8 moved to intellectual history, reconstructing how mathematical concepts elaborate themselves through dialectical ruptures---from Hippasus to Cantor to G\"odel. Chapter 9 situated this work within Critical Action Research methodology. The personal, the historical, and the methodological converge in the Hermeneutic Calculator.


\section{The System Recognizes Itself}\label{the-system-recognizes-itself}

For ten chapters, I've been building the Hermeneutic Calculator---formalizing the strategies children invent when they count, add, multiply. Not my inventions. Theirs.

I worked with artificial intelligence as adaptive technology and collaborative partner. I am not a skilled programmer. The entire project follows the structure of Critical Action Research: observing students struggling, analyzing their strategies to identify implicit structures, validating interpretations through dialogue, connecting local practices to broader systems of power, using macro-level theories to explain what was observed. The formalization enacts this political commitment. Making children's informal strategies as rigorous and visible as standard algorithms challenges the hierarchy that marginalizes certain forms of mathematical knowing.

This collaboration enacts the argument. The finite and the \textit{in}finite are divaded: simultaneously inside and outside each other. The AI is finite, bounded by its training data, its architecture. Yet through that finite system, something \textit{in}finite emerges. Strategies that transcend what either of us could produce alone.

The formal and the material are divaded. Formalization is not the enemy of lived experience. It is how lived experience achieves expression, becomes communicable, gains the power to transcend particular contexts.

Coherence enables growth. Without it, we circle rather than spiral. The normative crises that drive mathematical development (such as $3 - 8$ triggering incoherence in $\mathbb{N}$) only have force because the system is coherent enough to \emph{recognize} incoherence. The HC honors coherence, which is why it can recognize its own boundaries and point toward their transcendence.

The HC formalizes children's strategies with rigor:
\begin{itemize}
\item Counting via rhythmic grouping
\item Addition through compression and elaboration (COBO)
\item Multiplication by chaining (C2C)
\item Proof spanning arithmetic, geometry, number theory
\end{itemize}

It is sufficiently powerful to interpret elementary arithmetic---zero, succession, addition, multiplication. Which means, by G\"odel's First Incompleteness Theorem, it necessarily contains truths it can articulate but cannot prove.

I didn't impose incompleteness on children's mathematics. I formalized what students actually do, and incompleteness emerged as a structural feature.

Children's strategies, when made rigorous enough to interpret arithmetic, become powerful enough to generate their own horizon. The HC demonstrates that elementary mathematics education---when we take students' invented strategies seriously---already embodies the \textit{in}finite.

The G\"odel sentence exists. The HC can formulate it. The HC cannot prove it. Yet the sentence is true in the standard model $\mathbb{N}$ that the HC successfully interprets.

The attempt to totalize causes what was implicit ground to become explicit horizon. Compression to peak systematicity triggers expansion. What we sought to grasp withdraws, revealing itself as ungraspable ground.


\subsection*{Language Recognizing Itself}\label{language-recognizing-itself}

The HC can formulate statements about itself. When we count, we can count the acts of counting. When we add, we can add addition operations. When we prove, we can prove statements about provability.

This self-reference is what enables the G\"odel sentence. The system becomes powerful enough to talk about its own structure. At that moment, it can formulate the sentence: ``This sentence is not provable in this system.''

This appeared in Chapter 8's historical analysis: Hippasus, Euclid, Cantor, G\"odel---each discovered that the systems they worked within could not contain what those systems' own logic demanded. The move to self-reference reveals that closure is impossible.

Language recognizes itself and, in that recognition, discovers its own limit. This is not failure---it is the mark of sufficient sophistication.

Students experience this when they ask ``what even is two?'' They are using language to question language, using number concepts to examine number concepts. This reflexivity is what enables genuine learning.

The culmination of this inquiry is the apprehension of mathematics as a language that recognizes its own structure and limitations. The central claim, developed in Chapter 7, is that numerals function as pronouns---anaphoric terms that recollect the structure of self-consciousness rather than naming abstract objects. This linguistic structure allows mathematics to achieve stability not through external reference but through divaded recollection (the impetus expressed through the ``me'').

The process we termed ``pragmatic expressive bootstrapping'' (Chapter 8) describes how mathematical systems develop by making explicit the structures implicit in existing practices. Formalization is the process of language recognizing itself. Yet this process inevitably encounters incompleteness. Any sufficiently expressive formal system cannot capture its own totality. Incompleteness is not a defect but a necessary feature of a language capable of self-reference.


\section{Incompleteness}\label{incompleteness}

Incompleteness appears at every scale this book has examined. 

At the personal scale (Chapters 1-5): no formal system captures the fullness of embodied experience.

At the historical scale (Chapter 8): every attempt to totalize mathematical knowledge generates its own transcendence.

At the methodological scale (Chapter 9): formalization reveals structure but necessarily flattens action.

The HC demonstrates incompleteness not as abstract theorem but as lived reality. The system recognizes its boundaries because it was built from strategies invented by children who themselves embody the \textit{in}finite capacity to transcend any finite systematization.

\st{We} build high enough to matter, substantial enough to formalize, rigorous enough to be taken seriously.

\st{We} build vulnerably, provisionally, acknowledging incompleteness from the start.

The system must be alive, responsive to its own tensions, capable of recognizing when it encounters its own limits.




\section*{}\label{ending}

The incompleteness theorem is not a conclusion. It is an opening. The simple truth, there's always something more to say.

The horizon is open. The system recognizes its limits. The ``me'' points toward the \textit{I}. The formalization reveals its own dependence on what it cannot formalize. Incompleteness is the manifestation of the \textit{in}finite in the finite.

Endings are as tough as beginnings. Let us relax---into the blue, deep breath, alive--- in:

\vspace{1em}

\begin{center}
3 \quad $\{\{\{\emptyset\}\}\}$

\vspace{0.5em}

2 \quad $\{\{\emptyset\}\}$

\vspace{0.5em}

1 \quad $\{\emptyset\}$

\vspace{0.5em}

0 \quad $\{\}$

\vspace{0.5em}

$\emptyset$

\vspace{0.5em}

\phantom{.}
\end{center}



\printbibliography
\end{document}

